%!TEX program = xelatex
\documentclass[11pt,a4paper]{article}
\usepackage[margin=2.2cm]{geometry}
\usepackage{amsmath,amssymb,amsthm,mathtools}
\usepackage{bm}
\usepackage{enumitem}
\usepackage{booktabs}
\usepackage{array}
\usepackage{xcolor}
\usepackage{tikz}
\usetikzlibrary{arrows.meta,positioning}
\newcommand{\mserr}[1]{\textcolor{red}{#1}}
\newcommand{\errpanel}[1]{%
\begin{center}
\fcolorbox{red!75!black}{red!8}{%
\begin{minipage}{0.94\linewidth}
\setlength{\parskip}{0.45em plus 0.1em}
#1
\end{minipage}}
\end{center}
}
\usepackage{hyperref}
\usepackage{fancyhdr}
\usepackage{xeCJK}

\hypersetup{
    colorlinks=true,
    linkcolor=blue!60!black,
    citecolor=blue!60!black,
    urlcolor=blue!60!black
}

\pagestyle{fancy}
\fancyhf{}
\rhead{\small Qiuzhen QE Computational \& Applied Math --- Spring 2026}
\lhead{\small Official Qualifying Exam}
\rfoot{\small Page \thepage}

\DeclareMathOperator*{\argmax}{arg\,max}
\DeclareMathOperator*{\argmin}{arg\,min}
\DeclareMathOperator{\rank}{rank}
\DeclareMathOperator{\tr}{tr}

\setlist[itemize]{topsep=3pt,itemsep=2pt,parsep=0pt}
\setlist[enumerate]{topsep=3pt,itemsep=2pt,parsep=0pt}

\newcommand{\examstatement}[1]{%
  \par\addvspace{0.85em}%
  \noindent{\bfseries #1}\par\nobreak\addvspace{0.28em}%
}

% Shared amsthm note/remark/theorem styles
% Shared math extras for qe-review.
% Requires already-loaded: amsmath, amssymb.
%
% Classic pitfall: AMS blackboard-bold fonts have no digit glyphs, so
% \mathbb{1} renders as overlapping garbage. Map the digit 1 to dsfont's
% \mathds{1}; leave \mathbb{R}, \mathbb{P}, etc. on AMS.
%
% Usage (path relative to the main .tex being compiled):
%   notes:  % Shared math extras for qe-review.
% Requires already-loaded: amsmath, amssymb.
%
% Classic pitfall: AMS blackboard-bold fonts have no digit glyphs, so
% \mathbb{1} renders as overlapping garbage. Map the digit 1 to dsfont's
% \mathds{1}; leave \mathbb{R}, \mathbb{P}, etc. on AMS.
%
% Usage (path relative to the main .tex being compiled):
%   notes:  % Shared math extras for qe-review.
% Requires already-loaded: amsmath, amssymb.
%
% Classic pitfall: AMS blackboard-bold fonts have no digit glyphs, so
% \mathbb{1} renders as overlapping garbage. Map the digit 1 to dsfont's
% \mathds{1}; leave \mathbb{R}, \mathbb{P}, etc. on AMS.
%
% Usage (path relative to the main .tex being compiled):
%   notes:  \input{../../common/qe-math-extras.tex}
%   exams:  \input{../../../common/qe-math-extras.tex}

\usepackage{dsfont}
\usepackage{pdftexcmds}

\makeatletter
\let\qe@amsmmathbb\mathbb
\renewcommand{\mathbb}[1]{%
  \ifnum\pdf@strcmp{\unexpanded{#1}}{1}=\z@
    \mathds{1}%
  \else
    \qe@amsmmathbb{#1}%
  \fi
}
\makeatother

% Preferred indicator macros (either style is safe after the remap above).
\newcommand{\bbone}{\mathds{1}}
\newcommand{\ind}{\mathbf{1}}

%   exams:  % Shared math extras for qe-review.
% Requires already-loaded: amsmath, amssymb.
%
% Classic pitfall: AMS blackboard-bold fonts have no digit glyphs, so
% \mathbb{1} renders as overlapping garbage. Map the digit 1 to dsfont's
% \mathds{1}; leave \mathbb{R}, \mathbb{P}, etc. on AMS.
%
% Usage (path relative to the main .tex being compiled):
%   notes:  \input{../../common/qe-math-extras.tex}
%   exams:  \input{../../../common/qe-math-extras.tex}

\usepackage{dsfont}
\usepackage{pdftexcmds}

\makeatletter
\let\qe@amsmmathbb\mathbb
\renewcommand{\mathbb}[1]{%
  \ifnum\pdf@strcmp{\unexpanded{#1}}{1}=\z@
    \mathds{1}%
  \else
    \qe@amsmmathbb{#1}%
  \fi
}
\makeatother

% Preferred indicator macros (either style is safe after the remap above).
\newcommand{\bbone}{\mathds{1}}
\newcommand{\ind}{\mathbf{1}}


\usepackage{dsfont}
\usepackage{pdftexcmds}

\makeatletter
\let\qe@amsmmathbb\mathbb
\renewcommand{\mathbb}[1]{%
  \ifnum\pdf@strcmp{\unexpanded{#1}}{1}=\z@
    \mathds{1}%
  \else
    \qe@amsmmathbb{#1}%
  \fi
}
\makeatother

% Preferred indicator macros (either style is safe after the remap above).
\newcommand{\bbone}{\mathds{1}}
\newcommand{\ind}{\mathbf{1}}

%   exams:  % Shared math extras for qe-review.
% Requires already-loaded: amsmath, amssymb.
%
% Classic pitfall: AMS blackboard-bold fonts have no digit glyphs, so
% \mathbb{1} renders as overlapping garbage. Map the digit 1 to dsfont's
% \mathds{1}; leave \mathbb{R}, \mathbb{P}, etc. on AMS.
%
% Usage (path relative to the main .tex being compiled):
%   notes:  % Shared math extras for qe-review.
% Requires already-loaded: amsmath, amssymb.
%
% Classic pitfall: AMS blackboard-bold fonts have no digit glyphs, so
% \mathbb{1} renders as overlapping garbage. Map the digit 1 to dsfont's
% \mathds{1}; leave \mathbb{R}, \mathbb{P}, etc. on AMS.
%
% Usage (path relative to the main .tex being compiled):
%   notes:  \input{../../common/qe-math-extras.tex}
%   exams:  \input{../../../common/qe-math-extras.tex}

\usepackage{dsfont}
\usepackage{pdftexcmds}

\makeatletter
\let\qe@amsmmathbb\mathbb
\renewcommand{\mathbb}[1]{%
  \ifnum\pdf@strcmp{\unexpanded{#1}}{1}=\z@
    \mathds{1}%
  \else
    \qe@amsmmathbb{#1}%
  \fi
}
\makeatother

% Preferred indicator macros (either style is safe after the remap above).
\newcommand{\bbone}{\mathds{1}}
\newcommand{\ind}{\mathbf{1}}

%   exams:  % Shared math extras for qe-review.
% Requires already-loaded: amsmath, amssymb.
%
% Classic pitfall: AMS blackboard-bold fonts have no digit glyphs, so
% \mathbb{1} renders as overlapping garbage. Map the digit 1 to dsfont's
% \mathds{1}; leave \mathbb{R}, \mathbb{P}, etc. on AMS.
%
% Usage (path relative to the main .tex being compiled):
%   notes:  \input{../../common/qe-math-extras.tex}
%   exams:  \input{../../../common/qe-math-extras.tex}

\usepackage{dsfont}
\usepackage{pdftexcmds}

\makeatletter
\let\qe@amsmmathbb\mathbb
\renewcommand{\mathbb}[1]{%
  \ifnum\pdf@strcmp{\unexpanded{#1}}{1}=\z@
    \mathds{1}%
  \else
    \qe@amsmmathbb{#1}%
  \fi
}
\makeatother

% Preferred indicator macros (either style is safe after the remap above).
\newcommand{\bbone}{\mathds{1}}
\newcommand{\ind}{\mathbf{1}}


\usepackage{dsfont}
\usepackage{pdftexcmds}

\makeatletter
\let\qe@amsmmathbb\mathbb
\renewcommand{\mathbb}[1]{%
  \ifnum\pdf@strcmp{\unexpanded{#1}}{1}=\z@
    \mathds{1}%
  \else
    \qe@amsmmathbb{#1}%
  \fi
}
\makeatother

% Preferred indicator macros (either style is safe after the remap above).
\newcommand{\bbone}{\mathds{1}}
\newcommand{\ind}{\mathbf{1}}


\usepackage{dsfont}
\usepackage{pdftexcmds}

\makeatletter
\let\qe@amsmmathbb\mathbb
\renewcommand{\mathbb}[1]{%
  \ifnum\pdf@strcmp{\unexpanded{#1}}{1}=\z@
    \mathds{1}%
  \else
    \qe@amsmmathbb{#1}%
  \fi
}
\makeatother

% Preferred indicator macros (either style is safe after the remap above).
\newcommand{\bbone}{\mathds{1}}
\newcommand{\ind}{\mathbf{1}}

% Shared amsthm environments for qe-review (minimal).
% Requires already-loaded: amsthm.
% Safe to \input after \usepackage{amsmath,amssymb,amsthm,...}.
%
% Shared numbering uses aliascnt so cleveref keys off the environment
% name (Lemma, Proposition, Exercise, ...) rather than the parent
% counter (Theorem, Definition, Remark). Without aliascnt, \cref{lem:...}
% prints ``Theorem 9.13'' instead of ``Lemma 9.13''.

\usepackage{aliascnt}

\theoremstyle{plain}
\newtheorem{theorem}{Theorem}[section]

\newaliascnt{lemma}{theorem}
\newtheorem{lemma}[lemma]{Lemma}
\aliascntresetthe{lemma}

\newaliascnt{proposition}{theorem}
\newtheorem{proposition}[proposition]{Proposition}
\aliascntresetthe{proposition}

\newaliascnt{corollary}{theorem}
\newtheorem{corollary}[corollary]{Corollary}
\aliascntresetthe{corollary}

\newaliascnt{claim}{theorem}
\newtheorem{claim}[claim]{Claim}
\aliascntresetthe{claim}

\newaliascnt{fact}{theorem}
\newtheorem{fact}[fact]{Fact}
\aliascntresetthe{fact}

\theoremstyle{definition}
\newtheorem{definition}{Definition}[section]

\newaliascnt{exercise}{definition}
\newtheorem{exercise}[exercise]{Exercise}
\aliascntresetthe{exercise}

\newaliascnt{assumption}{definition}
\newtheorem{assumption}[assumption]{Assumption}
\aliascntresetthe{assumption}

\newaliascnt{notation}{definition}
\newtheorem{notation}[notation]{Notation}
\aliascntresetthe{notation}

\newaliascnt{construction}{definition}
\newtheorem{construction}[construction]{Construction}
\aliascntresetthe{construction}

% Example: document-wide numbering (not per section / chapter)
\newtheorem{example}{Example}

\theoremstyle{remark}
\newtheorem{remark}{Remark}[section]
\newtheorem*{remark*}{Remark}

\newaliascnt{heuristic}{remark}
\newtheorem{heuristic}[heuristic]{Heuristic}
\aliascntresetthe{heuristic}

\newaliascnt{convention}{remark}
\newtheorem{convention}[convention]{Convention}
\aliascntresetthe{convention}

\newaliascnt{warning}{remark}
\newtheorem{warning}[warning]{Warning}
\aliascntresetthe{warning}

\newaliascnt{proofsketch}{remark}
\newtheorem{proofsketch}[proofsketch]{Proof sketch}
\aliascntresetthe{proofsketch}

\newaliascnt{checkpoint}{remark}
\newtheorem{checkpoint}[checkpoint]{Checkpoint}
\aliascntresetthe{checkpoint}

% Note: document-wide numbering (not per section)
\newtheorem{note}{Note}
\newtheorem*{note*}{Note}

\newaliascnt{examnote}{note}
\newtheorem{examnote}[examnote]{Exam note}
\aliascntresetthe{examnote}

\newaliascnt{study}{note}
\newtheorem{study}[study]{Study note}
\aliascntresetthe{study}

\newaliascnt{summary}{note}
\newtheorem{summary}[summary]{Summary}
\aliascntresetthe{summary}

\newtheorem{examproblem}{Problem}[section]
\newtheorem*{examproblem*}{Problem}

% Bilingual aliases
\newenvironment{dingli}{\begin{theorem}}{\end{theorem}}
\newenvironment{yinli}{\begin{lemma}}{\end{lemma}}
\newenvironment{mingti}{\begin{proposition}}{\end{proposition}}
\newenvironment{tuilun}{\begin{corollary}}{\end{corollary}}
\newenvironment{dingyi}{\begin{definition}}{\end{definition}}
\newenvironment{li}{\begin{example}}{\end{example}}
\newenvironment{zhu}{\begin{remark}}{\end{remark}}
\newenvironment{biji}{\begin{note}}{\end{note}}
\newenvironment{kaoshi}{\begin{examnote}}{\end{examnote}}

\renewcommand{\qedsymbol}{\ensuremath{\square}}

% Shared cleveref setup for qe-review.
% Load AFTER: hyperref, and AFTER all \newtheorem / amsthm environments.
% Shared theorem counters are aliased in qe-amsthm-envs.tex (aliascnt),
% otherwise \cref on a lemma/proposition prints ``Theorem''.
%
% Usage (path relative to the main .tex being compiled):
%   notes:  % Shared cleveref setup for qe-review.
% Load AFTER: hyperref, and AFTER all \newtheorem / amsthm environments.
% Shared theorem counters are aliased in qe-amsthm-envs.tex (aliascnt),
% otherwise \cref on a lemma/proposition prints ``Theorem''.
%
% Usage (path relative to the main .tex being compiled):
%   notes:  % Shared cleveref setup for qe-review.
% Load AFTER: hyperref, and AFTER all \newtheorem / amsthm environments.
% Shared theorem counters are aliased in qe-amsthm-envs.tex (aliascnt),
% otherwise \cref on a lemma/proposition prints ``Theorem''.
%
% Usage (path relative to the main .tex being compiled):
%   notes:  \input{../../common/qe-cleveref.tex}
%   exams:  \input{../../../common/qe-cleveref.tex}
%
% Prefer \cref / \Cref over \ref. Use \Cref at sentence start.
% Unnumbered sectioning (e.g. \paragraph with default secnumdepth)
% produces an empty cleveref type; map that to "Paragraph".

\usepackage[nameinlink,noabbrev]{cleveref}

% Fallback for unnumbered \paragraph / \subsubsection labels
\crefname{}{Paragraph}{Paragraphs}
\Crefname{}{Paragraph}{Paragraphs}

% Numbered theorem-like (plain)
\crefname{theorem}{Theorem}{Theorems}
\Crefname{theorem}{Theorem}{Theorems}
\crefname{lemma}{Lemma}{Lemmas}
\Crefname{lemma}{Lemma}{Lemmas}
\crefname{proposition}{Proposition}{Propositions}
\Crefname{proposition}{Proposition}{Propositions}
\crefname{corollary}{Corollary}{Corollaries}
\Crefname{corollary}{Corollary}{Corollaries}
\crefname{claim}{Claim}{Claims}
\Crefname{claim}{Claim}{Claims}
\crefname{fact}{Fact}{Facts}
\Crefname{fact}{Fact}{Facts}

% Definition-like
\crefname{definition}{Definition}{Definitions}
\Crefname{definition}{Definition}{Definitions}
\crefname{example}{Example}{Examples}
\Crefname{example}{Example}{Examples}
\crefname{exercise}{Exercise}{Exercises}
\Crefname{exercise}{Exercise}{Exercises}
\crefname{assumption}{Assumption}{Assumptions}
\Crefname{assumption}{Assumption}{Assumptions}
\crefname{notation}{Notation}{Notations}
\Crefname{notation}{Notation}{Notations}
\crefname{construction}{Construction}{Constructions}
\Crefname{construction}{Construction}{Constructions}

% Remark-like
\crefname{remark}{Remark}{Remarks}
\Crefname{remark}{Remark}{Remarks}
\crefname{heuristic}{Heuristic}{Heuristics}
\Crefname{heuristic}{Heuristic}{Heuristics}
\crefname{convention}{Convention}{Conventions}
\Crefname{convention}{Convention}{Conventions}
\crefname{warning}{Warning}{Warnings}
\Crefname{warning}{Warning}{Warnings}
\crefname{proofsketch}{Proof sketch}{Proof sketches}
\Crefname{proofsketch}{Proof sketch}{Proof sketches}
\crefname{checkpoint}{Checkpoint}{Checkpoints}
\Crefname{checkpoint}{Checkpoint}{Checkpoints}

% Document-wide notes / exam problems
\crefname{note}{Note}{Notes}
\Crefname{note}{Note}{Notes}
\crefname{examnote}{Exam note}{Exam notes}
\Crefname{examnote}{Exam note}{Exam notes}
\crefname{study}{Study note}{Study notes}
\Crefname{study}{Study note}{Study notes}
\crefname{summary}{Summary}{Summaries}
\Crefname{summary}{Summary}{Summaries}
\crefname{examproblem}{Problem}{Problems}
\Crefname{examproblem}{Problem}{Problems}

% Sectioning / floats (explicit, noabbrev-friendly)
\crefname{section}{Section}{Sections}
\Crefname{section}{Section}{Sections}
\crefname{subsection}{Subsection}{Subsections}
\Crefname{subsection}{Subsection}{Subsections}
\crefname{subsubsection}{Subsubsection}{Subsubsections}
\Crefname{subsubsection}{Subsubsection}{Subsubsections}
\crefname{paragraph}{Paragraph}{Paragraphs}
\Crefname{paragraph}{Paragraph}{Paragraphs}
\crefname{equation}{Equation}{Equations}
\Crefname{equation}{Equation}{Equations}
\crefname{figure}{Figure}{Figures}
\Crefname{figure}{Figure}{Figures}
\crefname{table}{Table}{Tables}
\Crefname{table}{Table}{Tables}
\crefname{enumi}{Item}{Items}
\Crefname{enumi}{Item}{Items}

%   exams:  % Shared cleveref setup for qe-review.
% Load AFTER: hyperref, and AFTER all \newtheorem / amsthm environments.
% Shared theorem counters are aliased in qe-amsthm-envs.tex (aliascnt),
% otherwise \cref on a lemma/proposition prints ``Theorem''.
%
% Usage (path relative to the main .tex being compiled):
%   notes:  \input{../../common/qe-cleveref.tex}
%   exams:  \input{../../../common/qe-cleveref.tex}
%
% Prefer \cref / \Cref over \ref. Use \Cref at sentence start.
% Unnumbered sectioning (e.g. \paragraph with default secnumdepth)
% produces an empty cleveref type; map that to "Paragraph".

\usepackage[nameinlink,noabbrev]{cleveref}

% Fallback for unnumbered \paragraph / \subsubsection labels
\crefname{}{Paragraph}{Paragraphs}
\Crefname{}{Paragraph}{Paragraphs}

% Numbered theorem-like (plain)
\crefname{theorem}{Theorem}{Theorems}
\Crefname{theorem}{Theorem}{Theorems}
\crefname{lemma}{Lemma}{Lemmas}
\Crefname{lemma}{Lemma}{Lemmas}
\crefname{proposition}{Proposition}{Propositions}
\Crefname{proposition}{Proposition}{Propositions}
\crefname{corollary}{Corollary}{Corollaries}
\Crefname{corollary}{Corollary}{Corollaries}
\crefname{claim}{Claim}{Claims}
\Crefname{claim}{Claim}{Claims}
\crefname{fact}{Fact}{Facts}
\Crefname{fact}{Fact}{Facts}

% Definition-like
\crefname{definition}{Definition}{Definitions}
\Crefname{definition}{Definition}{Definitions}
\crefname{example}{Example}{Examples}
\Crefname{example}{Example}{Examples}
\crefname{exercise}{Exercise}{Exercises}
\Crefname{exercise}{Exercise}{Exercises}
\crefname{assumption}{Assumption}{Assumptions}
\Crefname{assumption}{Assumption}{Assumptions}
\crefname{notation}{Notation}{Notations}
\Crefname{notation}{Notation}{Notations}
\crefname{construction}{Construction}{Constructions}
\Crefname{construction}{Construction}{Constructions}

% Remark-like
\crefname{remark}{Remark}{Remarks}
\Crefname{remark}{Remark}{Remarks}
\crefname{heuristic}{Heuristic}{Heuristics}
\Crefname{heuristic}{Heuristic}{Heuristics}
\crefname{convention}{Convention}{Conventions}
\Crefname{convention}{Convention}{Conventions}
\crefname{warning}{Warning}{Warnings}
\Crefname{warning}{Warning}{Warnings}
\crefname{proofsketch}{Proof sketch}{Proof sketches}
\Crefname{proofsketch}{Proof sketch}{Proof sketches}
\crefname{checkpoint}{Checkpoint}{Checkpoints}
\Crefname{checkpoint}{Checkpoint}{Checkpoints}

% Document-wide notes / exam problems
\crefname{note}{Note}{Notes}
\Crefname{note}{Note}{Notes}
\crefname{examnote}{Exam note}{Exam notes}
\Crefname{examnote}{Exam note}{Exam notes}
\crefname{study}{Study note}{Study notes}
\Crefname{study}{Study note}{Study notes}
\crefname{summary}{Summary}{Summaries}
\Crefname{summary}{Summary}{Summaries}
\crefname{examproblem}{Problem}{Problems}
\Crefname{examproblem}{Problem}{Problems}

% Sectioning / floats (explicit, noabbrev-friendly)
\crefname{section}{Section}{Sections}
\Crefname{section}{Section}{Sections}
\crefname{subsection}{Subsection}{Subsections}
\Crefname{subsection}{Subsection}{Subsections}
\crefname{subsubsection}{Subsubsection}{Subsubsections}
\Crefname{subsubsection}{Subsubsection}{Subsubsections}
\crefname{paragraph}{Paragraph}{Paragraphs}
\Crefname{paragraph}{Paragraph}{Paragraphs}
\crefname{equation}{Equation}{Equations}
\Crefname{equation}{Equation}{Equations}
\crefname{figure}{Figure}{Figures}
\Crefname{figure}{Figure}{Figures}
\crefname{table}{Table}{Tables}
\Crefname{table}{Table}{Tables}
\crefname{enumi}{Item}{Items}
\Crefname{enumi}{Item}{Items}

%
% Prefer \cref / \Cref over \ref. Use \Cref at sentence start.
% Unnumbered sectioning (e.g. \paragraph with default secnumdepth)
% produces an empty cleveref type; map that to "Paragraph".

\usepackage[nameinlink,noabbrev]{cleveref}

% Fallback for unnumbered \paragraph / \subsubsection labels
\crefname{}{Paragraph}{Paragraphs}
\Crefname{}{Paragraph}{Paragraphs}

% Numbered theorem-like (plain)
\crefname{theorem}{Theorem}{Theorems}
\Crefname{theorem}{Theorem}{Theorems}
\crefname{lemma}{Lemma}{Lemmas}
\Crefname{lemma}{Lemma}{Lemmas}
\crefname{proposition}{Proposition}{Propositions}
\Crefname{proposition}{Proposition}{Propositions}
\crefname{corollary}{Corollary}{Corollaries}
\Crefname{corollary}{Corollary}{Corollaries}
\crefname{claim}{Claim}{Claims}
\Crefname{claim}{Claim}{Claims}
\crefname{fact}{Fact}{Facts}
\Crefname{fact}{Fact}{Facts}

% Definition-like
\crefname{definition}{Definition}{Definitions}
\Crefname{definition}{Definition}{Definitions}
\crefname{example}{Example}{Examples}
\Crefname{example}{Example}{Examples}
\crefname{exercise}{Exercise}{Exercises}
\Crefname{exercise}{Exercise}{Exercises}
\crefname{assumption}{Assumption}{Assumptions}
\Crefname{assumption}{Assumption}{Assumptions}
\crefname{notation}{Notation}{Notations}
\Crefname{notation}{Notation}{Notations}
\crefname{construction}{Construction}{Constructions}
\Crefname{construction}{Construction}{Constructions}

% Remark-like
\crefname{remark}{Remark}{Remarks}
\Crefname{remark}{Remark}{Remarks}
\crefname{heuristic}{Heuristic}{Heuristics}
\Crefname{heuristic}{Heuristic}{Heuristics}
\crefname{convention}{Convention}{Conventions}
\Crefname{convention}{Convention}{Conventions}
\crefname{warning}{Warning}{Warnings}
\Crefname{warning}{Warning}{Warnings}
\crefname{proofsketch}{Proof sketch}{Proof sketches}
\Crefname{proofsketch}{Proof sketch}{Proof sketches}
\crefname{checkpoint}{Checkpoint}{Checkpoints}
\Crefname{checkpoint}{Checkpoint}{Checkpoints}

% Document-wide notes / exam problems
\crefname{note}{Note}{Notes}
\Crefname{note}{Note}{Notes}
\crefname{examnote}{Exam note}{Exam notes}
\Crefname{examnote}{Exam note}{Exam notes}
\crefname{study}{Study note}{Study notes}
\Crefname{study}{Study note}{Study notes}
\crefname{summary}{Summary}{Summaries}
\Crefname{summary}{Summary}{Summaries}
\crefname{examproblem}{Problem}{Problems}
\Crefname{examproblem}{Problem}{Problems}

% Sectioning / floats (explicit, noabbrev-friendly)
\crefname{section}{Section}{Sections}
\Crefname{section}{Section}{Sections}
\crefname{subsection}{Subsection}{Subsections}
\Crefname{subsection}{Subsection}{Subsections}
\crefname{subsubsection}{Subsubsection}{Subsubsections}
\Crefname{subsubsection}{Subsubsection}{Subsubsections}
\crefname{paragraph}{Paragraph}{Paragraphs}
\Crefname{paragraph}{Paragraph}{Paragraphs}
\crefname{equation}{Equation}{Equations}
\Crefname{equation}{Equation}{Equations}
\crefname{figure}{Figure}{Figures}
\Crefname{figure}{Figure}{Figures}
\crefname{table}{Table}{Tables}
\Crefname{table}{Table}{Tables}
\crefname{enumi}{Item}{Items}
\Crefname{enumi}{Item}{Items}

%   exams:  % Shared cleveref setup for qe-review.
% Load AFTER: hyperref, and AFTER all \newtheorem / amsthm environments.
% Shared theorem counters are aliased in qe-amsthm-envs.tex (aliascnt),
% otherwise \cref on a lemma/proposition prints ``Theorem''.
%
% Usage (path relative to the main .tex being compiled):
%   notes:  % Shared cleveref setup for qe-review.
% Load AFTER: hyperref, and AFTER all \newtheorem / amsthm environments.
% Shared theorem counters are aliased in qe-amsthm-envs.tex (aliascnt),
% otherwise \cref on a lemma/proposition prints ``Theorem''.
%
% Usage (path relative to the main .tex being compiled):
%   notes:  \input{../../common/qe-cleveref.tex}
%   exams:  \input{../../../common/qe-cleveref.tex}
%
% Prefer \cref / \Cref over \ref. Use \Cref at sentence start.
% Unnumbered sectioning (e.g. \paragraph with default secnumdepth)
% produces an empty cleveref type; map that to "Paragraph".

\usepackage[nameinlink,noabbrev]{cleveref}

% Fallback for unnumbered \paragraph / \subsubsection labels
\crefname{}{Paragraph}{Paragraphs}
\Crefname{}{Paragraph}{Paragraphs}

% Numbered theorem-like (plain)
\crefname{theorem}{Theorem}{Theorems}
\Crefname{theorem}{Theorem}{Theorems}
\crefname{lemma}{Lemma}{Lemmas}
\Crefname{lemma}{Lemma}{Lemmas}
\crefname{proposition}{Proposition}{Propositions}
\Crefname{proposition}{Proposition}{Propositions}
\crefname{corollary}{Corollary}{Corollaries}
\Crefname{corollary}{Corollary}{Corollaries}
\crefname{claim}{Claim}{Claims}
\Crefname{claim}{Claim}{Claims}
\crefname{fact}{Fact}{Facts}
\Crefname{fact}{Fact}{Facts}

% Definition-like
\crefname{definition}{Definition}{Definitions}
\Crefname{definition}{Definition}{Definitions}
\crefname{example}{Example}{Examples}
\Crefname{example}{Example}{Examples}
\crefname{exercise}{Exercise}{Exercises}
\Crefname{exercise}{Exercise}{Exercises}
\crefname{assumption}{Assumption}{Assumptions}
\Crefname{assumption}{Assumption}{Assumptions}
\crefname{notation}{Notation}{Notations}
\Crefname{notation}{Notation}{Notations}
\crefname{construction}{Construction}{Constructions}
\Crefname{construction}{Construction}{Constructions}

% Remark-like
\crefname{remark}{Remark}{Remarks}
\Crefname{remark}{Remark}{Remarks}
\crefname{heuristic}{Heuristic}{Heuristics}
\Crefname{heuristic}{Heuristic}{Heuristics}
\crefname{convention}{Convention}{Conventions}
\Crefname{convention}{Convention}{Conventions}
\crefname{warning}{Warning}{Warnings}
\Crefname{warning}{Warning}{Warnings}
\crefname{proofsketch}{Proof sketch}{Proof sketches}
\Crefname{proofsketch}{Proof sketch}{Proof sketches}
\crefname{checkpoint}{Checkpoint}{Checkpoints}
\Crefname{checkpoint}{Checkpoint}{Checkpoints}

% Document-wide notes / exam problems
\crefname{note}{Note}{Notes}
\Crefname{note}{Note}{Notes}
\crefname{examnote}{Exam note}{Exam notes}
\Crefname{examnote}{Exam note}{Exam notes}
\crefname{study}{Study note}{Study notes}
\Crefname{study}{Study note}{Study notes}
\crefname{summary}{Summary}{Summaries}
\Crefname{summary}{Summary}{Summaries}
\crefname{examproblem}{Problem}{Problems}
\Crefname{examproblem}{Problem}{Problems}

% Sectioning / floats (explicit, noabbrev-friendly)
\crefname{section}{Section}{Sections}
\Crefname{section}{Section}{Sections}
\crefname{subsection}{Subsection}{Subsections}
\Crefname{subsection}{Subsection}{Subsections}
\crefname{subsubsection}{Subsubsection}{Subsubsections}
\Crefname{subsubsection}{Subsubsection}{Subsubsections}
\crefname{paragraph}{Paragraph}{Paragraphs}
\Crefname{paragraph}{Paragraph}{Paragraphs}
\crefname{equation}{Equation}{Equations}
\Crefname{equation}{Equation}{Equations}
\crefname{figure}{Figure}{Figures}
\Crefname{figure}{Figure}{Figures}
\crefname{table}{Table}{Tables}
\Crefname{table}{Table}{Tables}
\crefname{enumi}{Item}{Items}
\Crefname{enumi}{Item}{Items}

%   exams:  % Shared cleveref setup for qe-review.
% Load AFTER: hyperref, and AFTER all \newtheorem / amsthm environments.
% Shared theorem counters are aliased in qe-amsthm-envs.tex (aliascnt),
% otherwise \cref on a lemma/proposition prints ``Theorem''.
%
% Usage (path relative to the main .tex being compiled):
%   notes:  \input{../../common/qe-cleveref.tex}
%   exams:  \input{../../../common/qe-cleveref.tex}
%
% Prefer \cref / \Cref over \ref. Use \Cref at sentence start.
% Unnumbered sectioning (e.g. \paragraph with default secnumdepth)
% produces an empty cleveref type; map that to "Paragraph".

\usepackage[nameinlink,noabbrev]{cleveref}

% Fallback for unnumbered \paragraph / \subsubsection labels
\crefname{}{Paragraph}{Paragraphs}
\Crefname{}{Paragraph}{Paragraphs}

% Numbered theorem-like (plain)
\crefname{theorem}{Theorem}{Theorems}
\Crefname{theorem}{Theorem}{Theorems}
\crefname{lemma}{Lemma}{Lemmas}
\Crefname{lemma}{Lemma}{Lemmas}
\crefname{proposition}{Proposition}{Propositions}
\Crefname{proposition}{Proposition}{Propositions}
\crefname{corollary}{Corollary}{Corollaries}
\Crefname{corollary}{Corollary}{Corollaries}
\crefname{claim}{Claim}{Claims}
\Crefname{claim}{Claim}{Claims}
\crefname{fact}{Fact}{Facts}
\Crefname{fact}{Fact}{Facts}

% Definition-like
\crefname{definition}{Definition}{Definitions}
\Crefname{definition}{Definition}{Definitions}
\crefname{example}{Example}{Examples}
\Crefname{example}{Example}{Examples}
\crefname{exercise}{Exercise}{Exercises}
\Crefname{exercise}{Exercise}{Exercises}
\crefname{assumption}{Assumption}{Assumptions}
\Crefname{assumption}{Assumption}{Assumptions}
\crefname{notation}{Notation}{Notations}
\Crefname{notation}{Notation}{Notations}
\crefname{construction}{Construction}{Constructions}
\Crefname{construction}{Construction}{Constructions}

% Remark-like
\crefname{remark}{Remark}{Remarks}
\Crefname{remark}{Remark}{Remarks}
\crefname{heuristic}{Heuristic}{Heuristics}
\Crefname{heuristic}{Heuristic}{Heuristics}
\crefname{convention}{Convention}{Conventions}
\Crefname{convention}{Convention}{Conventions}
\crefname{warning}{Warning}{Warnings}
\Crefname{warning}{Warning}{Warnings}
\crefname{proofsketch}{Proof sketch}{Proof sketches}
\Crefname{proofsketch}{Proof sketch}{Proof sketches}
\crefname{checkpoint}{Checkpoint}{Checkpoints}
\Crefname{checkpoint}{Checkpoint}{Checkpoints}

% Document-wide notes / exam problems
\crefname{note}{Note}{Notes}
\Crefname{note}{Note}{Notes}
\crefname{examnote}{Exam note}{Exam notes}
\Crefname{examnote}{Exam note}{Exam notes}
\crefname{study}{Study note}{Study notes}
\Crefname{study}{Study note}{Study notes}
\crefname{summary}{Summary}{Summaries}
\Crefname{summary}{Summary}{Summaries}
\crefname{examproblem}{Problem}{Problems}
\Crefname{examproblem}{Problem}{Problems}

% Sectioning / floats (explicit, noabbrev-friendly)
\crefname{section}{Section}{Sections}
\Crefname{section}{Section}{Sections}
\crefname{subsection}{Subsection}{Subsections}
\Crefname{subsection}{Subsection}{Subsections}
\crefname{subsubsection}{Subsubsection}{Subsubsections}
\Crefname{subsubsection}{Subsubsection}{Subsubsections}
\crefname{paragraph}{Paragraph}{Paragraphs}
\Crefname{paragraph}{Paragraph}{Paragraphs}
\crefname{equation}{Equation}{Equations}
\Crefname{equation}{Equation}{Equations}
\crefname{figure}{Figure}{Figures}
\Crefname{figure}{Figure}{Figures}
\crefname{table}{Table}{Tables}
\Crefname{table}{Table}{Tables}
\crefname{enumi}{Item}{Items}
\Crefname{enumi}{Item}{Items}

%
% Prefer \cref / \Cref over \ref. Use \Cref at sentence start.
% Unnumbered sectioning (e.g. \paragraph with default secnumdepth)
% produces an empty cleveref type; map that to "Paragraph".

\usepackage[nameinlink,noabbrev]{cleveref}

% Fallback for unnumbered \paragraph / \subsubsection labels
\crefname{}{Paragraph}{Paragraphs}
\Crefname{}{Paragraph}{Paragraphs}

% Numbered theorem-like (plain)
\crefname{theorem}{Theorem}{Theorems}
\Crefname{theorem}{Theorem}{Theorems}
\crefname{lemma}{Lemma}{Lemmas}
\Crefname{lemma}{Lemma}{Lemmas}
\crefname{proposition}{Proposition}{Propositions}
\Crefname{proposition}{Proposition}{Propositions}
\crefname{corollary}{Corollary}{Corollaries}
\Crefname{corollary}{Corollary}{Corollaries}
\crefname{claim}{Claim}{Claims}
\Crefname{claim}{Claim}{Claims}
\crefname{fact}{Fact}{Facts}
\Crefname{fact}{Fact}{Facts}

% Definition-like
\crefname{definition}{Definition}{Definitions}
\Crefname{definition}{Definition}{Definitions}
\crefname{example}{Example}{Examples}
\Crefname{example}{Example}{Examples}
\crefname{exercise}{Exercise}{Exercises}
\Crefname{exercise}{Exercise}{Exercises}
\crefname{assumption}{Assumption}{Assumptions}
\Crefname{assumption}{Assumption}{Assumptions}
\crefname{notation}{Notation}{Notations}
\Crefname{notation}{Notation}{Notations}
\crefname{construction}{Construction}{Constructions}
\Crefname{construction}{Construction}{Constructions}

% Remark-like
\crefname{remark}{Remark}{Remarks}
\Crefname{remark}{Remark}{Remarks}
\crefname{heuristic}{Heuristic}{Heuristics}
\Crefname{heuristic}{Heuristic}{Heuristics}
\crefname{convention}{Convention}{Conventions}
\Crefname{convention}{Convention}{Conventions}
\crefname{warning}{Warning}{Warnings}
\Crefname{warning}{Warning}{Warnings}
\crefname{proofsketch}{Proof sketch}{Proof sketches}
\Crefname{proofsketch}{Proof sketch}{Proof sketches}
\crefname{checkpoint}{Checkpoint}{Checkpoints}
\Crefname{checkpoint}{Checkpoint}{Checkpoints}

% Document-wide notes / exam problems
\crefname{note}{Note}{Notes}
\Crefname{note}{Note}{Notes}
\crefname{examnote}{Exam note}{Exam notes}
\Crefname{examnote}{Exam note}{Exam notes}
\crefname{study}{Study note}{Study notes}
\Crefname{study}{Study note}{Study notes}
\crefname{summary}{Summary}{Summaries}
\Crefname{summary}{Summary}{Summaries}
\crefname{examproblem}{Problem}{Problems}
\Crefname{examproblem}{Problem}{Problems}

% Sectioning / floats (explicit, noabbrev-friendly)
\crefname{section}{Section}{Sections}
\Crefname{section}{Section}{Sections}
\crefname{subsection}{Subsection}{Subsections}
\Crefname{subsection}{Subsection}{Subsections}
\crefname{subsubsection}{Subsubsection}{Subsubsections}
\Crefname{subsubsection}{Subsubsection}{Subsubsections}
\crefname{paragraph}{Paragraph}{Paragraphs}
\Crefname{paragraph}{Paragraph}{Paragraphs}
\crefname{equation}{Equation}{Equations}
\Crefname{equation}{Equation}{Equations}
\crefname{figure}{Figure}{Figures}
\Crefname{figure}{Figure}{Figures}
\crefname{table}{Table}{Tables}
\Crefname{table}{Table}{Tables}
\crefname{enumi}{Item}{Items}
\Crefname{enumi}{Item}{Items}

%
% Prefer \cref / \Cref over \ref. Use \Cref at sentence start.
% Unnumbered sectioning (e.g. \paragraph with default secnumdepth)
% produces an empty cleveref type; map that to "Paragraph".

\usepackage[nameinlink,noabbrev]{cleveref}

% Fallback for unnumbered \paragraph / \subsubsection labels
\crefname{}{Paragraph}{Paragraphs}
\Crefname{}{Paragraph}{Paragraphs}

% Numbered theorem-like (plain)
\crefname{theorem}{Theorem}{Theorems}
\Crefname{theorem}{Theorem}{Theorems}
\crefname{lemma}{Lemma}{Lemmas}
\Crefname{lemma}{Lemma}{Lemmas}
\crefname{proposition}{Proposition}{Propositions}
\Crefname{proposition}{Proposition}{Propositions}
\crefname{corollary}{Corollary}{Corollaries}
\Crefname{corollary}{Corollary}{Corollaries}
\crefname{claim}{Claim}{Claims}
\Crefname{claim}{Claim}{Claims}
\crefname{fact}{Fact}{Facts}
\Crefname{fact}{Fact}{Facts}

% Definition-like
\crefname{definition}{Definition}{Definitions}
\Crefname{definition}{Definition}{Definitions}
\crefname{example}{Example}{Examples}
\Crefname{example}{Example}{Examples}
\crefname{exercise}{Exercise}{Exercises}
\Crefname{exercise}{Exercise}{Exercises}
\crefname{assumption}{Assumption}{Assumptions}
\Crefname{assumption}{Assumption}{Assumptions}
\crefname{notation}{Notation}{Notations}
\Crefname{notation}{Notation}{Notations}
\crefname{construction}{Construction}{Constructions}
\Crefname{construction}{Construction}{Constructions}

% Remark-like
\crefname{remark}{Remark}{Remarks}
\Crefname{remark}{Remark}{Remarks}
\crefname{heuristic}{Heuristic}{Heuristics}
\Crefname{heuristic}{Heuristic}{Heuristics}
\crefname{convention}{Convention}{Conventions}
\Crefname{convention}{Convention}{Conventions}
\crefname{warning}{Warning}{Warnings}
\Crefname{warning}{Warning}{Warnings}
\crefname{proofsketch}{Proof sketch}{Proof sketches}
\Crefname{proofsketch}{Proof sketch}{Proof sketches}
\crefname{checkpoint}{Checkpoint}{Checkpoints}
\Crefname{checkpoint}{Checkpoint}{Checkpoints}

% Document-wide notes / exam problems
\crefname{note}{Note}{Notes}
\Crefname{note}{Note}{Notes}
\crefname{examnote}{Exam note}{Exam notes}
\Crefname{examnote}{Exam note}{Exam notes}
\crefname{study}{Study note}{Study notes}
\Crefname{study}{Study note}{Study notes}
\crefname{summary}{Summary}{Summaries}
\Crefname{summary}{Summary}{Summaries}
\crefname{examproblem}{Problem}{Problems}
\Crefname{examproblem}{Problem}{Problems}

% Sectioning / floats (explicit, noabbrev-friendly)
\crefname{section}{Section}{Sections}
\Crefname{section}{Section}{Sections}
\crefname{subsection}{Subsection}{Subsections}
\Crefname{subsection}{Subsection}{Subsections}
\crefname{subsubsection}{Subsubsection}{Subsubsections}
\Crefname{subsubsection}{Subsubsection}{Subsubsections}
\crefname{paragraph}{Paragraph}{Paragraphs}
\Crefname{paragraph}{Paragraph}{Paragraphs}
\crefname{equation}{Equation}{Equations}
\Crefname{equation}{Equation}{Equations}
\crefname{figure}{Figure}{Figures}
\Crefname{figure}{Figure}{Figures}
\crefname{table}{Table}{Tables}
\Crefname{table}{Table}{Tables}
\crefname{enumi}{Item}{Items}
\Crefname{enumi}{Item}{Items}

% PDF sidebar outline + printed TOC for QE exam transcripts.
% Load in the preamble AFTER hyperref and AFTER \newcommand{\examstatement}.
\hypersetup{
  unicode=true,
  bookmarks=true,
  bookmarksopen=true,
  bookmarksopenlevel=3,
  bookmarksnumbered=false,
  bookmarksdepth=3
}
\IfFileExists{bookmark.sty}{%
  \usepackage{bookmark}%
  \bookmarksetup{open,openlevel=3,numbered=false,depth=3}%
}{}
% Unnumbered in print, but still written to the TOC / PDF outline
% down to \subsubsection. Starred headings create neither.
\setcounter{secnumdepth}{-1}
\setcounter{tocdepth}{3}
\makeatletter
\@ifundefined{examstatement}{}{%
  \let\qe@origexamstatement\examstatement
  \renewcommand{\examstatement}[1]{%
    \par\phantomsection
    \addcontentsline{toc}{subsection}{#1}%
    \qe@origexamstatement{#1}%
  }%
}
\makeatother


\title{\textbf{QUALIFYING EXAM: APPLIED MATHEMATICS}\\[0.5em]\Large SPRING 2026}
\date{}
\author{}

\begin{document}

\maketitle
\vspace{-3em}

\tableofcontents

\section{Statement of the problems}

\examstatement{Problem 1. [10 points]}
Assume that locally the function \(f \in C^{2}(\mathbb{R})\) has a unique zero \(x_{\star}\) with \(f'(x_{\star}) \neq 0\). Consider the iteration

\[
x_{k+1} = x_{k} - \frac{[f(x_{k})]^{2}}{f(x_{k} + f(x_{k})) - f(x_{k})},\qquad k = 0, 1, \dots.
\]

Prove that \(x_{k} \to x_{\star}\), and estimate its order of local convergence.

\examstatement{Problem 2. [20 points]}
For any continuous function \(f\), define \(\langle f, g \rangle = \int_{-1}^{1} f(x)g(x)\,\mathrm{d}x\) and \(\|f\|_2 = \sqrt{\langle f, f \rangle}\). Write
\[
a_{n}=\frac{(2n)!}{2^{n}(n!)^{2}}.
\]
Define the \(n\)-th monic Legendre polynomial by: \(P_{0}(x)=1\), \(P_{1}(x)=x\), \(P_{n+1}(x)=xP_{n}(x)-\frac{n^{2}}{4n^{2}-1}P_{n-1}(x)\).
\begin{enumerate}[label=(\alph*)]
    \item Compute \(\langle P_{n}, P_{m} \rangle\), for \(m,n\in\mathbb{N}\).
    \item Prove that \(P_{n}\) solves
    \[
    \min_{\substack{\text{polynomial }p\\ \deg p\le n,\, p^{(n)}(1)=n!}}\|p\|_{2}
    \]
    and compute the minimal value.
    \item For \(f\in C^{n+1}[-1,1]\), choose the interpolation nodes \(x_{0},x_{1},\ldots,x_{n}\) as the \(n+1\) zeros of \(P_{n+1}(x)\), and denote its degree-\(n\) interpolation polynomial by \(L_{n}\). Prove
    \[
    \|f-L_{n}\|_{2}\le\frac{2\|f^{(n+1)}\|_{2}}{(2n+3)a_{n+1}(n+1)!}.
    \]
    \item For \(f\in C^{2n+2}[-1,1]\), derive the Gauss quadrature formula of \(I(f)=\int_{-1}^{1}f(x)\,\mathrm{d}x\) for \(n+1\) quadrature nodes, written by \(I_{n+1}(f)\), and estimate its error \(|I(f)-I_{n+1}(f)|\).
\end{enumerate}

\examstatement{Problem 3. [20 points]}
Given \(A \in \mathbb{R}^{l \times m}\), \(B \in \mathbb{R}^{n \times p}\), \(C \in \mathbb{R}^{l \times p}\). Consider the problem \(\min_{X}\|AXB - C\|_{F}\), where \(\|\cdot\|_{F}\) is the Frobenius norm of a matrix.
\begin{enumerate}[label=(\alph*)]
    \item Prove that the minimal value is \(0\), if and only if \(\rank(A) = \rank\bigl(\begin{bmatrix} A & C \end{bmatrix}\bigr)\), \(\rank(B) = \rank\begin{pmatrix} B \\ C \end{pmatrix}\).
    \item Show that there exists a unique minimal point \(X_{\star}\), if and only if \(\rank(A) = m\), \(\rank(B) = n\).
    \item Propose a numerical method to compute the unique \(X_{\star}\) in (b).
    \item Does \(X_{\star}\) also minimize \(\|AXB - C\|_{2}\)? Why? Here \(\|\cdot\|_{2}\) is the spectral norm of a matrix.
\end{enumerate}

\examstatement{Problem 4. [15 points]}
For the advection equation \(u_{t} + u_{x} = 0\), analyze the stability of the implicit schemes
\begin{enumerate}[label=(\alph*)]
    \item \(\dfrac{u_{j}^{n+1} - u_{j}^{n}}{\tau} + \dfrac{u_{j}^{n+1} - u_{j-1}^{n+1}}{h} = 0\);
    \item \(\dfrac{u_{j}^{n+1} - u_{j}^{n}}{\tau} + \dfrac{u_{j+1}^{n+1} - u_{j}^{n+1}}{h} = 0\);
    \item \(\dfrac{u_{j}^{n+1} - u_{j}^{n}}{\tau} + \dfrac{u_{j+1}^{n+1} - u_{j-1}^{n+1}}{2h} = 0\).
\end{enumerate}

\examstatement{Problem 5. [15 points]}
Write and prove the maximum principle of the centered finite difference scheme for discretizing the equation

\[
u_{xx} + u_{yy} + d u_{x} + e u_{y} + f u = 0, \quad f < 0,
\]

under some suitable assumptions. Apply the maximum principle to show the stability of the backward Euler scheme for the parabolic equation

\[
u_{t} - (u_{xx} + u_{yy} + d u_{x} + e u_{y} + f u) = 0, \quad (x, y) \in (0, 1) \times (0, 1),
\]

with homogeneous Dirichlet boundary conditions.

\noindent\textbf{You can choose either 6 or 7 to answer. The points will be decided as \(\max(6, 7)\).}

\examstatement{Problem 6. [20 points]}
Consider the following system

\[
\ddot{x} - 2x - x^{2} + x^{3} = 0.
\]

\begin{enumerate}[label=(\alph*)]
    \item (6 points) Find the static solutions to the system.
    \item (14 points) For sufficiently small but nonzero \(\epsilon\), show that the orbit starting at \((x(0), \dot{x}(0)) = (2 + \epsilon, 0)\) is closed in the phase space, i.e., \(x - \dot{x}\) plane. Find the approximation of the period of this orbit up to \(O(\epsilon^{2})\).
\end{enumerate}

\examstatement{Problem 7. [20 points]}
The polar of \(C \subseteq \mathbb{R}^{n}\) is defined as the set

\[
C^{\circ} = \bigl\{ y \in \mathbb{R}^{n} \mid y^{T} x \le 1 \text{ for all } x \in C \bigr\}.
\]

\begin{enumerate}[label=(\alph*)]
    \item (2 points) Show that \(C^{\circ}\) is convex (even if \(C\) is not).
    \item (3 points) What is the polar of a cone?
    \item (3 points) What is the polar of the unit ball for a norm \(\|\cdot\|\)?
    \item (4 points) What is the polar of the set \(C = \{ x \mid \mathbf{1}^{T} x = 1,\, x \succeq 0 \}\)?
    \item (8 points) Show that if \(C\) is closed and convex, with \(0 \in C\), then \((C^{\circ})^{\circ} = C\).
\end{enumerate}

\noindent\rule{\linewidth}{0.4pt}

\section{Statement and solutions}

\subsection{Problem 1. [10 points]}
Assume that locally the function \(f \in C^{2}(\mathbb{R})\) has a unique zero \(x_{\star}\) with \(f'(x_{\star}) \neq 0\). Consider the iteration

\[
x_{k+1} = x_{k} - \frac{[f(x_{k})]^{2}}{f(x_{k} + f(x_{k})) - f(x_{k})},\qquad k = 0, 1, \dots.
\]

Prove that \(x_{k} \to x_{\star}\), and estimate its order of local convergence.

\setcounter{section}{1}

\subsubsection{Theorems used in Problem 1}

\begin{theorem}[Taylor formula with Peano remainder]
\label{thm:peano}
Let \(n\ge 1\) and let \(F\) be \(n\) times differentiable at \(a\in\mathbb{R}\) (one-sided derivatives suffice). Write \(T_{m}(h)\coloneqq\sum_{k=0}^{m}\frac{F^{(k)}(a)}{k!}\,h^{k}\). Then, as \(h\to 0\),

\begin{enumerate}[label=\textup{(\roman*)}]
    \item
    \(F(a+h)=T_{n}(h)+o(h^{n})\).

    \item
    \(F(a+h)=T_{n-1}(h)+O(h^{n})\).
\end{enumerate}
\end{theorem}

\begin{proof}
For (i), the claim is

\[
\lim_{h\to 0}
\frac{F(a+h)-T_{n}(h)}{h^{n}}
=
0.
\]

The numerator and its first \(n-1\) derivatives vanish at \(h=0\), so L'H\^{o}pital applies \(n-1\) times and reduces the limit to

\[
\lim_{h\to 0}
\frac{F^{(n-1)}(a+h)-F^{(n-1)}(a)-F^{(n)}(a)\,h}{n!\,h}.
\]

The last step is the definition of \(F^{(n)}(a)\), not another application of L'H\^{o}pital.

For (ii),

\[
\biggl|
\frac{F(a+h)-T_{n-1}(h)}{h^{n}}
\biggr|
\le
\biggl|
\frac{F(a+h)-T_{n}(h)}{h^{n}}
\biggr|
+
\frac{\bigl|F^{(n)}(a)\bigr|}{n!}.
\]

The first summand tends to \(0\) by (i), hence is locally bounded. Thus the left-hand side is locally bounded, which is (ii).
\end{proof}

\begin{corollary}[\(C^{2}\) expansions]
\label{cor:c2-peano}
If \(F\) is twice differentiable at \(a\), then as \(h\to 0\),

\begin{equation}
\label{eq:c2-peano}
F(a+h)
=
F(a)+F'(a)h+\tfrac12 F''(a)\,h^{2}+o(h^{2})
=
F(a)+F'(a)h+O(h^{2}).
\end{equation}

In particular both identities hold for every \(F\in C^{2}\) in a neighborhood of \(a\).
\end{corollary}

\begin{corollary}[When the remainder is \(O(h^{n+1})\)]
\label{cor:peano-not-onplus1}
The expansion \(F(a+h)=T_{n}(h)+O(h^{n+1})\) is the case \(n+1\) of Theorem~\ref{thm:peano}(ii). Its hypothesis is existence of \(F^{(n+1)}(a)\), not mere \(n\)-fold differentiability at \(a\).

Thus \(F\in C^{2}\) yields \eqref{eq:c2-peano}, but does \emph{not} in general yield

\[
F(a+h)
=
F(a)+F'(a)h+\tfrac12 F''(a)\,h^{2}+O(h^{3}).
\]

If in addition \(F''\) is Hölder-\(\alpha\) at \(a\) for some \(\alpha\in(0,1]\), the integral remainder

\[
R_{2}(h)
=
\int_{0}^{h}(h-t)\bigl(F''(a+t)-F''(a)\bigr)\,\mathrm{d}t
\]

satisfies \(R_{2}(h)=O(|h|^{2+\alpha})\). The endpoint \(\alpha=1\) is Lipschitz continuity of \(F''\) (equivalently \(F\in C^{2,1}\)), which is the natural condition for \(O(h^{3})\) short of assuming \(F'''(a)\).
\end{corollary}

\begin{example}
\label{ex:c2-not-O3}
The germ \(F(h)=|h|^{5/2}\) is \(C^{2}\) on \(\mathbb{R}\), with \(F(0)=F'(0)=F''(0)=0\) and \(F''(h)=\frac{15}{4}|h|^{1/2}\) for \(h\neq 0\), but \(F'''(0)\) does not exist. Then \(F(h)/h^{2}=|h|^{1/2}\to 0\), while \(F(h)/h^{3}=|h|^{-1/2}\) is unbounded as \(h\to 0\). Hence \(F(h)=o(h^{2})\) and \(F(h)\neq O(h^{3})\).
\end{example}

\subsubsection{Manuscript proof idea}

The calculation below follows the handwritten notes, in the same order and with the same expansions. Write \(h_{k}\coloneqq x_{k}-x_{*}\) and \(a\coloneqq f'(x_{*})\), \(b\coloneqq f''(x_{*})\). \mserr{Red} marks a false claim or a missing justification.

The notes open with \mserr{``for \(k\) large enough''}. \mserr{Gap: this assumes \(x_{k}\to x_{*}\), which is half of what the exam asks to prove. The legitimate hypothesis is that \(x_{0}\) is sufficiently close to \(x_{*}\).} The error recurrence is correctly rewritten as

\[
h_{k+1}
=
h_{k}
-
\frac{\bigl[f(x_{k})\bigr]^{2}}{f\bigl(x_{k}+f(x_{k})\bigr)-f(x_{k})}.
\]

\mserr{Gap: well-definedness of the denominator is not checked.}

\paragraph{First expansion (order \(2\)).}
By \cref{cor:c2-peano}, the notes' first Taylor line is legal:

\[
f(x_{k})
=
a h_{k}+O(h_{k}^{2}).
\]

The composite expansion is likewise the \(n=2\) case of \cref{thm:peano}(ii):

\[
\begin{aligned}
f\bigl(x_{k}+f(x_{k})\bigr)
&=
a\bigl(h_{k}+f(x_{k})\bigr)
+
O\bigl((h_{k}+f(x_{k}))^{2}\bigr)
\\
&=
a\bigl(h_{k}+a h_{k}+O(h_{k}^{2})\bigr)
+
O(h_{k}^{2})
\\
&=
a(1+a)h_{k}+O(h_{k}^{2}).
\end{aligned}
\]

The \(O((h_{k}+f(x_{k}))^{2})=O(h_{k}^{2})\) substitution is legitimate because \(f(x_{k})=O(h_{k})\). Then

\[
f\bigl(x_{k}+f(x_{k})\bigr)-f(x_{k})
=
a^{2} h_{k}+O(h_{k}^{2}),
\qquad
\bigl[f(x_{k})\bigr]^{2}
=
a^{2} h_{k}^{2}+O(h_{k}^{3}),
\]

and, provided the denominator is \(a^{2}h_{k}(1+O(h_{k}))\) with \(a\neq 0\),

\[
h_{k+1}
=
h_{k}
-
\frac{h_{k}+O(h_{k}^{2})}{1+O(h_{k})}
=
h_{k}
-
\bigl(h_{k}+O(h_{k}^{2})\bigr)\bigl(1+O(h_{k})\bigr)
=
O(h_{k}^{2}).
\]

The inversion \((1+O(h_{k}))^{-1}=1+O(h_{k})\) is the same \(O\)-class identity, and needs \(|h_{k}|\) small. From here the notes conclude \(|h_{k+1}|\le M|h_{k}|^{2}\), hence quadratic convergence. \mserr{Gap: an \(O(\,\cdot\,)\) as \(h_{k}\to 0\) produces such an \(M\) only on a neighborhood of \(x_{*}\). Quadratic attraction then requires the additional restriction \(|h_{0}|<1/(2M)\), otherwise \(M|h_{0}|^{2}\) need not be smaller than \(|h_{0}|\). The uniqueness of the zero is not used.}

\paragraph{Second expansion (asymptotic constant).}
The notes restart with

\[
f(x_{k})
=
a h_{k}+\tfrac12 b h_{k}^{2}
+
\mserr{O(h_{k}^{3})}.
\]

\mserr{Illegal remainder: \(f\in C^{2}\) gives \(T_{2}+o(h^{2})\) by \cref{cor:c2-peano}, not \(T_{2}+O(h^{3})\). The latter is \cref{cor:peano-not-onplus1}, and needs \(f'''(x_{*})\). See \cref{ex:c2-not-O3}.} The same overstatement is repeated for the composite:

\[
f\bigl(x_{k}+f(x_{k})\bigr)
=
a\bigl(h_{k}+f(x_{k})\bigr)
+
\tfrac12 b\bigl(h_{k}+f(x_{k})\bigr)^{2}
+
\mserr{O(h_{k}^{3})}.
\]

Substituting the (already illegal) expansion of \(f(x_{k})\) and writing \(\delta_{k}\coloneqq h_{k}+f(x_{k})=(1+a)h_{k}+\tfrac12 b h_{k}^{2}+\mserr{O(h_{k}^{3})}\), the notes keep only

\[
f\bigl(x_{k}+f(x_{k})\bigr)
=
a(1+a)h_{k}
+
\mserr{\bigl(\tfrac12 b+ab\bigr)h_{k}^{2}}
+
\mserr{O(h_{k}^{3})}.
\]

\mserr{Algebra error, independent of the remainder class: \(\tfrac12 b\,\delta_{k}^{2}\) contributes \(\tfrac12 b(1+a)^{2}h_{k}^{2}\) at leading order, and \(a\cdot\tfrac12 b h_{k}^{2}\) contributes \(\tfrac12 ab\,h_{k}^{2}\). The \(h_{k}^{2}\) coefficient is \(\tfrac12 ab+\tfrac12 b(1+a)^{2}=\tfrac12 b+\tfrac32 ab+\tfrac12 a^{2}b\), not \(\tfrac12 b+ab\).} Subtracting \(f(x_{k})\) then produces the notes' denominator

\[
f\bigl(x_{k}+f(x_{k})\bigr)-f(x_{k})
=
a^{2}h_{k}
+
\mserr{ab\,h_{k}^{2}}
+
\mserr{O(h_{k}^{3})},
\]

whereas the missing \((1+a)^{2}\) term changes the \(h_{k}^{2}\) coefficient from \(ab\) to \(\tfrac12 ab(3+a)\). The numerator

\[
\bigl[f(x_{k})\bigr]^{2}
=
a^{2}h_{k}^{2}+ab\,h_{k}^{3}+\mserr{O(h_{k}^{4})}
\]

would be \(a^{2}h_{k}^{2}+ab\,h_{k}^{3}+o(h_{k}^{3})\) under the legal remainder \(o(h_{k}^{2})\). The notes continue

\[
h_{k+1}
=
h_{k}
-
\frac{h_{k}+\frac{b}{a}h_{k}^{2}+\mserr{O(h_{k}^{3})}}{1+\frac{b}{a}h_{k}+\mserr{O(h_{k}^{2})}}
=
h_{k}
-
\Bigl(h_{k}+\frac{b}{a}h_{k}^{2}+\mserr{O(h_{k}^{3})}\Bigr)
\Bigl(1-\frac{b}{a}h_{k}+\mserr{O(h_{k}^{2})}\Bigr).
\]

Even internally, the two displayed \(\frac{b}{a}h_{k}^{2}\) contributions cancel, leaving \(h_{k+1}=\mserr{O(h_{k}^{3})}\) at this (incorrect) truncation. \mserr{The final claimed leading term \(h_{k+1}=-\frac{b}{a}h_{k}^{2}(1+o(1))\) is inconsistent with that cancellation, and is the wrong constant: the legal limit is \(\frac{b(1+a)}{2a}\), see the correct proof.}

\errpanel{
The first block is the right proof of local order \(2\), once ``\(k\) large'' is replaced by ``\(x_{0}\) close'', the denominator is shown to be nonzero, and \(O(h_{k}^{2})\) is converted into a neighborhood bound \(|h_{k+1}|\le M|h_{k}|^{2}\) with \(|h_{0}|<1/(2M)\).

The second block cannot be repaired by keeping \(O(h^{3})\). Under \(C^{2}\) one must write \(o(h^{2})\) (or Lagrange of order \(1\)). The \(h^{2}\) coefficient of \(f(x+f(x))\) must retain \(\tfrac12 f''(x_{*})(1+f'(x_{*}))^{2}\). The asymptotic constant is \(\frac{f''(x_{*})(1+f'(x_{*}))}{2f'(x_{*})}\), not \(-\frac{f''(x_{*})}{f'(x_{*})}\).
}

\subsubsection{Solution of Problem 1}

\paragraph{Answer.}
If \(x_{0}\) is sufficiently close to \(x_{*}\), the iteration is well-defined in a neighborhood of \(x_{*}\), \(x_{k}\to x_{*}\), and the local \(Q\)-order is \(2\): there exist \(M,\delta>0\) such that \(|x_{k}-x_{*}|<\delta\) implies

\[
\bigl|x_{k+1}-x_{*}\bigr|
\le
M\bigl|x_{k}-x_{*}\bigr|^{2}.
\]

If in addition \(x_{k}\neq x_{*}\) for all \(k\), then

\[
\lim_{k\to\infty}
\frac{x_{k+1}-x_{*}}{(x_{k}-x_{*})^{2}}
=
\frac{f''(x_{*})\bigl(1+f'(x_{*})\bigr)}{2f'(x_{*})}.
\]

\begin{proof}[Correct proof]
Write \(a\coloneqq f'(x_{*})\), \(b\coloneqq f''(x_{*})\), \(h_{k}\coloneqq x_{k}-x_{*}\), and

\[
\varphi(x)
\coloneqq
x
-
\frac{f(x)^{2}}{f\bigl(x+f(x)\bigr)-f(x)},
\]

so that \(x_{k+1}=\varphi(x_{k})\) whenever the denominator is nonzero. Since \(f\in C^{2}(\mathbb{R})\) one has \(a\neq 0\) and \(f'\) continuous, and \cref{cor:c2-peano} supplies

\begin{equation}
\label{eq:p1-f-exp}
f(x_{*}+h)
=
ah+\tfrac12 b h^{2}+o(h^{2})
=
ah+O(h^{2}),
\qquad h\to 0.
\end{equation}

\textbf{Step 1 (denominator).}
Let \(h=x-x_{*}\) and \(\delta\coloneqq h+f(x_{*}+h)\). Then \(\delta=(1+a)h+O(h^{2})=O(h)\), and a second application of \eqref{eq:p1-f-exp} gives

\[
f(x_{*}+\delta)
=
a\delta+O(\delta^{2})
=
a(1+a)h+O(h^{2}).
\]

Subtracting \(f(x_{*}+h)=ah+O(h^{2})\) yields

\[
f\bigl(x+f(x)\bigr)-f(x)
=
a^{2}h+O(h^{2}).
\]

Choose \(\delta_{0}>0\) so that \(|h|<\delta_{0}\) implies \(\bigl|O(h^{2})\bigr|\le \tfrac12 a^{2}|h|\). Then for \(0<|h|<\delta_{0}\) the denominator is nonzero (it has the sign of \(a^{2}h\)), and \(\varphi\) is well-defined on the punctured ball. At \(h=0\) one has \(f(x_{*})=0\), so \(x_{*}\) is a fixed point of any continuous extension of \(\varphi\).

\textbf{Step 2 (quadratic bound).}
On the same ball,

\[
f(x)^{2}
=
\bigl(ah+O(h^{2})\bigr)^{2}
=
a^{2}h^{2}+O(h^{3}),
\]

hence

\[
\varphi(x)-x_{*}
=
h
-
\frac{a^{2}h^{2}+O(h^{3})}{a^{2}h+O(h^{2})}
=
h
-
h\cdot\frac{1+O(h)}{1+O(h)}
=
O(h^{2}).
\]

Thus there exist \(M>0\) and \(\delta\in(0,\delta_{0}]\) with

\[
\bigl|\varphi(x)-x_{*}\bigr|
\le
M\bigl|x-x_{*}\bigr|^{2}
\qquad\text{whenever}\qquad
|x-x_{*}|<\delta.
\]

\textbf{Step 3 (local attraction).}
Shrink \(\delta\) further so that \(M\delta\le\tfrac12\). If \(|h_{0}|<\delta\), then \(|h_{1}|\le M|h_{0}|^{2}\le\tfrac12|h_{0}|<\delta\), and by induction the orbit remains in the ball, stays well-defined, and satisfies \(|h_{k+1}|\le\tfrac12|h_{k}|\). Therefore \(h_{k}\to 0\), i.e.\ \(x_{k}\to x_{*}\). The same bound \(|h_{k+1}|\le M|h_{k}|^{2}\) is the local \(Q\)-order \(2\) estimate. Uniqueness of the zero is used only to identify the limit: if a well-defined orbit converged to some \(L\), passing to the limit in the recurrence would force \(f(L)=0\), hence \(L=x_{*}\).

\textbf{Step 4 (asymptotic constant).}
Return to \eqref{eq:p1-f-exp} and keep the quadratic term. Write \(r(h)=o(h^{2})\) for the Peano remainder, so \(f(x_{*}+h)=ah+\tfrac12 bh^{2}+r(h)\). Then

\[
\delta
=
h+f(x_{*}+h)
=
(1+a)h+\tfrac12 bh^{2}+r(h).
\]

In all cases \(\delta=O(h)\), so \(r(\delta)=o(\delta^{2})=o(h^{2})\). Expanding \(\delta^{2}=(1+a)^{2}h^{2}+o(h^{2})\) (valid even if \(a=-1\), in which case \(\delta=o(h)\) and \(\delta^{2}=o(h^{2})\)) gives

\[
f(x_{*}+\delta)
=
a\delta+\tfrac12 b\delta^{2}+r(\delta)
=
a(1+a)h
+
\bigl(\tfrac12 ab+\tfrac12 b(1+a)^{2}\bigr)h^{2}
+
o(h^{2}).
\]

The denominator is therefore

\[
\begin{aligned}
D(h)
&\coloneqq
f(x_{*}+\delta)-f(x_{*}+h)
\\
&=
a^{2}h
+
\bigl(\tfrac12 ab+\tfrac12 b(1+a)^{2}-\tfrac12 b\bigr)h^{2}
+
o(h^{2})
\\
&=
a^{2}h+\tfrac12 ab(3+a)\,h^{2}+o(h^{2}).
\end{aligned}
\]

The numerator is

\[
N(h)
\coloneqq
f(x_{*}+h)^{2}
=
\bigl(ah+\tfrac12 bh^{2}+r(h)\bigr)^{2}
=
a^{2}h^{2}+ab\,h^{3}+o(h^{3}),
\]

where \(2ah\cdot r(h)=o(h^{3})\) uses only \(r(h)=o(h^{2})\), not an \(O(h^{3})\) remainder. For \(h\neq 0\) small,

\[
\frac{N(h)}{D(h)}
=
h\cdot
\frac{1+\frac{b}{a}h+o(h)}{1+\frac{b}{2a}(3+a)h+o(h)}
=
h
+
\Bigl(\frac{b}{a}-\frac{b(3+a)}{2a}\Bigr)h^{2}
+
o(h^{2}).
\]

The coefficient simplifies to \(-\frac{b(1+a)}{2a}\), hence

\[
\varphi(x_{*}+h)-x_{*}
=
h-\frac{N(h)}{D(h)}
=
\frac{b(1+a)}{2a}\,h^{2}+o(h^{2}).
\]

Along the locally convergent orbit of Step~3 one may divide by \(h_{k}^{2}\neq 0\) and pass to the limit.

If \(a=-1\) or \(b=0\), the displayed constant vanishes and the same identity yields only \(h_{k+1}=o(h_{k}^{2})\); the \(Q\)-order is then at least \(2\), and a higher order would require more regularity than \(C^{2}\).
\end{proof}

\begin{remark}[Do not start from \(x_{k}\) or \(h_{k}\)]
\label{rem:p1-h-first}
The manuscript order is the obstruction, not a missing identity. Peano's theorem (\cref{thm:peano}, \cref{cor:c2-peano}) is a statement about a dummy increment: as \(h\to 0\),
\[
f(x_{*}+h)=ah+O(h^{2}).
\]
It is not a statement about the orbit \(\{x_{k}\}\). Writing \(x_{k}\) (or \(h_{k}=x_{k}-x_{*}\)) first, and then quoting \(O(h_{k}^{2})\), silently assumes that \(h_{k}\) is already small --- which is half of what the exam asks to prove, and is exactly the illegitimate ``for \(k\) large enough'' at the start of the notes.

The legal order is the one in the proof above.

\begin{enumerate}[label=\textup{\arabic*.}]
    \item Introduce a dummy \(h=x-x_{*}\), not the sequence \(h_{k}\). All Taylor/Peano expansions, the denominator \(a^{2}h+O(h^{2})\), and the inversion \((1+O(h))^{-1}=1+O(h)\) are functions of this \(h\).
    \item Restrict \(h\) to a neighborhood of \(0\): choose \(\delta_{0}\) so that the \(O(h^{2})\) remainder is at most \(\tfrac12 a^{2}|h|\) whenever \(|h|<\delta_{0}\). Only inside that ball are the expansions valid and \(\varphi\) well-defined. (The range is a punctured interval around \(0\), not a statement about \(k\).)
    \item Convert the \(O(\,\cdot\,)\) as \(h\to 0\) into a uniform bound \(|\varphi(x)-x_{*}|\le M|h|^{2}\) on a possibly smaller ball \(|h|<\delta\), then shrink further to \(M\delta\le\tfrac12\).
    \item Only now specialize to the orbit: if \(|h_{0}|<\delta\), induction keeps every \(h_{k}\) inside the same ball, so the Peano hypotheses remain in force along the sequence, and \(h_{k}\to 0\).
\end{enumerate}

Step~4 of the proof (the asymptotic constant) is the same pattern: the identity
\[
\varphi(x_{*}+h)-x_{*}
=
\frac{b(1+a)}{2a}\,h^{2}+o(h^{2})
\]
is first proved for dummy \(h\to 0\); dividing by \(h_{k}^{2}\) is legal only after Step~3 has already put \(h_{k}\to 0\). Reversing the order --- expand in \(x_{k}\) first, box \(h_{k}\) later --- is why the notes cannot close the argument.
\end{remark}

\subsection{Problem 2. [20 points]}
For any continuous function \(f\), define \(\langle f, g \rangle = \int_{-1}^{1} f(x)g(x)\,\mathrm{d}x\) and \(\|f\|_2 = \sqrt{\langle f, f \rangle}\). Write
\[
a_{n}=\frac{(2n)!}{2^{n}(n!)^{2}}.
\]
Define the \(n\)-th monic Legendre polynomial by: \(P_{0}(x)=1\), \(P_{1}(x)=x\), \(P_{n+1}(x)=xP_{n}(x)-\frac{n^{2}}{4n^{2}-1}P_{n-1}(x)\).
\begin{enumerate}[label=(\alph*)]
    \item Compute \(\langle P_{n}, P_{m} \rangle\), for \(m,n\in\mathbb{N}\).
    \item Prove that \(P_{n}\) solves
    \[
    \min_{\substack{\text{polynomial }p\\ \deg p\le n,\, p^{(n)}(1)=n!}}\|p\|_{2}
    \]
    and compute the minimal value.
    \item For \(f\in C^{n+1}[-1,1]\), choose the interpolation nodes \(x_{0},x_{1},\ldots,x_{n}\) as the \(n+1\) zeros of \(P_{n+1}(x)\), and denote its degree-\(n\) interpolation polynomial by \(L_{n}\). Prove
    \[
    \|f-L_{n}\|_{2}\le\frac{2\|f^{(n+1)}\|_{2}}{(2n+3)a_{n+1}(n+1)!}.
    \]
    \item For \(f\in C^{2n+2}[-1,1]\), derive the Gauss quadrature formula of \(I(f)=\int_{-1}^{1}f(x)\,\mathrm{d}x\) for \(n+1\) quadrature nodes, written by \(I_{n+1}(f)\), and estimate its error \(|I(f)-I_{n+1}(f)|\).
\end{enumerate}

\subsubsection{Hints for Problem 2}

The family \(\{P_n\}\) is the monic orthogonal polynomial sequence for \(\langle\cdot,\cdot\rangle\) on \([-1,1]\). The given three-term recurrence is the monic form of Bonnet's formula. Write \(Q_n\) for the \emph{standard} Legendre polynomial (Rodrigues / \(Q_n(1)=1\)); then \(a_n\) is its leading coefficient, \(Q_n=a_n P_n\), and \(|Q_n|\le 1\) on \([-1,1]\). (The exam's \(L_n\) is the interpolant in (c), not \(Q_n\).) Do not quote \(\langle Q_n,Q_m\rangle=\frac{2}{2n+1}\delta_{nm}\) as an unproved oracle unless you derive it.

\begin{enumerate}[label=\textup{(\alph*)}]
    \item
    Orthogonality first, then the squared norm. A naive induction on the given recurrence does \emph{not} close: it only yields \(\langle P_{n+1},P_{n-1}\rangle=\|P_n\|_2^2-\gamma_n\|P_{n-1}\|_2^2\), which is circular until the norm ratio is known. Identify \(\{P_n\}\) with the monic Rodrigues sequence, whose three-term coefficient is \(\beta_n=\|P_n\|_2^2/\|P_{n-1}\|_2^2=\gamma_n\). Then \(\langle xP_n,P_{n-1}\rangle\) equals both \(\|P_n\|_2^2\) and \(\gamma_n\|P_{n-1}\|_2^2\) without invoking \(P_{n-2}\). Start from \(\|P_0\|_2^2=2\) and telescope. Equivalent closed form, once \(Q_n=a_n P_n\) is in hand:

    \[
    \langle P_n,P_m\rangle
    =
    \frac{2}{(2n+1)a_n^2}\,\delta_{nm}.
    \]

    \item
    The constraint \(p^{(n)}(1)=n!\) is exactly ``\(p\) is monic of degree \(n\)'': \(p^{(n)}\equiv n!\,\mathrm{lc}(p)\). Write \(p=P_n+q\) with \(q\in\mathcal{P}_{n-1}\). Expand \(\|P_n+q\|_2^2\) and kill the cross term by (a). Uniqueness is the same identity. The minimal value is \(\|P_n\|_2\), not \(\|P_n\|_2^2\).

    \item
    The nodes are the zeros of \(P_{n+1}\), and \(P_{n+1}\) is already monic, so the nodal polynomial is \(\omega=P_{n+1}\). The pointwise formula \(f(x)-L_n(x)=\frac{f^{(n+1)}(\xi_x)}{(n+1)!}\omega(x)\) is the wrong shape for an \(L^2\) bound on \(f^{(n+1)}\): \(\xi_x\) is not a function of a dummy integration variable. The printed \(L^2\)--\(L^2\) constant is moreover too small for \(n\ge 1\) (the monomial \(x^{n+1}/(n+1)!\) is a counter-example). What one can prove is
    \[
    \|f-L_n\|_2
    \le
    \frac{\|P_{n+1}\|_2}{(n+1)!}\,\|f^{(n+1)}\|_\infty.
    \]
    The Peano integral remainder is the correct vehicle if one insists on controlling \(\|f^{(n+1)}\|_2\), but Hilbert--Schmidt on the kernel does not recover the printed factor. A crude substitute \(|Q_{n+1}|\le 1\) only yields the weaker factor \(1/a_{n+1}\).

    \item
    Take the same nodes as in (c). The interpolatory rule \(I_{n+1}(f)=\sum_{j=0}^n w_j f(x_j)\) with \(w_j=\int_{-1}^1 \ell_j\) is Gauss--Legendre. Algebraic precision \(2n+1\) (not \(n\)): for \(r\in\mathcal{P}_{2n+1}\) write \(r=qP_{n+1}+s\) with \(q,s\in\mathcal{P}_n\), then \(\int qP_{n+1}=0\) by (a) and \(I_{n+1}(r)=I(s)=I(r)\). The error for \(f\in C^{2n+2}\) is the monic Gauss remainder

    \[
    I(f)-I_{n+1}(f)
    =
    \frac{f^{(2n+2)}(\xi)}{(2n+2)!}\,\|P_{n+1}\|_2^2,
    \]

    so the bound is \(\frac{2}{(2n+3)a_{n+1}^2(2n+2)!}\|f^{(2n+2)}\|_\infty\). Note the extra \(a_{n+1}\) relative to (c), and the jump from order \(n+1\) to \(2n+2\).
\end{enumerate}

\setcounter{section}{2}

\subsubsection{Theorems used in Problem 2}

Write \(\mathcal{P}_{k}\) for the space of polynomials of degree at most \(k\), and \(\gamma_{n}\coloneqq\frac{n^{2}}{4n^{2}-1}\) for \(n\ge 1\). The given recurrence is \(P_{n+1}=xP_{n}-\gamma_{n}P_{n-1}\) for \(n\ge 1\), with \(P_{0}=1\) and \(P_{1}=x\). Each \(P_{n}\) is monic of degree \(n\) by a trivial induction. The inner product is symmetric and satisfies \(\langle xf,g\rangle=\langle f,xg\rangle\).

\begin{lemma}[Uniqueness of monic orthogonal polynomials]
\label{lem:monic-op-unique}
For each \(n\ge 0\) there is at most one monic \(p\in\mathcal{P}_{n}\) with \(\langle p,q\rangle=0\) for every \(q\in\mathcal{P}_{n-1}\) (the condition is vacuous for \(n=0\)).
\end{lemma}

\begin{proof}
If \(p\) and \(r\) are two such polynomials, then \(p-r\in\mathcal{P}_{n-1}\) and
\[
\|p-r\|_{2}^{2}
=
\langle p,p-r\rangle-\langle r,p-r\rangle
=
0,
\]
hence \(p=r\).
\end{proof}

\begin{lemma}[Rodrigues representation]
\label{lem:rodrigues}
For \(n\ge 0\) define
\[
R_{n}(x)
\coloneqq
\frac{n!}{(2n)!}\,
\frac{\mathrm{d}^{n}}{\mathrm{d}x^{n}}(x^{2}-1)^{n}.
\]
Then \(R_{n}\) is monic of degree \(n\), \(\langle R_{n},q\rangle=0\) for all \(q\in\mathcal{P}_{n-1}\), and the coefficient of \(x^{n-2}\) in \(R_{n}\) (for \(n\ge 1\)) is
\[
c_{n}
=
-\frac{n(n-1)}{2(2n-1)}.
\]
In particular \(R_{0}=P_{0}\) and \(R_{1}=P_{1}\).
\end{lemma}

\begin{proof}
Expand \((x^{2}-1)^{n}=x^{2n}-n x^{2n-2}+O(x^{2n-4})\). Differentiating \(n\) times,
\[
\frac{\mathrm{d}^{n}}{\mathrm{d}x^{n}}x^{2n}
=
\frac{(2n)!}{n!}\,x^{n},
\qquad
\frac{\mathrm{d}^{n}}{\mathrm{d}x^{n}}x^{2n-2}
=
\frac{(2n-2)!}{(n-2)!}\,x^{n-2}
\quad(n\ge 2),
\]
and the \(n=1\) case is \(R_{1}(x)=\frac12\cdot 2x=x\), so \(c_{1}=0\). For \(n\ge 2\),
\[
R_{n}(x)
=
x^{n}
-
\frac{n!\cdot n\cdot(2n-2)!}{(2n)!\,(n-2)!}\,x^{n-2}
+
O(x^{n-4})
=
x^{n}
-
\frac{n(n-1)}{2(2n-1)}\,x^{n-2}
+
O(x^{n-4}).
\]
The leading term is \(x^{n}\) in every degree (including \(R_{0}=1\)), so \(R_{n}\) is monic. The displayed formula for \(c_{n}\) also holds at \(n=1\).

For orthogonality, write \(u_{n}(x)=(x^{2}-1)^{n}=(x-1)^{n}(x+1)^{n}\). Then \(u_{n}\) has zeros of order \(n\) at \(x=\pm 1\), hence \(u_{n}^{(j)}(\pm 1)=0\) for \(0\le j\le n-1\). Integrating by parts \(n\) times against any \(q\in\mathcal{P}_{n-1}\) therefore produces vanishing boundary terms and
\[
\langle R_{n},q\rangle
=
\frac{n!}{(2n)!}
\int_{-1}^{1} u_{n}^{(n)}(x)\,q(x)\,\mathrm{d}x
=
\frac{n!}{(2n)!}\,(-1)^{n}
\int_{-1}^{1} u_{n}(x)\,q^{(n)}(x)\,\mathrm{d}x
=
0,
\]
since \(q^{(n)}\equiv 0\).
\end{proof}

\begin{lemma}[Three-term recurrence of monic orthogonal polynomials]
\label{lem:monic-3term}
Let \(\{\pi_{n}\}_{n\ge 0}\) be the unique sequence of monic polynomials with \(\pi_{n}\) of degree \(n\) and \(\pi_{n}\perp\mathcal{P}_{n-1}\). Then \(\pi_{0}=1\), \(\pi_{1}=x\), and
\[
\pi_{n+1}(x)
=
x\,\pi_{n}(x)
-
\beta_{n}\,\pi_{n-1}(x),
\qquad
\beta_{n}
=
\frac{\|\pi_{n}\|_{2}^{2}}{\|\pi_{n-1}\|_{2}^{2}}
\quad(n\ge 1).
\]
Moreover \(\pi_{n}(-x)=(-1)^{n}\pi_{n}(x)\).
\end{lemma}

\begin{proof}
Uniqueness is \cref{lem:monic-op-unique}; existence is \cref{lem:rodrigues}. The polynomial \(\widetilde{\pi}_{n}(x)\coloneqq(-1)^{n}\pi_{n}(-x)\) is likewise monic of degree \(n\). For \(q\in\mathcal{P}_{n-1}\) the substitution \(y=-x\) gives
\[
\langle\widetilde{\pi}_{n},q\rangle
=
(-1)^{n}
\int_{-1}^{1}\pi_{n}(-x)\,q(x)\,\mathrm{d}x
=
(-1)^{n}
\langle\pi_{n},q(-\,\cdot\,)\rangle
=
0,
\]
so \(\widetilde{\pi}_{n}=\pi_{n}\) by uniqueness. In particular \(x\pi_{n}(x)^{2}\) is odd, hence \(\langle x\pi_{n},\pi_{n}\rangle=0\).

Now \(x\pi_{n}-\pi_{n+1}\) has degree at most \(n\). Expanding in the basis \(\{\pi_{j}\}_{0\le j\le n}\) and taking inner products with \(\pi_{k}\) for \(k\le n-2\) (an empty range when \(n=1\)) yields
\[
\langle x\pi_{n}-\pi_{n+1},\pi_{k}\rangle
=
\langle\pi_{n},x\pi_{k}\rangle
=
0,
\]
because \(\deg(x\pi_{k})\le n-1\) and \(\pi_{n+1}\perp\mathcal{P}_{n}\). The coefficient of \(\pi_{n}\) is
\[
\frac{\langle x\pi_{n},\pi_{n}\rangle}{\|\pi_{n}\|_{2}^{2}}
=
0.
\]
The remaining coefficient is
\[
\frac{\langle x\pi_{n},\pi_{n-1}\rangle}{\|\pi_{n-1}\|_{2}^{2}}
=
\frac{\langle\pi_{n},x\pi_{n-1}\rangle}{\|\pi_{n-1}\|_{2}^{2}}
=
\frac{\|\pi_{n}\|_{2}^{2}}{\|\pi_{n-1}\|_{2}^{2}},
\]
where the last step uses that \(x\pi_{n-1}-\pi_{n}\) has degree at most \(n-1\). This is the claimed recurrence. Finally \(\pi_{0}=1\), and \(\pi_{1}(x)=x-c\) with \(\langle\pi_{1},1\rangle=0\) forces \(c=0\).
\end{proof}

\begin{lemma}[Simple zeros in \((-1,1)\)]
\label{lem:op-zeros}
Let \(\{\pi_{n}\}\) be the monic orthogonal sequence of \cref{lem:monic-3term}. For every \(n\ge 1\), the polynomial \(\pi_{n}\) has \(n\) distinct roots in \((-1,1)\).
\end{lemma}

\begin{proof}
Let \(r_{1},\ldots,r_{k}\) be the points of \((-1,1)\) at which \(\pi_{n}\) changes sign, and write \(\pi(x)\coloneqq\prod_{j=1}^{k}(x-r_{j})\) (the product is \(1\) if \(k=0\)). Then \(\pi_{n}\pi\) does not change sign on \([-1,1]\), and is a non-zero polynomial, hence \(\langle\pi_{n},\pi\rangle\neq 0\). But \(\pi_{n}\perp\mathcal{P}_{n-1}\), so necessarily \(\deg\pi=k\ge n\). Combined with \(k\le n\) this yields \(k=n\), and every root is simple (each contributes a sign change).
\end{proof}

\begin{lemma}[Standard Legendre polynomial]
\label{lem:std-legendre-bound}
Write \(u(x)\coloneqq(x^{2}-1)^{n}\) and
\[
Q_{n}(x)
\coloneqq
\frac{1}{2^{n}n!}\,u^{(n)}(x)
=
a_{n}R_{n}(x).
\]
Then \(Q_{n}\) satisfies Legendre's equation
\[
(1-x^{2})Q_{n}''(x)-2x Q_{n}'(x)+n(n+1)Q_{n}(x)=0,
\]
one has \(Q_{n}(1)=1\) and \(Q_{n}(-1)=(-1)^{n}\), and \(\lvert Q_{n}(x)\rvert\le 1\) for all \(x\in[-1,1]\). After the identification \(P_{n}=R_{n}\) in (a), this is \(Q_{n}=a_{n}P_{n}\).
\end{lemma}

\begin{proof}
The identity \((x^{2}-1)u'=2nxu\) is immediate. Differentiate it \(n+1\) times and write \(v\coloneqq u^{(n)}\). Leibniz's rule on the left retains only derivatives of \(x^{2}-1\) of order at most \(2\):
\[
(x^{2}-1)v''+2(n+1)xv'+\binom{n+1}{2}\cdot 2\,v
=
2n\bigl(xv'+(n+1)v\bigr).
\]
Since \(\binom{n+1}{2}\cdot 2=n(n+1)\), rearranging gives \((x^{2}-1)v''+2xv'-n(n+1)v=0\), or equivalently the claimed equation for \(Q_{n}=v/(2^{n}n!)\).

At \(x=1\), write \(u(x)=(x-1)^{n}(x+1)^{n}\). Leibniz's rule retains only the term that differentiates \((x-1)^{n}\) exactly \(n\) times, giving \(u^{(n)}(1)=n!\,2^{n}\). Thus \(Q_{n}(1)=1\). The function \(u\) is even or odd with \(n\), so \(Q_{n}\) has the same parity and \(Q_{n}(-1)=(-1)^{n}\).

If \(n=0\) then \(Q_{0}\equiv 1\). If \(n\ge 1\), set
\[
E(x)
\coloneqq
(1-x^{2})\bigl(Q_{n}'(x)\bigr)^{2}+n(n+1)Q_{n}(x)^{2}.
\]
Differentiating and substituting \((1-x^{2})Q_{n}''+n(n+1)Q_{n}=2x Q_{n}'\) produces
\[
E'(x)
=
2Q_{n}'(x)\Bigl(-x Q_{n}'(x)+(1-x^{2})Q_{n}''(x)+n(n+1)Q_{n}(x)\Bigr)
=
2x\bigl(Q_{n}'(x)\bigr)^{2}.
\]
Thus \(E\) is decreasing on \([-1,0]\) and increasing on \([0,1]\), so its maximum on \([-1,1]\) occurs at an endpoint. There \(E(\pm 1)=n(n+1)\). A fortiori \(n(n+1)Q_{n}(x)^{2}\le E(x)\le n(n+1)\), hence \(\lvert Q_{n}(x)\rvert\le 1\).
\end{proof}

\begin{proposition}[Taylor remainders by repeated integration by parts]
\label{prop:taylor-ibp}
Let \(k\ge 0\) and let \(I\) be a compact interval. Suppose \(f^{(k+1)}\) is Riemann integrable on \(I\) (in particular this holds if \(f\in C^{k+1}(I)\), and also if \(f^{(k+1)}\in L^{1}(I)\) with \(f^{(k)}\) absolutely continuous). Write
\[
T_{k}(x;a)
\coloneqq
\sum_{j=0}^{k}\frac{f^{(j)}(a)}{j!}(x-a)^{j}
\]
for the Taylor polynomial of degree \(k\) at \(a\in I\). Then for every \(x\in I\),
\begin{equation}
\label{eq:taylor-integral}
f(x)-T_{k}(x;a)
=
\frac{1}{k!}
\int_{a}^{x}(x-t)^{k}f^{(k+1)}(t)\,\mathrm{d}t.
\end{equation}
If in addition \(f^{(k+1)}\) is continuous on \(I\), there exists \(\xi\) strictly between \(a\) and \(x\) (or \(\xi=a\) if \(x=a\)) such that
\begin{equation}
\label{eq:taylor-lagrange}
f(x)-T_{k}(x;a)
=
\frac{f^{(k+1)}(\xi)}{(k+1)!}(x-a)^{k+1}.
\end{equation}
These are Qingshu's (10.4) and (10.3) respectively \cite{qingshu11}: the integral remainder is obtained by integrating by parts \(k\) times from the fundamental theorem; the Lagrange remainder is the same integral after the weighted mean value theorem.
\end{proposition}

\begin{proof}
The case \(x=a\) is \(0=0\). Fix \(x\neq a\) in \(I\). The fundamental theorem gives
\[
f(x)-f(a)
=
\int_{a}^{x}f'(t)\,\mathrm{d}t.
\]
Integrate by parts with \(u=f'(t)\) and \(\mathrm{d}v=\mathrm{d}t\), using the primitive \(v=-(x-t)\) (so that the boundary term at \(t=x\) vanishes):
\[
\int_{a}^{x}f'(t)\,\mathrm{d}t
=
\Bigl[-f'(t)(x-t)\Bigr]_{a}^{x}
+
\int_{a}^{x}(x-t)f''(t)\,\mathrm{d}t
=
f'(a)(x-a)
+
\int_{a}^{x}(x-t)f''(t)\,\mathrm{d}t.
\]
Repeating with \(u=f''(t)\) and \(\mathrm{d}v=(x-t)\,\mathrm{d}t\), now \(v=-(x-t)^{2}/2\), produces
\[
\int_{a}^{x}(x-t)f''(t)\,\mathrm{d}t
=
f''(a)\frac{(x-a)^{2}}{2}
+
\frac{1}{2}
\int_{a}^{x}(x-t)^{2}f'''(t)\,\mathrm{d}t.
\]
After \(k\) steps the pattern is
\[
\int_{a}^{x}\frac{(x-t)^{j}}{j!}f^{(j+1)}(t)\,\mathrm{d}t
=
\frac{f^{(j+1)}(a)}{(j+1)!}(x-a)^{j+1}
+
\int_{a}^{x}\frac{(x-t)^{j+1}}{(j+1)!}f^{(j+2)}(t)\,\mathrm{d}t,
\]
which is integration by parts with \(u=f^{(j+1)}(t)\) and \(v=-(x-t)^{j+1}/(j+1)!\). Starting from \(j=0\) and stopping at \(j=k-1\) yields \eqref{eq:taylor-integral}. (If one only assumes \(f^{(k+1)}\in L^{1}\), the last step is the definition of the absolutely continuous primitive \(f^{(k)}\).)

Now assume \(f^{(k+1)}\) is continuous. On the segment joining \(a\) and \(x\), the kernel \(t\mapsto(x-t)^{k}\) does not change sign. The weighted integral mean value theorem therefore supplies \(\xi\) between \(a\) and \(x\) with
\[
\int_{a}^{x}(x-t)^{k}f^{(k+1)}(t)\,\mathrm{d}t
=
f^{(k+1)}(\xi)
\int_{a}^{x}(x-t)^{k}\,\mathrm{d}t
=
f^{(k+1)}(\xi)\frac{(x-a)^{k+1}}{k+1}.
\]
(The last evaluation is \(\bigl[-(x-t)^{k+1}/(k+1)\bigr]_{a}^{x}\), and holds whether \(x>a\) or \(x<a\).) Substituting into \eqref{eq:taylor-integral} gives \eqref{eq:taylor-lagrange}.
\end{proof}

\begin{remark}[Truncated powers as the kernel of \(D^{k+1}\)]
\label{rem:truncated-power}
The monomials \((x-t)^{k}\) in \eqref{eq:taylor-integral} are, for \(t\) between \(a\) and \(x\), the restriction of the truncated powers
\[
(x-t)_{+}^{k}
\coloneqq
\bigl(\max\{x-t,0\}\bigr)^{k}.
\]
Distributionally, \(t\mapsto (x-t)_{+}^{k}/k!\) is a fundamental solution of \(\frac{\mathrm{d}^{k+1}}{\mathrm{d}t^{k+1}}\) (up to sign, depending on whether one differentiates in \(t\) or in \(x\)): each derivative lowers the exponent, and after \(k+1\) steps one obtains a Heaviside function and then a Dirac mass at \(t=x\). Thus \eqref{eq:taylor-integral} is the statement that the Taylor remainder is the pairing of \(f^{(k+1)}\) against this Green kernel with all jets at the single node \(a\) killed. Equivalently, on the oriented interval from \(a\) to \(x\),
\[
K_{\mathrm{Taylor}}(x,t)
=
(x-t)_{+}^{k}
\qquad
\text{for \(t\) between \(a\) and \(x\)}.
\]

If one interpolates at several nodes rather than matching a jet at one point, the same construction produces Peano's kernel: the remainder functional \(R_{x}\) annihilates \(\mathcal{P}_{k}\), hence
\[
R_{x}(f)
=
\frac{1}{k!}
\int K(x,t)\,f^{(k+1)}(t)\,\mathrm{d}t,
\qquad
K(x,t)
=
R_{x}\bigl((\,\cdot\,-t)_{+}^{k}\bigr).
\]
This is \cref{thm:peano-lagrange} below (\(k=n\)). The truncated power is still the Green function of \(D^{k+1}\); the interpolation conditions merely subtract the unique element of \(\mathcal{P}_{k}\) that matches the same nodal data, so \(K(\,\cdot\,,t)\) vanishes at the nodes.

Qingshu never writes the formula \(K(x,t)=R_{x}((\,\cdot\,-t)_{+}^{n})\). What the notes do write, in Theorem~10.2 (\S10.1.4, printed pp.~175--177), is the two-endpoint Green function \(k(x,y)\) for the second-order two-point problem \(u''=f\) on \([a,b]\) with \(u(a)=u(b)=0\): start from the Heaviside kernel of \(D\), integrate once more to get a kernel of \(D^{2}\), and impose \(k(x,a)=k(x,b)=0\). That \(k(x,y)\) is the \(n=1\), two-node special case of the same idea (Lagrange interpolant of degree \(1\) at the endpoints); it is not already the kernel of the degree-\(n\) interpolant on \([-1,1]\). Finally, the scalar \(K=K(x)\) in \S10.1.3 is Rolle's constant in the \(K\)-method, unrelated to either kernel.
\end{remark}

\begin{theorem}[Lagrange remainder]
\label{thm:lagrange-remainder}
Let \(x_{0},\ldots,x_{n}\) be distinct points of \([-1,1]\), write \(\omega(x)\coloneqq\prod_{j=0}^{n}(x-x_{j})\), and let \(L_{n}(f)\) be the unique element of \(\mathcal{P}_{n}\) interpolating \(f\in C^{n+1}[-1,1]\) at these nodes. Then for every \(x\) there is \(\theta_{x}\) in the convex hull of \(\{x,x_{0},\ldots,x_{n}\}\) such that
\[
f(x)-L_{n}(f)(x)
=
\frac{f^{(n+1)}(\theta_{x})}{(n+1)!}\,\omega(x).
\]
\end{theorem}

\begin{proof}
This is Qingshu's Theorem~10.1 with \(s_{j}=0\) for all \(j\). If \(x\) is a node the identity is \(0=0\). Otherwise choose \(K\) so that
\[
F(y)
\coloneqq
f(y)-L_{n}(f)(y)-K\,\omega(y)
\]
vanishes at \(y=x\). Then \(F\) has \(n+2\) zeros, Rolle's theorem \(n+1\) times produces \(\theta_{x}\) with \(F^{(n+1)}(\theta_{x})=0\), and \(\deg L_{n}(f)\le n\) with \(\omega\) monic of degree \(n+1\) forces \(K=f^{(n+1)}(\theta_{x})/(n+1)!\).
\end{proof}

\begin{theorem}[Peano remainder for Lagrange interpolation]
\label{thm:peano-lagrange}
Let \(x_{0},\ldots,x_{n}\) be distinct points of \([-1,1]\), let \(L_{n}(f)\) be the unique element of \(\mathcal{P}_{n}\) interpolating \(f\in C^{n+1}[-1,1]\) at these nodes, and write \(R_{x}(f)\coloneqq f(x)-L_{n}(f)(x)\). Then
\[
R_{x}(f)
=
\frac{1}{n!}
\int_{-1}^{1} K(x,t)\,f^{(n+1)}(t)\,\mathrm{d}t,
\qquad
K(x,t)
=
R_{x}\bigl((\,\cdot\,-t)_{+}^{n}\bigr),
\]
where \(y_{+}\coloneqq\max\{y,0\}\). In particular \(R_{x}\) vanishes on \(\mathcal{P}_{n}\).
\end{theorem}

\begin{proof}
The functional \(R_{x}\) is linear in \(f\) and vanishes on \(\mathcal{P}_{n}\) by uniqueness of interpolation. As in \cref{rem:truncated-power}, Taylor's formula with integral remainder at the left endpoint (\cref{prop:taylor-ibp}, \(a=-1\) and \(k=n\)) rewrites the remainder as a pairing against truncated powers:
\[
f(y)
=
\sum_{k=0}^{n}\frac{f^{(k)}(-1)}{k!}(y+1)^{k}
+
\frac{1}{n!}
\int_{-1}^{y}(y-t)^{n}f^{(n+1)}(t)\,\mathrm{d}t
=
T_{n}(y)
+
\frac{1}{n!}
\int_{-1}^{1}(y-t)_{+}^{n}f^{(n+1)}(t)\,\mathrm{d}t.
\]
The sum \(T_{n}\) lies in \(\mathcal{P}_{n}\), so \(R_{x}(T_{n})=0\). Applying \(R_{x}\) under the integral (justified by uniform convergence of Riemann sums, \(K(x,\cdot)\) being piecewise polynomial) gives the claim.
\end{proof}

\begin{theorem}[Hermite remainder at double nodes]
\label{thm:hermite-double}
Let \(x_{0},\ldots,x_{n}\) be distinct, let \(\omega(x)\coloneqq\prod_{j=0}^{n}(x-x_{j})\), and let \(H\in\mathcal{P}_{2n+1}\) interpolate \(f\in C^{2n+2}[-1,1]\) in the Hermite sense \(H(x_{j})=f(x_{j})\) and \(H'(x_{j})=f'(x_{j})\) for each \(j\). Then for every \(x\) there is \(\theta_{x}\) between \(x\) and the nodes such that
\[
f(x)-H(x)
=
\frac{f^{(2n+2)}(\theta_{x})}{(2n+2)!}\,\omega(x)^{2}.
\]
\end{theorem}

\begin{proof}
This is Qingshu's Theorem~10.1 with \(s_{j}=1\) for all \(j\), so \(\sum(s_{j}+1)=2n+2\) and the nodal factor is \(\omega^{2}\). For a self-contained Rolle argument: freeze \(x\) not a node, choose \(K\) so that
\[
F(y)
\coloneqq
f(y)-H(y)-K\,\omega(y)^{2}
\]
vanishes at \(y=x\). Then \(F\) has zeros of order at least \(2\) at each \(x_{j}\) and a further zero at \(x\), hence at least \(2n+3\) zeros counting multiplicity. Rolle's theorem \(2n+2\) times produces \(\theta_{x}\) with \(F^{(2n+2)}(\theta_{x})=0\). Since \(\deg H\le 2n+1\) and \(\omega^{2}\) is monic of degree \(2n+2\), this forces \(K=f^{(2n+2)}(\theta_{x})/(2n+2)!\).
\end{proof}

\begin{lemma}[Gauss quadrature precision]
\label{lem:gauss-precision}
Let \(\{\pi_{k}\}\) be as in \cref{lem:monic-3term}. Write \(x_{0},\ldots,x_{n}\) for the zeros of \(\pi_{n+1}\), which are simple in \((-1,1)\) by \cref{lem:op-zeros}, and write \(\ell_{j}\) for the associated Lagrange basis. Define
\[
I_{n+1}(f)
\coloneqq
\sum_{j=0}^{n}w_{j}f(x_{j}),
\qquad
w_{j}
\coloneqq
\int_{-1}^{1}\ell_{j}(x)\,\mathrm{d}x.
\]
Then \(w_{j}>0\) and \(I_{n+1}(p)=\int_{-1}^{1}p\) for every \(p\in\mathcal{P}_{2n+1}\).
\end{lemma}

\begin{proof}
Each \(\ell_{j}\) has degree \(n\), so \(\ell_{j}^{2}-\ell_{j}\) has degree at most \(2n\) and vanishes at every node (at \(x_{j}\) one has \(1-1=0\); at \(x_{k}\) for \(k\neq j\) one has \(0-0=0\)). Thus \(\ell_{j}^{2}-\ell_{j}=\pi_{n+1}q_{j}\) for some \(q_{j}\in\mathcal{P}_{n-1}\). Orthogonality of \(\pi_{n+1}\) to \(\mathcal{P}_{n}\) yields \(\int_{-1}^{1}(\ell_{j}^{2}-\ell_{j})=0\), hence \(w_{j}=\int_{-1}^{1}\ell_{j}^{2}\). The integrand is continuous and not identically zero, so \(w_{j}>0\).

If \(p\in\mathcal{P}_{n}\) then \(p=\sum_{j=0}^{n}p(x_{j})\ell_{j}\), so \(I_{n+1}(p)=\int_{-1}^{1}p\) already. If \(p\in\mathcal{P}_{2n+1}\), divide as \(p=q\pi_{n+1}+s\) with \(q,s\in\mathcal{P}_{n}\). Then \(\int_{-1}^{1}q\pi_{n+1}=0\) by orthogonality, while \(p(x_{j})=s(x_{j})\), so \(I_{n+1}(p)=I_{n+1}(s)=\int_{-1}^{1}s=\int_{-1}^{1}p\).
\end{proof}

\subsubsection{Solution of Problem 2(a)}

\paragraph{Answer.}
For all \(m,n\in\mathbb{N}\) (including \(0\)),
\[
\langle P_{n},P_{m}\rangle
=
\frac{2}{(2n+1)a_{n}^{2}}\,\delta_{nm}
=
\frac{2^{2n+1}(n!)^{4}}{(2n)!\,(2n+1)!}\,\delta_{nm}.
\]

\begin{proof}
By \cref{lem:rodrigues,lem:monic-op-unique}, the Rodrigues polynomials \(\{R_{n}\}\) are precisely the monic orthogonal sequence of \cref{lem:monic-3term}. That lemma therefore supplies a recurrence \(R_{n+1}=xR_{n}-\beta_{n}R_{n-1}\). The expansion \(R_{n}(x)=x^{n}+c_{n}x^{n-2}+O(x^{n-4})\) from \cref{lem:rodrigues} (with no \(x^{n-1}\) term, by parity) gives
\[
R_{n+1}(x)
=
x^{n+1}+c_{n+1}x^{n-1}+O(x^{n-3}),
\qquad
xR_{n}(x)-\beta_{n}R_{n-1}(x)
=
x^{n+1}+(c_{n}-\beta_{n})x^{n-1}+O(x^{n-3}).
\]
Matching coefficients of \(x^{n-1}\) yields \(\beta_{n}=c_{n}-c_{n+1}\), hence
\[
\beta_{n}
=
c_{n}-c_{n+1}
=
-\frac{n(n-1)}{2(2n-1)}
+
\frac{(n+1)n}{2(2n+1)}
=
\frac{n}{2}
\left(
\frac{n+1}{2n+1}
-
\frac{n-1}{2n-1}
\right)
=
\frac{n}{2}\cdot\frac{2n}{4n^{2}-1}
=
\frac{n^{2}}{4n^{2}-1}
=
\gamma_{n}.
\]
(The elementary identity \((n+1)(2n-1)-(n-1)(2n+1)=2n\) was used.) Thus \(\{R_{n}\}\) satisfies the same initial conditions and the same three-term recurrence as \(\{P_{n}\}\), so \(P_{n}=R_{n}\) for every \(n\). In particular \(\langle P_{n},P_{m}\rangle=0\) whenever \(n\neq m\).

It remains to evaluate the diagonal. From \cref{lem:monic-3term} one has \(\|P_{n}\|_{2}^{2}=\gamma_{n}\|P_{n-1}\|_{2}^{2}\) for \(n\ge 1\). Equivalently, for every \(n\ge 1\),
\[
\langle xP_{n},P_{n-1}\rangle
=
\langle P_{n+1}+\gamma_{n}P_{n-1},P_{n-1}\rangle
=
\gamma_{n}\|P_{n-1}\|_{2}^{2},
\]
while \(xP_{n-1}-P_{n}\) has degree at most \(n-1\), so orthogonality gives
\[
\langle xP_{n},P_{n-1}\rangle
=
\langle P_{n},xP_{n-1}\rangle
=
\langle P_{n},P_{n}+(xP_{n-1}-P_{n})\rangle
=
\|P_{n}\|_{2}^{2}.
\]
(The second identity does not expand \(xP_{n-1}\) by the recurrence, so it is valid at \(n=1\) with no \(P_{-1}\).) Starting from \(\|P_{0}\|_{2}^{2}=\int_{-1}^{1}1\,\mathrm{d}x=2\) and telescoping,
\[
\|P_{n}\|_{2}^{2}
=
2\prod_{k=1}^{n}\gamma_{k}
=
2\prod_{k=1}^{n}\frac{k^{2}}{(2k-1)(2k+1)}
=
2\cdot
\frac{(n!)^{2}}{\displaystyle\prod_{k=1}^{n}(2k-1)\cdot\prod_{k=1}^{n}(2k+1)}.
\]
The odd factorials are
\[
\prod_{k=1}^{n}(2k-1)
=
\frac{(2n)!}{2^{n}n!},
\qquad
\prod_{k=1}^{n}(2k+1)
=
\frac{(2n+1)!}{2^{n}n!},
\]
where the second identity uses \(1\cdot 3\cdots(2n+1)=\frac{(2n+1)!}{2^{n}n!}\) and the observation that the extra factor \(1\) does not change the product. Therefore
\[
\|P_{n}\|_{2}^{2}
=
2\cdot(n!)^{2}\cdot\frac{2^{n}n!}{(2n)!}\cdot\frac{2^{n}n!}{(2n+1)!}
=
\frac{2^{2n+1}(n!)^{4}}{(2n)!\,(2n+1)!}.
\]
On the other hand \(a_{n}=\frac{(2n)!}{2^{n}(n!)^{2}}\) gives
\[
\frac{2}{(2n+1)a_{n}^{2}}
=
\frac{2}{2n+1}\cdot\frac{2^{2n}(n!)^{4}}{\bigl[(2n)!\bigr]^{2}}
=
\frac{2^{2n+1}(n!)^{4}}{(2n)!\,(2n+1)!},
\]
which is the same quantity. The empty product for \(n=0\) is \(1\), and \(a_{0}=1\), so the formula also covers \(\|P_{0}\|_{2}^{2}=2\).
\end{proof}

\subsubsection{Solution of Problem 2(b)}

\paragraph{Answer.}
The admissible polynomials are precisely the monic polynomials of degree \(n\). The unique minimiser is \(P_{n}\), and the minimal value is
\[
\|P_{n}\|_{2}
=
\frac{\sqrt{2}}{a_{n}\sqrt{2n+1}}
=
\frac{2^{n}(n!)^{2}}{\sqrt{(2n)!\,(2n+1)!}}.
\]

\begin{proof}
If \(\deg p\le n\), then \(p^{(n)}\) is the constant \(n!\,\mathrm{lc}(p)\), where the leading coefficient is interpreted as \(0\) when \(\deg p<n\). The constraint \(p^{(n)}(1)=n!\) therefore forces \(\mathrm{lc}(p)=1\), hence \(\deg p=n\) and \(p\) is monic.

Write \(p=P_{n}+q\) with \(q\in\mathcal{P}_{n-1}\). This is a bijection between monic polynomials of degree \(n\) and \(\mathcal{P}_{n-1}\). Expand
\[
\|p\|_{2}^{2}
=
\|P_{n}+q\|_{2}^{2}
=
\|P_{n}\|_{2}^{2}+2\langle P_{n},q\rangle+\|q\|_{2}^{2}.
\]
The cross term vanishes by (a), so \(\|p\|_{2}^{2}=\|P_{n}\|_{2}^{2}+\|q\|_{2}^{2}\ge\|P_{n}\|_{2}^{2}\), with equality if and only if \(q=0\). Thus \(P_{n}\) is the unique minimiser, and the minimal value is \(\|P_{n}\|_{2}\) (not \(\|P_{n}\|_{2}^{2}\)). The explicit formula is the square root of the diagonal computed in (a).
\end{proof}

\begin{remark}[Interpolation remainder, after Qingshu]
\label{rem:p2-qingshu-interp}
Parts (c)--(d) are Qingshu's Chapter~10 interpolation package \cite{qingshu11}, specialised to the inner product of this problem. Rolle plus an auxiliary function is Hermite interpolation (Qingshu, Theorem~10.1); the same identity has an integral remainder whose kernel is the Peano function built from truncated powers, as in \cref{rem:truncated-power} (Qingshu writes only the two-endpoint Green function \(k(x,y)\) in Theorem~10.2 / \S10.1.4).

\emph{Pointwise form (Qingshu, Theorem~10.1).}
Let \(x_{0}<\cdots<x_{n}\) be nodes in \([-1,1]\) and write \(\omega(x)\coloneqq\prod_{j=0}^{n}(x-x_{j})\). If \(L_{n}\) is the unique element of \(\mathcal{P}_{n}\) interpolating \(f\in C^{n+1}[-1,1]\) at these nodes, then for each \(x\) there is \(\theta_{x}\) between \(x\) and the nodes such that

\begin{equation}
\label{eq:qingshu-lagrange}
f(x)-L_{n}(x)
=
\frac{f^{(n+1)}(\theta_{x})}{(n+1)!}\,\omega(x).
\end{equation}

This is the \(K\)-method: freeze \(x\), choose \(K\) so that \(F\coloneqq f-L_{n}-K\omega\) has \(n+2\) zeros, and apply Rolle \(n+1\) times. One node of multiplicity \(k+1\) recovers Taylor with Lagrange remainder (Qingshu, (10.3)), which is also \eqref{eq:taylor-lagrange} obtained by integrating by parts rather than by Rolle; see \cref{prop:taylor-ibp}. Two simple nodes recover the mean value theorem. Hermite interpolation (matching jets \(s_{j}\geqslant 1\)) is the same statement with \(\omega(x)=\prod_{j}(x-x_{j})^{s_{j}+1}\).

\emph{Why (c) cannot stop at \eqref{eq:qingshu-lagrange}.}
The factor \(\theta_{x}\) is not a function of a dummy integration variable. Integrating \(\lvert f-L_{n}\rvert^{2}\) therefore does not produce \(\lVert f^{(n+1)}\rVert_{2}\) by Cauchy--Schwarz. Qingshu already meets this obstruction when integrating a pointwise remainder (notes after (10.14)): the identity may be integrated, but \(\theta(x)\) cannot be treated as an integration variable. That is why (c) wants the Peano form.

\emph{Integral form (Qingshu, \S10.1.4).}
The one-node case is already \eqref{eq:taylor-integral} (Qingshu, (10.4)). For several nodes the same interpolation rewrites as

\begin{equation}
\label{eq:qingshu-peano}
f(x)-L_{n}(x)
=
\frac{1}{n!}
\int_{-1}^{1} K(x,t)\,f^{(n+1)}(t)\,\mathrm{d}t,
\end{equation}

where \(K(x,t)=R_{x}((\,\cdot\,-t)_{+}^{n})\) as in \cref{rem:truncated-power,thm:peano-lagrange}. Qingshu does not write this truncated-power formula; the notes construct the kernel by integrating the Heaviside and imposing nodal conditions, and write the two-endpoint case as Theorem~10.2 with Green kernel \(k(x,y)\). Hilbert--Schmidt on \(K\) does convert \eqref{eq:qingshu-peano} into a genuine \(L^{2}\)--\(L^{2}\) bound, with constant \(\lVert K\rVert_{L^{2}([-1,1]^{2})}/n!\); this constant is \emph{not} the printed factor \(2/((2n+3)a_{n+1}(n+1)!)\). The pointwise form \eqref{eq:qingshu-lagrange} does convert into an \(L^{2}\) bound controlled by \(\lVert\omega\rVert_{2}\), provided one uses \(\lVert f^{(n+1)}\rVert_{\infty}\). For the nodes of (c), \(\omega=P_{n+1}\) (already monic), so that factor is the diagonal of (a).

\emph{What this does for the rest of Problem~2.}
Part (b) is the \(L^{2}\) analogue of choosing Chebyshev nodes: among monic polynomials of degree \(n\), \(\lVert\omega\rVert_{2}\) is minimised precisely at \(\omega=P_{n}\). Qingshu's Exercises~10.4--10.5 integrate the remainder and apply the integral mean value theorem to obtain the midpoint and trapezoid identities; Gauss quadrature in (d) is the same move at \(n+1\) Legendre zeros. Algebraic precision jumps from \(n\) to \(2n+1\) because \(\omega=P_{n+1}\) is orthogonal to \(\mathcal{P}_{n}\) by (a). Equivalently, in Qingshu's Hermite language, the quadrature error is the remainder for interpolation of type \(s_{j}=1\) at those zeros, hence involves \(P_{n+1}^{2}\) and a derivative of order \(2n+2\):

\[
I(f)-I_{n+1}(f)
=
\frac{f^{(2n+2)}(\xi)}{(2n+2)!}\,\lVert P_{n+1}\rVert_{2}^{2}.
\]
\end{remark}

\subsubsection{Solution of Problem 2(c)}

\paragraph{Answer.}
The polynomial \(P_{n+1}\) has \(n+1\) distinct zeros in \((-1,1)\), so the interpolation problem is well-posed. The printed \(L^{2}\)--\(L^{2}\) bound is false for every \(n\ge 1\). What holds is
\[
\|f-L_{n}\|_{2}
\le
\frac{\|P_{n+1}\|_{2}}{(n+1)!}\,\|f^{(n+1)}\|_{\infty}
=
\frac{\sqrt{2}}{a_{n+1}\sqrt{2n+3}\,(n+1)!}\,\|f^{(n+1)}\|_{\infty}.
\]
A coarser substitute \(\lvert P_{n+1}\rvert\le 1/a_{n+1}\) (from \(\lvert Q_{n+1}\rvert\le 1\) in \cref{lem:std-legendre-bound}) only replaces \(\|P_{n+1}\|_{2}\) by \(\sqrt{2}/a_{n+1}\).

\begin{proof}
By (a) and \cref{lem:monic-3term}, \(P_{n+1}\) is the monic orthogonal polynomial of degree \(n+1\), so \cref{lem:op-zeros} supplies \(n+1\) distinct zeros \(x_{0},\ldots,x_{n}\) in \((-1,1)\). Let \(L_{n}\) be the unique element of \(\mathcal{P}_{n}\) interpolating \(f\) at these nodes. The nodal polynomial is already monic of degree \(n+1\) and vanishes at the nodes, hence coincides with \(P_{n+1}\). \Cref{thm:lagrange-remainder} therefore yields
\[
f(x)-L_{n}(x)
=
\frac{f^{(n+1)}(\theta_{x})}{(n+1)!}\,P_{n+1}(x)
\]
for some \(\theta_{x}\in[-1,1]\). Taking absolute values and the \(L^{2}\) norm,
\[
\|f-L_{n}\|_{2}
\le
\frac{\|f^{(n+1)}\|_{\infty}}{(n+1)!}\,\|P_{n+1}\|_{2}.
\]
The explicit factor is the square root of the diagonal in (a) at index \(n+1\).

The same remainder also has the Peano form of \cref{thm:peano-lagrange}. That representation is the correct vehicle if one insists on an \(L^{2}\) control of \(f^{(n+1)}\): Cauchy--Schwarz in \(t\) followed by \(\|\cdot\|_{2}\) in \(x\) gives
\[
\|f-L_{n}\|_{2}
\le
\frac{\|K\|_{L^{2}([-1,1]^{2})}}{n!}\,\|f^{(n+1)}\|_{2}.
\]
The kernel \(K(x,\cdot)=R_{x}((\,\cdot\,-t)_{+}^{n})\) is not a scalar multiple of \(P_{n+1}(x)\), and the Hilbert--Schmidt constant is not \(2/((2n+3)a_{n+1}(n+1)!)\).

It remains to record that the printed inequality cannot be salvaged. Take \(f(x)=x^{n+1}/(n+1)!\). Then \(f^{(n+1)}\equiv 1\) and \(x^{n+1}=P_{n+1}+q\) with \(q\in\mathcal{P}_{n}\), so the interpolant of \(f\) is \(q/(n+1)!\) and \(f-L_{n}=P_{n+1}/(n+1)!\). The ratio of the two \(L^{2}\) norms is therefore
\[
\frac{\|f-L_{n}\|_{2}}{\|f^{(n+1)}\|_{2}}
=
\frac{\|P_{n+1}\|_{2}}{(n+1)!\,\sqrt{2}}.
\]
The printed constant is \(a_{n+1}\|P_{n+1}\|_{2}^{2}/(n+1)!\). The monomial would obey the printed bound if and only if \(1/\sqrt{2}\le a_{n+1}\|P_{n+1}\|_{2}=\sqrt{2/(2n+3)}\), i.e.\ if and only if \(n=0\). For \(n\ge 1\) the printed inequality already fails on this single polynomial. Explicitly, for \(n=1\) one has \(P_{2}(x)=x^{2}-1/3\), \(a_{2}=3/2\), \(\|P_{2}\|_{2}^{2}=8/45\), and \(f(x)=x^{2}/2\) gives
\[
\frac{\|f-L_{1}\|_{2}}{\|f''\|_{2}}
=
\frac{1}{3\sqrt{5}}
>
\frac{2}{15}
=
\frac{2}{(2\cdot 1+3)a_{2}\cdot 2!}.
\]

The pointwise factor \(\theta_{x}\) is not a function of a dummy integration variable, so one cannot replace \(\|f^{(n+1)}\|_{\infty}\) by \(\|f^{(n+1)}\|_{2}\) inside \cref{thm:lagrange-remainder} and still apply Cauchy--Schwarz. That is the obstruction already flagged after Qingshu's (10.14).
\end{proof}

\subsubsection{Solution of Problem 2(d)}

\paragraph{Answer.}
Let \(x_{0},\ldots,x_{n}\) be the zeros of \(P_{n+1}\), and let \(\ell_{j}\) be the associated Lagrange basis. The \((n+1)\)-node Gauss--Legendre rule is
\[
I_{n+1}(f)
=
\sum_{j=0}^{n}w_{j}f(x_{j}),
\qquad
w_{j}
=
\int_{-1}^{1}\ell_{j}(x)\,\mathrm{d}x
=
\int_{-1}^{1}\ell_{j}(x)^{2}\,\mathrm{d}x>0.
\]
It is exact on \(\mathcal{P}_{2n+1}\). For \(f\in C^{2n+2}[-1,1]\) there exists \(\xi\in(-1,1)\) such that
\[
I(f)-I_{n+1}(f)
=
\frac{f^{(2n+2)}(\xi)}{(2n+2)!}\,\|P_{n+1}\|_{2}^{2},
\]
and therefore
\[
\bigl\lvert I(f)-I_{n+1}(f)\bigr\rvert
\le
\frac{2}{(2n+3)a_{n+1}^{2}(2n+2)!}\,\|f^{(2n+2)}\|_{\infty}.
\]

\begin{proof}
The nodes exist by (c). \Cref{lem:gauss-precision} applied to \(\pi_{n+1}=P_{n+1}\) gives the stated quadrature rule, the positivity of the weights, and exactness on \(\mathcal{P}_{2n+1}\).

To estimate the error, let \(H\in\mathcal{P}_{2n+1}\) be the Hermite interpolant of \(f\) at the double nodes \((x_{j},x_{j})\), so \(H(x_{j})=f(x_{j})\) and \(H'(x_{j})=f'(x_{j})\). \Cref{thm:hermite-double} yields
\[
f(x)-H(x)
=
\frac{f^{(2n+2)}(\theta_{x})}{(2n+2)!}\,P_{n+1}(x)^{2}.
\]
Since \(f-H\) vanishes to order at least \(2\) at each simple zero of \(P_{n+1}\), the quotient \(\varphi\coloneqq(f-H)/P_{n+1}^{2}\) extends continuously to \([-1,1]\). Off the nodes one has \(\varphi(x)=f^{(2n+2)}(\theta_{x})/(2n+2)!\), so the image of \(\varphi\) lies in the interval with endpoints \(\min f^{(2n+2)}/(2n+2)!\) and \(\max f^{(2n+2)}/(2n+2)!\). Write \(I(g)\coloneqq\int_{-1}^{1}g\). Then \(I(f)-I(H)=I(\varphi P_{n+1}^{2})\). Because \(P_{n+1}^{2}\ge 0\) is not identically zero and \(\varphi\) is continuous, the integral mean value theorem supplies \(\eta\) with \(I(\varphi P_{n+1}^{2})=\varphi(\eta)\,\|P_{n+1}\|_{2}^{2}\). The intermediate value theorem for the continuous function \(f^{(2n+2)}\) then produces \(\xi\in[-1,1]\) with \(\varphi(\eta)=f^{(2n+2)}(\xi)/(2n+2)!\), hence
\[
I(f)-I(H)
=
\frac{f^{(2n+2)}(\xi)}{(2n+2)!}\,\|P_{n+1}\|_{2}^{2}.
\]
(The selection \(x\mapsto\theta_{x}\) need not be continuous; it is never integrated.) Exactness on \(\mathcal{P}_{2n+1}\) gives \(I(H)=I_{n+1}(H)\). But \(H\) and \(f\) agree at the nodes, so \(I_{n+1}(H)=I_{n+1}(f)\). This is the remainder identity. The bound follows by taking absolute values and inserting the diagonal of (a) at index \(n+1\); alternatively, \(\lvert I(\varphi P_{n+1}^{2})\rvert\le\|f^{(2n+2)}\|_{\infty}/(2n+2)!\,\|P_{n+1}\|_{2}^{2}\) uses only the range of \(\varphi\), with no mean value theorem.

The bound is sharp in the \(C^{2n+2}\) scale: for \(f(x)=x^{2n+2}/(2n+2)!\) one has \(f^{(2n+2)}\equiv 1\) and \(f-H=P_{n+1}^{2}/(2n+2)!\) (the Hermite interpolant of a monic polynomial of degree \(2n+2\) at the double zeros of \(P_{n+1}\) is the remainder of the division by \(P_{n+1}^{2}\)), so the error equals \(\|P_{n+1}\|_{2}^{2}/(2n+2)!\).
\end{proof}

\subsection{Problem 3. [20 points]}
Given \(A \in \mathbb{R}^{l \times m}\), \(B \in \mathbb{R}^{n \times p}\), \(C \in \mathbb{R}^{l \times p}\). Consider the problem \(\min_{X}\|AXB - C\|_{F}\), where \(\|\cdot\|_{F}\) is the Frobenius norm of a matrix.
\begin{enumerate}[label=(\alph*)]
    \item Prove that the minimal value is \(0\), if and only if \(\rank(A) = \rank\bigl(\begin{bmatrix} A & C \end{bmatrix}\bigr)\), \(\rank(B) = \rank\begin{pmatrix} B \\ C \end{pmatrix}\).
    \item Show that there exists a unique minimal point \(X_{\star}\), if and only if \(\rank(A) = m\), \(\rank(B) = n\).
    \item Propose a numerical method to compute the unique \(X_{\star}\) in (b).
    \item Does \(X_{\star}\) also minimize \(\|AXB - C\|_{2}\)? Why? Here \(\|\cdot\|_{2}\) is the spectral norm of a matrix.
\end{enumerate}

\setcounter{section}{3}

\subsubsection{Definitions used in Problem 3}

The unknown is \(X\in\mathbb{R}^{m\times n}\), so that \(AXB\) has the size of \(C\). Write \(\Phi(X)\coloneqq AXB\). Then \(\Phi:\mathbb{R}^{m\times n}\to\mathbb{R}^{l\times p}\) is linear, and the exam is ordinary least squares for \(\Phi\) once \(\mathbb{R}^{l\times p}\) carries the Frobenius inner product \(\langle M,N\rangle_{F}=\operatorname{tr}(M^{T}N)\). The gradient in (b) is the directional derivative of \(\tfrac12\operatorname{tr}(R^{T}R)\) in the matrix variable \(X\); one never has to pass through \(\operatorname{vec}\) or convert \(\|\cdot\|_{F}\) into an Euclidean \(\ell^{2}\) vector. The Kronecker form is an isometric coordinate chart on the same inner-product space (Ding's OLS package after stacking \cite{ding24}), useful for (c) as a warning not to assemble \(B^{T}\otimes A\). Parts (a)--(c) are the range, the kernel, and a constructive inverse of \(\Phi\). Part (d) asks whether the same \(X_{\star}\) is also a minimizer for the spectral norm; that is a different projection, because \(\|\cdot\|_{2}\) does not come from an inner product.

\begin{definition}[Right multiplication as column action, left as row action]
\label{def:left-right-mult}
Let \(M\) have columns \(m_{1},\ldots,m_{k}\). If \(B\) has as many rows as \(M\) has columns, the \(j\)-th column of the product is
\[
(MB)_{\cdot j}
=
M b_{\cdot j}
=
\sum_{i} b_{ij}\, m_{i}.
\]
Right multiplication therefore replaces the columns of \(M\) by linear combinations of those columns. The three elementary column operations (swap two columns, scale a column by a nonzero scalar, add a multiple of one column to another) are the special case in which \(B\) is an invertible elementary matrix. In this problem the right-hand factor is in general rectangular, so the action is an arbitrary column combination, not an elementary transformation.

Dually, left multiplication \(AM\) replaces the rows of \(M\) by linear combinations of those rows. Elementary row operations are left multiplication by invertible elementary matrices.

In particular, for the data \(A,B,C\) of the problem:
\begin{enumerate}[label=\textup{(\roman*)}]
    \item
    There exists \(Y\in\mathbb{R}^{m\times p}\) with \(C=AY\) if and only if every column of \(C\) is a linear combination of columns of \(A\), if and only if \(\operatorname{range}(C)\subseteq\operatorname{range}(A)\), if and only if \(\rank\begin{bmatrix} A & C \end{bmatrix}=\rank(A)\).

    \item
    There exists \(Z\in\mathbb{R}^{l\times n}\) with \(C=ZB\) if and only if every row of \(C\) is a linear combination of rows of \(B\), if and only if \(\operatorname{range}(C^{T})\subseteq\operatorname{range}(B^{T})\), if and only if \(\rank\begin{bmatrix} B \\ C \end{bmatrix}=\rank(B)\). Equivalently, \(\ker(B)\subseteq\ker(C)\).

    \item
    There exists \(X\in\mathbb{R}^{m\times n}\) with \(Y=XB\) if and only if every row of \(Y\) is a linear combination of rows of \(B\), if and only if \(\ker(B)\subseteq\ker(Y)\).
\end{enumerate}
\end{definition}

\begin{lemma}[Full column rank \(\Leftrightarrow\) left inverse; full row rank \(\Leftrightarrow\) right inverse]
\label{lem:one-sided-inverse}
Let \(M\in\mathbb{R}^{q\times r}\). The following are equivalent.
\begin{enumerate}[label=\textup{(\roman*)}]
    \item
    \(\rank(M)=r\) (independent columns; \(M\) is injective).
    \item
    \(M^{T}M\) is SPD, hence invertible.
    \item
    There exists a left inverse \(Q\in\mathbb{R}^{r\times q}\) with \(QM=I_{r}\). One such matrix is \(Q=(M^{T}M)^{-1}M^{T}\).
\end{enumerate}
Dually, \(\rank(M)=q\) (independent rows; \(M\) is surjective) if and only if \(MM^{T}\) is SPD, if and only if there exists a right inverse \(R\in\mathbb{R}^{r\times q}\) with \(MR=I_{q}\). One such matrix is \(R=M^{T}(MM^{T})^{-1}\).
\end{lemma}

\begin{proof}
A left inverse forces injectivity: \(Mv=0\) implies \(v=QMv=0\), so the columns are independent and \(\rank(M)=r\). Conversely, independent columns mean \(\ker(M)=\{0\}\), hence
\[
v^{T}(M^{T}M)v
=
\|Mv\|_{2}^{2}
=
0
\quad\Longrightarrow\quad
Mv=0
\quad\Longrightarrow\quad
v=0,
\]
so \(M^{T}M\) is SPD. Then \(Q=(M^{T}M)^{-1}M^{T}\) satisfies \(QM=I_{r}\).

The same \(Q\) has a column-space reading, which is the construction to keep in hand. The columns \(m_{1},\ldots,m_{r}\) of \(M\) are a basis of \(\operatorname{range}(M)\). A left inverse is any linear map \(Q:\mathbb{R}^{q}\to\mathbb{R}^{r}\) with \(Q m_{j}=e_{j}\). Extend those columns to a basis of \(\mathbb{R}^{q}\) and send the extra vectors to \(0\); the Gram formula above is the special choice that vanishes on \(\operatorname{range}(M)^{\perp}\) and inverts \(M\) on its column space. In the language of \cref{def:left-right-mult}, left multiplication by \(Q\) undoes the column family of \(M\).

The row statements are the column statements applied to \(M^{T}\): a right inverse of \(M\) is a left inverse of \(M^{T}\), and \(MR=I_{q}\) means the rows of \(M\) are sent back to the standard basis of \(\mathbb{R}^{q}\).
\end{proof}

\begin{definition}[Frobenius inner product and Frobenius norm]
\label{def:frobenius}
Let \(M=(m_{ij})\) and \(N=(n_{ij})\) be real matrices of the same size. The \emph{Frobenius inner product} is
\[
\langle M,N\rangle_{F}
\coloneqq
\tr(M^{T}N)
=
\sum_{i,j} m_{ij}n_{ij},
\]
and the \emph{Frobenius norm} is
\[
\|M\|_{F}
\coloneqq
\sqrt{\langle M,M\rangle_{F}}
=
\Bigl(\sum_{i,j} m_{ij}^{2}\Bigr)^{1/2}.
\]
If \(\sigma_{1}(M)\ge\cdots\ge\sigma_{r}(M)>0\) are the positive singular values of \(M\), then
\[
\|M\|_{F}
=
\Bigl(\sum_{k=1}^{r}\sigma_{k}(M)^{2}\Bigr)^{1/2}.
\]
Equivalently, \(\|M\|_{F}=\|\operatorname{vec}(M)\|_{2}\), where \(\operatorname{vec}\) stacks columns (\cref{def:vec-kronecker}). This inner product makes every \(\mathbb{R}^{q\times r}\) a Euclidean space, so \(\min_{X}\|AXB-C\|_{F}\) is the orthogonal projection of \(C\) onto \(\operatorname{range}(\Phi)\).
\end{definition}

\begin{definition}[Spectral / operator \(2\)-norm]
\label{def:spectral}
For \(M\in\mathbb{R}^{q\times r}\), the \emph{spectral norm} (operator norm induced by the vector Euclidean norm) is
\[
\|M\|_{2}
\coloneqq
\sup_{\|x\|_{2}=1}\|Mx\|_{2}
=
\sigma_{\max}(M)
=
\sqrt{\lambda_{\max}(M^{T}M)}.
\]
This is \emph{not} the entrywise \(\ell^{2}\) norm: the latter is \(\|\cdot\|_{F}\). Orthogonal invariance: if \(Q\) and \(W\) are orthogonal of the matching sizes, then \(\|QMW\|_{2}=\|M\|_{2}\) and \(\|QMW\|_{F}=\|M\|_{F}\).
\end{definition}

\begin{lemma}[Comparison of \(\|\cdot\|_{2}\) and \(\|\cdot\|_{F}\)]
\label{lem:norm-compare}
For every \(M\),
\[
\|M\|_{2}
\le
\|M\|_{F}
\le
\sqrt{\rank(M)}\,\|M\|_{2}.
\]
In particular the two norms are equivalent on each fixed matrix space, but they are not scalar multiples, and they induce different best-approximation problems onto a proper subspace.
\end{lemma}

\begin{proof}
Both norms are orthogonally invariant, so it is enough to take \(M=\operatorname{diag}(\sigma_{1},\ldots,\sigma_{r})\) in compact SVD coordinates, with \(\sigma_{1}\ge\cdots\ge\sigma_{r}>0\). Then \(\|M\|_{2}=\sigma_{1}\) and \(\|M\|_{F}^{2}=\sum_{k}\sigma_{k}^{2}\), and \(\sigma_{1}^{2}\le\sum_{k}\sigma_{k}^{2}\le r\sigma_{1}^{2}\).
\end{proof}

\begin{remark}[Inner-product norms versus operator norms]
\label{rem:inner-product-norms}
A norm on a real vector space \(V\) is an \emph{inner-product norm} (Hilbert / Euclidean) when \(\|x\|=\sqrt{\langle x,x\rangle}\) for some inner product on \(V\). Jordan--von Neumann: this happens if and only if the parallelogram law
\[
\|x+y\|^{2}+\|x-y\|^{2}
=
2\bigl(\|x\|^{2}+\|y\|^{2}\bigr)
\]
holds for all \(x,y\in V\). Three labels that collide:
\begin{itemize}[leftmargin=1.4em]
    \item
    \emph{Inner-product / Hilbert / Euclidean}: the norm of the space in which one differentiates or projects (here \(\mathbb{R}^{m\times n}\) or \(\mathbb{R}^{l\times p}\)).
    \item
    \emph{Induced / operator / subordinate} matrix norm: \(\|M\|=\sup_{\|x\|=1}\|Mx\|\), made from a \emph{vector} norm. Being induced by the vector inner product \(\ell^{2}\) does \emph{not} make \(\|\cdot\|_{2}\) an inner-product norm on matrices.
    \item
    \emph{Schatten-\(p\)}: \(\bigl(\sum_{k}\sigma_{k}(M)^{p}\bigr)^{1/p}\). Only \(p=2\) (Frobenius) is an inner-product norm; \(p=1\) is nuclear, \(p=\infty\) is spectral.
\end{itemize}

{\footnotesize
\begin{center}
\begin{tabular}{@{}>{\raggedright\arraybackslash}p{2.55cm}
                   >{\raggedright\arraybackslash}p{3.55cm}
                   >{\centering\arraybackslash}p{1.85cm}
                   >{\raggedright\arraybackslash}p{5.35cm}@{}}
\toprule
Name & Definition & Inner-product norm? & What it is / is not \\
\midrule
vector \(\|\cdot\|_{2}\)
&
\(\sqrt{x^{T}x}\)
&
yes, on \(\mathbb{R}^{d}\)
&
The Euclidean structure that induces the \emph{operator} \(\|\cdot\|_{2}\) below; not a matrix norm.
\\[0.35em]
vector \(\|\cdot\|_{p}\), \(p\neq 2\)
&
\(\bigl(\sum_{i}\lvert x_{i}\rvert^{p}\bigr)^{1/p}\)
&
no
&
Fails parallelogram. Induces the operator \(p\)-norms in the next two rows.
\\[0.35em]
Frobenius \(\|\cdot\|_{F}\)
&
\(\sqrt{\operatorname{tr}(M^{T}N)}\) at \(N=M\); equivalently \(\bigl(\sum_{k}\sigma_{k}^{2}\bigr)^{1/2}\)
&
yes, on matrices
&
Schatten-\(2\), Hilbert--Schmidt, entrywise \(\ell^{2}\). Problem~3(a)--(c) lives here: \(\nabla\) and orthogonal projection exist. Not an operator norm except on \(1\times 1\).
\\[0.35em]
spectral \(\|\cdot\|_{2}\)
&
\(\sup_{\|x\|_{2}=1}\|Mx\|_{2}=\sigma_{\max}(M)\)
&
no, on matrices
&
Operator / induced / Schatten-\(\infty\). Induced by vector \(\ell^{2}\), but parallelogram fails on \(\mathbb{R}^{q\times r}\). Problem~3(d).
\\[0.35em]
operator \(\|\cdot\|_{1}\)
&
\(\sup_{\|x\|_{1}=1}\|Mx\|_{1}=\max_{j}\sum_{i}\lvert m_{ij}\rvert\)
&
no
&
Max column sum. Induced by vector \(\ell^{1}\), not an inner product.
\\[0.35em]
operator \(\|\cdot\|_{\infty}\)
&
\(\sup_{\|x\|_{\infty}=1}\|Mx\|_{\infty}=\max_{i}\sum_{j}\lvert m_{ij}\rvert\)
&
no
&
Max row sum. Induced by vector \(\ell^{\infty}\).
\\[0.35em]
nuclear \(\|\cdot\|_{*}\)
&
\(\sum_{k}\sigma_{k}(M)\)
&
no
&
Schatten-\(1\), trace class. Dual to spectral; convex relaxation of rank. Not used in this exam.
\\[0.35em]
entrywise \(\ell^{p}\), \(p\neq 2\)
&
\(\bigl(\sum_{ij}\lvert m_{ij}\rvert^{p}\bigr)^{1/p}\)
&
no
&
Not unitarily invariant (except \(p=2\)). The case \(p=\infty\) is \(\max_{ij}\lvert m_{ij}\rvert\), distinct from operator \(\infty\).
\\
\bottomrule
\end{tabular}
\end{center}
}

Parallelogram for spectral, at \(A=e_{1}e_{1}^{T}\) and \(B=e_{2}e_{2}^{T}\) in \(\mathbb{R}^{2\times 2}\): \(\|A\|_{2}=\|B\|_{2}=1\), while \(A+B=I_{2}\) and \(A-B=\operatorname{diag}(1,-1)\) both have spectral norm \(1\), so the left-hand side is \(2\) and the right-hand side is \(4\). The same pair is orthogonal of unit Frobenius length, and \(\|\cdot\|_{F}\) does satisfy the law. Consequently \(\min\|AXB-C\|_{F}\) is an orthogonal projection and \(\min\|AXB-C\|_{2}\) is not.
\end{remark}

\begin{definition}[Vectorization and Kronecker product]
\label{def:vec-kronecker}
For \(M\in\mathbb{R}^{q\times r}\) with columns \(m_{1},\ldots,m_{r}\), set \(\operatorname{vec}(M)\coloneqq(m_{1}^{T},\ldots,m_{r}^{T})^{T}\in\mathbb{R}^{qr}\) (column-major stacking). For \(P=(p_{ij})\in\mathbb{R}^{a\times b}\) and \(Q\in\mathbb{R}^{c\times d}\), the \emph{Kronecker product} \(P\otimes Q\) is the \(ac\times bd\) block matrix whose \((i,j)\)-block is the scalar multiple \(p_{ij}Q\). The left factor \(P\) therefore determines the coarse block grid; the right factor \(Q\) is copied, scaled, into every block. It is associative, and \((P\otimes Q)^{T}=P^{T}\otimes Q^{T}\). Mixed-product rule: \((P\otimes Q)(R\otimes S)=PR\otimes QS\) whenever the ordinary products exist. This is the MATLAB / standard convention that makes \cref{lem:vec-axb} hold; the opposite block convention would replace \(B^{T}\otimes A\) by \(A\otimes B^{T}\).
\end{definition}

\begin{example}[Two \(2\times 2\) Kronecker products]
\label{ex:kronecker}
Write
\[
P
=
\begin{pmatrix}
1 & 2 \\
3 & 4
\end{pmatrix},
\qquad
Q
=
\begin{pmatrix}
5 & 6 \\
7 & 8
\end{pmatrix}.
\]
The block recipe \(P\otimes Q\) is the \(4\times 4\) matrix
\[
P\otimes Q
=
\begin{pmatrix}
1\,Q & 2\,Q \\
3\,Q & 4\,Q
\end{pmatrix}
=
\begin{pmatrix}
5 & 6 & 10 & 12 \\
7 & 8 & 14 & 16 \\
15 & 18 & 20 & 24 \\
21 & 24 & 28 & 32
\end{pmatrix}.
\]
The opposite order \(Q\otimes P\) is a different \(4\times 4\) matrix (its top-left block is \(5P\), not \(1\,Q\)). The mixed-product rule is already visible at this size: \((P\otimes I_{2})(I_{2}\otimes Q)=P\otimes Q\).

The same recipe is what linearises left multiplication after \(\operatorname{vec}\). Take
\[
A
=
\begin{pmatrix}
1 & 2 \\
0 & 1
\end{pmatrix},
\qquad
X
=
\begin{pmatrix}
3 & 4 \\
5 & 6
\end{pmatrix}.
\]
Then
\[
AX
=
\begin{pmatrix}
13 & 16 \\
5 & 6
\end{pmatrix},
\qquad
\operatorname{vec}(X)
=
\begin{pmatrix}
3 \\ 5 \\ 4 \\ 6
\end{pmatrix},
\qquad
\operatorname{vec}(AX)
=
\begin{pmatrix}
13 \\ 5 \\ 16 \\ 6
\end{pmatrix},
\]
and
\[
I_{2}\otimes A
=
\begin{pmatrix}
A & 0 \\
0 & A
\end{pmatrix}
=
\begin{pmatrix}
1 & 2 & 0 & 0 \\
0 & 1 & 0 & 0 \\
0 & 0 & 1 & 2 \\
0 & 0 & 0 & 1
\end{pmatrix}.
\]
Matrix--vector multiplication recovers the stacked product: \((I_{2}\otimes A)\operatorname{vec}(X)=\operatorname{vec}(AX)\). Right multiplication by a \(2\times 2\) matrix \(B\) is the other factor \(B^{T}\otimes I_{2}\); composing the two is \cref{lem:vec-axb}.
\end{example}

\begin{lemma}[\(\operatorname{vec}(AXB)=(B^{T}\otimes A)\operatorname{vec}(X)\)]
\label{lem:vec-axb}
For \(A\in\mathbb{R}^{l\times m}\), \(X\in\mathbb{R}^{m\times n}\), \(B\in\mathbb{R}^{n\times p}\),
\[
\operatorname{vec}(AXB)
=
(B^{T}\otimes A)\,\operatorname{vec}(X).
\]
In particular \(\operatorname{vec}(A^{T}MB^{T})=(B\otimes A^{T})\,\operatorname{vec}(M)\). Consequently
\[
\|AXB-C\|_{F}
=
\bigl\|(B^{T}\otimes A)\operatorname{vec}(X)-\operatorname{vec}(C)\bigr\|_{2},
\]
so Problem~3 is ordinary least squares for the design matrix \(B^{T}\otimes A\in\mathbb{R}^{lp\times mn}\), which is the setting of Ding's Chapter~3 \cite{ding24} after stacking.
\end{lemma}

\begin{proof}
If \(B\) has columns \(b_{1},\ldots,b_{p}\), then \(XB\) has columns \(Xb_{j}\) and \(AXB\) has columns \(AXb_{j}\). Stacking gives
\[
\operatorname{vec}(AXB)
=
(I_{p}\otimes A)\,\operatorname{vec}(XB).
\]
The same identity with \(A=I_{m}\) yields \(\operatorname{vec}(XB)=(B^{T}\otimes I_{m})\operatorname{vec}(X)\). Compose and apply the mixed-product rule:
\[
(I_{p}\otimes A)(B^{T}\otimes I_{m})
=
B^{T}\otimes A.
\]
The identity \(\operatorname{vec}(PMQ)=(Q^{T}\otimes P)\operatorname{vec}(M)\) is the same statement, with the triple \((A,X,B)\) renamed \((P,M,Q)\). The adjoint form follows by taking \(P=A^{T}\) and \(Q=B^{T}\). The norm identity is \(\|M\|_{F}=\|\operatorname{vec}(M)\|_{2}\) from \cref{def:frobenius}.
\end{proof}

\begin{lemma}[Vector calculus: linear and quadratic forms]
\label{lem:ding-vec-calc}
Let \(x\in\mathbb{R}^{p}\). For \(a\in\mathbb{R}^{p}\) and a symmetric \(G\in\mathbb{R}^{p\times p}\),
\[
\frac{\partial}{\partial x}(a^{T}x)
=
a,
\qquad
\frac{\partial}{\partial x}(x^{T}Gx)
=
2Gx.
\]
Consequently, for \(K\in\mathbb{R}^{q\times p}\) and \(y\in\mathbb{R}^{q}\),
\[
\frac{\partial}{\partial x}\Bigl(\tfrac12\|Kx-y\|_{2}^{2}\Bigr)
=
K^{T}(Kx-y).
\]
The Hessian of the right-hand side (as a function of \(x\)) is the Gram matrix \(K^{T}K\).
\end{lemma}

\begin{proof}
The two displays are the vector-calculus identities in Appendix~A of Ding \cite{ding24} (linear form; quadratic form with symmetric Gram matrix). For the least-squares function expand
\[
\tfrac12\|Kx-y\|_{2}^{2}
=
\tfrac12 x^{T}(K^{T}K)x-(K^{T}y)^{T}x+\tfrac12\|y\|_{2}^{2}.
\]
The matrix \(K^{T}K\) is symmetric, so the two identities give \(\nabla_{x}=K^{T}Kx-K^{T}y\). This is the same first-order condition as the OLS normal equation \(K^{T}(y-Kx)=0\) in Ding, Chapter~3.
\end{proof}

\begin{lemma}[Gradient of the squared Frobenius residual]
\label{lem:frob-grad}
Let \(f(X)=\tfrac12\|AXB-C\|_{F}^{2}\). Then, with respect to \(\langle\cdot,\cdot\rangle_{F}\),
\[
\nabla f(X)
=
A^{T}(AXB-C)B^{T}.
\]
The Hessian quadratic form is \(\mathrm{d}^{2}f(H,H)=\|AHB\|_{F}^{2}\ge 0\), so \(f\) is convex, and strictly convex if and only if \(AHB=0\) forces \(H=0\).
\end{lemma}

\begin{proof}
Write \(R\coloneqq AXB-C\), so that
\[
f(X)
=
\tfrac12\|R\|_{F}^{2}
=
\tfrac12\operatorname{tr}(R^{T}R).
\]
The map \(X\mapsto R\) is affine. The directional derivative of \(f\) at \(X\) in the direction \(H\in\mathbb{R}^{m\times n}\) is therefore the first variation of \(\tfrac12\|R\|_{F}^{2}\):
\[
\mathrm{d}f(X)[H]
=
\langle R,AHB\rangle_{F}
=
\operatorname{tr}\bigl(R^{T}AHB\bigr).
\]
Cycle the trace (\(\operatorname{tr}(UVW)=\operatorname{tr}(VWU)\), here sending \(B\) to the front) to obtain
\[
\operatorname{tr}\bigl(R^{T}AHB\bigr)
=
\operatorname{tr}\bigl(BR^{T}AH\bigr)
=
\langle A^{T}RB^{T},H\rangle_{F}.
\]
The Riesz representative of \(\mathrm{d}f(X)\) in the Frobenius inner product is therefore \(\nabla f(X)=A^{T}RB^{T}=A^{T}(AXB-C)B^{T}\). Differentiating once more, \(\mathrm{d}^{2}f(H,H)=\langle AHB,AHB\rangle_{F}=\|AHB\|_{F}^{2}\).

The \(\operatorname{vec}\) route is the same computation in coordinates, not a different derivative. By \cref{lem:vec-axb}, \(f(X)=\tfrac12\|Kx-c\|_{2}^{2}\) with \(x=\operatorname{vec}(X)\), \(c=\operatorname{vec}(C)\), and \(K=B^{T}\otimes A\). \Cref{lem:ding-vec-calc} then gives \(\nabla_{x}f=K^{T}(Kx-c)=(B\otimes A^{T})\operatorname{vec}(R)\), which reshapes to the same matrix \(A^{T}RB^{T}\).
\end{proof}

\begin{remark}[Trace calculus versus \(\operatorname{vec}\)]
\label{rem:p3-trace}
The identity \(\|M\|_{F}=\|\operatorname{vec}(M)\|_{2}\) is an isometry of Euclidean spaces; it does not change the critical-point equation. For the exam it is shorter to stay on matrices: differentiate \(\tfrac12\operatorname{tr}(R^{T}R)\) as above, read off \(\nabla f=A^{T}RB^{T}\), and write the normal equation \(A^{T}AXBB^{T}=A^{T}CB^{T}\). Kronecker products are needed only if one insists on quoting Ding's vector OLS package, or as a size warning in (c).
\end{remark}

\begin{definition}[Hadamard / entrywise product]
\label{def:hadamard}
For matrices \(M,N\) of the same size, the \emph{Hadamard product} \(M\circ N\) is the matrix with entries \((M\circ N)_{ij}=M_{ij}N_{ij}\). It does not enlarge the array (unlike Kronecker). If \(D=\operatorname{diag}(\sigma_{1},\ldots,\sigma_{m})\) and \(E=\operatorname{diag}(\tau_{1},\ldots,\tau_{n})\), then
\[
(D\tilde{X}E)_{ij}
=
\sigma_{i}\tau_{j}\,\tilde{X}_{ij}
=
(\sigma\tau^{T}\circ\tilde{X})_{ij}.
\]
Left-and-right multiplication by diagonals is therefore Hadamard multiplication against the rank-one matrix of products \(\sigma_{i}\tau_{j}\). Inverting that action, when every product is nonzero, is entrywise division: \(\tilde{X}_{ij}=\tilde{C}_{ij}/(\sigma_{i}\tau_{j})\).
\end{definition}

\begin{lemma}[The Gram matrix squares the condition number]
\label{lem:gram-cond}
Let \(M\) have full column rank, with \(2\)-norm condition number \(\kappa_{2}(M)=\sigma_{\max}(M)/\sigma_{\min}(M)\). The singular values of \(M^{T}M\) are the squares of those of \(M\), so
\[
\kappa_{2}(M^{T}M)
=
\kappa_{2}(M)^{2}.
\]
In floating-point arithmetic, forming \(M^{T}M\) (or inverting it) therefore squares the sensitivity of \(M\). Householder QR or SVD applied to \(M\) itself does not \cite{trefethen97,golub96}.
\end{lemma}

\begin{proof}
If \(M=U\Sigma V^{T}\) is a thin SVD, then \(M^{T}M=V\Sigma^{2}V^{T}\). The ratio of extreme eigenvalues of \(M^{T}M\) is \((\sigma_{\max}/\sigma_{\min})^{2}\).
\end{proof}

\begin{definition}[Moore--Penrose inverse, via SVD]
\label{def:pinv}
Let \(M\in\mathbb{R}^{q\times r}\) have compact SVD \(M=U\Sigma V^{T}\) with \(\Sigma=\operatorname{diag}(\sigma_{1},\ldots,\sigma_{\rho})\), \(\sigma_{k}>0\), and \(\rho=\rank(M)\). The \emph{Moore--Penrose inverse} is \(M^{+}\coloneqq V\Sigma^{-1}U^{T}\). Then \(MM^{+}\) is the orthogonal projector onto \(\operatorname{range}(M)\), and \(M^{+}M\) is the orthogonal projector onto \(\operatorname{range}(M^{T})\). If \(M\) has full column rank, \(M^{+}=(M^{T}M)^{-1}M^{T}\); if \(M\) has full row rank, \(M^{+}=M^{T}(MM^{T})^{-1}\).
\end{definition}

\begin{lemma}[Range of \(X\mapsto AXB\)]
\label{lem:phi-range}
Let \(\Phi(X)=AXB\). For \(C\in\mathbb{R}^{l\times p}\), the following are equivalent.
\begin{enumerate}[label=\textup{(\roman*)}]
    \item \(C\in\operatorname{range}(\Phi)\), i.e.\ \(AXB=C\) for some \(X\).
    \item \(AA^{+}CB^{+}B=C\).
    \item \(\operatorname{range}(C)\subseteq\operatorname{range}(A)\) and \(\operatorname{range}(C^{T})\subseteq\operatorname{range}(B^{T})\).
    \item \(\rank(A)=\rank\begin{bmatrix} A & C \end{bmatrix}\) and \(\rank(B)=\rank\begin{bmatrix} B \\ C \end{bmatrix}\).
\end{enumerate}
In that case one solution is \(X=A^{+}CB^{+}\). This is part~(a): the minimal value of \(\|AXB-C\|_{F}\) is \(0\) if and only if these rank identities hold.
\end{lemma}

\begin{proof}
Always \(AA^{+}CB^{+}B=\Phi(A^{+}CB^{+})\), so (ii) implies (i), and the displayed \(X\) works. Conversely, if \(C=AXB\), then \(\operatorname{range}(C)\subseteq\operatorname{range}(A)\) and \(\operatorname{range}(C^{T})\subseteq\operatorname{range}(B^{T}X^{T})\subseteq\operatorname{range}(B^{T})\), which is (iii), and then \(AA^{+}C=C\) and \(CB^{+}B=C\), hence (ii).

The identity \(AA^{+}C=C\) says that the orthogonal projector onto \(\operatorname{range}(A)\) fixes every column of \(C\), which is \(\operatorname{range}(C)\subseteq\operatorname{range}(A)\), equivalently \(\rank\begin{bmatrix} A & C \end{bmatrix}=\rank(A)\). Likewise \(CB^{+}B=C\) is \(\operatorname{range}(C^{T})\subseteq\operatorname{range}(B^{T})\), equivalently \(\rank\begin{bmatrix} B \\ C \end{bmatrix}=\rank(B)\). Thus (ii), (iii), and (iv) coincide.

Finally, \(\min_{X}\|AXB-C\|_{F}=0\) if and only if \(C\) lies in the (closed) subspace \(\operatorname{range}(\Phi)\).
\end{proof}

\begin{lemma}[Injectivity of \(X\mapsto AXB\)]
\label{lem:phi-injective}
The map \(\Phi\) is injective if and only if \(\rank(A)=m\) and \(\rank(B)=n\). In that case, for every \(C\) there is a unique Frobenius minimizer
\[
X_{\star}
=
A^{+}CB^{+}
=
(A^{T}A)^{-1}A^{T}\,C\,B^{T}(BB^{T})^{-1}.
\]
This is part~(b). Existence of a minimizer is free (finite-dimensional least squares); uniqueness is injectivity of \(\Phi\).
\end{lemma}

\begin{proof}
If \(\rank(A)<m\), pick \(v\neq 0\) in \(\operatorname{ker}(A)\) and \(w\neq 0\) in \(\mathbb{R}^{n}\). Then \(X=vw^{T}\neq 0\) and \(\Phi(X)=(Av)(w^{T}B)=0\). If \(\rank(B)<n\), pick \(w\neq 0\) in \(\operatorname{ker}(B^{T})\) and \(v\neq 0\) in \(\mathbb{R}^{m}\). Then \(X=vw^{T}\neq 0\) and \(\Phi(X)=(Av)(w^{T}B)=0\).

Conversely, if \(\rank(A)=m\) and \(\rank(B)=n\), then \(\operatorname{ker}(A)=\{0\}\) and \(B\) is surjective. From \(AXB=0\) we get \(XB=0\), hence \(X=XBB^{+}=0\). So \(\Phi\) is injective, the normal equation \((B^{T}\otimes A)^{T}(B^{T}\otimes A)\operatorname{vec}(X)=(B^{T}\otimes A)^{T}\operatorname{vec}(C)\) has a unique solution, and substituting the full-rank formulae for \(A^{+}\) and \(B^{+}\) from \cref{def:pinv} gives the displayed \(X_{\star}\). Directly: \(\Phi(A^{+}CB^{+})=AA^{+}CB^{+}B=C\) only when \(C\) is in the range; in general \(AA^{+}CB^{+}B\) is the orthogonal projection of \(C\) onto \(\operatorname{range}(\Phi)\), because \(AA^{+}\) and \(B^{+}B\) are the orthogonal projectors onto \(\operatorname{range}(A)\) and \(\operatorname{range}(B^{T})\).
\end{proof}

\begin{remark}[Parts (c)--(d)]
\label{rem:p3-cd}
\emph{Numerical method for the unique \(X_{\star}\).}
Part~(c) is not asking for another closed form. The formula of~(b) already computes \(X_{\star}\) in exact arithmetic; a \emph{numerical method} is an algorithm that evaluates that map in floating point, with a cost and a backward-error story. Forming \(A^{T}A\) and \(BB^{T}\) squares the condition numbers (\cref{lem:gram-cond}). Forming \(B^{T}\otimes A\) is algebraically the right least-squares matrix and computationally the wrong one: it is \(lp\times mn\). The usable algorithms factor \(A\) and \(B\) separately (thin Householder QR, or thin SVD). After SVD the remaining solve is Hadamard division by \(\sigma_{i}(A)\sigma_{j}(B)\) (\cref{def:hadamard}). After QR it is two triangular substitutions. Both run under the rank hypotheses of~(b).

\emph{Spectral versus Frobenius.}
The unique \(X_{\star}\) of (b) is the orthogonal projection onto \(\operatorname{range}(\Phi)\) in \(\langle\cdot,\cdot\rangle_{F}\). Minimizing \(\|AXB-C\|_{2}\) is Chebyshev approximation in the operator norm, whose metric projection onto a proper subspace is typically non-unique and typically a different point. \Cref{lem:norm-compare} only compares the two numbers \(\|M\|_{2}\) and \(\|M\|_{F}\) at a fixed \(M\); it does not identify argminima. \Cref{rem:inner-product-norms} records the terminology: \(\|\cdot\|_{F}\) is an inner-product norm on matrices, while the spectral \(\|\cdot\|_{2}\) is an induced operator norm and is not. The next example is a case of (b) in which \(X_{\star}\) does \emph{not} minimize the spectral residual.
\end{remark}

\begin{example}[Unique Frobenius minimizer that is not spectral]
\label{ex:p3-spectral}
Take \(l=p=2\), \(m=n=1\), \(A=\begin{pmatrix}1\\0\end{pmatrix}\), \(B=\begin{pmatrix}1&1\end{pmatrix}\), \(C=I_{2}\). Then \(\rank(A)=1=m\), \(\rank(B)=1=n\), and \(X=x\in\mathbb{R}\) is a scalar, with
\[
AXB
=
x
\begin{pmatrix}1&1\\0&0\end{pmatrix}.
\]
The Frobenius residual is
\[
\|AXB-I_{2}\|_{F}^{2}
=
(x-1)^{2}+x^{2}+1,
\]
uniquely minimized at \(x_{\star}=\tfrac12\). The spectral residual at this point is the matrix \(\begin{pmatrix}-\tfrac12&\tfrac12\\0&-1\end{pmatrix}\), whose squared singular values are the eigenvalues of
\[
\begin{pmatrix}1/4&-1/4\\-1/4&5/4\end{pmatrix},
\]
namely \(\frac{3\pm\sqrt{5}}{4}\). Hence
\[
\|AX_{\star}B-C\|_{2}
=
\sqrt{\frac{3+\sqrt{5}}{4}}
>
1.
\]
At \(x=0\), one has \(AXB-C=-I_{2}\) and \(\|-I_{2}\|_{2}=1\), which is strictly smaller. Thus the unique Frobenius minimizer of (b) need not minimize \(\|AXB-C\|_{2}\).
\end{example}

\subsubsection{Solution of Problem 3(a)}

\paragraph{Answer.}
The minimal value is \(0\) if and only if both rank identities hold.

\begin{proof}
The objective \(\|AXB-C\|_{F}\) is nonnegative, and vanishes if and only if \(AXB=C\). Write \(\Phi(X)=AXB\).

\emph{Necessity.}
Suppose \(AXB=C\). By \cref{def:left-right-mult}, this is a right multiplication of \(A\) by \(Y\coloneqq XB\): the columns of \(C\) are linear combinations of the columns of \(A\), hence
\[
\rank\begin{bmatrix} A & C \end{bmatrix}
=
\rank(A).
\]
Dually, it is a left multiplication of \(B\) by \(Z\coloneqq AX\): the rows of \(C\) are linear combinations of the rows of \(B\), hence
\[
\rank\begin{bmatrix} B \\ C \end{bmatrix}
=
\rank(B).
\]

\emph{Sufficiency.}
The identity \(\rank\begin{bmatrix} B \\ C \end{bmatrix}=\rank(B)\) is a statement about columns of a stacked matrix. Make those columns visible by elementary column operations on \(B\), and perform the same operations on \(C\).

Concretely: there is an invertible \(P\in\mathbb{R}^{p\times p}\) (a product of elementary matrices) such that
\[
BP
=
\begin{bmatrix} B_{1} & 0 \end{bmatrix},
\]
where \(B_{1}\in\mathbb{R}^{n\times r}\) has full column rank \(r=\rank(B)\). The zero block is \(\ker(B)\) written in the new coordinates. Apply the same right multiplication to \(C\),
\[
CP
=
\begin{bmatrix} C_{1} & C_{2} \end{bmatrix},
\]
and the stacked matrix becomes
\[
\begin{bmatrix} B \\ C \end{bmatrix}P
=
\begin{bmatrix} B_{1} & 0 \\ C_{1} & C_{2} \end{bmatrix}.
\]
Column operations do not change rank, so the right-hand side still has rank \(r=\rank(B_{1})\). The last block column is \(\bigl(\begin{smallmatrix}0\\ C_{2}\end{smallmatrix}\bigr)\). If some column \(c\) of \(C_{2}\) were nonzero, that stacked column would be nonzero, yet it cannot lie in the span of the first \(r\) stacked columns: a combination with coefficients \(\alpha\in\mathbb{R}^{r}\) produces \(\bigl(\begin{smallmatrix}B_{1}\alpha\\ C_{1}\alpha\end{smallmatrix}\bigr)\), and \(B_{1}\alpha=0\) forces \(\alpha=0\) because \(B_{1}\) has independent columns. Hence \(C_{2}=0\), and
\[
CP
=
\begin{bmatrix} C_{1} & 0 \end{bmatrix}.
\]
(The second rank condition is exactly this: after simplifying \(B\), there are no leftover columns in \(C\).)

The first rank condition is unchanged by the same operations, because \(CP\) and \(C\) have the same column space. Thus \(\operatorname{range}(C_{1})\subseteq\operatorname{range}(C)\subseteq\operatorname{range}(A)\), so \cref{def:left-right-mult}(i) supplies \(Y_{1}\in\mathbb{R}^{m\times r}\) with \(C_{1}=AY_{1}\).

It remains only to solve \(XB_{1}=Y_{1}\). The matrix \(B_{1}\) has independent columns, so \cref{lem:one-sided-inverse} supplies a left inverse \(Q\) with \(QB_{1}=I_{r}\) (for instance \(Q=(B_{1}^{T}B_{1})^{-1}B_{1}^{T}\)). Set \(X=Y_{1}Q\). Then \(XB_{1}=Y_{1}QB_{1}=Y_{1}\), and
\[
AX(BP)
=
AX\begin{bmatrix} B_{1} & 0 \end{bmatrix}
=
\begin{bmatrix} AY_{1} & 0 \end{bmatrix}
=
\begin{bmatrix} C_{1} & 0 \end{bmatrix}
=
CP.
\]
Since \(P\) is invertible, \(AXB=C\). The residual vanishes, and the minimal value is \(0\).
\end{proof}

The argument never invents an auxiliary matrix on \(\ker(B)\). Column operations put \(\ker(B)\) into the explicit zero block of \(BP\); the stacked rank condition then forces the matching block of \(CP\) to vanish as well. What remains is a full-column-rank right factor \(B_{1}\), through which the column coordinates \(Y_{1}\) of \(C_{1}\) factor uniquely on the left.

\subsubsection{Solution of Problem 3(b)}

\paragraph{Answer.}
A unique minimizer exists if and only if \(\rank(A)=m\) and \(\rank(B)=n\). In that case
\[
X_{\star}
=
(A^{T}A)^{-1}A^{T}\,C\,B^{T}(BB^{T})^{-1}.
\]

\begin{proof}
Minimizing \(\|AXB-C\|_{F}\) is the same as minimizing the square. Set
\[
f(X)
\coloneqq
\tfrac12\|AXB-C\|_{F}^{2}
=
\tfrac12\operatorname{tr}\bigl((AXB-C)^{T}(AXB-C)\bigr)
\]
on the space of matrices \(X\in\mathbb{R}^{m\times n}\), with gradient taken in \(\langle\cdot,\cdot\rangle_{F}\). Differentiating the trace as in \cref{lem:frob-grad} (see also \cref{rem:p3-trace}) gives
\[
\nabla f(X)
=
A^{T}(AXB-C)B^{T},
\]
with no intermediate conversion of \(\|\cdot\|_{F}\) into a vector \(\ell^{2}\) norm. (Equivalently, after \(\operatorname{vec}\) this is OLS with design matrix \(B^{T}\otimes A\), \cref{lem:vec-axb}.) Critical points are the solutions of the normal equation
\[
A^{T}AX_{\star}BB^{T}
=
A^{T}CB^{T}.
\]
The second variation \(\mathrm{d}^{2}f(H,H)=\|AHB\|_{F}^{2}\) is independent of \(X\) and nonnegative, so \(f\) is a convex quadratic: every critical point is a global minimizer. The equation \(\nabla f(X)=0\) is \(\Phi^{*}\Phi X=\Phi^{*}C\) with \(\Phi^{*}(Y)=A^{T}YB^{T}\). This linear system is always solvable, because \(\operatorname{range}(\Phi^{*}\Phi)=\operatorname{range}(\Phi^{*})\). The solution is unique if and only if \(\Phi^{*}\Phi\) is invertible, if and only if the Hessian quadratic form of \cref{lem:frob-grad} is positive definite, if and only if \(\|AHB\|_{F}=0\) forces \(H=0\), if and only if \(\Phi\) is injective. If \(\Phi\) is not injective then \(f(X+H)=f(X)\) for every \(H\in\ker\Phi\), so the minimizers form a positive-dimensional affine space.

\emph{Injectivity of \(\Phi\).}
If \(\rank(A)<m\), pick \(v\neq 0\) in \(\ker(A)\) and \(w\neq 0\) in \(\mathbb{R}^{n}\), and set \(H=vw^{T}\neq 0\). Then \(AHB=(Av)(w^{T}B)=0\). If \(\rank(B)<n\), pick \(w\neq 0\) in \(\ker(B^{T})\) and \(v\neq 0\) in \(\mathbb{R}^{m}\), and set \(H=vw^{T}\neq 0\). Then \(w^{T}B=0\), so \(AHB=0\). Conversely, if \(\rank(A)=m\) and \(\rank(B)=n\), left multiplication by \(A\) kills no nonzero column, and \cref{lem:one-sided-inverse} supplies a right inverse \(R\) of \(B\) with \(BR=I_{n}\). From \(AHB=0\) one gets \(HB=0\), then \(H=HBR=0\).

\emph{The unique critical point.}
Under \(\rank(A)=m\) and \(\rank(B)=n\), \cref{lem:one-sided-inverse} says that \(A\) has the left inverse \((A^{T}A)^{-1}A^{T}\) and \(B\) has the right inverse \(B^{T}(BB^{T})^{-1}\). The normal equation therefore solves as the displayed \(X_{\star}\). Equivalently, \(\nabla f(X_{\star})=0\) is \(A^{T}(AX_{\star}B-C)B^{T}=0\), which is the chain-rule factor \(A^{T}(\,\cdot\,)B^{T}\) from \cref{lem:frob-grad}.
\end{proof}

\subsubsection{Solution of Problem 3(c)}

\paragraph{Answer.}
Thin Householder QR of \(A\) and of \(B^{T}\), then two triangular solves for \(R_{A}X_{\star}R_{B}^{T}=Q_{A}^{T}CQ_{B}\). Equivalently: thin SVD of each factor, then Hadamard division \(\tilde{X}_{ij}=\tilde{C}_{ij}/(\sigma_{i}(A)\sigma_{j}(B))\).

\begin{proof}
The formula of~(b),
\[
X_{\star}
=
(A^{T}A)^{-1}A^{T}\,C\,B^{T}(BB^{T})^{-1},
\]
is an exact-arithmetic expression, not a numerical method. ``Numerical'' here is the Trefethen--Bau / Golub--Van Loan distinction \cite{trefethen97,golub96}: an algorithm that realises the same map in floating point, together with its cost and its backward error. Three procedures compute the same \(X_{\star}\) in exact arithmetic under \(\rank(A)=m\) and \(\rank(B)=n\); they are not interchangeable in floating point.

\emph{What fails as a method, though it is algebraically correct.}
Forming the Kronecker matrix \(K=B^{T}\otimes A\in\mathbb{R}^{lp\times mn}\) and solving \(\min_{x}\|Kx-\operatorname{vec}(C)\|_{2}\) is the vectorised problem of \cref{lem:vec-axb}. It works, and it is the wrong algorithm: storage and work scale as a single dense least-squares problem of size \(lp\times mn\), and the Kronecker structure is thrown away. Forming the Gram matrices in the formula of~(b) (or Cholesky on \(A^{T}A\) and \(BB^{T}\)) likewise works in exact arithmetic, but \cref{lem:gram-cond} says \(\kappa_{2}(A^{T}A)=\kappa_{2}(A)^{2}\), so every digit of \(\kappa_{2}(A)\) is doubled before one even touches \(C\). That is the standard reason one does not compute a thin QR of \(A\) by first forming \(A^{T}A\).

\emph{QR, which is the default exam answer.}
Because \(\rank(A)=m\), a thin Householder QR \(A=Q_{A}R_{A}\) exists with \(Q_{A}\in\mathbb{R}^{l\times m}\) orthonormal columns and \(R_{A}\in\mathbb{R}^{m\times m}\) invertible (upper triangular, nonzero diagonal). Because \(\rank(B)=n\), a thin QR \(B^{T}=Q_{B}R_{B}\) exists with \(R_{B}\in\mathbb{R}^{n\times n}\) invertible. Then \(A^{+}=R_{A}^{-1}Q_{A}^{T}\) and \(B^{+}=Q_{B}R_{B}^{-T}\), and the formula of~(b) rearranges as
\[
R_{A}X_{\star}R_{B}^{T}
=
Q_{A}^{T}CQ_{B}.
\]
The right-hand side is two products with orthonormal factors (stable). The left-hand side is a triangular matrix equation: substitute along columns of \(R_{A}\), then along rows of \(R_{B}^{T}\). Classical Gram--Schmidt is \emph{not} a substitute for this QR: it produces the same factors in exact arithmetic and is unstable. The method is Householder (or a rank-revealing QR). It runs precisely because~(b) makes \(R_{A}\) and \(R_{B}\) invertible; if a diagonal entry of \(R_{A}\) vanished, \(\Phi\) would not be injective and \(X_{\star}\) would not be unique.

\emph{SVD, and where the Hadamard product sits.}
The same construction in singular-vector coordinates makes the core solve diagonal, hence entrywise. Let \(A=U_{A}\Sigma_{A}V_{A}^{T}\) and \(B=U_{B}\Sigma_{B}V_{B}^{T}\) be thin SVDs. Under~(b) every diagonal entry of \(\Sigma_{A}=\operatorname{diag}(\sigma_{1},\ldots,\sigma_{m})\) and of \(\Sigma_{B}=\operatorname{diag}(\tau_{1},\ldots,\tau_{n})\) is positive. Orthogonal invariance of \(\|\cdot\|_{F}\) gives
\[
\|AXB-C\|_{F}
=
\bigl\|\Sigma_{A}\tilde{X}\Sigma_{B}-\tilde{C}\bigr\|_{F},
\]
where \(\tilde{X}=V_{A}^{T}XU_{B}\) and \(\tilde{C}=U_{A}^{T}CV_{B}\). By \cref{def:hadamard},
\[
(\Sigma_{A}\tilde{X}\Sigma_{B})_{ij}
=
\sigma_{i}\tau_{j}\,\tilde{X}_{ij}
=
(\sigma\tau^{T}\circ\tilde{X})_{ij}.
\]
The unique minimizer in these coordinates is therefore the Hadamard (entrywise) division
\[
\tilde{X}_{ij}
=
\frac{\tilde{C}_{ij}}{\sigma_{i}\tau_{j}},
\qquad
X_{\star}
=
V_{A}\tilde{X}\,U_{B}^{T}
=
V_{A}\Sigma_{A}^{-1}(U_{A}^{T}CV_{B})\Sigma_{B}^{-1}U_{B}^{T}.
\]
No linear system remains: after two SVDs one rescales the \(m\times n\) array \(\tilde{C}\) entry by entry. This is why an SVD write-up mentions a Hadamard product. It is not a different problem; it is \(\Sigma_{A}(\,\cdot\,)\Sigma_{B}\) written without summation. The algorithm is backward stable and additionally exhibits the small singular values if \(A\) or \(B\) is ill-conditioned but still full rank.

\emph{Cost.}
Both usable methods factor \(A\) (\(l\times m\)) and \(B\) (\(n\times p\)) separately, then do \(O(lpm+mnp)\) work on \(C\) plus \(m\times n\) triangular or diagonal algebra. They never assemble \(B^{T}\otimes A\).
\end{proof}

\subsubsection{Solution of Problem 3(d)}

\paragraph{Answer.}
No. The unique Frobenius minimizer \(X_{\star}\) of (b) need not minimize \(\|AXB-C\|_{2}\).

\begin{proof}
By \cref{def:frobenius}, \(\|\cdot\|_{F}\) is the Euclidean norm of the stacked columns, equivalently the \(\ell^{2}\) norm of the singular values. The associated projection onto \(\operatorname{range}(\Phi)\) is orthogonal for \(\langle\cdot,\cdot\rangle_{F}\). By \cref{def:spectral},
\[
\|M\|_{2}
=
\sup_{\|x\|_{2}=1}\|Mx\|_{2}
\]
is the maximal Euclidean stretch of a column combination \(Mx\). That is again a right-multiplication quantity, but it records only the worst column combination, not the \(\ell^{2}\) sum over all entries. The two geometries on \(\mathbb{R}^{l\times p}\) are equivalent as norms (\cref{lem:norm-compare}) and nevertheless induce different best-approximation problems onto a proper subspace.

A concrete case of (b) is \cref{ex:p3-spectral}: \(A=\begin{pmatrix}1\\0\end{pmatrix}\), \(B=\begin{pmatrix}1&1\end{pmatrix}\), \(C=I_{2}\). Then \(\rank(A)=1=m\), \(\rank(B)=1=n\), so (b) supplies a unique Frobenius minimizer \(x_{\star}=\tfrac12\), with
\[
\|AX_{\star}B-C\|_{2}
=
\sqrt{\frac{3+\sqrt{5}}{4}}
>
1
=
\|-I_{2}\|_{2}
=
\|A\cdot 0\cdot B-C\|_{2}.
\]
Thus \(X_{\star}\) does not minimize the spectral residual.
\end{proof}

\subsection{Problem 4. [15 points]}
For the advection equation \(u_{t} + u_{x} = 0\), analyze the stability of the implicit schemes
\begin{enumerate}[label=(\alph*)]
    \item \(\dfrac{u_{j}^{n+1} - u_{j}^{n}}{\tau} + \dfrac{u_{j}^{n+1} - u_{j-1}^{n+1}}{h} = 0\);
    \item \(\dfrac{u_{j}^{n+1} - u_{j}^{n}}{\tau} + \dfrac{u_{j+1}^{n+1} - u_{j}^{n+1}}{h} = 0\);
    \item \(\dfrac{u_{j}^{n+1} - u_{j}^{n}}{\tau} + \dfrac{u_{j+1}^{n+1} - u_{j-1}^{n+1}}{2h} = 0\).
\end{enumerate}

\setcounter{section}{4}
\setcounter{theorem}{0}
\setcounter{definition}{0}
\setcounter{remark}{0}
\setcounter{note}{0}

\subsubsection{Preliminaries: what ``stability'' means, and how one computes it}

The question is not asking whether the PDE \(u_t+u_x=0\)\footnote{The (linear, constant-speed) \emph{advection} equation, also called transport or one-way wave, is \(u_t+a u_x=0\). Characteristics are the lines \(x-at=\mathrm{const}\), and the solution is the shift \(u(x,t)=u_0(x-at)\). This problem is the case \(a=1\).} is well-posed. That Cauchy problem is well-posed: the solution is a pure translation \(u(x,t)=u_0(x-t)\), so every \(L^p\) norm is conserved. The question is whether a \emph{difference scheme} can manufacture exponentially growing grid functions that the PDE does not have. Implicit time-stepping does not, by itself, prevent that; scheme (b) is the standard counter-example.

The official CAM syllabus names this circle as finite-difference methods, stability and convergence, and Lax equivalence; the standard references are Gustafsson--Kreiss--Oliger \cite{gko95}, LeVeque \cite{leveque07}, and Strikwerda \cite{strikwerda04}. Finite elements \cite{brennerscott10} are a different discretization (weak form, Sobolev spaces) and are not the language of Problems~4--5. The three schemes below only force von Neumann algebra. Consistency as a Taylor residual, convergence as a mesh limit, Lax equivalence as a theorem, the CFL obstruction, the Du~Fort--Frankel trap, and the FEM/spectral contrast are recorded after the solution (\cref{thm:lax-eq,def:convergence,exer:dufort,ex:fem-poisson}).

\paragraph{The one-page map.}

A linear evolution PDE is turned into a marching scheme in three moves: put unknowns on a grid; replace each derivative by a difference quotient; decide which time level the spatial differences sit on (explicit versus implicit). Three properties of the resulting scheme are then asked, in a fixed order.

\begin{center}
\begin{tikzpicture}[
  box/.style={draw, rounded corners=2pt, align=center, inner sep=5pt, font=\small},
  arr/.style={-{Latex}, thick}
]
  \node[box] (pde) at (0,2.4) {PDE \(u_t+Lu=0\)};
  \node[box] (sch) at (4.3,2.4) {difference scheme};
  \node[box] (con) at (0,0.6) {consistency\\(truncation \(\to 0\))};
  \node[box] (sta) at (4.3,0.6) {stability\\(no spurious growth)};
  \node[box] (cvg) at (2.15,-1.1) {convergence\\(\(U\to u\) as \(h,\tau\to 0\))};
  \draw[arr] (pde) -- (sch);
  \draw[arr] (sch) -- (con);
  \draw[arr] (sch) -- (sta);
  \draw[arr] (con) -- (cvg);
  \draw[arr] (sta) -- (cvg);
\end{tikzpicture}
\end{center}

Lax equivalence is the arrow into the bottom box: for a linear well-posed IVP, a consistent scheme converges if and only if it is stable \cite{leveque07}. Qiuzhen almost never asks you to quote the theorem by name. It asks you to \emph{compute} stability, usually by von Neumann (Problem~4, and most years' Problem~4/5) or by a discrete maximum principle (Problem~5).

\paragraph{Qiuzhen's actual exam language, 2023--2026.}

The PDE slot is extremely stereotyped. Learn two computational machines, not a survey of schemes.

{\footnotesize
\begin{center}
\begin{tabular}{@{}>{\raggedright\arraybackslash}p{2.35cm}
                   >{\raggedright\arraybackslash}p{4.55cm}
                   >{\raggedright\arraybackslash}p{3.35cm}
                   >{\raggedright\arraybackslash}p{3.55cm}@{}}
\toprule
Paper & Typical ask & Machine & Relation to 2026 \\
\midrule
2023 Fall P5 & Lax--Friedrichs for \(u_t+a u_{xxx}=0\) & von Neumann \(+\) truncation & same \(g(\xi)\) algebra as P4 \\
2024 Spring P4 & \(\theta\)-method for heat & von Neumann \(+\) truncation & implicit in time, like P4 \\
2024 Spring P5 & 1D monotone 3-point operator & max principle & 1D warm-up for P5 \\
2024 Fall P4 & first-order system, Lax--Friedrichs flavour & von Neumann & Fourier on a vector unknown \\
2024 Fall P5 & 2D centered convection--diffusion, \(f<0\) & max principle & P5 without backward Euler \\
2025 Spring P5 & Lax--Wendroff advection & \(\ell_2\) energy identity & alternative to Fourier \\
2025 Fall P4 & Du Fort--Frankel, 2D heat & consistency \(+\) von Neumann & \(|g|\le 1\) is not enough if inconsistent \\
2025 Fall P5 & implicit wave scheme & von Neumann / Jury & two-step \(g\), quadratic symbol \\
2026 Spring P4 & three implicit advection schemes & von Neumann & this problem \\
2026 Spring P5 & 2024 Fall P5 \(+\) backward Euler & max principle \(\Rightarrow\ell_\infty\) & next problem \\
\bottomrule
\end{tabular}
\end{center}
}

\begin{notation}[Grid, Courant number]
\label{not:grid-p4}
Fix mesh size \(h>0\) and time step \(\tau>0\). Write \(x_j=jh\) and \(t_n=n\tau\). The unknown \(u_j^n\) is a grid function intended to approximate \(u(x_j,t_n)\). For advection with speed \(1\), the mesh Courant number is
\[
r
\coloneqq
\frac{\tau}{h}.
\]
A scheme is \emph{explicit} if \(u^{n+1}\) is given by a formula in \(u^n\); it is \emph{implicit} if computing \(u^{n+1}\) requires solving a linear system coupling neighbouring nodes at time \(t_{n+1}\). All three schemes in the exam are implicit (backward Euler in time). The same spatial operator, written both ways, is \cref{ex:expl-impl}.
\end{notation}

\begin{example}[Explicit versus implicit, same upwind stencil]
\label{ex:expl-impl}
For \(u_t+u_x=0\), replace \(u_x\) by the backward difference \((u_j-u_{j-1})/h\). If that difference is taken at the old time \(t_n\),
\[
\frac{u_j^{n+1}-u_j^n}{\tau}
+
\frac{u_j^n-u_{j-1}^n}{h}
=0
\qquad\Longleftrightarrow\qquad
u_j^{n+1}
=
(1-r)\,u_j^n
+
r\,u_{j-1}^n.
\]
The right-hand side is a formula in already-known values: the scheme is explicit. If the same difference is taken at the new time \(t_{n+1}\), one obtains exam scheme~(a),
\[
\frac{u_j^{n+1}-u_j^n}{\tau}
+
\frac{u_j^{n+1}-u_{j-1}^{n+1}}{h}
=0
\qquad\Longleftrightarrow\qquad
-r\,u_{j-1}^{n+1}
+
(1+r)\,u_j^{n+1}
=
u_j^n.
\]
The unknowns at time \(t_{n+1}\) are coupled. On a periodic interval this is a cyclic bidiagonal system for the vector \(u^{n+1}\); that is what ``implicit'' means. (On the infinite line one can still march in the index \(j\), but the spatial operator still sits at the new time.) Schemes~(b) and~(c) are the same backward-Euler clock with the forward difference and the centered difference, respectively, in place of the backward one. Stability of the explicit member is computed in \cref{ex:explicit-upwind}; implicit does not by itself imply unconditional stability, which is the point of~(b).
\end{example}

\begin{definition}[Standard difference quotients]
\label{def:quotients}
For a grid function \(v_j\),
\begin{align*}
D^+v_j
&\coloneqq
\frac{v_{j+1}-v_j}{h},
\qquad
D^-v_j
\coloneqq
\frac{v_j-v_{j-1}}{h},
\qquad
D^0v_j
\coloneqq
\frac{v_{j+1}-v_{j-1}}{2h}.
\end{align*}
If \(v(x)\) is smooth, Taylor expansion at \(x_j\) gives
\[
D^+v=v_x+\tfrac12 h v_{xx}+O(h^2),\qquad
D^-v=v_x-\tfrac12 h v_{xx}+O(h^2),\qquad
D^0v=v_x+O(h^2).
\]
Thus \(D^\pm\) are first-order one-sided approximations of \(\partial_x\), and \(D^0\) is second-order centered. Replacing \(u_t\) by \((u_j^{n+1}-u_j^n)/\tau\) is likewise first-order in time (backward Euler).
\end{definition}

The three exam schemes are exactly backward Euler in time, combined with \(D^-\), \(D^+\), and \(D^0\) in space:
\begin{enumerate}[label=(\alph*)]
    \item implicit upwind (the one-sided difference looks \emph{backward} in \(x\), which is the inflow side when information travels to the right);
    \item implicit downwind (the opposite one-sided difference);
    \item implicit centered.
\end{enumerate}

\begin{definition}[Local truncation error / consistency]
\label{def:truncation}
Write a scheme in \emph{PDE form}: an operator \(\mathcal L_{h,\tau}\) acting on grid functions, set equal to zero. For exam scheme~(a) that operator is
\[
\bigl(\mathcal L_{h,\tau} U\bigr)_j^{n+1}
\coloneqq
\frac{U_j^{n+1}-U_j^n}{\tau}
+
\frac{U_j^{n+1}-U_{j-1}^{n+1}}{h}.
\]
The computed unknown \(U\) is defined by \(\mathcal L_{h,\tau}U=0\). It is not the same object as the PDE solution \(u\). Let \(u\) be a smooth solution of \(u_t+u_x=0\), and write \(u_j^n\coloneqq u(x_j,t_n)\) for its samples on the grid. The \emph{local truncation error} (or residual) is what you get by feeding those samples to the same operator:
\[
\mathcal T_j^{n+1}
\coloneqq
\bigl(\mathcal L_{h,\tau} u\bigr)_j^{n+1}
=
\frac{u(x_j,t_{n+1})-u(x_j,t_n)}{\tau}
+
\frac{u(x_j,t_{n+1})-u(x_j-h,t_{n+1})}{h}.
\]
This is an explicit number once \(u\) is given: every term on the right is a value of a known function. The scheme is \emph{consistent} with the PDE when this residual vanishes under mesh refinement; the formal statement is \cref{def:consistency}. The order is the largest pair \((p,q)\) for which \(\mathcal T=O(\tau^p+h^q)\).

Three objects that are easy to conflate:
\begin{enumerate}[label=\textup{(\roman*)}]
    \item \(U_j^n\): the grid function the scheme actually computes, \(\mathcal L_{h,\tau}U=0\).
    \item \(u_j^n=u(x_j,t_n)\): samples of a PDE solution. In general \(\mathcal L_{h,\tau}u\neq 0\).
    \item \(e_j^n=U_j^n-u_j^n\): the global error. Consistency is a statement about \(\mathcal T\), not about \(e\). Under stability, \(e=O(\mathcal T)\) on a fixed time interval (\cref{def:convergence,thm:lax-eq}).
\end{enumerate}
If the scheme is rewritten in update form \(U^{n+1}=S_{h,\tau}U^n\), the exact samples satisfy \(u^{n+1}=S_{h,\tau}u^n+\tau\mathcal T^{n+1}\), so \(\mathcal T\) as defined above is a residual \emph{per unit time}, comparable to \(u_t+u_x\).
\end{definition}

\begin{example}[Scheme (a): the residual as a Taylor polynomial]
\label{ex:trunc-a}
Expand at \((x_j,t_{n+1})\). The backward Euler quotient is
\[
\frac{u(x_j,t_{n+1})-u(x_j,t_n)}{\tau}
=
u_t
-
\tfrac12\tau\,u_{tt}
+
O(\tau^2),
\]
and \(D^-u(\,\cdot\,,t_{n+1})\) at \(x_j\) is \(u_x-\tfrac12 h u_{xx}+O(h^2)\) by \cref{def:quotients}. Adding,
\[
\mathcal T_j^{n+1}
=
\bigl(u_t+u_x\bigr)
-
\tfrac12\tau\,u_{tt}
-
\tfrac12 h\,u_{xx}
+
O(\tau^2+h^2).
\]
On solutions of \(u_t+u_x=0\) the first parenthesis vanishes, so
\[
\mathcal T_j^{n+1}
=
-\tfrac12\tau\,u_{tt}(x_j,t_{n+1})
-
\tfrac12 h\,u_{xx}(x_j,t_{n+1})
+
O(\tau^2+h^2)
=
O(\tau+h).
\]
The scheme is consistent of order \(O(\tau+h)\). Scheme~(b) is the same computation with \(D^+\), which flips the sign of the \(O(h)\) term:
\[
\mathcal T^{(b)}
=
-\tfrac12\tau\,u_{tt}
+
\tfrac12 h\,u_{xx}
+
O(\tau^2+h^2).
\]
\end{example}

\begin{example}[Scheme (c): the \(O(h)\) term cancels]
\label{ex:trunc-c}
Keep backward Euler in time and replace \(D^-\) by \(D^0\). Then \(D^0u=u_x+\tfrac16 h^2 u_{xxx}+O(h^4)\), and
\[
\mathcal T_j^{n+1}
=
-\tfrac12\tau\,u_{tt}(x_j,t_{n+1})
+
\tfrac16 h^2 u_{xxx}(x_j,t_{n+1})
+
O(\tau^2+h^4)
=
O(\tau+h^2).
\]
\end{example}

\begin{example}[A travelling wave, so every derivative is explicit]
\label{ex:trunc-sine}
The function \(u(x,t)=\sin(x-t)\) solves \(u_t+u_x=0\), with \(u_{tt}=u_{xx}=-\sin(x-t)\) and \(u_{xxx}=-\cos(x-t)\). Substituting into \cref{ex:trunc-a,ex:trunc-c} at \((x_j,t_{n+1})\) gives the concrete residuals
\begin{align*}
\mathcal T^{(a)}
&=
\tfrac12(\tau+h)\sin(x_j-t_{n+1})+O(\tau^2+h^2),
\\
\mathcal T^{(c)}
&=
\tfrac12\tau\sin(x_j-t_{n+1})
-
\tfrac16 h^2\cos(x_j-t_{n+1})
+O(\tau^2+h^4).
\end{align*}
No unknown grid function appears: these are ordinary Taylor remainders of a known smooth function. As \(h,\tau\to 0\) both tend to \(0\), so both schemes are consistent; (c) is one order higher in space.
\end{example}

\begin{example}[Explicit upwind, expansion at the old time]
\label{ex:trunc-expl-upwind}
The explicit scheme of \cref{ex:expl-impl} has
\[
\bigl(\mathcal L_{h,\tau}U\bigr)_j^n
=
\frac{U_j^{n+1}-U_j^n}{\tau}
+
\frac{U_j^n-U_{j-1}^n}{h}.
\]
Expanding at \((x_j,t_n)\) now produces a \emph{forward} Euler quotient \(u_t+\tfrac12\tau u_{tt}+O(\tau^2)\), and \(D^-\) at time \(t_n\), hence
\[
\mathcal T_j^n
=
\tfrac12\tau\,u_{tt}(x_j,t_n)
-
\tfrac12 h\,u_{xx}(x_j,t_n)
+
O(\tau^2+h^2).
\]
Same order \(O(\tau+h)\) as implicit upwind; the time-level of the spatial difference only changes the sign of the \(O(\tau)\) term. Explicit centered (\cref{ex:centered-unstable}) is \(O(\tau+h^2)\) and still useless, because consistency does not prevent \(|g|>1\).
\end{example}

Consistency is \emph{not} asked in 2026 Problem~4, but it is asked in 2023 Fall, 2024 Spring/Fall, and 2025 Fall. The three exam schemes are \cref{ex:trunc-a,ex:trunc-c}: (a) and (b) are \(O(\tau+h)\), (c) is \(O(\tau+h^2)\).

\begin{definition}[Consistency]
\label{def:consistency}
Let \(\mathcal L_{h,\tau}\) be a linear scheme written in PDE form, as in \cref{def:truncation}, and let \(\mathcal T_j^n=(\mathcal L_{h,\tau}u)_j^n\) be the local truncation error obtained by feeding a smooth PDE solution \(u\) to that operator. The scheme is \emph{consistent} with the PDE if, at every fixed \((x,t)\) in the space--time domain,
\[
\mathcal T_j^n
\to
0
\qquad
\text{as }
(h,\tau)\to(0,0),
\]
for every sufficiently smooth solution \(u\). Equivalently, in a discrete norm, \(\|\mathcal T^n\|_h\to 0\) uniformly for \(t_n\le T\). The scheme has order \((p,q)\) when
\[
\mathcal T
=
O(\tau^p+h^q)
\]
as \((h,\tau)\to(0,0)\). Consistency is the statement that some positive \(p,q\) exist; the pair \((p,q)\) is the rate.

The limit is taken along the family of meshes actually used. If a CFL or parabolic restriction is part of that family (for instance \(\tau/h\le 1\), or \(\tau=O(h^2)\)), consistency is only claimed under that restriction. Du~Fort--Frankel is the standard trap: it is unconditionally von Neumann stable, but consistent with the heat equation if and only if \(\tau/h\to 0\) (\cref{exer:dufort}).

Consistency is a statement about \(\mathcal T\), not about the global error \(U-u\), and not about \(\lvert g(\xi)\rvert\). Lax equivalence (\cref{thm:lax-eq}) is the pairing: for a linear well-posed IVP, consistency plus stability yields convergence of the same order as \(\mathcal T\).
\end{definition}

\begin{definition}[Lax--Richtmyer stability]
\label{def:lax-richtmyer}
Consider a linear scheme written as \(U^{n+1}=S_{h,\tau}U^n\) on grid functions, with a discrete norm \(\|\cdot\|_h\) (usually \(\ell_2\) or \(\ell_\infty\)). The scheme is \emph{stable} on a time interval \([0,T]\) if there exist \(C,\alpha\), independent of \(h\) and \(\tau\) (for all sufficiently small steps, possibly subject to a restriction such as \(r\le 1\)), such that
\[
\|U^n\|_h
\le
C\,e^{\alpha t_n}\,\|U^0\|_h,
\qquad
0\le t_n\le T.
\]
In particular, if \(\|S_{h,\tau}\|\le 1\), one may take \(C=1\) and \(\alpha=0\): the scheme is non-expansive, which is stronger than needed but typical for conservative or dissipative advection schemes.
\end{definition}

Two remarks that prevent the usual confusion with ODE ``absolute stability''.

\begin{enumerate}[label=\textup{(\roman*)}]
    \item The constant \(C e^{\alpha T}\) is allowed to grow in \(T\), because the PDE itself may grow (heat equation with a positive zeroth-order term, for instance). What is forbidden is growth like \((1/h)^n\) or \(g^n\) with \(|g|>1\) independent of \(\tau\): that would explode as the mesh is refined with \(t_n\) fixed.
    \item For pure advection or heat on the line, the PDE does not grow, and the practical test is \(\|U^n\|_h\le\|U^0\|_h\), or at least \(\|U^n\|_h\le C\|U^0\|_h\) with \(C\) independent of \(h,\tau\).
\end{enumerate}

\begin{definition}[von Neumann / Fourier criterion]
\label{def:von-neumann}
On the infinite line (or a periodic interval), a constant-coefficient linear scheme admits Fourier modes
\[
u_j^n
=
g(\xi)^n\,e^{\mathrm i j\xi},
\qquad
\xi\in\mathbb R.
\]
Here \(\xi\) is the dimensionless frequency (one may write \(\xi=\kappa h\) with \(\kappa\) the dual variable to \(x\)). Substituting produces a scalar (or a matrix, for systems / multistep schemes) \emph{amplification factor} \(g(\xi)\). The scheme is von Neumann stable if there is \(K\), independent of \(\xi,h,\tau\), such that
\[
\bigl|g(\xi)\bigr|
\le
1+K\tau
\qquad
\text{for all }\xi.
\]
For the advection equation, and for every scheme in this problem, \(g\) depends on \((h,\tau)\) only through \(r=\tau/h\), so the criterion collapses to
\[
\bigl|g(\xi)\bigr|\le 1
\qquad
\text{for all }\xi\in[0,2\pi].
\]
\end{definition}

Fourier substitution is the method that produces \(g(\xi)\). It is not a second definition of stability, and it is not a Fourier transform in time. It is the statement that a constant-coefficient scheme on \(\mathbb Z\) (or on a periodic grid) is diagonalized by spatial exponentials, so \(\ell_2\) stability reduces to a scalar inequality for each frequency.

\paragraph{What the method is for.}

Lax--Richtmyer stability (\cref{def:lax-richtmyer}) asks for a bound on \(\|U^n\|_h\) for \emph{every} initial grid function. That looks infinite-dimensional. For a linear scheme whose coefficients do not depend on \(j\), the shift \(U_j\mapsto U_{j\pm 1}\) is a normal operator on \(\ell_2(\mathbb Z)\), and the complex exponentials \(j\mapsto e^{\mathrm i j\xi}\) are its eigenfunctions. Any such scheme is therefore a Fourier multiplier: in frequency space, one step of the scheme is ordinary multiplication by a number \(g(\xi)\). Checking that no frequency grows is then necessary and sufficient for an \(\ell_2\) bound. That is the whole point. Truncation error asks a different question (how well the stencil matches the PDE). Fourier substitution never sees the PDE solution; it only sees the homogeneous scheme.

The continuous analogue is the same idea used on the PDE itself. A mode \(u(x,t)=e^{\mathrm i\kappa x}y(t)\) in \(u_t+u_x=0\) yields \(y'=-\mathrm i\kappa y\), hence \(y(t)=e^{-\mathrm i\kappa t}y(0)\) with \(\lvert e^{-\mathrm i\kappa t}\rvert=1\). The PDE does not amplify any frequency. A difference scheme is allowed to dissipate (\(|g|<1\)) but not to amplify (\(|g|>1\)).

\begin{proposition}[Plancherel \(\Rightarrow\) von Neumann is \(\ell_2\) stability]
\label{prop:vn-plancherel}
Let \(U=(U_j)_{j\in\mathbb Z}\) be square-summable, and write
\[
\hat U(\xi)
=
\sum_{j\in\mathbb Z}U_j e^{-\mathrm i j\xi},
\qquad
\xi\in[0,2\pi],
\]
so that Parseval's identity reads
\[
\sum_{j\in\mathbb Z}\lvert U_j\rvert^2
=
\frac{1}{2\pi}\int_0^{2\pi}\lvert\hat U(\xi)\rvert^2\,\mathrm d\xi.
\]
Suppose a one-step linear scheme with \(j\)-independent coefficients acts as \(\widehat{SU}(\xi)=g(\xi)\hat U(\xi)\). Then
\[
\|S^n U\|_{\ell_2}
=
\Bigl(\frac{1}{2\pi}\int_0^{2\pi}\lvert g(\xi)\rvert^{2n}\lvert\hat U(\xi)\rvert^2\,\mathrm d\xi\Bigr)^{1/2}
\le
\bigl(\operatorname{ess\,sup}_{\xi}\lvert g(\xi)\rvert\bigr)^n
\|U\|_{\ell_2}.
\]
Thus \(\lvert g(\xi)\rvert\le 1\) for a.e.\ \(\xi\) if and only if \(\|U^n\|_2\le\|U^0\|_2\) for every \(\ell_2\) initial datum. If \(\lvert g(\xi_0)\rvert=1+\delta>1\) on a set of positive measure, some modes grow like \((1+\delta)^n=(1+\delta)^{t_n/\tau}\), which is not a bound of the form \(C e^{\alpha T}\) as \(\tau\to 0\). The GKO relaxation \(\lvert g\rvert\le 1+K\tau\) is the same statement with PDE growth \(e^{KT}\) allowed \cite{gko95,leveque07,strikwerda04}.
\end{proposition}

On a periodic grid of \(N\) nodes the same argument is a finite DFT: the matrix of the scheme is circulant, hence diagonalized by the discrete Fourier modes \(e^{2\pi\mathrm i jk/N}\), and \(\lvert g\rvert\) is the spectral radius of one time step. The infinite-line formulae are the \(N\to\infty\) version. Either way, it is enough to test the scheme on the family \(u_j=e^{\mathrm i j\xi}\): those functions already span the space in which \(\ell_2\) stability is measured.

The ansatz \(u_j^n=g^n e^{\mathrm i j\xi}\) is the statement that a pure frequency remains a pure frequency, with amplitude multiplied by \(g\) at each step. The sampled physical wave \(e^{\mathrm i\kappa x_j}=e^{\mathrm i\kappa jh}\) is exactly this family with \(\xi=\kappa h\). On the grid, \(\xi\) and \(\xi+2\pi\) are indistinguishable (aliasing), so one may restrict to \(\xi\in[0,2\pi]\).

\paragraph{The substitution, written as algebra.}

Fix one frequency \(\xi\). Every grid value of the mode is a multiple of the same complex number \(e^{\mathrm i j\xi}\), and a time step only changes the prefactor:

\begin{center}
{\small
\begin{tabular}{@{}lll@{}}
\toprule
grid value & substitution & relative to \(u_j^n=g^n e^{\mathrm i j\xi}\) \\
\midrule
\(u_j^n\) & \(g^n e^{\mathrm i j\xi}\) & \(1\) \\
\(u_{j\pm 1}^n\) & \(g^n e^{\mathrm i(j\pm 1)\xi}\) & \(e^{\pm\mathrm i\xi}\) \\
\(u_j^{n+1}\) & \(g^{n+1}e^{\mathrm i j\xi}\) & \(g\) \\
\(u_{j\pm 1}^{n+1}\) & \(g^{n+1}e^{\mathrm i(j\pm 1)\xi}\) & \(g\,e^{\pm\mathrm i\xi}\) \\
\(u_j^{n-1}\) & \(g^{n-1}e^{\mathrm i j\xi}\) & \(g^{-1}\) \\
\bottomrule
\end{tabular}
}
\end{center}

The last row is for two-step schemes. After substitution, every term contains the common factor \(g^n e^{\mathrm i j\xi}\neq 0\). Cancel it; what remains is a scalar (or polynomial) equation for \(g\). Then estimate \(\lvert g\rvert\).

\begin{example}[Fourier substitution, written out for explicit upwind]
\label{ex:fourier-subst}
Take the explicit scheme of \cref{ex:expl-impl},
\[
u_j^{n+1}
=
(1-r)u_j^n
+
r\,u_{j-1}^n.
\]
Insert \(u_j^n=g^n e^{\mathrm i j\xi}\). The left-hand side is \(g^{n+1}e^{\mathrm i j\xi}\). On the right,
\[
u_{j-1}^n
=
g^n e^{\mathrm i(j-1)\xi}
=
e^{-\mathrm i\xi}\,g^n e^{\mathrm i j\xi}.
\]
The equation becomes
\[
g^{n+1}e^{\mathrm i j\xi}
=
(1-r)\,g^n e^{\mathrm i j\xi}
+
r e^{-\mathrm i\xi}\,g^n e^{\mathrm i j\xi}.
\]
Divide by \(g^n e^{\mathrm i j\xi}\):
\[
g(\xi)
=
1-r+r e^{-\mathrm i\xi}
=
1-r+r\cos\xi
-
\mathrm i r\sin\xi.
\]
This is the amplification factor: one time step multiplies the frequency-\(\xi\) component of the grid function by this complex number. The \(\ell_2\) question is whether \(\lvert g(\xi)\rvert\le 1\) for every \(\xi\). Write \(g=A+\mathrm i B\) with \(A=1-r+r\cos\xi\) and \(B=-r\sin\xi\), and use \(\lvert g\rvert^2=A^2+B^2\):
\begin{align*}
\lvert g(\xi)\rvert^2
&=
(1-r+r\cos\xi)^2
+
r^2\sin^2\xi
=
(1-r)^2
+
2r(1-r)\cos\xi
+
r^2
\\
&=
1-2r(1-r)(1-\cos\xi).
\end{align*}
Since \(1-\cos\xi\ge 0\), the correction \(-2r(1-r)(1-\cos\xi)\) is nonpositive for all \(\xi\) if and only if \(r(1-r)\ge 0\), i.e. \(0\le r\le 1\). The same identity rules out \(r<0\): then \(r(1-r)<0\), so the correction is nonnegative and \(\lvert g(\xi)\rvert^2>1\) whenever \(\xi\not\equiv 0\pmod{2\pi}\). At the checkerboard \(\xi=\pi\) (where \(e^{\mathrm i j\pi}=(-1)^j\)) one has \(g=1-2r\) in any case: if \(r>1\) then \(\lvert g\rvert=2r-1>1\), and if \(r<0\) then \(\lvert g\rvert=1-2r>1\). In this problem \(h,\tau>0\), so \(r=\tau/h>0\) by \cref{not:grid-p4}; the remaining restriction is the upper bound. Allowing \(r<0\) anyway (negative speed, or this left-looking stencil used as downwind) is unconditionally unstable. Thus the scheme is von Neumann stable if and only if \(0\le r\le 1\).

Two sanity checks that do not require the full \(\lvert g\rvert^2\) expansion. At \(\xi=0\) (constants), \(g=1\): a constant grid function is translated without change, as the PDE conserves constants. At \(\xi=\pi\), the scheme reduces to the scalar recurrence \(v^{n+1}=(1-2r)v^n\) on \(u_j^n=v^n(-1)^j\); this is the mode that detects both \(r>1\) and \(r<0\). Fourier substitution is exactly ``run the scheme on every such mode, including the ones you would not think to test by hand''.
\end{example}

The geometric view of the same \(g\) is sometimes faster on a script: for \(0\le r\le 1\), \(g\) is a convex combination of \(1\) and \(e^{-\mathrm i\xi}\), both of which lie on the unit circle, hence \(\lvert g\rvert\le 1\). The \(\lvert g\rvert^2\) computation is the one that also handles implicit schemes, whose \(g\) is a quotient rather than a convex combination.

\paragraph{The mechanical recipe (write this on the script).}

\begin{enumerate}[label=\textup{\arabic*.}]
    \item Insert \(u_j^n=g^n e^{\mathrm i j\xi}\) into the scheme. Every shift \(j\mapsto j\pm 1\) multiplies by \(e^{\pm\mathrm i\xi}\); every time step multiplies by \(g\), as in the table above.
    \item Cancel the common factor \(g^n e^{\mathrm i j\xi}\) and solve for \(g=g(\xi;r)\).
    \item Compute \(\lvert g\rvert^2\) by writing the denominator (or the whole of \(g\)) as \(A+\mathrm i B\) and using \(\lvert A+\mathrm i B\rvert^2=A^2+B^2\).
    \item Decide whether \(\lvert g(\xi)\rvert\le 1\) for \emph{all} \(\xi\in[0,2\pi]\), and if not, find the restriction on \(r\) (or report unconditional instability). The dangerous frequencies are typically \(\xi=\pi\) or \(\xi=\pi/2\); they are worth testing first as a sanity check, but they do not replace the all-\(\xi\) estimate.
\end{enumerate}

For a two-step scheme the same substitution is used, but \(u^{n+1}\), \(u^n\), and \(u^{n-1}\) all appear, so cancelling \(g^n e^{\mathrm i j\xi}\) leaves a \emph{quadratic polynomial} in the unknown \(g\). That polynomial, and the Jury test that keeps both roots in the unit disk, are computed in \cref{ex:theta-heat} rather than quoted as a slogan.

\paragraph{A dictionary, before the examples.}

The next calculations all answer the same pair of exam questions --- is the scheme consistent, and for which meshes is it von Neumann stable --- but they use words that are easy to mix. Here they are as objects, not as slogans.

\begin{definition}[Conditional / unconditional von Neumann stability]
\label{def:cond-uncond}
Fix a family of schemes indexed by the steps \(h,\tau>0\). After Fourier substitution one has an amplification factor \(g(\xi;h,\tau)\).
\begin{itemize}
    \item The scheme is \emph{von Neumann stable} for a given \((h,\tau)\) if \(\lvert g(\xi;h,\tau)\rvert\le 1\) for every \(\xi\in[0,2\pi]\) (or \(\lvert g\rvert\le 1+K\tau\), in the GKO form).
    \item It is \emph{conditionally stable} if that inequality holds only when \((h,\tau)\) obeys an extra restriction, typically a bound on \(r=\tau/h\) or on \(\mu=\tau/h^2\). Explicit upwind is the model: stable iff \(0<r\le 1\).
    \item It is \emph{unconditionally stable} if the inequality holds for every \(h,\tau>0\), with no restriction on \(r\) or \(\mu\). Backward Euler on the heat equation is the model.
    \item It is \emph{unconditionally unstable} if, for every \(h,\tau>0\) however small, there exists some \(\xi\) with \(\lvert g(\xi)\rvert>1\). Explicit centered advection is the model. ``Unconditional'' here modifies the instability: there is no mesh on which the scheme becomes stable.
\end{itemize}
Consistency (\cref{def:truncation}) is a different question: \(\mathcal T\to 0\) as \(h,\tau\to 0\). A scheme may be consistent and unconditionally unstable (\cref{ex:centered-unstable}), or stable and inconsistent (\cref{exer:dufort}). Neither adjective is implied by ``implicit''.
\end{definition}

\begin{definition}[One-step versus two-step / three-level]
\label{def:one-two-step}
A scheme is \emph{one-step} if the vector \(U^{n+1}\) is determined by \(U^n\) alone: \(U^{n+1}=S_{h,\tau}U^n\). Every scheme in 2026 Problem~4 is one-step (backward Euler in time). Fourier substitution then produces a \emph{linear} equation for \(g\), hence a unique amplification factor.

A scheme is \emph{two-step} (also called \emph{three-level}, or, in the 2024 Spring wording, a ``three-step method'') if \(U^{n+1}\) depends on both \(U^n\) and \(U^{n-1}\). Starting the march requires two initial time levels. Fourier substitution then produces a quadratic
\[
A(\xi)\,g^2+B(\xi)\,g+C(\xi)=0
\]
with two roots \(g_\pm(\xi)\). Both must lie in the closed unit disk: one root is the physical mode (the one that approximates \(e^{\lambda\tau}\)), the other is a parasitic computational mode created by the extra time level. The polynomial in \(g\) is not the truncation error. Truncation of order \(O(\tau^2)\) is sometimes called a quadratic error in \(\tau\); that is a statement about \(\mathcal T\), not about this quadratic.
\end{definition}

\begin{proposition}[Jury's criterion for a real quadratic]
\label{prop:jury}
Let \(p(g)=Ag^2+Bg+C\) with \(A>0\) and \(A,B,C\in\mathbb R\). Both roots satisfy \(\lvert g\rvert\le 1\) if and only if
\[
p(1)\ge 0,
\qquad
p(-1)\ge 0,
\qquad
\lvert C\rvert\le A.
\]
(If one wants both roots \emph{strictly} inside the open disk, replace \(\ge,\le\) by \(>,<\). Roots on the unit circle are permitted for von Neumann stability, and occur for non-dissipative schemes.) This is the discrete analogue of the Routh--Hurwitz criterion: Routh--Hurwitz tests \(\operatorname{Re}\lambda\le 0\) for \(y'=\lambda y\); Jury tests \(\lvert g\rvert\le 1\) for \(y^{n+1}=g y^n\).
\end{proposition}

\begin{example}[The ODE test equation, and what ``symbol'' means]
\label{ex:ode-test}
Fourier substitution on a PDE scheme is the same calculation one already knows for ODEs, applied one frequency at a time. The ODE test equation is \(y'=\lambda y\) with \(\lambda\in\mathbb C\). Forward Euler \(y^{n+1}=y^n+\tau\lambda y^n\) has \(g=1+\tau\lambda\); backward Euler
\[
\frac{y^{n+1}-y^n}{\tau}
=
\lambda y^{n+1}
\qquad\Longrightarrow\qquad
g=\frac{1}{1-\tau\lambda}.
\]
The method is \emph{absolutely stable} at this \(\lambda\) if \(\lvert g(\tau\lambda)\rvert\le 1\). Backward Euler is \(A\)-stable: \(\lvert g\rvert\le 1\) whenever \(\operatorname{Re}\lambda\le 0\), i.e.\ on the whole left half-plane, which is exactly the region where the ODE itself does not grow.

A constant-coefficient PDE is a superposition of such scalar ODEs. Insert \(u(x,t)=e^{\mathrm i\kappa x}y(t)\) into \(u_t+u_x=0\):
\[
y'+\mathrm i\kappa y=0,
\qquad
\lambda=-\mathrm i\kappa,
\qquad
\operatorname{Re}\lambda=0.
\]
The exact amplification over a time step \(\tau\) is \(e^{\lambda\tau}=e^{-\mathrm i\kappa\tau}\), of modulus \(1\). A difference scheme replaces \(\partial_x\) by a difference quotient, hence replaces \(\lambda=-\mathrm i\kappa\) by a \emph{discrete symbol} \(\lambda_h(\xi)\), the eigenvalue of that difference operator on the grid mode \(e^{\mathrm i j\xi}\). Explicitly, on \(v_j=e^{\mathrm i j\xi}\),
\begin{align*}
D^-v_j
&=
\frac{1-e^{-\mathrm i\xi}}{h}\,v_j,
&
D^+v_j
&=
\frac{e^{\mathrm i\xi}-1}{h}\,v_j,
&
D^0v_j
&=
\frac{\mathrm i\sin\xi}{h}\,v_j.
\end{align*}
The scheme \(u_t+u_x=0\) discretized as \((U^{n+1}-U^n)/\tau+D U^{n+1}=0\) is backward Euler on \(y'=\lambda_h y\) with \(\lambda_h\) equal to minus the spatial operator on that mode:
\begin{align*}
\lambda_h^-
&=
-\frac{1-e^{-\mathrm i\xi}}{h}
=
-\frac{1-\cos\xi}{h}-\mathrm i\frac{\sin\xi}{h},
\\
\lambda_h^+
&=
-\frac{e^{\mathrm i\xi}-1}{h}
=
\frac{1-\cos\xi}{h}-\mathrm i\frac{\sin\xi}{h},
\\
\lambda_h^0
&=
-\mathrm i\frac{\sin\xi}{h}.
\end{align*}
Upwind has \(\operatorname{Re}\lambda_h^-\le 0\) (dissipative, like heat). Centered has \(\operatorname{Re}\lambda_h^0=0\) (purely imaginary, like the PDE). Downwind has \(\operatorname{Re}\lambda_h^+\ge 0\) (anti-dissipative: the scalar ODE itself grows). Backward Euler damps the first two unconditionally, because they lie in the closed left half-plane; it damps the third only when \(\lvert\tau\lambda_h^+\rvert\) is large enough that \(g=1/(1-\tau\lambda_h^+)\) falls back into the unit disk, which is the restriction \(r\ge 1\) in scheme~(b). That is the only fact needed to see why ``implicit \(\Rightarrow\) unconditionally stable'' is true for heat and false for downwind advection.
\end{example}

What Fourier substitution does \emph{not} do: it does not treat variable coefficients (except by freezing them, which is heuristic). It does not treat boundary conditions (that is GKS theory in \cite{gko95}), and it does not give \(\ell_\infty\) stability. Problem~5 uses a discrete maximum principle for that reason.

This is the computation the exam wants. It is necessary for Lax--Richtmyer stability in \(\ell_2\), and for constant-coefficient Cauchy problems it is essentially sufficient \cite{gko95,leveque07}. The three schemes of Problem~4 are the same substitution, with the spatial difference evaluated at time \(t_{n+1}\) so that \(g\) appears in the denominator.

\begin{example}[Explicit upwind, the comparison scheme]
\label{ex:explicit-upwind}
The scheme, the substitution, and the identity \(\lvert g\rvert^2=1-2r(1-r)(1-\cos\xi)\) are \cref{ex:fourier-subst}: the scheme is \emph{conditionally} von Neumann stable, the condition being \(0<r\le 1\). What remains is the name of that condition.

CFL is Courant--Friedrichs--Lewy, a statement about domains of dependence, not about \(\lvert g\rvert\). The characteristic of \(u_t+u_x=0\) through the node \((x_j,t_{n+1})\) hits time \(t_n\) at \(x_j-\tau\). An explicit scheme that only reads \(\{u_{j-1}^n,u_j^n\}\) can see at most distance \(h\) to the left, so a necessary condition for the numerical value to depend on the data the PDE actually uses is \(\tau\le h\), i.e. \(r\le 1\). For this particular scheme the CFL restriction and the von Neumann restriction coincide. They need not: explicit centered satisfies a CFL bound of the same form and is nevertheless unconditionally unstable (\cref{exer:cfl}). Fourier substitution is the sufficient \(\ell_2\) test; CFL is only an obstruction.
\end{example}

\begin{example}[Explicit centered advection is unconditionally unstable]
\label{ex:centered-unstable}
The scheme is forward Euler in time and centered in space,
\[
\frac{U_j^{n+1}-U_j^n}{\tau}
+
\frac{U_{j+1}^n-U_{j-1}^n}{2h}
=0
\qquad\Longleftrightarrow\qquad
U_j^{n+1}
=
U_j^n
-
\tfrac12 r\bigl(U_{j+1}^n-U_{j-1}^n\bigr).
\]
\emph{What we compute.} Two independent things: the truncation \(\mathcal T\) (consistency) and the amplification \(g\) (stability).

\emph{Consistency.} Expand at \((x_j,t_n)\). The forward Euler quotient is \(u_t+\tfrac12\tau u_{tt}+O(\tau^2)\) and \(D^0u=u_x+\tfrac16 h^2 u_{xxx}+O(h^4)\), so on solutions of \(u_t+u_x=0\)
\[
\mathcal T_j^n
=
\tfrac12\tau\,u_{tt}
+
\tfrac16 h^2 u_{xxx}
+
O(\tau^2+h^4)
=
O(\tau+h^2).
\]
The scheme is consistent of order \(O(\tau+h^2)\). If stability held, Lax equivalence would give convergence of that order. Stability fails, so this order is irrelevant.

\emph{Stability.} The substitution table gives \(U_j^{n+1}=g U_j^n\) and \(U_{j\pm 1}^n=e^{\pm\mathrm i\xi}U_j^n\). Cancelling \(U_j^n\) produces
\[
\frac{g-1}{\tau}
+
\frac{e^{\mathrm i\xi}-e^{-\mathrm i\xi}}{2h}
=0
\qquad\Longrightarrow\qquad
g
=
1-\mathrm i r\sin\xi.
\]
Hence \(\lvert g\rvert^2=1+r^2\sin^2\xi\). For every \(r>0\) one has \(\lvert g(\pi/2)\rvert^2=1+r^2>1\). There is no restriction on \(r\) that saves the scheme: it is unconditionally unstable. Equivalently, the discrete symbol of \(-D^0\) is purely imaginary, and forward Euler applied to \(y'=\mathrm i\omega y\) yields \(g=1+\mathrm i\tau\omega\) with \(\lvert g\rvert>1\) as soon as \(\omega\neq 0\). Centered differences are the right spatial operator for an implicit time stepper (\cref{ex:trunc-c} and exam scheme~(c)); they are the wrong operator for explicit first-order advection.
\end{example}

\begin{example}[Backward Euler for heat, unconditional \(\ell_2\) stability]
\label{ex:be-heat}
The PDE is now \(u_t=u_{xx}\), not advection. The scheme is backward Euler in time and the standard second difference in space:
\[
\frac{U_j^{n+1}-U_j^n}{\tau}
=
\frac{U_{j+1}^{n+1}-2U_j^{n+1}+U_{j-1}^{n+1}}{h^2}.
\]
\emph{What we compute.} The amplification \(g(\xi;\mu)\) as a function of the parabolic mesh ratio
\[
\mu
\coloneqq
\frac{\tau}{h^2}.
\]
(The hyperbolic ratio \(r=\tau/h\) is the wrong quantity here: heat has a second space derivative, so the dimensionless combination is \(\tau/h^2\).)

\emph{Substitution.} \(U_j^{n+1}=g U_j^n\) and \(U_{j\pm 1}^{n+1}=g e^{\pm\mathrm i\xi}U_j^n\), so the second difference at time \(t_{n+1}\) is
\[
\frac{g\bigl(e^{\mathrm i\xi}-2+e^{-\mathrm i\xi}\bigr)}{h^2}U_j^n
=
-\frac{4g\sin^2(\xi/2)}{h^2}U_j^n,
\]
using \(1-\cos\xi=2\sin^2(\xi/2)\). The scheme becomes
\[
\frac{g-1}{\tau}
=
-\frac{4\mu\sin^2(\xi/2)}{\tau}\,g
\qquad\Longrightarrow\qquad
g
=
\frac{1}{1+\eta},
\qquad
\eta
\coloneqq
4\mu\sin^2(\xi/2)\ge 0.
\]
Thus \(0<g(\xi)\le 1\) for every \(\mu>0\) and every \(\xi\). The scheme is unconditionally von Neumann stable. Equivalently: the discrete symbol of \(\delta_{xx}\) on this mode is \(\lambda_h=-4\sin^2(\xi/2)/h^2\le 0\), and backward Euler is \(A\)-stable on the negative axis.

\emph{Consistency, for the record.} Expanding at \(t_{n+1}\) gives \(\mathcal T=O(\tau+h^2)\). Combined with the stability just proved, Lax equivalence yields convergence of order \(O(\tau+h^2)\) with no mesh restriction. The explicit sibling (forward Euler plus the same second difference) has the same truncation but is only conditionally stable, \(\mu\le 1/2\) (\cref{exer:ftcs}). This pair is the model of ``implicit \(\Rightarrow\) unconditionally stable'' that scheme~(b) of the exam will break: (b) is implicit and only conditionally stable.
\end{example}

\begin{example}[2024 Spring \(\theta\)-method: truncation, a quadratic in \(g\), and Jury]
\label{ex:theta-heat}
The 2024 Spring Problem~4 is the typical ``analyse truncation error and stability'' prompt. The PDE is the heat equation \(u_t-a^2 u_{xx}=0\). The scheme, with time step \(k>0\) and a parameter \(\theta\ge 0\), is
\[
(1+\theta)\frac{U_j^{n+1}-U_j^n}{k}
-
\theta\frac{U_j^n-U_j^{n-1}}{k}
=
a^2\frac{U_{j+1}^{n+1}-2U_j^{n+1}+U_{j-1}^{n+1}}{h^2}.
\]
Three time levels appear, so this is a two-step implicit method in the sense of \cref{def:one-two-step} (the exam called it a three-step method). For \(\theta=0\) it collapses to backward Euler (\cref{ex:be-heat}). The unknowns at time \(t_{n+1}\) still occur linearly on both sides, so each step is an implicit linear solve; the extra ingredient is that the scheme also looks at \(U^{n-1}\).

\emph{What we are asked for.} The local truncation \(\mathcal T\) as a function of \(\theta\), the value of \(\theta\) at which the order is highest, and the von Neumann condition on \((k,h,\theta)\).

\emph{Truncation.} Let \(u\) be a smooth solution of \(u_t=a^2 u_{xx}\), and define \(\mathcal T\) by substituting samples of \(u\) into the scheme written in PDE form (left-hand side minus right-hand side). The spatial operator sits at time \(t_{n+1}\), so expand at \((x_j,t_{n+1})\). The two time quotients, expanded at \(t_{n+1}\), begin
\begin{align*}
\frac{u^{n+1}-u^n}{k}
&=
u_t-\tfrac12 k\,u_{tt}+O(k^2),
\\
\frac{u^n-u^{n-1}}{k}
&=
u_t-\tfrac32 k\,u_{tt}+O(k^2),
\end{align*}
hence the weighted left-hand side is
\[
u_t+(\theta-\tfrac12)k\,u_{tt}+O(k^2).
\]
The right-hand side is \(a^2 u_{xx}+\tfrac{1}{12}a^2 h^2 u_{xxxx}+O(h^4)\). Using \(u_t=a^2 u_{xx}\),
\[
\mathcal T
=
(\theta-\tfrac12)k\,u_{tt}
-
\tfrac{1}{12}a^2 h^2 u_{xxxx}
+
O(k^2+h^4).
\]
If \(\theta\neq\tfrac12\), this is \(O(k+h^2)\). If \(\theta=\tfrac12\), the \(O(k)\) term cancels and \(\mathcal T=O(k^2+h^2)\). That is the highest order; it is a quadratic error in the time step, meaning \(\mathcal T=O(k^2+\cdots)\), not meaning the amplification polynomial is quadratic. The two uses of ``quadratic'' are unrelated.

\emph{Amplification polynomial.} Insert \(U_j^n=g^n e^{\mathrm i j\xi}\). Set \(\mu=a^2 k/h^2\) and \(\eta=4\mu\sin^2(\xi/2)\ge 0\), as in \cref{ex:be-heat}. The identity \(e^{\mathrm i\xi}-2+e^{-\mathrm i\xi}=-4\sin^2(\xi/2)\) turns the scheme into
\[
(1+\theta)(g-1)-\theta\bigl(1-g^{-1}\bigr)+\eta g=0.
\]
Multiply through by \(g\) (legitimate for \(g\neq 0\), and \(g=0\) is not a root unless \(\theta=0\), which already reduces to one-step):
\[
(1+\theta+\eta)\,g^2-(1+2\theta)\,g+\theta=0.
\]
This is the quadratic in \(g\). Write \(A=1+\theta+\eta>0\), \(B=-(1+2\theta)\), \(C=\theta\). The two roots \(g_\pm(\xi)\) are the two amplification factors. Von Neumann stability requires \(\lvert g_\pm(\xi)\rvert\le 1\) for every \(\xi\), equivalently Jury's criterion (\cref{prop:jury}) for every \(\eta\ge 0\):
\begin{align*}
p(1)
&=
A+B+C
=
\eta
\ge 0,
\\
p(-1)
&=
A-B+C
=
2+4\theta+\eta
>0,
\\
\lvert C\rvert\le A
&\qquad
\theta\le 1+\theta+\eta,
\end{align*}
all three of which hold for every \(\theta\ge 0\) and every \(\eta\ge 0\). Thus the scheme is unconditionally von Neumann stable for every \(\theta\ge 0\). The parameter \(\theta\) changes the truncation order; it does not create a stability restriction.

The 2023--2025 papers that say ``analyse truncation error and stability'' want exactly these two computations: a Taylor residual \(\mathcal T\), then either a linear \(g\) with \(\lvert g\rvert^2=A^2+B^2\) or a quadratic in \(g\) with Jury. 2026 Problem~4 asks only for the second half, and only in the linear-\(g\) case.
\end{example}


\begin{examnote}[What to write in fifteen minutes]
\label{exnote:p4-script}
Problem~4 does not ask for truncation error. Write \cref{def:von-neumann,thm:vn-decide}, insert the mode, compute \(|g|^2\), and state the condition on \(r=\tau/h\). Do not say ``implicit, hence stable'': that sentence loses (b).
\end{examnote}

\subsubsection{How one decides ``stable'' or ``unstable''}

The sentence ``hence \(\lvert g_a(\xi)\rvert\le 1\), so (a) is unconditionally von Neumann stable'' is not a new principle invented for this scheme. It is the following test, applied to a number \(g(\xi)\) produced by Fourier substitution. The definitions that feed it are \cref{def:lax-richtmyer} (what stability \emph{means}), \cref{def:von-neumann} (the Fourier form), \cref{prop:vn-plancherel} (why \(\lvert g\rvert\) controls \(\ell_2\)), and \cref{def:cond-uncond} (conditional versus unconditional). What was missing is the decision lemma that turns \(\lvert g\rvert\le 1\) or \(\lvert g\rvert>1\) into a verdict.

\begin{theorem}[Von Neumann decision lemma]
\label{thm:vn-decide}
Let a linear one-step scheme with \(j\)-independent coefficients, on \(\mathbb Z\) or a periodic grid, have amplification factor \(g(\xi)\) obtained by substituting \(U_j^n=g^n e^{\mathrm i j\xi}\).
\begin{enumerate}[label=\textup{(\roman*)}]
    \item \emph{Stable.} If \(\lvert g(\xi)\rvert\le 1\) for every \(\xi\in[0,2\pi]\), then
    \[
    \|U^n\|_2
    \le
    \|U^0\|_2
    \qquad
    (n\tau\le T).
    \]
    This is Lax--Richtmyer stability in \(\ell_2\) with \(C=1\) and \(\alpha=0\) in \cref{def:lax-richtmyer}.
    \item \emph{Unstable.} If there exists even one frequency \(\xi_\star\) with \(\lvert g(\xi_\star)\rvert=1+\delta>1\), the pure mode \(U_j^0=e^{\mathrm i j\xi_\star}\) grows as \(\lvert g\rvert^n=(1+\delta)^{t_n/\tau}\). No constants \(C,\alpha\) independent of the mesh can bound \(\|U^n\|_2/\|U^0\|_2\) as \(h,\tau\to 0\) with \(t_n\le T\) fixed. The scheme is not Lax--Richtmyer stable.
    \item \emph{GKO form.} If the PDE itself may grow, replace (i) by \(\lvert g(\xi)\rvert\le 1+K\tau\). For \(u_t+u_x=0\) one may take \(K=0\), so (i) is the whole test used below.
    \item \emph{Unconditional / conditional.} If (i) holds for every \(r=\tau/h>0\), the scheme is unconditionally von Neumann stable. If (i) holds only for \(r\) in a nonempty proper subset of \((0,\infty)\), it is conditionally stable. If (ii) occurs for every \(r>0\), it is unconditionally unstable.
\end{enumerate}
For a two-step scheme, replace the single \(g\) by both roots of the amplification quadratic and apply \cref{prop:jury}: both roots in the closed unit disk is the analogue of (i).
\end{theorem}

The algebra in (a)--(c) is exclusively (i) versus (ii), and then (iv) to name the \(r\)-dependence. Truncation never enters. Several other tests exist in the same circle; they are not interchangeable. The next table is the method map, in the same layout as the Grade~2 ODE table in Problem~6.

{\footnotesize
\begin{center}
\begin{tabular}{@{}>{\raggedright\arraybackslash}p{2.55cm}
                   >{\raggedright\arraybackslash}p{3.85cm}
                   >{\raggedright\arraybackslash}p{3.85cm}
                   >{\raggedright\arraybackslash}p{3.55cm}@{}}
\toprule
Criterion & What it asks & How one computes it & Compared with the others \\
\midrule
Lax--Richtmyer
(\cref{def:lax-richtmyer})
&
Does \(\|U^n\|_h\le C e^{\alpha t_n}\|U^0\|_h\) hold for \emph{every} grid initial datum, with \(C,\alpha\) independent of \(h,\tau\)?
&
This is the definition, not a computation. One proves it by a test below.
&
The official meaning of ``stable''. Von Neumann, energy, and the maximum principle are machines that establish it in a specific norm.
\\
\addlinespace
von Neumann, \(\lvert g\rvert\le 1\)
(\cref{thm:vn-decide}(i)--(ii))
&
Does any Fourier mode grow? Equivalent to Lax--Richtmyer in \(\ell_2\), for constant-coefficient Cauchy / periodic problems (\cref{prop:vn-plancherel}).
&
Substitute \(U_j^n=g^n e^{\mathrm i j\xi}\), cancel, compute \(\lvert g\rvert^2=A^2+B^2\). Stable iff \(\lvert g(\xi)\rvert\le 1\) for all \(\xi\); unstable if some \(\xi\) has \(\lvert g\rvert>1\).
&
\textbf{The 2026 P4 test.} Does not see boundaries, variable coefficients, or \(\ell_\infty\). Does not see truncation: \(\lvert g\rvert\le 1\) is not consistency.
\\
\addlinespace
GKO form \(\lvert g\rvert\le 1+K\tau\)
&
Same, but the PDE is allowed to grow like \(e^{KT}\).
&
Identical substitution; the threshold is \(e^{K\tau}\) per step rather than \(1\).
&
For advection and heat, \(K=0\), so this collapses to \(\lvert g\rvert\le 1\). Needed when a zeroth-order term \(+cu\) with \(c>0\) is present.
\\
\addlinespace
Jury
(\cref{prop:jury})
&
Both roots of a two-step amplification quadratic lie in the closed unit disk.
&
Form \(Ag^2+Bg+C=0\), check \(p(1)\ge 0\), \(p(-1)\ge 0\), \(\lvert C\rvert\le A\).
&
The von Neumann test when \(U^{n+1}\) depends on \(U^n\) and \(U^{n-1}\) (2024 Spring \(\theta\)-method, 2025 Fall wave / Du~Fort--Frankel). Not used in 2026 P4 (one-step).
\\
\addlinespace
CFL
&
Does the numerical domain of dependence contain the characteristic domain?
&
Compare stencil width to characteristic speed: for explicit \(u_t+u_x=0\), necessary condition \(r\le 1\).
&
Necessary for explicit hyperbolic schemes, never sufficient (\cref{exer:cfl}). Vacuous for implicit schemes (the linear system couples the whole line). Not a replacement for \(\lvert g\rvert\).
\\
\addlinespace
\(\ell_2\) energy identity
&
Is \(\|U^{n+1}\|_2\le\|U^n\|_2\) by an algebraic inner-product identity, without Fourier?
&
Multiply the scheme by \(U_j^{n+1}\) (or \(U_j^{n+1}-U_j^n\)) and sum. Completing the square produces a telescoping \(\ell_2\) budget.
&
Sufficient for Lax--Richtmyer in \(\ell_2\); often equivalent to \(\lvert g\rvert\le 1\) when both apply. Optional on the 2026 P4 script (written out for (a) below). The 2025 Spring Lax--Wendroff question used this instead of Fourier.
\\
\addlinespace
discrete maximum principle
&
Can an interior grid value exceed the boundary / previous-time maximum?
&
Sign pattern of a stencil (nonnegative off-diagonal, strictly diagonally dominant, or mesh P\'eclet). Problem~5.
&
\(\ell_\infty\) stability, including variable coefficients and Dirichlet data. Not equivalent to \(\lvert g\rvert\le 1\). The other CAM machine.
\\
\addlinespace
ODE absolute stability / \(A\)-stability
(\cref{ex:ode-test})
&
For \(y'=\lambda y\), is \(\lvert g(\tau\lambda)\rvert\le 1\)?
&
Write the time stepper on the test equation. Backward Euler: \(g=1/(1-\tau\lambda)\), \(A\)-stable on \(\operatorname{Re}\lambda\le 0\).
&
Explains \emph{why} a given \(g(\xi)\) stays in the unit disk, once the spatial operator has been replaced by its symbol \(\lambda_h(\xi)\). It is the same von Neumann number, not a different criterion.
\\
\bottomrule
\end{tabular}
\end{center}
}

A shorter comparison of the schemes already computed, against the same tests:

{\footnotesize
\begin{center}
\begin{tabular}{@{}>{\raggedright\arraybackslash}p{3.3cm}
                   >{\raggedright\arraybackslash}p{3.7cm}
                   >{\raggedright\arraybackslash}p{3.5cm}
                   >{\raggedright\arraybackslash}p{3.3cm}@{}}
\toprule
Scheme & von Neumann (\cref{thm:vn-decide}) & CFL & Energy / other \\
\midrule
explicit upwind
&
conditionally stable, \(r\le 1\)
(\cref{ex:fourier-subst})
&
same restriction, and here it coincides with \(\lvert g\rvert\le 1\)
&
convex combination of \(1\) and \(e^{-\mathrm i\xi}\)
\\
explicit centered
&
unconditionally unstable: \(\lvert g(\pi/2)\rvert^2=1+r^2>1\) for every \(r>0\)
&
\(r\le 1\) still necessary, and useless
&
consistent \(O(\tau+h^2)\); Lax forbids convergence
\\
BE heat
&
unconditionally stable, \(g=1/(1+\eta)\in(0,1]\)
&
not a hyperbolic obstruction
&
model of ``implicit \(\Rightarrow\) unconditional'', which (b) breaks
\\
exam (a) implicit \(D^-\)
&
unconditional: \(\lvert g_a\rvert\le 1\) for all \(r>0\)
&
vacuous (implicit)
&
energy identity below, same \(\ell_2\) bound
\\
exam (b) implicit \(D^+\)
&
conditional: stable iff \(r\ge 1\) (the reverse CFL)
&
vacuous as an obstruction, but the wrong wind needs heavy implicit damping
&
\(\operatorname{Re}\lambda_h^+\ge 0\); \(A\)-stability is not free
\\
exam (c) implicit \(D^0\)
&
unconditional: \(\lvert g_c\rvert^2=1/(1+r^2\sin^2\xi)\)
&
vacuous
&
purely imaginary symbol, backward Euler on the imaginary axis
\\
\bottomrule
\end{tabular}
\end{center}
}

\subsubsection{Solution of Problem 4}

Throughout, \(r=\tau/h>0\) and \(u_j^n=g^n e^{\mathrm i j\xi}\). The substitution is the one of \cref{ex:fourier-subst}; the only change is that the spatial difference sits at time \(t_{n+1}\), so \(g\) appears on both the time difference and the spatial difference.

\paragraph{(a). Implicit upwind.}

The scheme is
\[
\frac{u_j^{n+1}-u_j^n}{\tau}
+
\frac{u_j^{n+1}-u_{j-1}^{n+1}}{h}
=0.
\]
The table in the Fourier-substitution paragraph gives \(u_j^{n+1}=g\cdot u_j^n\) and \(u_{j-1}^{n+1}=g e^{-\mathrm i\xi}u_j^n\). Inserting and cancelling \(u_j^n=g^n e^{\mathrm i j\xi}\) produces
\[
\frac{g-1}{\tau}
+
\frac{g\bigl(1-e^{-\mathrm i\xi}\bigr)}{h}
=0,
\]
hence
\[
g-1+r g\bigl(1-e^{-\mathrm i\xi}\bigr)=0,
\qquad
g_a(\xi)
=
\frac{1}{1+r\bigl(1-e^{-\mathrm i\xi}\bigr)}.
\]
The denominator is \(1+r-r\cos\xi+\mathrm i r\sin\xi\), so
\begin{align*}
\bigl|1+r(1-e^{-\mathrm i\xi})\bigr|^2
&=
(1+r-r\cos\xi)^2+(r\sin\xi)^2
=
(1+r)^2-2r(1+r)\cos\xi+r^2
\\
&=
1+2r(1+r)(1-\cos\xi)
\ge
1.
\end{align*}
Hence \(\lvert g_a(\xi)\rvert\le 1\) for every \(r>0\) and every \(\xi\). By \cref{thm:vn-decide}(i) and~(iv), scheme (a) is unconditionally von Neumann stable, which is Lax--Richtmyer stability in \(\ell_2\).

The same conclusion has an \(\ell_2\) energy form, which is optional on the script. Multiply the scheme by \(u_j^{n+1}\) and sum in \(j\) (periodic or \(\ell_2(\mathbb Z)\)). The algebraic identities
\[
\bigl(u_j^{n+1}-u_j^n\bigr)u_j^{n+1}
=
\tfrac12\Bigl(\lvert u_j^{n+1}\rvert^2-\lvert u_j^n\rvert^2+\lvert u_j^{n+1}-u_j^n\rvert^2\Bigr)
\]
and
\[
\bigl(u_j^{n+1}-u_{j-1}^{n+1}\bigr)u_j^{n+1}
=
\tfrac12\Bigl(\lvert u_j^{n+1}\rvert^2-\lvert u_{j-1}^{n+1}\rvert^2+\lvert u_j^{n+1}-u_{j-1}^{n+1}\rvert^2\Bigr)
\]
produce, after summing,
\[
\|u^{n+1}\|_2^2
+
\|u^{n+1}-u^n\|_2^2
+
r\|\delta^-u^{n+1}\|_2^2
=
\|u^n\|_2^2,
\]
so \(\|u^{n+1}\|_2\le\|u^n\|_2\). This is the energy row of the criterion table: the same \(\ell_2\) bound as \(\lvert g_a\rvert\le 1\), obtained without Fourier. It is optional on the script.

\paragraph{(b). Implicit downwind.}

Now
\[
\frac{u_j^{n+1}-u_j^n}{\tau}
+
\frac{u_{j+1}^{n+1}-u_j^{n+1}}{h}
=0,
\]
hence
\[
g_b(\xi)
=
\frac{1}{1+r\bigl(e^{\mathrm i\xi}-1\bigr)}.
\]
The denominator is \(1-r+r\cos\xi+\mathrm i r\sin\xi\), and
\begin{align*}
\bigl|1+r(e^{\mathrm i\xi}-1)\bigr|^2
&=
(1-r+r\cos\xi)^2+(r\sin\xi)^2
=
1-2r(1-r)(1-\cos\xi).
\end{align*}
The identity \(1-\cos\xi\ge 0\) decides the sign of the correction.

\begin{itemize}
    \item If \(0<r<1\), then \(1-r>0\), so the right-hand side is strictly less than \(1\) for \(\xi\not\equiv 0\bmod 2\pi\). Thus \(|g_b|>1\) on a set of frequencies, and the scheme is unstable.
    \item If \(r=1\), the right-hand side equals \(1\), so \(|g_b|\equiv 1\).
    \item If \(r>1\), then \(1-r<0\), so the correction is nonnegative and the right-hand side is at least \(1\). Thus \(|g_b|\le 1\).
\end{itemize}

Scheme (b) is von Neumann stable if and only if \(r\ge 1\). In the language of \cref{thm:vn-decide}: (ii) when \(0<r<1\), (i) when \(r\ge 1\), hence (iv) conditional, with the reverse of the explicit CFL restriction. The spatial difference looks the wrong way relative to the characteristics \(x-t=\mathrm{const}\); when \(r\) is small the scheme behaves like explicit downwind (unstable), and only a sufficiently heavy implicit step damps that instability.

\paragraph{(c). Implicit centered.}

The scheme is
\[
\frac{u_j^{n+1}-u_j^n}{\tau}
+
\frac{u_{j+1}^{n+1}-u_{j-1}^{n+1}}{2h}
=0,
\]
hence
\[
g_c(\xi)
=
\frac{1}{1+\mathrm i r\sin\xi},
\qquad
\bigl|g_c(\xi)\bigr|^2
=
\frac{1}{1+r^2\sin^2\xi}
\le 1.
\]
Unconditionally stable by \cref{thm:vn-decide}(i) and~(iv), and strictly dissipative except at \(\xi\in\pi\mathbb Z\). Equivalently, the centered symbol of \(\partial_x\) is \(\mathrm i(\sin\xi)/h\), purely imaginary, and backward Euler is \(A\)-stable on the imaginary axis.

\begin{summary}[Problem 4]
\label{sum:p4}
Let \(r=\tau/h\).
{\footnotesize
\begin{center}
\begin{tabular}{@{}>{\raggedright\arraybackslash}p{4.2cm}
                   >{\raggedright\arraybackslash}p{6.4cm}
                   >{\raggedright\arraybackslash}p{3.6cm}@{}}
\toprule
Scheme & Amplification & Stability \\
\midrule
(a) implicit \(D^-\) (upwind) & \(g=1/\bigl(1+r(1-e^{-\mathrm i\xi})\bigr)\) & unconditional \\
(b) implicit \(D^+\) (downwind) & \(g=1/\bigl(1+r(e^{\mathrm i\xi}-1)\bigr)\) & stable iff \(r\ge 1\) \\
(c) implicit \(D^0\) (centered) & \(g=1/(1+\mathrm i r\sin\xi)\) & unconditional \\
\bottomrule
\end{tabular}
\end{center}
}
\end{summary}

\subsubsection{Syllabus objects that (a)--(c) do not compute}

Problem~4 asks only for von Neumann stability of three implicit advection schemes. The CAM PDE-numerics line is larger: finite-difference methods, stability and convergence, Lax equivalence; also finite elements and spectral methods \cite{gko95,leveque07,strikwerda04,brennerscott10,trefethen00}. The objects below are the ones named in the preliminaries that the three amplification factors do not force you to write. They are the usual next questions on a later paper (truncation, CFL, Du~Fort--Frankel, a weak-form comparison).

\paragraph{Finite-difference methods.}

A finite-difference method is a recipe, not a single scheme. Put the unknown on a grid (\cref{not:grid-p4}); replace each derivative by a difference quotient (\cref{def:quotients}); decide the time level of the spatial stencil (explicit, implicit, or multistep). The output is a linear marching rule \(U^{n+1}=S_{h,\tau}U^n\), or, for a stationary elliptic problem, a sparse linear system \(A_h U=F\) (Problem~5). Every Qiuzhen difference question is an instance of this recipe with a named stencil: upwind / centered / Lax--Friedrichs / Lax--Wendroff / Du~Fort--Frankel / five-point Laplacian.

The local accuracy of the recipe is the truncation error of \cref{def:truncation}: feed the sampled PDE solution to \(\mathcal L_{h,\tau}\) and Taylor-expand. The residual \(\mathcal T\) is computed in \cref{ex:trunc-a,ex:trunc-c,ex:trunc-sine}; it is not the global error \(U-u\).

\begin{exercise}[Truncation of scheme (b), as a check]
\label{exer:trunc-a}
Using only \cref{def:quotients} and a Taylor expansion at \((x_j,t_{n+1})\), recover
\[
\bigl(\mathcal T^{(b)}\bigr)_j^{n+1}
=
-\tfrac12\tau\,u_{tt}
+
\tfrac12 h\,u_{xx}
+
O(\tau^2+h^2)
\]
on solutions of \(u_t+u_x=0\). For \(u=\sin(x-t)\), show that this is \(\tfrac12(\tau-h)\sin(x_j-t_{n+1})+O(\tau^2+h^2)\).
\end{exercise}

\begin{proof}[Solution]
Backward Euler at \(t_{n+1}\) is again \(u_t-\tfrac12\tau u_{tt}+O(\tau^2)\). The spatial piece is \(D^+u=u_x+\tfrac12 h u_{xx}+O(h^2)\). Adding and using \(u_t+u_x=0\) leaves the displayed residual. On the travelling wave, \(u_{tt}=u_{xx}=-\sin(x-t)\), so \(-\tfrac12\tau u_{tt}+\tfrac12 h u_{xx}=\tfrac12(\tau-h)\sin(x_j-t_{n+1})\).
\end{proof}

\paragraph{Convergence.}

Stability controls the homogeneous scheme. Convergence compares the computed grid function to the PDE solution, in a discrete norm, in the joint limit \(h,\tau\to 0\) with time held in a fixed interval.

\begin{definition}[Convergence]
\label{def:convergence}
Let \(u\) be a smooth solution of a linear IVP, and let \(U^n\) be the grid function produced by a scheme from the sampled initial datum \(U^0_j=u(x_j,0)\) (plus consistent boundary data, if any). Write \(u^n_j\coloneqq u(x_j,t_n)\) and \(e^n\coloneqq U^n-u^n\). The scheme is \emph{convergent} of order \((p,q)\) in the discrete norm \(\|\cdot\|_h\) if, whenever \(t_n\le T\) and \((h,\tau)\to(0,0)\) through any mesh restriction under which the scheme is stable,
\[
\|e^n\|_h
\le
C_T\bigl(\tau^p+h^q\bigr)
\to 0.
\]
Consistency is \(\|\mathcal T^n\|_h\to 0\); convergence is \(\|e^n\|_h\to 0\). The two limits are not the same: \(\mathcal T\) is the residual of the exact solution in the stencil, while \(e\) is the error of the computed solution.
\end{definition}

The error itself satisfies a forced scheme. If \(U^{n+1}=S_{h,\tau}U^n\) and the exact samples obey \(u^{n+1}=S_{h,\tau}u^n+\tau\mathcal T^{n+1}\) (the convention that \(\mathcal T\) is a residual per unit time), then
\[
e^{n+1}
=
S_{h,\tau}e^n
+
\tau\mathcal T^{n+1}.
\]
Unrolling,
\begin{equation}
\label{eq:discrete-duhamel}
e^n
=
S_{h,\tau}^n e^0
+
\tau\sum_{k=1}^{n} S_{h,\tau}^{n-k}\mathcal T^{k}.
\end{equation}
A uniform bound on \(\|S_{h,\tau}^m\|\) for \(m\tau\le T\) turns a small truncation into a small global error. That bound is precisely Lax--Richtmyer stability (\cref{def:lax-richtmyer}). Without it, \eqref{eq:discrete-duhamel} can amplify a consistent residual (or even \(e^0=0\) plus roundoff) into an \(O(1)\) or exponentially large grid function as the mesh is refined.

\paragraph{Lax equivalence.}

This is the official theorem named on the CAM syllabus \cite{laxrichtmyer56}. Textbook statements are in LeVeque \cite{leveque07}, Strikwerda \cite{strikwerda04}, and GKO \cite{gko95}. Qiuzhen almost never asks you to quote it. It is the reason one is allowed to spend the whole of Problem~4 on \(|g|\le 1\).

\begin{theorem}[Lax equivalence]
\label{thm:lax-eq}
Consider a linear well-posed evolutionary IVP and a consistent finite-difference approximation \(U^{n+1}=S_{h,\tau}U^n\), in a discrete norm in which the PDE itself is well-posed. The scheme is convergent if and only if it is stable.
\end{theorem}

\begin{proof}[Proof idea]
\emph{Stability \(+\) consistency \(\Rightarrow\) convergence.} Bound \eqref{eq:discrete-duhamel}. Stability supplies \(\|S_{h,\tau}^m\|_h\le C e^{\alpha T}\) for \(m\tau\le T\). If \(\|e^0\|_h=0\) (exact sampling) and \(\|\mathcal T^k\|_h\le\varepsilon(h,\tau)\to 0\), then
\[
\|e^n\|_h
\le
C T e^{\alpha T}\,\varepsilon(h,\tau)
\to 0.
\]
If the truncation is \(O(\tau^p+h^q)\), the same estimate gives the rate in \cref{def:convergence}. The constant \(C e^{\alpha T}\) is allowed to grow in \(T\), matching \cref{def:lax-richtmyer} and the GKO form \(|g|\le 1+K\tau\) of von Neumann \cite{gko95}: a factor \(1+K\tau\) per step is \(e^{KT}\) at time \(T\), which is PDE growth, not mesh blow-up.

\emph{Convergence \(\Rightarrow\) stability.} If the family \(\{S_{h,\tau}^n:n\tau\le T\}\) is not uniformly bounded, Banach--Steinhaus produces initial grid functions of unit norm whose images become arbitrarily large as \(h,\tau\to 0\). Those data may be realized (or approximated) by smooth PDE initial values, and the computed solution cannot tend to the well-posed PDE solution. Thus a convergent consistent scheme is stable.

\emph{Convergence \(\Rightarrow\) consistency} is the remaining corner: if a stable scheme converges for all smooth solutions of a well-posed linear IVP, the truncation must vanish (apply the scheme to the exact solution and pass to the limit). In exam writing one checks consistency by Taylor expansion and stability by von Neumann or a maximum principle; Lax then gives convergence for free.
\end{proof}

The hypotheses are not decoration. Linearity is used to write a marching operator \(S_{h,\tau}\) independent of the solution; well-posedness of the PDE is used so that the exact solution is a legitimate comparison object; the same discrete norm must be used for stability and for convergence. Von Neumann is the \(\ell_2\) machine. A discrete maximum principle is the \(\ell_\infty\) machine (Problem~5). An \(\ell_2\)-stable scheme need not be \(\ell_\infty\)-stable.

Two standard ways the theorem can fail if a hypothesis is dropped: consistent but unstable, and stable but inconsistent. Both occur as named schemes.

\begin{example}[Consistent and unstable, hence not convergent]
\label{ex:lax-centered}
The explicit centered scheme of \cref{ex:centered-unstable} is consistent of order \(O(\tau+h^2)\) and has \(|g(\xi)|^2=1+r^2\sin^2\xi\). On a mode \(\xi=\pi/2\),
\[
\bigl|g\bigr|^n
=
\bigl(1+r^2\bigr)^{n/2}
=
\bigl(1+r^2\bigr)^{t_n/(2\tau)}.
\]
With \(r=\tau/h\) held fixed, this is \(e^{c t_n/h}\) as \(h\to 0\), which is not a bound of the form \(C e^{\alpha T}\). The scheme is unstable, so Lax equivalence forbids convergence: high-frequency components of a smooth initial datum are amplified without bound under mesh refinement, even though each individual truncation residual tends to zero.
\end{example}

\begin{exercise}[CFL is necessary, not sufficient]
\label{exer:cfl}
For \(u_t+u_x=0\), the characteristic through \((x_j,t_{n+1})\) hits \(t_n\) at \(x_j-\tau\). An explicit three-point stencil at time \(t_n\) has numerical domain of dependence \(\{x_{j-1},x_j,x_{j+1}\}\), hence numerical speed \(h/\tau=1/r\). Courant--Friedrichs--Lewy says that a necessary condition for convergence (hence, by \cref{thm:lax-eq}, for stability) is that the numerical domain contain the characteristic domain, i.e. \(r\le 1\). Show that explicit centered advection satisfies this CFL restriction whenever \(r\le 1\), and is nevertheless unconditionally unstable. What, then, is CFL doing?
\end{exercise}

\begin{proof}[Solution]
The stencil of \cref{ex:centered-unstable} uses both neighbours, so the CFL inequality is \(r\le 1\). Von Neumann nevertheless gives \(|g|^2=1+r^2\sin^2\xi>1\) for some \(\xi\) as soon as \(r>0\). CFL is a domain-of-dependence obstruction: if \(r>1\), information the PDE needs is invisible to the stencil, so the scheme cannot converge. It is not a stability certificate. Explicit upwind (\cref{ex:explicit-upwind}) is the scheme for which CFL \(r\le 1\) happens to coincide with \(|g|\le 1\); explicit centered is the scheme that shows the two conditions are distinct. Implicit schemes have an infinite numerical domain of dependence (a linear system couples the whole line), so CFL is automatically satisfied; that is why implicit downwind (b) can be stable for \(r\ge 1\) even though it looks the wrong way.
\end{proof}

\begin{exercise}[Du~Fort--Frankel: stable but inconsistent unless the mesh is parabolic]
\label{exer:dufort}
The Du~Fort--Frankel scheme for \(u_t=u_{xx}\) is the leapfrog time difference with the centered Laplacian, except that \(2u_j^n\) is replaced by \(u_j^{n+1}+u_j^{n-1}\):
\[
\frac{u_j^{n+1}-u_j^{n-1}}{2\tau}
=
\frac{u_{j+1}^n-u_j^{n+1}-u_j^{n-1}+u_{j-1}^n}{h^2}.
\]
Let \(\mu=\tau/h^2\).
\begin{enumerate}[label=\textup{(\alph*)}]
    \item Show that the local truncation error, expanded at \((x_j,t_n)\), contains a leading term \((\tau/h)^2 u_{tt}\). Conclude that the scheme is consistent with \(u_t=u_{xx}\) if and only if \(\tau/h\to 0\) as \((h,\tau)\to(0,0)\). In particular it is consistent under parabolic refinement \(\tau=O(h^2)\), and inconsistent under hyperbolic refinement \(\tau=O(h)\).
    \item Substitute \(u_j^n=g^n e^{\mathrm i j\xi}\) and obtain the quadratic
    \[
    (1+2\mu)\,g^2
    -
    4\mu\cos\xi\cdot g
    -
    (1-2\mu)
    =0.
    \]
    (The scheme is unconditionally von Neumann stable: both roots satisfy \(|g|\le 1\) for every \(\mu>0\) and every \(\xi\). The 2025 Fall paper asks exactly this pair, truncation plus \(|g|\).)
    \item Why does Lax equivalence not yield convergence for arbitrary \(\tau,h\to 0\)?
\end{enumerate}
\end{exercise}

\begin{proof}[Solution]
(a) Write the right-hand side as the usual second difference minus a second time difference:
\[
u_{j+1}^n-u_j^{n+1}-u_j^{n-1}+u_{j-1}^n
=
\bigl(u_{j+1}^n-2u_j^n+u_{j-1}^n\bigr)
-
\bigl(u_j^{n+1}-2u_j^n+u_j^{n-1}\bigr).
\]
Taylor at \((x_j,t_n)\) gives \(h^2 u_{xx}+O(h^4)\) for the first parenthesis and \(\tau^2 u_{tt}+O(\tau^4)\) for the second. Dividing by \(h^2\) produces \(u_{xx}-(\tau/h)^2 u_{tt}+O(h^2+\tau^4/h^2)\). The left-hand side is \(u_t+O(\tau^2)\). The residual of \(u_t-u_{xx}\) is therefore \((\tau/h)^2 u_{tt}\) plus higher-order terms. This tends to zero for all smooth \(u\) if and only if \(\tau/h\to 0\).

(b) Insert the mode and cancel \(g^{n-1}e^{\mathrm i j\xi}\):
\[
\frac{g^2-1}{2\tau}
=
\frac{2g\cos\xi-g^2-1}{h^2}.
\]
Clearing \(2\tau\) and writing \(\mu=\tau/h^2\) produces the displayed quadratic. (Unconditional stability is a computation with the discriminant \(1-4\mu^2\sin^2\xi\): when it is nonnegative both real roots lie in \([-1,1]\); when it is negative the roots are complex conjugates of modulus \(\sqrt{(2\mu-1)/(2\mu+1)}<1\) for \(\mu>1/2\), and a direct check covers \(\mu\le 1/2\).)

(c) The scheme is not consistent unless \(\tau/h\to 0\). Lax equivalence requires consistency. If one refines with \(\tau/h\) bounded away from zero, the scheme is consistent with a different PDE,
\[
u_t
+
\bigl(\tfrac{\tau}{h}\bigr)^2 u_{tt}
=
u_{xx},
\]
and converges (if at all) to solutions of that hyperbolic regularization, not of the heat equation. This is the standard warning that ``\(|g|\le 1\) for all \(\mu\)'' is not a convergence theorem.
\end{proof}

\begin{exercise}[FTCS heat, the model parabolic CFL]
\label{exer:ftcs}
The forward-time centered-space scheme for \(u_t=u_{xx}\) is
\[
\frac{u_j^{n+1}-u_j^n}{\tau}
=
\frac{u_{j+1}^n-2u_j^n+u_{j-1}^n}{h^2}.
\]
Set \(\mu=\tau/h^2\). Compute \(g(\xi)\) and prove von Neumann stability if and only if \(0<\mu\le 1/2\). Combined with the truncation \(O(\tau+h^2)\), Lax equivalence yields convergence of order \(O(\tau+h^2)\) under that mesh restriction. (Compare \cref{ex:be-heat}: backward Euler removes the restriction. Compare \cref{ex:theta-heat}: the \(\theta\)-method interpolates between the two.)
\end{exercise}

\begin{proof}[Solution]
Fourier substitution gives \(g=1-4\mu\sin^2(\xi/2)\). This is real, and \(g\le 1\) automatically. The remaining inequality \(g\ge -1\) is \(2-4\mu\sin^2(\xi/2)\ge 0\) for all \(\xi\), hence \(\mu\le 1/2\). At \(\mu=1/2\) one has \(g=\cos\xi\), so \(|g|\le 1\) still. Consistency is the standard Taylor computation: the time difference is \(u_t+\tfrac12\tau u_{tt}+O(\tau^2)\) and the second difference is \(u_{xx}+\tfrac{1}{12}h^2 u_{xxxx}+O(h^4)\).
\end{proof}

\paragraph{What the three books are for.}

LeVeque \cite{leveque07} is the first reading: difference quotients (Ch.~1), backward Euler and absolute stability for ODEs (Ch.~2), advection / CFL / von Neumann / Lax (Chs.~8--10). Strikwerda \cite{strikwerda04} is the elliptic companion (discrete maximum principle, Ch.~2) and a second Fourier source (Chs.~5--6). Gustafsson--Kreiss--Oliger \cite{gko95} is the syllabus item: Fourier multipliers, the \(1+K\tau\) form of von Neumann, and dissipativity. Problem~4 lives in the LeVeque/GKO hyperbolic chapter; Problem~5 lives in Strikwerda's elliptic chapter. None of the three is a finite-element book.

\paragraph{Finite elements, in one page.}

The same elliptic operator \(Lu=u_{xx}+u_{yy}+\cdots\) can be discretized without writing a difference quotient. Multiply by a test function, integrate by parts, and restrict the weak form to a finite-dimensional subspace. That is the language of Brenner--Scott \cite{brennerscott10}, and it is not what ``centered finite difference scheme'' means in Problem~5.

\begin{example}[1D Poisson, Galerkin, C\'ea]
\label{ex:fem-poisson}
On \((0,1)\) with \(u(0)=u(1)=0\), the problem \(-u''=f\) is equivalent to: find \(u\in V\coloneqq H_0^1(0,1)\) such that
\[
a(u,v)
\coloneqq
\int_0^1 u'v'\,\mathrm{d}x
=
\int_0^1 fv\,\mathrm{d}x
\qquad
\text{for all }v\in V.
\]
Let \(V_h\subset V\) be the space of continuous piecewise-linear functions on a uniform mesh of size \(h\), vanishing at the endpoints (hat-function basis). The Galerkin method seeks \(u_h\in V_h\) with \(a(u_h,v_h)=(f,v_h)\) for all \(v_h\in V_h\). Subtracting, \(a(u-u_h,v_h)=0\): the error is \(a\)-orthogonal to \(V_h\). Coercivity and continuity of \(a\) on \(V\) (here \(a(w,w)=\lvert w\rvert_{H^1}^2\)) give C\'ea's lemma \cite{brennerscott10}:
\[
\|u-u_h\|_V
\le
C\inf_{v_h\in V_h}\|u-v_h\|_V.
\]
The infimum is interpolation error. For \(u\in H^2\cap H_0^1\), the nodal interpolant \(I_h u\) satisfies \(\lvert u-I_h u\rvert_{H^1}=O(h)\) and \(\|u-I_h u\|_{L^2}=O(h^2)\), hence
\[
\lvert u-u_h\rvert_{H^1}=O(h),
\qquad
\|u-u_h\|_{L^2}=O(h^2)
\]
after an Aubin--Nitsche lift. The stiffness matrix of the hat basis is the tridiagonal matrix of the three-point Laplacian, up to the factor \(h\): on this one-dimensional uniform mesh, linear FEM and the standard second-difference scheme produce the same algebraic system. In several space dimensions, on unstructured triangulations, or with variable coefficients in divergence form, they do not. Stability is then coercivity of \(a\) on \(V_h\) (or an inf-sup condition), not an amplification factor \(g(\xi)\).
\end{example}

\begin{exercise}[Where FEM would change the exam]
\label{exer:fem-not-p4}
Suppose Problem~5 had said ``piecewise-linear finite elements on a uniform triangulation'' instead of ``centered finite difference scheme''. Which of the following would still be a correct first sentence, and which would be a category error: (i)~write a five-point stencil and check mesh P\'eclet conditions; (ii)~write the weak form, choose \(V_h\), and invoke a discrete maximum principle or C\'ea; (iii)~insert \(g^n e^{\mathrm i j\xi}\) and compute \(|g|\)?
\end{exercise}

\begin{proof}[Solution]
(i) is the finite-difference first sentence, hence a category error for an FEM prompt (the five-point stencil is a particular nodal arrangement of second differences, not a Galerkin statement). (ii) is the correct first sentence. (iii) is von Neumann, which requires a constant-coefficient marching scheme on the whole line or a torus; it is the language of Problem~4, not of a stationary elliptic FEM problem on the square, and not of Problem~5 even in the finite-difference formulation (variable \(d,e,f\), Dirichlet boundary). A discrete maximum principle can still be available for some FEM schemes (acute triangulations, lumped mass), but it is proved from the sign pattern of the stiffness matrix, not from \(g(\xi)\).
\end{proof}

\paragraph{Spectral methods.}

The remaining syllabus discretization is spectral \cite{trefethen00}. Replace the unknown by a global interpolant (trigonometric on a periodic interval, Chebyshev on a bounded interval) and collocate the PDE at the interpolation nodes, or equivalently apply a differentiation matrix to the vector of nodal values. For analytic data the error is typically \(O(\rho^{-N})\) in the number \(N\) of modes, against the algebraic \(O(h^p)\) of a \(p\)-th order difference scheme. Stability of a Fourier spectral method for constant-coefficient periodic evolution is again a statement about symbols: the collocation operator is a Fourier multiplier, and one checks that the resulting ODE \(\hat u_t=\lambda_k\hat u\) is integrated stably. That is the same von Neumann machine with \(\xi\) restricted to the retained modes. Chebyshev collocation on an interval is not a Fourier multiplier; the analogue of CFL becomes a restriction \(\tau=O(N^{-2})\) for explicit time-stepping, because the largest differentiation-matrix eigenvalue is \(O(N^2)\). Qiuzhen has not asked a spectral PDE scheme in the 2023--2026 CAM papers; the nearby object is a differentiation matrix / interpolation remainder, which lives in the approximation slot, not here.

\begin{examnote}[What of this belongs on a script]
\label{exnote:p4-syllabus}
For 2026 Problem~4, none of it. Write \(|g|^2\) as in \cref{exnote:p4-script}. If a later paper says ``truncation error and stability'', add \cref{ex:trunc-a,exer:ftcs}. If it says ``Du~Fort--Frankel'', add \cref{exer:dufort}. If it says ``finite element'', start from \cref{ex:fem-poisson}, not from \(D^\pm\). Lax equivalence is the warrant that the stability computation is the whole story; it is not itself an answer.
\end{examnote}


\subsection{Problem 5. [15 points]}
Write and prove the maximum principle of the centered finite difference scheme for discretizing the equation

\[
u_{xx} + u_{yy} + d u_{x} + e u_{y} + f u = 0, \quad f < 0,
\]

under some suitable assumptions. Apply the maximum principle to show the stability of the backward Euler scheme for the parabolic equation

\[
u_{t} - (u_{xx} + u_{yy} + d u_{x} + e u_{y} + f u) = 0, \quad (x, y) \in (0, 1) \times (0, 1),
\]

with homogeneous Dirichlet boundary conditions.

\setcounter{section}{5}
\setcounter{theorem}{0}
\setcounter{definition}{0}
\setcounter{remark}{0}
\setcounter{note}{0}

\subsubsection{Preliminaries: how one discretizes a linear PDE, and why a maximum principle is a stability theorem}

Problem~4 tested \(\ell_2\) stability of a marching scheme by Fourier analysis on the whole line. Problem~5 tests \(\ell_\infty\) stability of an elliptic stencil, then of a parabolic time-stepper, on a bounded square with Dirichlet data. The Fourier machine does not apply: the coefficients \(d,e,f\) need not be constant, and the boundary is essential. The replacement machine is a discrete maximum principle for a monotone difference operator.

This is the same 15-point question as 2024 Fall Problem~5, with one extra sentence: after the elliptic statement, apply it to backward Euler. The 1D ancestor is 2024 Spring Problem~5 (a three-point operator \(Lu_j=-(a_j u_{j-1}-b_j u_j+c_j u_{j+1})+q_j u_j\)). If those two papers are worked, 2026 Problem~5 is a transcription.

Finite elements \cite{brennerscott10} would discretize the same elliptic operator weakly, on a triangulation, with a stiffness matrix. That is not what ``centered finite difference scheme'' means here.

\paragraph{The discretization recipe (write the scheme before proving anything).}

The target class is a linear second-order operator in two space variables,
\[
Lu
\coloneqq
u_{xx}+u_{yy}+d\,u_x+e\,u_y+f u.
\]
The elliptic exam is \(Lu=0\) (with \(f<0\)). The parabolic exam is
\[
u_t-Lu=0
\quad\text{on }(0,1)^2,
\qquad
u=0\text{ on the spatial boundary}.
\]
A finite-difference scheme is produced as follows.

\begin{enumerate}[label=\textup{\arabic*.}]
    \item \textbf{Grid.} Choose mesh sizes \(h_x=1/M\) and \(h_y=1/N\), nodes \(x_i=ih_x\), \(y_j=jh_y\), and unknowns \(U_{ij}\approx u(x_i,y_j)\). Interior nodes are \(1\le i\le M-1\), \(1\le j\le N-1\); boundary nodes carry Dirichlet data (here \(0\)).
    \item \textbf{Replace each derivative by a difference quotient.} The complete catalogue --- forward, backward, and centered, for first and second derivatives --- is \cref{def:quotients-1d,def:quotients-2d,not:time-quotients}. ``Centered'' in the exam statement means the second-order column of that catalogue: \(\delta_x,\delta_y\) for first derivatives and \(\delta_{xx},\delta_{yy}\) for second derivatives, not a one-sided stencil and not \((D^0)^2\).
    \[
    \delta_x U_{ij}
    =
    \frac{U_{i+1,j}-U_{i-1,j}}{2h_x},
    \qquad
    \delta_{xx}U_{ij}
    =
    \frac{U_{i+1,j}-2U_{ij}+U_{i-1,j}}{h_x^2},
    \]
    and likewise \(\delta_y\), \(\delta_{yy}\).
    \item \textbf{Write the five-point stencil} by collecting coefficients of the four neighbours and the center.
    \item \textbf{For a parabolic equation, discretize time separately.} Forward Euler uses \(L_h U^n\); backward Euler uses \(L_h U^{n+1}\); Crank--Nicolson averages the two. Problem~5 asks for backward Euler.
    \item \textbf{Prove a comparison principle} for the spatial operator (maximum principle), then feed it into the time-stepper.
\end{enumerate}

The same five steps, with different stencils, generate every Qiuzhen finite-difference question: Lax--Friedrichs and Lax--Wendroff are step~2 with a specific averaging; Du Fort--Frankel replaces \(2u_j^n\) in the Laplacian by \(u_j^{n+1}+u_j^{n-1}\); the wave scheme of 2025 Fall P5 is a three-level average of \(\delta_{xx}\).

\paragraph{The difference-quotient catalogue (this problem uses only the centered column).}

Problem~4 already introduced the three first-derivative quotients \(D^\pm,D^0\) on a 1D grid (\cref{def:quotients}). Elliptic operators also need second derivatives, and the exam adjective ``centered'' is a choice among three columns. The next two definitions are the complete spatial catalogue; the time quotients are \cref{not:time-quotients}. None of the one-sided second differences appear in 2026 Problem~5, but they are the objects one would write if the exam had said ``forward'' or ``backward'' instead of ``centered''.

\begin{definition}[Forward, backward, and centered quotients, one variable]
\label{def:quotients-1d}
Let \(v_j=v(x_j)\) be a grid function on a uniform mesh of size \(h>0\). The three standard approximations of \(\partial_x\) are those of \cref{def:quotients}:
\begin{align*}
D^+v_j
&\coloneqq
\frac{v_{j+1}-v_j}{h},
\qquad
D^-v_j
\coloneqq
\frac{v_j-v_{j-1}}{h},
\qquad
D^0v_j
\coloneqq
\frac{v_{j+1}-v_{j-1}}{2h}.
\end{align*}
The three standard approximations of \(\partial_{xx}\) are the corresponding second differences:
\begin{align*}
(D^+)^2 v_j
&\coloneqq
\frac{v_{j+2}-2v_{j+1}+v_j}{h^2},
\\
(D^-)^2 v_j
&\coloneqq
\frac{v_j-2v_{j-1}+v_{j-2}}{h^2},
\\
\delta_{xx}v_j
&\coloneqq
\frac{v_{j+1}-2v_j+v_{j-1}}{h^2}
=
D^+D^-v_j
=
D^-D^+v_j.
\end{align*}
If \(v\) is smooth, Taylor expansion at \(x_j\) gives
\begin{align*}
D^+v
&=
v'+\tfrac12 h v''+O(h^2),
&
D^-v
&=
v'-\tfrac12 h v''+O(h^2),
&
D^0v
&=
v'+\tfrac16 h^2 v'''+O(h^4),
\\
(D^+)^2 v
&=
v''+h v'''+O(h^2),
&
(D^-)^2 v
&=
v''-h v'''+O(h^2),
&
\delta_{xx}v
&=
v''+\tfrac{1}{12}h^2 v^{(4)}+O(h^4).
\end{align*}
Thus \(D^\pm\) and \((D^\pm)^2\) are first-order accurate, while \(D^0\) and \(\delta_{xx}\) are second-order accurate. Two identifications that are easy to get wrong:
\begin{enumerate}[label=\textup{(\roman*)}]
    \item \(\delta_{xx}=D^+D^-=D^-D^+\), a compact three-point stencil. It is \emph{not} \((D^0)^2\):
    \[
    (D^0)^2 v_j
    =
    \frac{v_{j+2}-2v_j+v_{j-2}}{4h^2}
    =
    v''+\tfrac13 h^2 v^{(4)}+O(h^4),
    \]
    which is still second-order but uses a wider five-point molecule in 1D (nine-point in 2D) and is not the Laplacian of the exam.
    \item ``Second-order'' is overloaded. It may mean the \emph{derivative} being approximated (\(\partial_{xx}\)), or the \emph{accuracy} \(O(h^2)\). Centered \(\delta_{xx}\) is second-order in both senses; \((D^+)^2\) is a second derivative of only first-order accuracy.
\end{enumerate}
\end{definition}

\begin{definition}[The same operators in two variables]
\label{def:quotients-2d}
On the grid \((x_i,y_j)\) of the recipe, apply \cref{def:quotients-1d} independently in each coordinate, with steps \(h_x\) and \(h_y\). Write \(D_x^\pm\), \(D_x^0\), \(\delta_{xx}\) for the \(x\)-operators (the last already appeared as \(\delta_{xx}\) in the recipe), and \(D_y^\pm\), \(D_y^0\), \(\delta_{yy}\) analogously. The exam's centered scheme is the choice
\[
\partial_x
\;\mapsto\;
D_x^0=\delta_x,
\qquad
\partial_{xx}
\;\mapsto\;
\delta_{xx},
\qquad
\partial_y
\;\mapsto\;
\delta_y,
\qquad
\partial_{yy}
\;\mapsto\;
\delta_{yy}.
\]
A mixed derivative \(\partial_{xy}\) is not in this PDE, so the exam stencil never uses it. If it did appear, the catalogue product would be any \(D_x^\alpha D_y^\beta\); the centered choice is the four-corner stencil
\[
(\delta_x\delta_y)v_{i,j}
=
\frac{v_{i+1,j+1}-v_{i+1,j-1}-v_{i-1,j+1}+v_{i-1,j-1}}{4h_x h_y}.
\]
The other order \(\partial_{yx}\) is the same object for a \(C^{2}\) solution (Schwarz: \(\partial_x\partial_y=\partial_y\partial_x\)). On the product grid the difference operators in \(x\) and in \(y\) commute as well, so
\[
\delta_x\delta_y
=
\delta_y\delta_x,
\qquad
D_x^\alpha D_y^\beta
=
D_y^\beta D_x^\alpha,
\]
and there is no second mixed stencil to invent. One-sided mixed products such as \(D_x^+D_y^-\) remain well-defined; they are simply a different (first-order) approximation of the same \(\partial_{xy}=\partial_{yx}\).
\end{definition}

{\footnotesize
\begin{center}
\begin{tabular}{@{}>{\raggedright\arraybackslash}p{2.35cm}
                   >{\raggedright\arraybackslash}p{4.15cm}
                   >{\raggedright\arraybackslash}p{2.35cm}
                   >{\raggedright\arraybackslash}p{4.95cm}@{}}
\toprule
Operator & Formula at a node & Accuracy & Where it is used \\
\midrule
\(D^+\) (forward)
&
\((v_{j+1}-v_j)/h\)
&
\(O(h)\) for \(\partial_x\)
&
P4 scheme (b), downwind if the wind is to the right; explicit Euler in time
\\
\(D^-\) (backward)
&
\((v_j-v_{j-1})/h\)
&
\(O(h)\) for \(\partial_x\)
&
P4 scheme (a), upwind if the wind is to the right; backward Euler in time
\\
\(D^0\) (centered)
&
\((v_{j+1}-v_{j-1})/(2h)\)
&
\(O(h^2)\) for \(\partial_x\)
&
P4 scheme (c); P5 convection \(\delta_x,\delta_y\). Needs a mesh P\'eclet restriction for monotonicity
\\
\((D^+)^2\)
&
\((v_{j+2}-2v_{j+1}+v_j)/h^2\)
&
\(O(h)\) for \(\partial_{xx}\)
&
one-sided second derivative; not the five-point Laplacian
\\
\((D^-)^2\)
&
\((v_j-2v_{j-1}+v_{j-2})/h^2\)
&
\(O(h)\) for \(\partial_{xx}\)
&
the other one-sided second derivative
\\
\(\delta_{xx}=D^+D^-\)
&
\((v_{j+1}-2v_j+v_{j-1})/h^2\)
&
\(O(h^2)\) for \(\partial_{xx}\)
&
\textbf{this problem}, and every Qiuzhen five-point / three-point Laplacian
\\
\bottomrule
\end{tabular}
\end{center}
}

\begin{remark}[Upwind as the other spatial column, not used here]
\label{rem:upwind-p5}
If the mesh P\'eclet condition \(\lvert d\rvert h_x\le 2\) of \cref{ass:peclet} fails, centered \(\delta_x\) makes a neighbour coefficient negative and the maximum principle stops. The standard repair is to discretize the convection \(d u_x\) by a one-sided first derivative pointing \emph{against} the transport velocity of \(u_t=d u_x+\cdots\), i.e.\ \(D_x^+\) when \(d\ge 0\) and \(D_x^-\) when \(d<0\) (and likewise in \(y\)). Neighbour coefficients of the Laplacian then stay nonnegative with no restriction on \(h_x\), at the cost of reducing the convection term from \(O(h^2)\) to \(O(h)\). The diffusion \(\delta_{xx}\) is left centered: one does not replace \(\partial_{xx}\) by \((D^\pm)^2\) merely because convection was upwinded. The exam says ``centered'', so this remark is the alternative that one does \emph{not} write as the scheme, but that one names as the reason \cref{ass:peclet} is needed.
\end{remark}

\begin{notation}[Time quotients, the same three columns]
\label{not:time-quotients}
The identical forward / backward / centered choice exists in \(t\), independently of the spatial stencil. With time step \(\tau\) and \(U^n\) the grid function at \(t_n\),
\begin{align*}
\frac{U^{n+1}-U^n}{\tau}
&=
D^+_t U^n
=
D^-_t U^{n+1},
\qquad
\frac{U^{n+1}-U^{n-1}}{2\tau}
=
D^0_t U^n.
\end{align*}
Pairing with a spatial operator \(L_h\) produces the three standard time-steppers for \(u_t=Lu\):
\begin{enumerate}[label=\textup{(\roman*)}]
    \item \emph{Forward Euler} (explicit): \((U^{n+1}-U^n)/\tau=L_h U^n\). First-order in time; conditionally \(\ell_2\)-stable for heat (\(\mu\le 1/2\), \cref{exer:ftcs}); the discrete maximum principle typically imposes an even stricter \(\tau=O(h^2)\) restriction in \(\ell_\infty\).
    \item \emph{Backward Euler} (implicit): \((U^{n+1}-U^n)/\tau=L_h U^{n+1}\). First-order in time; unconditionally stable once \(L_h\) obeys a maximum principle. \textbf{This is what Problem~5 asks.}
    \item \emph{Centered in time:} leapfrog \((U^{n+1}-U^{n-1})/(2\tau)=L_h U^n\) is two-step and usually the wrong pairing with diffusion; Crank--Nicolson \((U^{n+1}-U^n)/\tau=L_h(U^{n+1}+U^n)/2\) is the second-order implicit average. Neither is asked here.
\end{enumerate}
Spatial ``centered'' and temporal ``backward Euler'' are independent choices. The exam combines them: centered five-point \(L_h\), backward Euler in \(t\).
\end{notation}

\begin{notation}[Centered discrete operator]
\label{not:Lh}
On the grid of the previous paragraph, define
\[
(L_h U)_{ij}
\coloneqq
\delta_{xx}U_{ij}+\delta_{yy}U_{ij}+d_{ij}\,\delta_x U_{ij}+e_{ij}\,\delta_y U_{ij}+f_{ij}\,U_{ij}.
\]
If \(d,e,f\) are constants, drop the indices. Expanding,
\begin{align*}
(L_h U)_{ij}
&=
a_E U_{i+1,j}+a_W U_{i-1,j}+a_N U_{i,j+1}+a_S U_{i,j-1}+a_C U_{ij},
\end{align*}
with
\begin{align*}
a_E
&=
\frac{1}{h_x^2}+\frac{d_{ij}}{2h_x},
&
a_W
&=
\frac{1}{h_x^2}-\frac{d_{ij}}{2h_x},
&
a_N
&=
\frac{1}{h_y^2}+\frac{e_{ij}}{2h_y},
&
a_S
&=
\frac{1}{h_y^2}-\frac{e_{ij}}{2h_y},
\end{align*}
and
\[
a_C
=
-\frac{2}{h_x^2}-\frac{2}{h_y^2}+f_{ij}.
\]
In particular \(a_E+a_W+a_N+a_S+a_C=f_{ij}\), which rearranges as the neighbour form used in every maximum-principle proof:
\begin{equation}
\label{eq:neighbour-form}
(L_h U)_{ij}
=
a_E(U_{i+1,j}-U_{ij})
+a_W(U_{i-1,j}-U_{ij})
+a_N(U_{i,j+1}-U_{ij})
+a_S(U_{i,j-1}-U_{ij})
+f_{ij}U_{ij}.
\end{equation}
\end{notation}

\begin{center}
\begin{tikzpicture}[
  nd/.style={circle, draw, minimum size=1.15cm, inner sep=1pt, font=\small, align=center},
  lab/.style={font=\scriptsize}
]
  \node[nd] (c) at (0,0) {\(U_{ij}\)\\[-1pt]{\scriptsize \(a_C\)}};
  \node[nd] (e) at (2.4,0) {\(U_{i+1,j}\)\\[-1pt]{\scriptsize \(a_E\)}};
  \node[nd] (w) at (-2.4,0) {\(U_{i-1,j}\)\\[-1pt]{\scriptsize \(a_W\)}};
  \node[nd] (n) at (0,2.15) {\(U_{i,j+1}\)\\[-1pt]{\scriptsize \(a_N\)}};
  \node[nd] (s) at (0,-2.15) {\(U_{i,j-1}\)\\[-1pt]{\scriptsize \(a_S\)}};
  \draw[thick] (w) -- (c) -- (e);
  \draw[thick] (s) -- (c) -- (n);
\end{tikzpicture}
\end{center}

\begin{assumption}[Mesh P\'eclet / monotonicity]
\label{ass:peclet}
The neighbour coefficients in \cref{not:Lh} are nonnegative, and the zeroth-order coefficient is strictly negative:
\[
\lvert d_{ij}\rvert h_x\le 2,
\qquad
\lvert e_{ij}\rvert h_y\le 2,
\qquad
f_{ij}<0.
\]
The first two inequalities are equivalent to \(a_E,a_W,a_N,a_S\ge 0\). They fail if the convection is strong and the mesh is coarse (cell P\'eclet number \(\lvert d\rvert h_x/2>1\)); then centered differences are the wrong spatial scheme and the maximum principle proof stops. Upwinding the first derivatives restores monotonicity without a mesh restriction, at the cost of first-order spatial accuracy. The exam says ``centered'' and ``suitable assumptions'', so one must write \cref{ass:peclet} explicitly.
\end{assumption}

\begin{theorem}[Discrete maximum principle]
\label{thm:dmp}
Assume \cref{ass:peclet}. If \(L_h U=0\) at every interior node, then \(U\) cannot attain a strictly positive maximum at an interior node. Consequently
\[
\max_{\text{all nodes}}U
\le
\max\bigl\{0,\,\max_{\text{boundary nodes}}U\bigr\}.
\]
Applying the same statement to \(-U\) yields the minimum principle, and therefore
\[
\max_{\text{all nodes}}\lvert U\rvert
\le
\max_{\text{boundary nodes}}\lvert U\rvert.
\]
\end{theorem}

\begin{proof}
Suppose, for a contradiction, that \(U_{ij}=M>0\) is an interior maximum. Then each neighbour difference in \eqref{eq:neighbour-form} is \(\le 0\). By \cref{ass:peclet} the four neighbour coefficients are \(\ge 0\), so those four terms are \(\le 0\), while \(f_{ij}M<0\). Hence \((L_h U)_{ij}<0\), contradicting \(L_h U=0\).

Thus any positive maximum is attained on the boundary (or else \(U\le 0\) everywhere). That is the displayed bound. The argument never used constancy of \(d,e,f\), only the sign conditions at the maximizing node.
\end{proof}

If the boundary data vanish and \(L_h U=0\), then \(U\le 0\) and \(-U\le 0\), so \(U\equiv 0\). Unique solvability of the discrete Dirichlet problem is therefore included.

\begin{example}[1D ancestor: 2024 Spring Problem~5]
\label{ex:1d-dmp}
On a 1D grid, the operator
\[
(Lu)_j=-(a_j u_{j-1}-b_j u_j+c_j u_{j+1})+q_j u_j
\]
with \(a_j,b_j,c_j>0\), \(q_j\ge 0\), and \(a_j+c_j\le b_j\), is the same monotonicity package. If \(u_r=M>0\) is an interior maximum and \(Lu\le 0\), the neighbour form forces \(Lu_r\ge(b_r-a_r-c_r+q_r)M\ge 0\), hence equality throughout and a constant state. The 2D proof above is that argument with four neighbours instead of two, and with \(f<0\) playing the role of a strictly positive \(q\).
\end{example}

\paragraph{Backward Euler, first as an ODE, then as a PDE scheme.}

For a linear ODE \(y'=Ay\), backward Euler is
\[
\frac{y^{n+1}-y^n}{\tau}
=
A y^{n+1},
\qquad
(I-\tau A)y^{n+1}=y^n.
\]
The one-step matrix is therefore \((I-\tau A)^{-1}\). The numerical-analysis bound one actually wants is not a Neumann series in \(\tau\|A\|\); it is a resolvent estimate under dissipativity.

\begin{lemma}[Resolvent bound for backward Euler]
\label{lem:be-resolvent}
Let \(A\) be a real square matrix and \(\tau>0\).
\begin{enumerate}[label=\textup{(\roman*)}]
    \item
    \emph{Inner product.} If \(\langle Ax,x\rangle\le 0\) for every \(x\) (in particular if \(A+A^{T}\) is negative semidefinite), then \(I-\tau A\) is invertible and
    \[
    \bigl\|(I-\tau A)^{-1}\bigr\|_{2}
    \le
    1.
    \]
    \item
    \emph{Logarithmic norm.} For any vector norm, write
    \[
    \mu[A]
    \coloneqq
    \lim_{\varepsilon\downarrow 0}\frac{\|I+\varepsilon A\|-1}{\varepsilon}.
    \]
    Then \(I-\tau A\) is invertible whenever \(\tau\mu[A]<1\), and
    \[
    \bigl\|(I-\tau A)^{-1}\bigr|
    \le
    \frac{1}{1-\tau\mu[A]}.
    \]
    In particular \(\mu[A]\le 0\) implies \(\bigl\|(I-\tau A)^{-1}\bigr\|\le 1\) for every \(\tau>0\). For the Euclidean and max norms,
    \[
    \mu_{2}[A]
    =
    \lambda_{\max}\Bigl(\tfrac{A+A^{T}}{2}\Bigr),
    \qquad
    \mu_{\infty}[A]
    =
    \max_{i}\Bigl(a_{ii}+\sum_{j\neq i}\lvert a_{ij}\rvert\Bigr).
    \]
\end{enumerate}
Thus \(\|y^{n+1}\|\le\|y^n\|\) with no restriction on \(\tau\), as soon as \(A\) is dissipative in the chosen norm. This is \(A\)-stability in matrix form. The Neumann bound \(\bigl\|(I-\tau A)^{-1}\bigr\|\le 1/(1-\tau\|A\|)\) (valid only for \(\tau\|A\|<1\)) is a different, small-step estimate; it does not give unconditional stability.
\end{lemma}

\begin{proof}
For (i), set \(x=(I-\tau A)y\). Then
\[
\langle x,y\rangle
=
\|y\|_{2}^{2}-\tau\langle Ay,y\rangle
\ge
\|y\|_{2}^{2},
\]
so \(\|y\|_{2}^{2}\le\langle x,y\rangle\le\|x\|_{2}\|y\|_{2}\). If \(y\neq 0\) this is \(\|y\|_{2}\le\|x\|_{2}\). Taking \(x=0\) gives \(y=0\), hence invertibility, and the operator-norm claim is the same inequality.

For (ii), the definition of \(\mu\) is equivalent to the resolvent bound \(\|(\zeta I-A)^{-1}\|\le 1/(\operatorname{Re}\zeta-\mu[A])\) for \(\operatorname{Re}\zeta>\mu[A]\) (the generator of the semigroup \(e^{tA}\) has growth bound \(\mu[A]\)). Substituting \(\zeta=1/\tau\) yields
\[
\bigl\|(I-\tau A)^{-1}\bigr\|
=
\tau^{-1}\bigl\|(\tau^{-1}I-A)^{-1}\bigr\|
\le
\frac{\tau^{-1}}{\tau^{-1}-\mu[A]}
=
\frac{1}{1-\tau\mu[A]}.
\]
The \(\ell_{\infty}\) formula is the first-order expansion of \(\|I+\varepsilon A\|_{\infty}=\max_{i}\bigl(1+\varepsilon a_{ii}+\varepsilon\sum_{j\neq i}\lvert a_{ij}\rvert+O(\varepsilon^{2})\bigr)\) for small \(\varepsilon>0\). If \(\mu_{\infty}[A]\le 0\), an elementary row argument avoids the limit: at a maximizing index \(k\) of \(y\), the identity \(((I-\tau A)y)_{k}=x_{k}\) and the row-sum hypothesis give \(\lvert y_{k}\rvert\le\|x\|_{\infty}\).
\end{proof}

If \(A\) is a scalar \(\lambda\) with \(\operatorname{Re}\lambda\le 0\), then \(g=1/(1-\tau\lambda)\) satisfies \(\lvert g\rvert\le 1\) with no restriction on \(\tau>0\): that is part~(i) in dimension one, already used in Problem~4(c). In several space dimensions the same formula is applied to the vector of interior unknowns, with \(A\) replaced by the matrix of \(L_{h}\) (including homogeneous Dirichlet by omitting boundary nodes). Under \cref{ass:peclet} one has \(\mu_{\infty}[L_{h}]\le 0\), so \cref{lem:be-resolvent} already yields \(\|U^{n+1}\|_{\infty}\le\|U^{n}\|_{\infty}\). The maximum-principle proof of \cref{thm:be-stable} is that row argument written without assembling the matrix.

On the grid, backward Euler for \(u_t=Lu\) reads
\begin{equation}
\label{eq:be}
\frac{U^{n+1}-U^n}{\tau}
=
L_h U^{n+1},
\qquad\text{i.e.}\qquad
U^n
=
U^{n+1}-\tau L_h U^{n+1}.
\end{equation}
Each time step is an elliptic solve for \(U^{n+1}\). The maximum principle controls that solve in \(\ell_\infty\), which is a stronger and more local statement than von Neumann.

\begin{theorem}[Backward Euler is unconditionally \(\ell_\infty\) stable]
\label{thm:be-stable}
Assume \cref{ass:peclet}, and impose homogeneous Dirichlet conditions at every time level. The scheme \eqref{eq:be} satisfies
\[
\|U^{n+1}\|_\infty
\le
\|U^n\|_\infty,
\]
with no restriction on \(\tau>0\). In particular \(\|U^n\|_\infty\le\|U^0\|_\infty\) for all \(n\).
\end{theorem}

\begin{proof}
Write \(M\coloneqq\max U^{n+1}\), the maximum being taken over all nodes, boundary included. The boundary values are \(0\), so either \(M\le 0\), or \(M>0\) is attained at an interior node \((i,j)\). In the latter case, the neighbour form \eqref{eq:neighbour-form} and \cref{ass:peclet} give \((L_h U^{n+1})_{ij}<0\), hence
\[
U_{ij}^n
=
U_{ij}^{n+1}-\tau(L_h U^{n+1})_{ij}
=
M-\tau\cdot(\text{a negative number})
>
M.
\]
Thus \(M<\max U^n\le\|U^n\|_\infty\). Combining with the case \(M\le 0\), one has \(\max U^{n+1}\le\|U^n\|_\infty\). The same argument applied to \(-U\) bounds \(-\min U^{n+1}\). Therefore \(\|U^{n+1}\|_\infty\le\|U^n\|_\infty\).
\end{proof}

The time step \(\tau\) never entered the sign of \(L_h U^{n+1}\). Unconditional stability is with respect to \(\tau\); the mesh restriction in \cref{ass:peclet} is on \(h_x,h_y\), not on \(\tau\).

\begin{examnote}[Two stability languages]
\label{exnote:two-languages}
Von Neumann (\cref{def:von-neumann}) is an \(\ell_2\) statement about symbols. The discrete maximum principle (\cref{thm:dmp}) is an \(\ell_\infty\) statement about signs of stencil coefficients. Qiuzhen uses the first for constant-coefficient Cauchy/periodic problems (advection, heat, waves), and the second for Dirichlet elliptic/parabolic problems with lower-order terms. Mixing the two on the script is a common way to lose the ``suitable assumptions''.
\end{examnote}

\subsubsection{Solution of Problem 5}

\paragraph{The scheme.}

On a uniform grid of \((0,1)^2\), the centered scheme for the elliptic equation is \(L_h U=0\) in the interior, with \(L_h\) as in \cref{not:Lh}. The suitable assumptions are \cref{ass:peclet}. (If the exam's \(d,e,f\) are constants, the same inequalities are imposed globally.)

\paragraph{Maximum principle.}

This is \cref{thm:dmp}. The one-line engine is \eqref{eq:neighbour-form}: at an interior positive maximum every neighbour difference is nonpositive, every neighbour coefficient is nonnegative, and \(fU<0\), so \(L_h U<0\).

\paragraph{Backward Euler.}

The parabolic equation is \(u_t=Lu\) with homogeneous Dirichlet data. Backward Euler is \eqref{eq:be}. \Cref{thm:be-stable} gives
\[
\|U^{n+1}\|_\infty\le\|U^n\|_\infty
\]
with no restriction on \(\tau>0\), which is \(\ell_\infty\) stability on \([0,T]\) in the sense of \cref{def:lax-richtmyer} (take \(C=1\) and \(\alpha=0\)).

If one wants a comparison form: the error \(E^n=U^n-u(\cdot,t_n)\) (restricted to the grid) satisfies the same scheme with a truncation forcing on the right-hand side. The maximum principle applied to \(E^{n+1}-C\tau\) with \(C\) large enough to dominate the residual yields the usual \(O(\tau+h_x^2+h_y^2)\) bound, but the exam only asked stability, not an error estimate.


\noindent\textbf{You can choose either 6 or 7 to answer. The points will be decided as \(\max(6, 7)\).}

\subsection{Problem 6. [20 points]}
Consider the following system

\[
\ddot{x} - 2x - x^{2} + x^{3} = 0.
\]

\begin{enumerate}[label=(\alph*)]
    \item (6 points) Find the static solutions to the system.
    \item (14 points) For sufficiently small but nonzero \(\epsilon\), show that the orbit starting at \((x(0), \dot{x}(0)) = (2 + \epsilon, 0)\) is closed in the phase space, i.e., \(x - \dot{x}\) plane. Find the approximation of the period of this orbit up to \(O(\epsilon^{2})\).
\end{enumerate}

\setcounter{section}{6}
\setcounter{theorem}{0}
\setcounter{definition}{0}
\setcounter{remark}{0}
\setcounter{note}{0}

\subsubsection{Preliminaries: ODE methods specialised to this problem}

The equation is autonomous, second-order, and conservative: there is no explicit \(t\), and no first-derivative damping. The Grade~2 ODE notes \cite{grade2ode} are a catalogue of elementary solution methods. Most of that catalogue does \emph{not} produce a closed-form \(x(t)\) here. Two items do: reduction of order when \(t\) is absent (energy), and the classification of equilibria in the phase plane. The three methods one actually reaches for on a second-order ODE --- separation of variables, the companion-system / matrix-exponential calculus, and the elementary tricks --- are written out after the energy/Lindstedt path that the script uses. The operator calculus is the general brute-force theory (rewrite as \(\dot{\mathbf{z}}=A\mathbf{z}+\mathbf{N}(\mathbf{z})\), integrate against \(e^{tA}\)); it does not evaluate in elementary functions on this cubic, which is why the script does not grind it. The closed-orbit uniqueness argument is the sibling note \cite{grade2period} (originally written for \(\ddot x+x+x^{3}=0\)). The \(O(\epsilon^{2})\) period is a Lindstedt--Poincar\'e correction on top of that energy picture, which the sophomore notes do not develop; it is the standard next tool once the orbit is known to be periodic \cite{nayfeh73}.

\paragraph{Method map (what the Grade~2 notes actually contain).}

{\footnotesize
\begin{center}
\begin{tabular}{@{}>{\raggedright\arraybackslash}p{3.5cm}
                   >{\raggedright\arraybackslash}p{5.3cm}
                   >{\raggedright\arraybackslash}p{5.0cm}@{}}
\toprule
Family in \cite{grade2ode} & Typical move & This problem \\
\midrule
separation of variables / homogeneous \(y/x\) / Bernoulli & first-order quadrature & not on \(\ddot x=f(x)\) directly; yes after energy (\cref{ex:sep-var}) \\
integrating factor / exact equation & \(P\,\mathrm{d}x+Q\,\mathrm{d}y=0\) & the energy identity \emph{is} an exact \(P\,\mathrm{d}x+Q\,\mathrm{d}\dot x=0\) \\
first-order, missing \(x\) or \(y\) & parameter \(p=y'\) & a different missing variable; see reduction of order below \\
Picard / existence--uniqueness & contraction on \(C^{0}\) & uniqueness of the IVP, used to close the orbit \\
higher-order linear / characteristic roots / Euler / variation of parameters / Laplace & linear superposition & only the linearization at an equilibrium \\
\textbf{reduction of order: \(t\) absent} & \(v=\dot x\), \(x\) as independent variable & \textbf{the energy integral} \\
power series / Bessel & Frobenius at a regular singular point & not this polynomial ODE \\
linear systems / matrix exponential & \(\dot{\mathbf{z}}=A\mathbf{z}+\mathbf{N}(\mathbf{z})\), Duhamel & general brute-force theory; does not close here (\cref{thm:duhamel,ex:duhamel-p6}) \\
Lyapunov \(V\) / linearization & sign of \(V\), eigenvalues of \(DF\) & \(V\) is the energy; linearization gives \(V''\) at equilibria \\
limit cycle & attracting cycle of a dissipative planar system & here the cycle is a \emph{level set of energy}, not a limit cycle \\
classification of equilibria & node / saddle / focus / center & (a) plus \(V''\gtrless 0\) \\
\bottomrule
\end{tabular}
\end{center}
}

\begin{definition}[Reduction of order when \(t\) is absent]
\label{def:reduce-autonomous}
\cite{grade2ode}, higher-order ODEs / reduction of order / independent variable \(t\) absent. For
\[
F\bigl(x,\dot x,\ddot x,\ldots,x^{(n)}\bigr)=0,
\]
set \(v=\dot x\) and treat \(x\) as the independent variable. Then
\[
\ddot x=v\frac{\mathrm{d}v}{\mathrm{d}x},\qquad
x'''=v\frac{\mathrm{d}}{\mathrm{d}x}\Bigl(v\frac{\mathrm{d}v}{\mathrm{d}x}\Bigr),\qquad
\ldots,
\]
and the order drops by one. For a second-order equation \(\ddot x=f(x)\) this is immediate: multiply by \(\dot x\) and integrate in \(t\),
\begin{equation}
\label{eq:energy}
\frac12\dot x^{2}+V(x)=E,
\qquad
V'(x)=-f(x).
\end{equation}
The first-order quadrature
\[
\mathrm{d}t=\frac{\mathrm{d}x}{\sqrt{2\bigl(E-V(x)\bigr)}}
\]
then produces the period of a turning-point orbit as
\begin{equation}
\label{eq:period-quad}
T
=
\sqrt{2}\int_{x_{\min}}^{x_{\max}}\frac{\mathrm{d}x}{\sqrt{E-V(x)}}.
\end{equation}
\end{definition}

That identity is the whole of part~(b)'s topology. The \(O(\epsilon^{2})\) expansion of \eqref{eq:period-quad} is messy (the integrand is singular at the endpoints). Lindstedt--Poincar\'e expands the same period by stretching time so that the solution stays \(2\pi\)-periodic order by order, and kills resonant forcing.

\begin{definition}[Static solution / equilibrium]
\label{def:static}
A \emph{static solution} of \(\ddot x=f(x)\) is a constant \(x(t)\equiv x_{\star}\) with \(f(x_{\star})=0\). In the phase plane \((x,\dot x)\) it is the equilibrium \((x_{\star},0)\). The linearization is \(\ddot y=f'(x_{\star})y\). If \(f'(x_{\star})<0\), equivalently \(V''(x_{\star})>0\), the equilibrium is a linear center; if \(f'(x_{\star})>0\), it is a saddle. For a conservative planar system the linear center is nonlinearly a genuine center: sublevel sets of \(E\) near \(x_{\star}\) are closed curves \cite{grade2ode}.
\end{definition}

\begin{proposition}[Energy \(\Rightarrow\) periodic orbit]
\label{prop:energy-periodic}
\cite{grade2period}. Let \(\ddot x=f(x)\) with \(f\) locally Lipschitz, and let \(V'=-f\), so that
\[
\frac12\dot x^{2}+V(x)=E
\]
is constant along solutions. Fix rest initial data \((x(0),\dot x(0))=(a,0)\), and write \(E=V(a)\). Suppose there is a second point \(b\neq a\) such that:
\begin{enumerate}[label=\textup{(\roman*)}]
    \item \(V(b)=V(a)\), and \(V(x)<V(a)\) for every \(x\) strictly between \(a\) and \(b\);
    \item \(V(x)>V(a)\) immediately outside the closed interval \(I\) with endpoints \(a\) and \(b\) (equivalently: \(I\) is a connected component of the sublevel set \(\{V\le V(a)\}\)).
\end{enumerate}
Then the orbit in the \((x,\dot x)\) plane is the simple closed curve \(\frac12\dot x^{2}+V(x)=V(a)\) over \(I\), and \(x(t)\) is periodic.

The two turning points are the only zeros of \(\dot x\) on that level set, because \(\frac12\dot x^{2}=V(a)-V(x)\) vanishes on \(I\) precisely at the endpoints. Energy conservation plus (ii) keeps the solution inside \(I\): leaving \(I\) would require \(V>V(a)\) with \(\frac12\dot x^{2}\ge 0\), which contradicts the conserved value \(E=V(a)\). After the second turning point the time-shifted solution shares Cauchy data with the original one, so uniqueness of the Lipschitz IVP identifies them and produces a period.
\end{proposition}

\begin{construction}[Lindstedt--Poincar\'e for a weakly nonlinear center]
\label{con:lindstedt}
After translating an equilibrium to the origin, write
\[
\ddot y+\omega_{0}^{2} y=g(y),\qquad g(y)=O(y^{2}),
\]
with initial data \(y(0)=\epsilon\), \(\dot y(0)=0\), and \(\epsilon\) small. A regular expansion \(y=\epsilon\cos(\omega_{0} t)+\cdots\) in the original time produces resonant forcing \(\cos(\omega_{0} t)\) at higher order, hence secular terms \(t\sin(\omega_{0} t)\). Those terms are an artefact of expanding at the linear frequency while the true period is slightly different. Lindstedt stretches the clock so that the unknown frequency is built into the independent variable: set \(s=\Omega t\) with \(\Omega\) itself expanded, and demand that every \(y_{k}(s)\) be \(2\pi\)-periodic. Then \(t\mapsto t+T\) with \(T=2\pi/\Omega\) is a period of \(y(t)\).

The quadratic interaction \(y_{1}^{2}\) Fourier-decomposes into a constant and \(\cos 2s\) only, so it does not resonate at order \(\epsilon^{2}\). The first solvability condition is therefore vacuous for a correction of size \(\epsilon\) in \(\Omega\), and one may take
\[
\Omega=\omega_{0}\bigl(1+\nu\epsilon^{2}+O(\epsilon^{4})\bigr),\qquad
y=\epsilon y_{1}(s)+\epsilon^{2} y_{2}(s)+\epsilon^{3} y_{3}(s)+\cdots,
\]
with \(y_{1}(0)=1\) and \(y_{k}(0)=y_{k}'(0)=0\) for \(k\ge 2\). At order \(\epsilon^{3}\) the coefficient of \(\cos s\) on the right-hand side of \(y_{3}''+y_{3}=\cdots\) is a linear function of \(\nu\); choosing \(\nu\) to cancel it removes the secular term and yields
\[
T=\frac{2\pi}{\Omega}=\frac{2\pi}{\omega_{0}}\bigl(1-\nu\epsilon^{2}\bigr)+O(\epsilon^{4}).
\]
One never needs the function \(y_{3}\) itself: the period is a solvability condition, not a quadrature of \(y_{3}\).
\end{construction}

\paragraph{Three methods one actually reaches for.}

A second-order scalar ODE is not a single algorithm. The three moves below are the ones that come to mind on an exam, in increasing generality. Separation of variables and the elementary tricks close first-order or linear problems. The companion-system calculus --- rewrite as \(\dot{\mathbf{z}}=A\mathbf{z}+\mathbf{N}(\mathbf{z})\), integrate against \(e^{At}\) --- is the brute-force theory that always produces an integral equation, and is what Yau applied / analysis papers ask when they write \(\mathbf{y}'=A\mathbf{y}+\mathbf{f}\). For \emph{this} problem the integral does not evaluate in elementary functions, which is why the script uses energy plus Lindstedt instead. The theory is still the thing one should be able to write.

\begin{example}[Separation of variables, the first-order move]
\label{ex:sep-var}
If a first-order equation can be written \(\mathrm{d}y/\mathrm{d}x=g(x)h(y)\) with \(h(y)\neq 0\), then
\[
\int\frac{\mathrm{d}y}{h(y)}
=
\int g(x)\,\mathrm{d}x+C,
\]
provided one also checks the equilibria \(h(y)=0\). The logistic equation \(\dot y=y(1-y)\) separates as
\[
\int\Bigl(\frac{1}{y}+\frac{1}{1-y}\Bigr)\mathrm{d}y
=
\int\mathrm{d}t
\qquad\Longrightarrow\qquad
y(t)
=
\frac{y_0 e^{t}}{1-y_0+y_0 e^{t}}.
\]
A homogeneous equation \(\mathrm{d}y/\mathrm{d}x=f(y/x)\) becomes separable after \(u=y/x\). None of this applies directly to \(\ddot x=2x+x^{2}-x^{3}\), which is second-order in \(t\). After the energy reduction \eqref{eq:energy}, however, the quadrature
\[
\mathrm{d}t
=
\frac{\mathrm{d}x}{\sqrt{2\bigl(E-V(x)\bigr)}}
\]
\emph{is} separation of variables. For the quartic \(V\) of this problem the integral is elliptic, not elementary, so separation produces the period as a definite integral and then stops. That is the right diagnosis, not a failed method.
\end{example}

\begin{definition}[Companion system]
\label{def:companion}
Every second-order equation \(\ddot x=f(x,\dot x,t)\) is a first-order system on the phase plane. Set
\[
\mathbf{z}
=
\begin{pmatrix} x \\ v \end{pmatrix},
\qquad
v=\dot x,
\qquad
\dot{\mathbf{z}}
=
\mathbf{F}(t,\mathbf{z})
=
\begin{pmatrix} v \\ f(x,v,t) \end{pmatrix}.
\]
For the exam equation \(f\) is independent of \(t\) and of \(v\), so \(\mathbf{F}\) is autonomous and conservative. Equilibria of \(\mathbf{F}\) are exactly the static solutions of part~(a). The linearization at an equilibrium \((x_{\star},0)\) is the constant-coefficient system \(\dot{\mathbf{w}}=A\mathbf{w}\) with
\[
A
=
D\mathbf{F}(x_{\star},0)
=
\begin{pmatrix}
0 & 1 \\
f'(x_{\star}) & 0
\end{pmatrix}.
\]
The characteristic polynomial is \(\lambda^{2}-f'(x_{\star})=0\). When \(f'(x_{\star})<0\) the eigenvalues are \(\pm i\omega\) with \(\omega=\sqrt{-f'(x_{\star})}\) (linear center); when \(f'(x_{\star})>0\) they are real and opposite (saddle). This recovers \cref{def:static}.
\end{definition}

\begin{definition}[Matrix exponential as a linear operator]
\label{def:matexp}
For \(A\in\mathrm{M}_n(\mathbb R)\) the operator \(e^{tA}\) is the norm-convergent series
\[
e^{tA}
\coloneqq
\sum_{k=0}^{\infty}\frac{(tA)^{k}}{k!},
\]
the unique solution of \(\dot\Phi=A\Phi\) with \(\Phi(0)=I\). Equivalently, if \(A=PJP^{-1}\) is a Jordan form, \(e^{tA}=P e^{tJ}P^{-1}\). The homogeneous IVP \(\dot{\mathbf{z}}=A\mathbf{z}\), \(\mathbf{z}(0)=\mathbf{z}_0\), is therefore
\[
\mathbf{z}(t)
=
e^{tA}\mathbf{z}_0.
\]
This is the general integration method for a constant-coefficient linear system: the matrix \(A\) is an operator on \(\mathbb R^n\), and time integration is the semigroup it generates. The same formula is the variation-of-constants kernel below. Computing \(e^{tA}\) is a finite algebraic problem once the Jordan form of \(A\) is known (eigenvalues, and a nilpotent block if there is a repeated root).
\end{definition}

\begin{example}[The linearization at \(x=2\), computed]
\label{ex:expA-center}
At the center \(x=2\) one has \(f'(2)=-V''(2)=-6\), so
\[
A
=
\begin{pmatrix} 0 & 1 \\ -6 & 0 \end{pmatrix}.
\]
The Cayley--Hamilton relation for this companion matrix is a scalar multiple of the identity: write \(\omega=\sqrt{6}\) and compute
\[
A^{2}
=
\begin{pmatrix} 0 & 1 \\ -\omega^{2} & 0 \end{pmatrix}
\begin{pmatrix} 0 & 1 \\ -\omega^{2} & 0 \end{pmatrix}
=
\begin{pmatrix} -\omega^{2} & 0 \\ 0 & -\omega^{2} \end{pmatrix}
=
-\omega^{2} I.
\]
Every higher power is therefore either a multiple of \(I\) or a multiple of \(A\). Inductively,
\[
A^{2m}=(-\omega^{2})^{m}I=(-1)^{m}\omega^{2m}I,
\qquad
A^{2m+1}=A^{2m}A=(-1)^{m}\omega^{2m}A,
\qquad m\ge 0.
\]
(The first few: \(A^{2}=-\omega^{2}I\), \(A^{3}=-\omega^{2}A\), \(A^{4}=\omega^{4}I\), \(A^{5}=\omega^{4}A\).) Split the defining series of \cref{def:matexp} into even and odd terms:
\begin{align*}
e^{tA}
&=
\sum_{k=0}^{\infty}\frac{t^{k}}{k!}A^{k}
=
\sum_{m=0}^{\infty}\frac{t^{2m}}{(2m)!}A^{2m}
+
\sum_{m=0}^{\infty}\frac{t^{2m+1}}{(2m+1)!}A^{2m+1}
\\
&=
\Biggl(\sum_{m=0}^{\infty}(-1)^{m}\frac{(\omega t)^{2m}}{(2m)!}\Biggr)I
+
\frac{1}{\omega}
\Biggl(\sum_{m=0}^{\infty}(-1)^{m}\frac{(\omega t)^{2m+1}}{(2m+1)!}\Biggr)A.
\end{align*}
The two scalar series are the Taylor series of cosine and sine, so
\[
e^{tA}
=
\cos(\omega t)\,I
+
\frac{\sin(\omega t)}{\omega}\,A.
\]
Written out to \(O(t^{5})\) before summing, the same identity is
\begin{align*}
e^{tA}
&=
I+tA+\frac{t^{2}}{2!}(-\omega^{2}I)+\frac{t^{3}}{3!}(-\omega^{2}A)+\frac{t^{4}}{4!}(\omega^{4}I)+\frac{t^{5}}{5!}(\omega^{4}A)+\cdots
\\
&=
\Bigl(1-\frac{(\omega t)^{2}}{2!}+\frac{(\omega t)^{4}}{4!}-\cdots\Bigr)I
+
\frac{1}{\omega}
\Bigl(\omega t-\frac{(\omega t)^{3}}{3!}+\frac{(\omega t)^{5}}{5!}-\cdots\Bigr)A.
\end{align*}
Substituting the matrices \(I\) and \(A\) gives the harmonic-oscillator fundamental matrix
\[
e^{tA}
=
\begin{pmatrix}
\cos\omega t & \omega^{-1}\sin\omega t \\
-\omega\sin\omega t & \cos\omega t
\end{pmatrix}.
\]
The linearized IVP \(\ddot y+6y=0\), \(y(0)=\epsilon\), \(\dot y(0)=0\) is therefore \(y(t)=\epsilon\cos(\omega t)\), period \(2\pi/\sqrt{6}\), which is the leading term of part~(b). This is the same calculation a Yau applied paper asks when it writes ``solve \(\dot{\mathbf{y}}=A\mathbf{y}\) by the matrix exponential''.
\end{example}

\begin{theorem}[Variation of constants / Duhamel convolution]
\label{thm:duhamel}
Let \(A\) be constant, and consider the linear nonhomogeneous system
\[
\dot{\mathbf{z}}
=
A\mathbf{z}+\mathbf{f}(t),
\qquad
\mathbf{z}(0)=\mathbf{z}_0.
\]
The unique solution is the convolution of the forcing against the semigroup \(e^{tA}\):
\begin{equation}
\label{eq:duhamel}
\mathbf{z}(t)
=
e^{tA}\mathbf{z}_0
+
\int_{0}^{t} e^{(t-s)A}\mathbf{f}(s)\,\mathrm{d}s.
\end{equation}
If instead \(A=A(t)\) is time-dependent and \(\Phi(t)\) is a fundamental matrix (\(\dot\Phi=A(t)\Phi\), \(\det\Phi\neq 0\)), the same identity reads
\[
\mathbf{z}(t)
=
\Phi(t)\Phi(0)^{-1}\mathbf{z}_0
+
\Phi(t)\int_{0}^{t}\Phi(s)^{-1}\mathbf{f}(s)\,\mathrm{d}s.
\]
This is the Grade~2 variation-of-constants formula \cite{grade2ode}. On a scalar \(n\)th-order linear equation it is the Wronskian formula for a particular solution. The kernel \(e^{(t-s)A}\) (or \(\Phi(t)\Phi(s)^{-1}\)) is the Green's function of the linear differential operator \(\frac{\mathrm{d}}{\mathrm{d}t}-A\). Integrating a linear ODE is convolution against that kernel.
\end{theorem}

\begin{proof}[Proof of \eqref{eq:duhamel}]
Set \(\mathbf{w}(t)=e^{-tA}\mathbf{z}(t)\). Then \(\dot{\mathbf{w}}=e^{-tA}\mathbf{f}(t)\), so \(\mathbf{w}(t)=\mathbf{z}_0+\int_0^t e^{-sA}\mathbf{f}(s)\,\mathrm{d}s\). Multiply by \(e^{tA}\). Equivalently, this is the integrating factor for the vector equation \(\dot{\mathbf{z}}-A\mathbf{z}=\mathbf{f}\), with integrating factor \(e^{-tA}\).
\end{proof}

\begin{definition}[Dirac delta]
\label{def:dirac}
The symbol \(\delta\) below is the Dirac unit impulse at \(t=0\), not a Kronecker \(\delta_{nm}\). As a distribution it is the evaluation functional
\[
\langle\delta,\varphi\rangle=\varphi(0)
\]
on test functions \(\varphi\). Informally: \(\delta(t)=0\) for \(t\neq 0\), and \(\int_{-\varepsilon}^{\varepsilon}\delta(t)\,\mathrm{d}t=1\) for every \(\varepsilon>0\). The ODE meaning is a unit kick at a single instant. A locally integrable function \(G\) (usually taken causal, \(G(t)=0\) for \(t<0\)) is a fundamental solution of a constant-coefficient operator \(L(D)=D^{n}+a_{n-1}D^{n-1}+\cdots+a_{0}\) when
\[
L(D)G=\delta
\]
in the distributional sense. Equivalently, \(G\) solves the homogeneous equation \(L(D)G=0\) for \(t>0\), with rest data at \(t=0^{+}\) except for a unit jump in the highest derivative that \(L\) sees:
\[
G(0^{+})=\dot G(0^{+})=\cdots=G^{(n-2)}(0^{+})=0,
\qquad
G^{(n-1)}(0^{+})=1.
\]
(The lower-order jet is continuous across \(t=0\); the \((n-1)\)-st derivative jumps by \(1\), which is what \(\delta\) records.) For the oscillator \(L(D)=D^{2}+\omega^{2}\) this is \(G(0^{+})=0\), \(\dot G(0^{+})=1\), hence \(G(t)=\omega^{-1}\sin(\omega t)\) for \(t>0\).
\end{definition}

\begin{remark}[Scalar Green's function]
\label{rem:scalar-green}
The same statement as \cref{thm:duhamel} for a constant-coefficient scalar operator \(L(D)\) is: pass to the companion system on \(\mathbf{z}=(y,y',\ldots,y^{(n-1)})^{\mathsf T}\), compute \(e^{tA}\), and read the first component. Equivalently, if \(G\) is the causal fundamental solution of \cref{def:dirac} and \(y_{h}\) solves the homogeneous equation with the same initial data, then
\[
y(t)
=
y_{h}(t)+(G*g)(t)
=
y_{h}(t)+\int_{0}^{t} G(t-s)\,g(s)\,\mathrm{d}s
\]
solves \(L(D)y=g\). The kernel \(G\) is the inverse Laplace transform of \(1/L(s)\). For the oscillator this \(G\) is the integral term of \eqref{eq:osc-duhamel} below.
\end{remark}

\begin{example}[Scalar form: the operator \((D^{2}+\omega^{2})^{-1}\) is a sine convolution]
\label{ex:green-osc}
The nonhomogeneous oscillator \(\ddot y+\omega^{2} y=g(t)\), \(y(0)=\epsilon\), \(\dot y(0)=0\), is the companion IVP
\[
\mathbf{z}
=
\begin{pmatrix} y \\ v \end{pmatrix},
\qquad
A
=
\begin{pmatrix} 0 & 1 \\ -\omega^{2} & 0 \end{pmatrix},
\qquad
\mathbf{f}(t)
=
\begin{pmatrix} 0 \\ g(t) \end{pmatrix},
\qquad
\mathbf{z}(0)
=
\begin{pmatrix} \epsilon \\ 0 \end{pmatrix}.
\]
The matrix \(A\) is the same as in \cref{ex:expA-center} with the present \(\omega\) in place of \(\sqrt{6}\), so
\[
e^{\tau A}
=
\begin{pmatrix}
\cos\omega\tau & \omega^{-1}\sin\omega\tau \\
-\omega\sin\omega\tau & \cos\omega\tau
\end{pmatrix}.
\]
Duhamel \eqref{eq:duhamel} is therefore
\[
\mathbf{z}(t)
=
e^{tA}\mathbf{z}(0)
+
\int_{0}^{t} e^{(t-s)A}\mathbf{f}(s)\,\mathrm{d}s.
\]
The homogeneous piece is the first column of \(e^{tA}\) times \(\epsilon\):
\[
e^{tA}\mathbf{z}(0)
=
\epsilon
\begin{pmatrix}
\cos\omega t \\
-\omega\sin\omega t
\end{pmatrix}.
\]
The forcing \(\mathbf{f}(s)\) is the second standard basis vector times \(g(s)\), so \(e^{(t-s)A}\mathbf{f}(s)\) is the second column of \(e^{(t-s)A}\) times \(g(s)\):
\[
e^{(t-s)A}\mathbf{f}(s)
=
g(s)
\begin{pmatrix}
\omega^{-1}\sin\bigl(\omega(t-s)\bigr) \\
\cos\bigl(\omega(t-s)\bigr)
\end{pmatrix}.
\]
The unknown \(y\) is the first component of \(\mathbf{z}\). Reading that row gives
\begin{equation}
\label{eq:osc-duhamel}
y(t)
=
\epsilon\cos\omega t
+
\frac{1}{\omega}\int_{0}^{t}\sin\bigl(\omega(t-s)\bigr)\,g(s)\,\mathrm{d}s.
\end{equation}
The kernel \(\omega^{-1}\sin(\omega(t-s))\) is exactly the \((1,2)\)-entry of \(e^{(t-s)A}\). In operator language, \(D^{2}+\omega^{2}\) is inverted by convolution against that causal Green's function of \cref{def:dirac}: it is the same object as the matrix exponential, written on the original second-order scalar equation.
\end{example}

\begin{example}[A linear convolution that does evaluate]
\label{ex:duhamel-linear}
Specialize \cref{ex:green-osc} to \(\ddot y+6y=1\) with rest data \(y(0)=\dot y(0)=0\). Then \(\omega=\sqrt{6}\), \(\epsilon=0\), and \(g\equiv 1\), so \(\mathbf{z}(0)=\mathbf{0}\) and \(\mathbf{f}(s)=(0,1)^{\mathsf T}\). The Duhamel integral collapses to the second column of \(e^{(t-s)A}\) alone:
\[
\mathbf{z}(t)
=
\int_{0}^{t} e^{(t-s)A}
\begin{pmatrix} 0 \\ 1 \end{pmatrix}
\mathrm{d}s
=
\int_{0}^{t}
\begin{pmatrix}
\omega^{-1}\sin\bigl(\omega(t-s)\bigr) \\
\cos\bigl(\omega(t-s)\bigr)
\end{pmatrix}
\mathrm{d}s.
\]
The first component is the definite integral
\[
y(t)
=
\frac{1}{\omega}\int_{0}^{t}\sin\bigl(\omega(t-s)\bigr)\,\mathrm{d}s.
\]
Substitute \(u=t-s\), so \(\mathrm{d}u=-\mathrm{d}s\) and the limits \(s=0\to t\) become \(u=t\to 0\):
\[
y(t)
=
\frac{1}{\omega}\int_{t}^{0}\sin(\omega u)\,(-\mathrm{d}u)
=
\frac{1}{\omega}\int_{0}^{t}\sin(\omega u)\,\mathrm{d}u
=
\frac{1-\cos\omega t}{\omega^{2}}
=
\frac{1-\cos(\sqrt{6}\,t)}{6}.
\]
The same integral against \(\cos\Omega t\) with \(\Omega\neq\omega\) produces the usual beats. This is the brute-force method working: \(A\) is a matrix, \(e^{tA}\) is a trigonometric polynomial, and the convolution is elementary. The obstruction on the exam is not the linear operator; it is that the forcing \(-5y^{2}-y^{3}\) still contains the unknown.
\end{example}

\begin{example}[The nonlinear remainder, as a Picard integral equation]
\label{ex:duhamel-p6}
After translating \(x=2+y\), the exam equation is \(\ddot y+6y=-5y^{2}-y^{3}\). In companion form,
\[
\dot{\mathbf{z}}
=
A\mathbf{z}+\mathbf{N}(\mathbf{z}),
\qquad
\mathbf{N}\bigl((y,v)^{\mathsf T}\bigr)
=
\begin{pmatrix} 0 \\ -5y^{2}-y^{3} \end{pmatrix},
\]
with the \(A\) of \cref{ex:expA-center}. Duhamel still applies, now with a state-dependent forcing \(\mathbf{f}(s)=\mathbf{N}(\mathbf{z}(s))\):
\[
\mathbf{z}(t)
=
e^{tA}\mathbf{z}(0)
+
\int_{0}^{t} e^{(t-s)A}\mathbf{N}\bigl(\mathbf{z}(s)\bigr)\,\mathrm{d}s.
\]
The scalar version is \eqref{eq:osc-duhamel} with \(g=-5y^{2}-y^{3}\) and \(\omega=\sqrt{6}\). This is an exact integral equation, not an explicit formula: the unknown sits inside the convolution. The general existence theory is Picard iteration on that integral,
\[
\mathbf{z}_{0}(t)=e^{tA}\mathbf{z}(0),
\qquad
\mathbf{z}_{k+1}(t)
=
e^{tA}\mathbf{z}(0)
+
\int_{0}^{t} e^{(t-s)A}\mathbf{N}\bigl(\mathbf{z}_{k}(s)\bigr)\,\mathrm{d}s,
\]
which converges locally by the usual contraction on a short time interval. Substituting the power series \(y=\epsilon\cos\omega t+O(\epsilon^{2})\) into the convolution and integrating term by term is a legitimate brute-force expansion --- and it produces secular terms \(t\sin\omega t\), because the quadratic and cubic interactions resonate with the homogeneous frequency. That is precisely why the expansion is not uniformly valid on times of order \(1/\epsilon^{2}\), and why Lindstedt stretches the time variable in \cref{con:lindstedt}. The operator calculus is the parent of the perturbation; it does not replace it.

For a polynomial nonlinearity the convolution integrals are elementary (products of sines and cosines), so one \emph{can} grind the expansion by hand to any finite order. One cannot sum it, and one cannot read off a closed orbit from the integral equation without a separate topological argument (energy, \cref{prop:energy-periodic}). That is the sense in which the method is general, and the sense in which it is not useful for the two sentences of part~(b).
\end{example}

\begin{remark}[What Yau papers do with the same calculus]
\label{rem:yau-eAt}
The companion-system / \(e^{tA}\) / Duhamel package is the standard written move when an applied or analysis paper asks for a linear system, a nonhomogeneous linear ODE, or a perturbative bound. Typical asks: compute \(e^{tA}\) from a Jordan form (stiff \(2\times 2\) examples); write the variation-of-constants formula and apply Gr\"onwall to a small integrable perturbation \(B(t)\) of a harmonic oscillator; invert \(D^{2}-\lambda^{2}\) or \(D^{2}+\omega^{2}\) by a sinh or sine kernel. Those are this paragraph, with \(\mathbf{N}\) linear in \(\mathbf{z}\) or independent of \(\mathbf{z}\). The present CAM question is the nonlinear conservative case, where the same formula is true and unhelpful as a closed form.
\end{remark}

\begin{example}[Elementary tricks that sometimes finish the problem]
\label{ex:ode-tricks}
The remaining Grade~2 devices are special-structure tests, not a general integration theory.
\begin{enumerate}[label=\textup{(\roman*)}]
    \item \emph{Exact energy / integrating factor in the phase plane.} For \(\ddot x=f(x)\), multiply by \(\dot x\): \(\frac{\mathrm{d}}{\mathrm{d}t}\bigl(\tfrac12\dot x^{2}+V(x)\bigr)=0\). This is the trick that \emph{does} solve the topology of part~(b). It is the statement that \(\mathbf{F}\) is Hamiltonian, equivalently that \(P\,\mathrm{d}x+Q\,\mathrm{d}v=0\) with \(P=f(x)\) and \(Q=-v\) is exact.
    \item \emph{Missing independent or dependent variable.} If \(\ddot x=f(\dot x,t)\) (no \(x\)), set \(p=\dot x\) and obtain a first-order equation for \(p(t)\). If \(\ddot x=f(x,\dot x)\) (no \(t\)), set \(p=\dot x\) and \(p\,\mathrm{d}p/\mathrm{d}x=f(x,p)\), which is \cref{def:reduce-autonomous}. The exam is the second case.
    \item \emph{Bernoulli / homogeneous / linear first-order.} After a reduction to first order, these are the quadrature catalogue of the method map. They do not fire on the original second-order polynomial.
    \item \emph{Undetermined coefficients / Laplace, for linear nonhomogeneous equations.} If one had \(\ddot y+6y=\cos 2t\), a particular solution is a multiple of \(\cos 2t\). The exam forcing is a function of \(y\) itself, so this test is empty until a perturbation expansion has already produced an explicit right-hand side (as in the \(\epsilon^{2}\) equation of the Lindstedt calculation).
    \item \emph{Reduction of order for linear equations when one solution is known.} If \(y_1\) solves a linear second-order equation, \(y=u y_1\) produces a first-order equation for \(u'\). Irrelevant here: the full equation is nonlinear.
\end{enumerate}
\end{example}


\begin{examnote}[What the script must contain]
\label{exnote:p6}
Part~(a) is \(f(x)=0\), three real roots, and (if one has time) \(V''\) at each. Part~(b) is two sentences of energy (the level \(V(2+\epsilon)\) lies in a potential well) plus the Lindstedt algebra through \(O(\epsilon^{3})\), including the explicit \(\cos s\) coefficient \(2\nu+49/108\). Do not quote a limit cycle theorem: there is no dissipation, so the orbit is a period annulus, not an attracting cycle. Do not grind the Duhamel convolution \eqref{eq:osc-duhamel} on the script: it is the general theory, and it does not evaluate.
\end{examnote}

\subsubsection{Solution of Problem 6}

\paragraph{(a). Static solutions.}

A constant solution satisfies \(f(x)=0\) with \(f(x)=2x+x^{2}-x^{3}\), i.e.
\[
x^{3}-x^{2}-2x=x(x^{2}-x-2)=x(x-2)(x+1)=0.
\]
The static solutions are
\[
x\equiv -1,\qquad x\equiv 0,\qquad x\equiv 2.
\]
The potential with \(\ddot x=-V'(x)\) is
\[
V(x)=-x^{2}-\frac{x^{3}}{3}+\frac{x^{4}}{4},
\qquad
V''(x)=-2-2x+3x^{2}.
\]
Thus \(V''(-1)=3>0\), \(V''(0)=-2<0\), \(V''(2)=6>0\): \(-1\) and \(2\) are centers, \(0\) is a saddle. In particular \(V(2)=-8/3\) and \(V(0)=0\).

\paragraph{(b). Closed orbit.}

Translate the center: \(x=2+y\). Then \(y(0)=\epsilon\), \(\dot y(0)=0\), and
\begin{equation}
\label{eq:p6-y}
\ddot y+6y+5y^{2}+y^{3}=0.
\end{equation}
The corresponding potential, normalized to vanish at the translated equilibrium, is
\[
W(y)
\coloneqq
V(2+y)-V(2)
=
3y^{2}+\frac53 y^{3}+\frac14 y^{4},
\]
so \(\ddot y=-W'(y)\) and the conserved energy of \eqref{eq:p6-y} is
\[
E
=
\frac12\dot y^{2}+W(y).
\]
(The original energy \(\frac12\dot x^{2}+V(x)\) differs from this \(E\) by the constant \(V(2)\); level sets are the same curves.) At the initial point one has \(\dot y(0)=0\), hence \(E=W(\epsilon)\). Taylor expansion uses \(W(0)=W'(0)=0\) and \(W''(0)=6\),
\[
E=W(\epsilon)=3\epsilon^{2}+O(\epsilon^{3}).
\]
In particular \(E>0\) for all sufficiently small \(\epsilon\neq 0\), and \(E\to 0\) as \(\epsilon\to 0\). This \(E\) is the energy of the orbit we are studying. It is not yet a statement about closedness: for that one needs a second turning point, which is a comparison with the energy of the adjacent saddle.

\emph{Saddle energy.} An equilibrium \((y_{\star},0)\) lies on the level set of energy \(W(y_{\star})\). That number is the \emph{saddle energy} (or center energy) of that equilibrium: it is simply \(E\) evaluated at a rest point, not an extra construction. Differentiating \(W\) gives the translated equilibria of part~(a),
\[
W'(y)=6y+5y^{2}+y^{3}=y(y+2)(y+3),
\]
so \(y=0\) (center, \(x=2\)), \(y=-2\) (saddle, \(x=0\)), and \(y=-3\) (center, \(x=-1\)). The saddle energy in these coordinates is
\[
W(-2)=V(0)-V(2)=8/3.
\]
In the original coordinates the same level is \(V(x)=V(0)=0\). The phase-plane curve \(\frac12\dot y^{2}+W(y)=8/3\) passes through the saddle \((-2,0)\). Because \(W\to+\infty\) as \(y\to\pm\infty\) and \(W\) has a second well at \(y=-3\), that level set is a figure-eight of two homoclinic loops sharing the saddle: one loop encloses \(y=0\), the other encloses \(y=-3\). Orbits with energy strictly below \(8/3\) cannot reach the saddle, and are trapped in one well or the other.

\emph{Sign chart of \(W'\), and the two turning points.} The factorization of \(W'\) gives
\[
W'<0\text{ on }(-\infty,-3),\qquad
W'>0\text{ on }(-3,-2),\qquad
W'<0\text{ on }(-2,0),\qquad
W'>0\text{ on }(0,+\infty).
\]
Thus \(W\) has a local minimum \(W(0)=0\) at the origin, a local maximum \(W(-2)=8/3\) at the saddle, and a local minimum \(W(-3)=9/4<8/3\) at the other center. Restrict attention to the well around the origin. On \([0,+\infty)\) one has \(W(0)=0\), \(W'>0\), and \(W(y)\to+\infty\), so for each \(E\in(0,8/3)\) there is a unique \(a(E)>0\) with \(W(a(E))=E\). On \([-2,0]\) one has \(W\) strictly decreasing from \(8/3\) to \(0\), so there is a unique \(b(E)\in(-2,0)\) with \(W(b(E))=E\). On the open interval \((b(E),a(E))\) one has \(W<E\) (decrease from \(E\) down to \(0\), then increase from \(0\) up to \(E\)). Immediately to the left of \(b(E)\) (still in \((-2,b(E))\)) and immediately to the right of \(a(E)\), one has \(W>E\).

\begin{center}
\begin{tikzpicture}[scale=0.92, >=Latex]
  \draw[->] (-4.15,0) -- (2.05,0) node[right, font=\small] {\(y\)};
  \draw[->] (0,-0.25) -- (0,3.15) node[above, font=\small] {\(W(y)\)};
  \begin{scope}
    \clip (-4.15,-0.25) rectangle (2.02,3.12);
    \draw[thick, domain=-3.95:0.95, samples=140]
      plot (\x, {3*\x*\x + (5.0/3.0)*\x*\x*\x + 0.25*\x*\x*\x*\x});
  \end{scope}
  \draw[dashed] (-4.05,2.667) -- (0.78,2.667);
  \node[font=\scriptsize, anchor=west] at (1.00,2.667) {\(W(-2)=\tfrac83\)};
  \draw[red!70!black] (-0.584,0.72) -- (0.437,0.72);
  \fill[red!70!black] (0.437,0.72) circle (1.5pt);
  \fill[red!70!black] (-0.584,0.72) circle (1.5pt);
  \node[red!70!black, font=\scriptsize, anchor=south west] at (0.50,0.80) {\(E=W(\epsilon)\)};
  \node[red!70!black, font=\scriptsize, below] at (0.437,-0.02) {\(\epsilon\)};
  \node[red!70!black, font=\scriptsize, below] at (-0.584,-0.02) {\(b\)};
  \fill (-2,2.667) circle (1.5pt);
  \fill (0,0) circle (1.5pt);
  \fill (-3,2.25) circle (1.5pt);
  \node[font=\scriptsize, below] at (-3.0,-0.12) {\(-3\)};
  \node[font=\scriptsize, below] at (-2.0,-0.12) {\(-2\)};
  \node[font=\scriptsize, anchor=north west] at (0.08,-0.06) {\(0\)};
\end{tikzpicture}
\end{center}

\emph{Small \(\epsilon\) places the data in this well.} Since \(W(\epsilon)=3\epsilon^{2}+O(\epsilon^{3})\) and \(W(0)=0\), there exists \(\delta\in(0,2)\) such that \(0<|\epsilon|<\delta\) implies \(0<W(\epsilon)<8/3\). For such \(\epsilon\) the conserved energy of the IVP satisfies \(0<E<8/3\). If \(\epsilon>0\) then \(a(E)=\epsilon\) and the second turning point is \(b=b(E)\in(-2,0)\); if \(\epsilon<0\) then \(b(E)=\epsilon\) and the second turning point is \(a(E)>0\). In either case the two endpoints of \(I=[b(E),a(E)]\) are the initial rest point and one other rest point of the same energy.

\emph{The hypotheses of \cref{prop:energy-periodic}, in the \(y\)-coordinate.} The force in \eqref{eq:p6-y} is polynomial, hence locally Lipschitz, and \(W'=-f\). The exam data are rest initial conditions at \(a=\epsilon\). The second turning point is the other endpoint of \(I=[b(E),a(E)]\): call it \(b\), so \(b=b(E)\) if \(\epsilon>0\) and \(b=a(E)\) if \(\epsilon<0\). Then:
\begin{enumerate}[label=\textup{(\roman*)}]
    \item \(W(b)=W(\epsilon)=E\), and \(W<E\) throughout the open interval between \(\epsilon\) and \(b\);
    \item \(W>E\) immediately outside \(I=[b(E),a(E)]\), so \(I\) is the connected component of \(\{W\le E\}\) containing the origin (compact, and strictly inside \((-2,+\infty)\)).
\end{enumerate}
\Cref{prop:energy-periodic} therefore applies: the level set \(\frac12\dot y^{2}+W(y)=E\) over \(I\) is a simple closed curve around the origin in the \((y,\dot y)\) plane, strictly inside the right-hand homoclinic loop of the saddle, and \(y(t)\) is periodic. Equivalently, in the original coordinates the orbit through \((2+\epsilon,0)\) is closed.

The comparison \(E<8/3\) is the only global ingredient. Without it, a large positive \(\epsilon\) could push \(W(\epsilon)\) past the saddle energy, the interval \(I\) would cease to be trapped in the well at \(0\), and the orbit would no longer be a small oscillation about \(x=2\). Linearization at a center does not by itself give a nonlinear closed orbit; the potential well below the saddle energy does.

\paragraph{(b). Period to \(O(\epsilon^{2})\).}

The linearization at \(y=0\) is \(\ddot y+6y=0\), period \(2\pi/\sqrt{6}\). The nonlinear terms shift that period by \(O(\epsilon^{2})\). Compute the shift by Lindstedt--Poincar\'e (\cref{con:lindstedt}), not by expanding the elliptic integral \eqref{eq:period-quad} (the integrand is singular at the turning points).

\emph{Change of clock.}
Set \(s=\Omega t\), so \(\mathrm{d}/\mathrm{d}t=\Omega\,\mathrm{d}/\mathrm{d}s\) and \(\mathrm{d}^{2}/\mathrm{d}t^{2}=\Omega^{2}\,\mathrm{d}^{2}/\mathrm{d}s^{2}\). Write \('=\mathrm{d}/\mathrm{d}s\). Then \eqref{eq:p6-y} becomes
\begin{equation}
\label{eq:p6-s}
\Omega^{2} y''+6y+5y^{2}+y^{3}=0.
\end{equation}
If \(y(s)\) is \(2\pi\)-periodic, the original period is \(T=2\pi/\Omega\). Expand both the unknown frequency and the unknown profile in the amplitude \(\epsilon=y(0)\):
\begin{equation}
\label{eq:p6-ansatz}
\Omega=\sqrt{6}\,\bigl(1+\nu\epsilon^{2}+O(\epsilon^{4})\bigr),
\qquad
y(s)=\epsilon y_{1}(s)+\epsilon^{2} y_{2}(s)+\epsilon^{3} y_{3}(s)+O(\epsilon^{4}),
\end{equation}
with each \(y_{k}\) \(2\pi\)-periodic. The initial conditions \(y(0)=\epsilon\) and \(\dot y(0)=0\) (hence \(y'(0)=0\)) become
\[
y_{1}(0)=1,\qquad y_{1}'(0)=0,
\qquad
y_{k}(0)=y_{k}'(0)=0\quad(k\ge 2).
\]
Squaring the frequency,
\[
\Omega^{2}
=
6\bigl(1+\nu\epsilon^{2}+O(\epsilon^{4})\bigr)^{2}
=
6\bigl(1+2\nu\epsilon^{2}+O(\epsilon^{4})\bigr).
\]
(The missing \(O(\epsilon)\) term in \(\Omega\) is justified after the \(\epsilon^{2}\) equation: it would be forced to vanish.) The nonlinearities start at quadratic order:
\begin{align*}
y^{2}
&=
\epsilon^{2} y_{1}^{2}+2\epsilon^{3} y_{1}y_{2}+O(\epsilon^{4}),
\\
y^{3}
&=
\epsilon^{3} y_{1}^{3}+O(\epsilon^{4}),
\\
\Omega^{2} y''
&=
6\bigl(1+2\nu\epsilon^{2}\bigr)\bigl(\epsilon y_{1}''+\epsilon^{2} y_{2}''+\epsilon^{3} y_{3}''\bigr)+O(\epsilon^{4})
\\
&=
6\bigl(\epsilon y_{1}''+\epsilon^{2} y_{2}''+\epsilon^{3} y_{3}''+2\nu\epsilon^{3} y_{1}''\bigr)+O(\epsilon^{4}).
\end{align*}
Substitute into \eqref{eq:p6-s} and equate coefficients of like powers of \(\epsilon\).

\emph{Order \(\epsilon\).}
\[
6y_{1}''+6y_{1}=0
\qquad\Longleftrightarrow\qquad
y_{1}''+y_{1}=0.
\]
The unique solution with \(y_{1}(0)=1\), \(y_{1}'(0)=0\) is
\[
y_{1}(s)=\cos s.
\]
This is the linear oscillator in the stretched time \(s\): one cycle in \(s\) is \(2\pi\), hence a period \(2\pi/\Omega\) in \(t\).

\emph{Order \(\epsilon^{2}\).}
The new terms are \(6y_{2}''+6y_{2}+5y_{1}^{2}\). Dividing by \(6\),
\[
y_{2}''+y_{2}
=
-\frac56 y_{1}^{2}.
\]
Use \(\cos^{2}s=\frac{1+\cos 2s}{2}\):
\[
y_{1}^{2}
=
\frac{1+\cos 2s}{2},
\qquad
y_{2}''+y_{2}
=
-\frac{5}{12}-\frac{5}{12}\cos 2s.
\]
No \(\cos s\) appears, so there is no solvability obstruction at this order. (If one had written \(\Omega=\sqrt{6}\,(1+\mu\epsilon+\nu\epsilon^{2}+\cdots)\), the extra term \(2\mu\cos s\) would appear here and force \(\mu=0\).) A particular solution is found by undetermined coefficients. A constant \(C\) satisfies \(C=-5/12\). A term \(B\cos 2s\) satisfies
\[
(B\cos 2s)''+B\cos 2s
=
-4B\cos 2s+B\cos 2s
=
-3B\cos 2s,
\]
so \(-3B=-5/12\) and \(B=5/36\). Thus a particular solution is \(-\frac{5}{12}+\frac{5}{36}\cos 2s\). The general \(2\pi\)-periodic solution adds the homogeneous kernel \(A\cos s+B\sin s\). Differentiating,
\[
y_{2}'(s)
=
-A\sin s+B\cos s-\frac{5}{18}\sin 2s,
\]
so \(y_{2}'(0)=B=0\). Then
\[
y_{2}(0)
=
A-\frac{5}{12}+\frac{5}{36}
=
0
\qquad\Longrightarrow\qquad
A
=
\frac{5}{12}-\frac{5}{36}
=
\frac{5}{18}.
\]
Hence
\begin{equation}
\label{eq:p6-y2}
y_{2}(s)
=
-\frac{5}{12}+\frac{5}{18}\cos s+\frac{5}{36}\cos 2s.
\end{equation}
The \(\cos s\) piece is not a secular term: it is a homogeneous correction fixing the initial condition. The frequency has not yet been determined.

\emph{Order \(\epsilon^{3}\), and the solvability condition.}
Collecting \(\epsilon^{3}\) in \eqref{eq:p6-s} gives
\[
6y_{3}''+6y_{3}+12\nu y_{1}''+10 y_{1}y_{2}+y_{1}^{3}=0.
\]
(The \(12\nu y_{1}''\) is \(6\cdot 2\nu y_{1}''\) from \(\Omega^{2}y''\); the \(10 y_{1}y_{2}\) is \(5\cdot 2 y_{1}y_{2}\) from the quadratic.) Divide by \(6\) and use \(y_{1}''=-\cos s\):
\begin{equation}
\label{eq:p6-y3}
y_{3}''+y_{3}
=
2\nu\cos s-\frac53 y_{1}y_{2}-\frac16 y_{1}^{3}.
\end{equation}
The homogeneous operator \(D^{2}+1\) annihilates \(\operatorname{span}\{\cos s,\sin s\}\). A \(2\pi\)-periodic \(y_{3}\) exists if and only if the right-hand side is orthogonal to that kernel, i.e.\ if the coefficients of \(\cos s\) and \(\sin s\) vanish. The problem is even in \(s\) (every forcing term is a cosine), so the \(\sin s\) coefficient is automatically zero. It remains to extract the coefficient of \(\cos s\).

The cubic identity \(\cos^{3}s=\frac{3\cos s+\cos 3s}{4}\) gives
\[
-\frac16 y_{1}^{3}
=
-\frac16\cdot\frac{3\cos s+\cos 3s}{4}
=
-\frac18\cos s-\frac{1}{24}\cos 3s.
\]
The quadratic interaction, using \eqref{eq:p6-y2} and \(\cos s\cos 2s=\frac12(\cos 3s+\cos s)\), is
\begin{align*}
y_{1}y_{2}
&=
\cos s\Bigl(-\frac{5}{12}+\frac{5}{18}\cos s+\frac{5}{36}\cos 2s\Bigr)
\\
&=
-\frac{5}{12}\cos s+\frac{5}{18}\cdot\frac{1+\cos 2s}{2}+\frac{5}{36}\cdot\frac{\cos 3s+\cos s}{2}
\\
&=
-\frac{5}{12}\cos s+\frac{5}{36}+\frac{5}{36}\cos 2s+\frac{5}{72}\cos 3s+\frac{5}{72}\cos s.
\end{align*}
The \(\cos s\) terms in \(y_{1}y_{2}\) are \(-\frac{5}{12}+\frac{5}{72}=-\frac{25}{72}\). Therefore the \(\cos s\) coefficient in \(-\frac53 y_{1}y_{2}\) is
\[
-\frac53\Bigl(-\frac{25}{72}\Bigr)
=
\frac{125}{216}.
\]
Adding the cubic contribution \(-\frac18=-\frac{27}{216}\),
\[
\frac{125}{216}-\frac{27}{216}
=
\frac{98}{216}
=
\frac{49}{108}.
\]
(The \(\frac{125}{216}\) comes from \(5y^{2}\) through \(y_{1}y_{2}\); the \(-\frac{27}{216}\) comes from \(y^{3}\). Dropping the quadratic term in \eqref{eq:p6-y} would replace \(\frac{49}{108}\) by \(-\frac18\) and give the wrong period.) The right-hand side of \eqref{eq:p6-y3} therefore contains
\[
\Bigl(2\nu+\frac{49}{108}\Bigr)\cos s
\]
plus non-resonant terms \(\mathrm{const}+\cos 2s+\cos 3s\). The solvability condition is
\[
2\nu+\frac{49}{108}=0
\qquad\Longrightarrow\qquad
\nu=-\frac{49}{216}.
\]
One can stop here: \(y_{3}\) exists and is \(2\pi\)-periodic, but it is not needed for \(T\).

\emph{Back to the original period.}
With \(\nu=-49/216\),
\[
\Omega
=
\sqrt{6}\Bigl(1+\nu\epsilon^{2}\Bigr)+O(\epsilon^{4})
=
\sqrt{6}\Bigl(1-\frac{49}{216}\epsilon^{2}\Bigr)+O(\epsilon^{4}).
\]
Inverting,
\[
T
=
\frac{2\pi}{\Omega}
=
\frac{2\pi}{\sqrt{6}}\bigl(1+\nu\epsilon^{2}\bigr)^{-1}+O(\epsilon^{4})
=
\frac{2\pi}{\sqrt{6}}\bigl(1-\nu\epsilon^{2}\bigr)+O(\epsilon^{4}).
\]
Since \(-\nu=49/216\), this is a slightly longer period than the linear value \(2\pi/\sqrt{6}\) (the well is anharmonic, and the quadratic term \(5y^{2}\) already lengthens the orbit on the side \(y>0\)):
\begin{equation}
\label{eq:p6-period}
T
=
\frac{2\pi}{\sqrt{6}}\Bigl(1+\frac{49}{216}\epsilon^{2}\Bigr)+O(\epsilon^{4}).
\end{equation}
The exam asks for an approximation up to \(O(\epsilon^{2})\), which is \eqref{eq:p6-period}. The remainder is \(O(\epsilon^{4})\) rather than \(O(\epsilon^{3})\) because the next frequency correction sits at \(\epsilon^{4}\).

\begin{summary}[Problem 6]
\label{sum:p6}
Static solutions \(x\equiv -1,0,2\). The point \(x=2\) is a center of \(V\). For small \(\epsilon\neq 0\) the energy \(V(2+\epsilon)\) sits in that well, so the orbit is closed. The period is \(\frac{2\pi}{\sqrt{6}}\bigl(1+\frac{49}{216}\epsilon^{2}\bigr)+O(\epsilon^{4})\).
\end{summary}

\subsubsection{Review: ODE exam points on past papers}
\label{sec:p6-ode-review}

The CAM syllabus has two ODE boxes. They are not the same course, and they have not been mixed on a single problem.

\begin{itemize}
    \item \emph{Mathematical models / asymptotics}: regular and singular perturbation, multiple scales, boundary layers, stationary phase \cite{bender99,nayfeh73}. This is the optional Problem~6 slot (against an optimization Problem~7). The present paper is the conservative-oscillator member of that family.
    \item \emph{Numerical solution of ODEs}: one-step and multistep methods, consistency, zero-stability, convergence, absolute stability, stiffness \cite{hairer93,hairer96,leveque07}. This has appeared as a standalone question only once (2023 Fall Problem~4, a Runge--Kutta tableau). Everywhere else it is the scalar test equation inside a PDE von Neumann computation (Problems~4--5 of the present paper, and the analogous slots in 2023--2025). The numerical half of this review is the ODE IVP syllabus, written at the same density as the rest of the notes; the PDE treatise stays in Problem~4.
\end{itemize}

The table is the inventory, 2023 Fall through 2026 Spring. ``P6-type'' means the optional applied-analysis question, not the problem number.

{\footnotesize
\begin{center}
\begin{tabular}{@{}>{\raggedright\arraybackslash}p{2.7cm}
                   >{\raggedright\arraybackslash}p{4.4cm}
                   >{\raggedright\arraybackslash}p{3.4cm}
                   >{\raggedright\arraybackslash}p{3.6cm}@{}}
\toprule
Paper & Question & Machine & Relation to 2026 P6(b) \\
\midrule
2023 Fall P4 & RK tableau, order, abs.\ stability & one-step consistency \(+\) \(R(z)\) & numerical ODE, not this orbit \\
2023 Fall P5 & Lax--Friedrichs for \(u_t+a u_{xxx}\) & PDE truncation \(+\) von Neumann & P4 of this paper \\
2023 Fall P7 & \(\epsilon\)-perturbation of a Sturm--Liouville \(\lambda\) & regular perturbation / Fredholm & different P6-type asymptotics \\
2024 Spring P4 & \(\theta\)-method for heat & truncation \(+\) Jury & P4 of this paper \\
2024 Spring & no oscillator / no RK & --- & P6/P7 were convexity and LP \\
2024 Fall P6 & \(\epsilon y''+\epsilon a(x)y'-y=g\), two BCs & matched inner/outer, width \(\sqrt{\epsilon}\) & singular perturbation, not Lindstedt \\
2025 Spring P6 & \(y''+y+\epsilon(y')^{3}=0\), large \(t\) & multiple scales, decaying amplitude & same oscillator class, damping \\
2025 Fall P6 & Rayleigh oscillator, large \(t\) & multiple scales, limit cycle & same class, weakly dissipative \\
2026 Spring P6 & conservative cubic, closed orbit \(+\) period & energy well \(+\) Lindstedt & this problem \\
2026 Spring P4--5 & implicit advection / max principle & PDE \(g(\xi)\), not ODE \(R(z)\) & uses the ODE test equation as a lemma \\
\bottomrule
\end{tabular}
\end{center}
}

Each exam point below is a theorem (or definition), then a worked example, then an exercise with a written solution. Stationary phase has not yet appeared as a P6-type question; it is still a syllabus item, so it is included in the same format.

{\footnotesize
\begin{center}
\begin{tabular}{@{}>{\raggedright\arraybackslash}p{3.3cm}
                   >{\raggedright\arraybackslash}p{3.7cm}
                   >{\raggedright\arraybackslash}p{3.5cm}
                   >{\raggedright\arraybackslash}p{3.5cm}@{}}
\toprule
Exam point & Theorem / definition & Example & Exercise \(+\) solution \\
\midrule
Energy \(\Rightarrow\) closed orbit & \cref{prop:energy-periodic} & this paper, P6(b) & --- (computed in the solution) \\
Lindstedt period & \cref{con:lindstedt,thm:fredholm-osc} & this paper, P6(b) & \cref{exer:duffing-period} \\
Multiple scales, damping & \cref{con:multiple-scales} & \cref{ex:ms-cubic-2025} & \cref{exer:ms-rayleigh} \\
Boundary layer / matching & \cref{thm:distinguished-balance} & \cref{ex:bl-model,ex:bl-2024fall} & \cref{exer:bl-other-end,exer:bl-2024-alpha0} \\
Eigenvalue perturbation & \cref{thm:eig-regular} & \cref{ex:eig-2023} & \cref{exer:eig-f1} \\
Stationary phase & \cref{thm:stat-phase} & \cref{ex:fresnel} & \cref{exer:bessel-phase} \\
One-step consistency / order & \cref{def:ode-consistency,thm:onestep-conv} & \cref{ex:fe-onestep,ex:heun-2023} & \cref{exer:trap-order} \\
Absolute / \(A\)-stability & \cref{def:abs-stab} & \cref{ex:heun-2023,ex:stab-functions} & \cref{exer:be-astable,exer:rk4-not-astable} \\
Linear multistep / zero-stab.\ & \cref{def:lmm,thm:root-condition,thm:dahlquist-eq} & \cref{ex:leapfrog-weak,ex:ab2,ex:bdf2} & \cref{exer:bdf2-order} \\
Stiffness & \cref{def:stiff,thm:dahlquist-barriers} & \cref{ex:stiff-proto} & \cref{exer:fe-stiff} \\
\bottomrule
\end{tabular}
\end{center}
}

\paragraph{Diagnosis: which perturbation machine.}

The first sentence of a P6-type script is the diagnosis, not the expansion.

\begin{enumerate}[label=\textup{(\roman*)}]
    \item No \(\epsilon\), second-order, no \(\dot x\) in the force: energy \(+\cref{prop:energy-periodic}\). Closed orbits are level sets. This is part~(b)'s topology.
    \item Small \(\epsilon\), conservative, a genuine periodic orbit whose period is asked: Lindstedt--Poincar\'e (\cref{con:lindstedt}). Stretch \(s=\Omega t\) and kill the resonant \(\cos s\) at the first order that produces one. That is part~(b)'s period.
    \item Small \(\epsilon\), damping or negative damping, the question says ``valid for large \(t\)'': two-time multiple scales (\cref{con:multiple-scales}). A regular expansion or a Lindstedt expansion in the original time is secular on \(t\sim 1/\epsilon\). Amplitude (and possibly phase) evolve on \(T=\epsilon t\). 2025 Spring and 2025 Fall.
    \item Small \(\epsilon\) multiplying the highest derivative, with boundary conditions at both ends: matched asymptotics (\cref{thm:distinguished-balance}). The reduced equation cannot meet both BCs. 2024 Fall.
    \item Small \(\epsilon\) in an eigenvalue: regular perturbation, solvability against the adjoint eigenfunction (\cref{thm:eig-regular}). 2023 Fall P7.
\end{enumerate}

\begin{warning}[Lindstedt is not multiple scales]
\label{warn:lindstedt-vs-ms}
Lindstedt and multiple scales both remove secular terms. They are not interchangeable. Lindstedt assumes a nearby periodic orbit and corrects its frequency; the amplitude is fixed by the initial energy. Multiple scales lets the amplitude run, and is the wrong first sentence on the present conservative problem (the orbit is closed, there is nothing to decay). Conversely, quoting a limit-cycle theorem on a conservative cubic is the error recorded in \cref{exnote:p6}.
\end{warning}

\paragraph{Fredholm solvability for a forced oscillator.}

Every frequency correction in Lindstedt, and every amplitude law in multiple scales, is the same linear algebra: a periodic solution of \(y''+\omega^{2}y=f\) exists if and only if \(f\) is orthogonal to the homogeneous kernel.

\begin{theorem}[Periodic solvability]
\label{thm:fredholm-osc}
Let \(\omega>0\) and let \(f\) be continuous and \(2\pi/\omega\)-periodic. The equation
\[
y''+\omega^{2} y=f(t)
\]
admits a \(2\pi/\omega\)-periodic solution if and only if
\[
\int_{0}^{2\pi/\omega}f(t)\cos(\omega t)\,\mathrm{d}t
=
\int_{0}^{2\pi/\omega}f(t)\sin(\omega t)\,\mathrm{d}t
=
0.
\]
If the two moments vanish, a periodic solution exists and is unique up to an element of the two-dimensional kernel \(\operatorname{span}\{\cos(\omega t),\sin(\omega t)\}\).
\end{theorem}

\begin{proof}
Variation of constants against \(\{\cos(\omega t),\sin(\omega t)\}\), or equivalently the Fourier coefficient of \(e^{\pm\mathrm{i}\omega t}\) in \(f\). A resonant forcing \(\cos(\omega t)\) or \(\sin(\omega t)\) produces a particular solution of the form \(t\sin(\omega t)\) or \(t\cos(\omega t)\), which is not periodic. Vanishing of those two coefficients is necessary and sufficient.
\end{proof}

\begin{remark}
\label{rem:fredholm-p6}
In the Lindstedt clock \(s=\Omega t\), the homogeneous operator is \(\mathrm{d}^{2}/\mathrm{d}s^{2}+1\), and \cref{thm:fredholm-osc} is the demand that the coefficients of \(\cos s\) and \(\sin s\) vanish at each order. That is how \(\nu\) is fixed in part~(b). In multiple scales the same obstruction, now with slowly varying amplitude, produces the modulation ODE for \(R(T)\).
\end{remark}

\paragraph{Conservative oscillator and Lindstedt.}

The topology of a closed orbit is \cref{prop:energy-periodic}; the period to \(O(\epsilon^{2})\) is \cref{con:lindstedt}. Both are computed in the solution of this paper's Problem~6. The drill below is the standard cubic oscillator with a frequency shift already at \(O(\epsilon)\), so that the solvability step is visible in one order rather than three.

\begin{example}[This paper, Problem~6(b)]
\label{ex:p6-lindstedt}
The equation is \(\ddot x-2x-x^{2}+x^{3}=0\). After \(x=2+y\), the potential well \(W(y)=3y^{2}+\frac53 y^{3}+\frac14 y^{4}\) has a centre at \(y=0\) and a saddle at \(y=-2\) of energy \(W(-2)=\frac83\). For \(0<W(\epsilon)<\frac83\) the level set is a simple closed curve (\cref{prop:energy-periodic}). Lindstedt with \(\omega_{0}=\sqrt{6}\) produces
\[
T
=
\frac{2\pi}{\sqrt{6}}\Bigl(1+\frac{49}{216}\epsilon^{2}\Bigr)+O(\epsilon^{4}).
\]
The quadratic term does not resonate; the first frequency correction is \(O(\epsilon^{2})\). See the solution of Problem~6(b) and \cref{con:lindstedt}.
\end{example}

\begin{exercise}[Duffing period to \(O(\epsilon)\)]
\label{exer:duffing-period}
Consider
\[
\ddot y+y+\epsilon y^{3}=0,
\qquad
y(0)=a,\quad \dot y(0)=0,
\]
with \(\epsilon\) small and \(a\) fixed. Using Lindstedt (\(s=\Omega t\), \(\Omega=1+\epsilon\omega_{1}+O(\epsilon^{2})\), \(y=a\cos s+\epsilon y_{1}(s)+O(\epsilon^{2})\), and \(y_{1}\) \(2\pi\)-periodic), show that
\[
\omega_{1}=\frac{3a^{2}}{8},
\qquad
T
=
2\pi\Bigl(1-\frac{3\epsilon a^{2}}{8}\Bigr)+O(\epsilon^{2}).
\]
\end{exercise}

\begin{proof}[Solution]
Write \(\mathrm{d}^{2}/\mathrm{d}t^{2}=\Omega^{2}\,\mathrm{d}^{2}/\mathrm{d}s^{2}\). Substituting the expansions into \(\Omega^{2} y''+y+\epsilon y^{3}=0\) and using \(\Omega^{2}=1+2\epsilon\omega_{1}+O(\epsilon^{2})\) gives, at order \(\epsilon^{0}\), \(y_{0}''+y_{0}=0\) with \(y_{0}(0)=a\), \(y_{0}'(0)=0\), hence \(y_{0}=a\cos s\). At order \(\epsilon^{1}\),
\[
y_{1}''+y_{1}
=
-2\omega_{1} y_{0}''-y_{0}^{3}
=
2\omega_{1} a\cos s-a^{3}\cos^{3}s.
\]
The identity \(\cos^{3}s=\frac{3\cos s+\cos 3s}{4}\) puts the right-hand side in the form
\[
\Bigl(2\omega_{1} a-\frac{3a^{3}}{4}\Bigr)\cos s-\frac{a^{3}}{4}\cos 3s.
\]
By \cref{thm:fredholm-osc} a \(2\pi\)-periodic \(y_{1}\) exists if and only if the coefficient of \(\cos s\) vanishes (the \(\sin s\) coefficient is automatically zero by evenness of the initial data), so \(\omega_{1}=3a^{2}/8\). Then \(\Omega=1+3\epsilon a^{2}/8+O(\epsilon^{2})\) and
\[
T=\frac{2\pi}{\Omega}=2\pi\Bigl(1-\frac{3\epsilon a^{2}}{8}\Bigr)+O(\epsilon^{2}).
\]
(The cubic Duffing shift is already \(O(\epsilon)\) because \(y_{0}^{3}\) contains a \(\cos s\) piece. In Problem~6 the first interaction is quadratic and that piece is absent, which is why the correction there is \(O(\epsilon^{2})\).)
\end{proof}

\paragraph{Multiple scales.}

\begin{construction}[Two-time expansion]
\label{con:multiple-scales}
For a weakly perturbed oscillator \(\ddot y+y=\epsilon g(y,\dot y)\) whose solutions are \emph{not} expected to stay at constant amplitude, introduce \(T=\epsilon t\) and expand
\[
y(t)=Y_{0}(t,T)+\epsilon Y_{1}(t,T)+\cdots.
\]
The chain rule is
\[
\frac{\mathrm{d}}{\mathrm{d}t}=\partial_{t}+\epsilon\partial_{T},
\qquad
\frac{\mathrm{d}^{2}}{\mathrm{d}t^{2}}=\partial_{tt}+2\epsilon\partial_{tT}+O(\epsilon^{2}).
\]
Order \(\epsilon^{0}\) is the harmonic oscillator in \(t\), with amplitude and phase depending on \(T\):
\[
Y_{0}(t,T)=R(T)\cos\bigl(t+\varphi(T)\bigr).
\]
Order \(\epsilon^{1}\) is an inhomogeneous oscillator. A solution \(Y_{1}\) that remains \(O(1)\) on intervals of length \(O(1/\epsilon)\) exists if and only if the coefficients of \(\sin\theta\) and \(\cos\theta\) vanish (\cref{thm:fredholm-osc}). Those two conditions are ODEs for \(R(T)\) and \(\varphi(T)\). The range of validity after one step is \(T=O(1)\), i.e.\ \(t=O(1/\epsilon)\): the expansion is not claimed for all \(t\ge 0\), and it is not a convergent series.
\end{construction}

\begin{example}[2025 Spring Problem~6, cubic damping]
\label{ex:ms-cubic-2025}
The 2025 Spring oscillator is this problem with a damping term instead of a potential well:
\[
\ddot y+y+\epsilon(\dot y)^{3}=0,
\qquad
y(0)=1,\quad \dot y(0)=0,
\qquad
\epsilon>0\text{ small}.
\]
The linearization is the same harmonic oscillator. The cubic is odd in velocity, so it dissipates energy:
\[
\frac{\mathrm{d}}{\mathrm{d}t}\Bigl(\tfrac12\dot y^{2}+\tfrac12 y^{2}\Bigr)=-\epsilon(\dot y)^{4}\le 0.
\]
There is no closed orbit to feed to Lindstedt. A regular expansion \(y=\cos t+\epsilon y_{1}(t)+\cdots\) produces a resonant forcing \(\sin t\) at order \(\epsilon\), hence a secular term \(t\cos t\), which is \(O(1)\) already at \(t\sim 1/\epsilon\).

Apply \cref{con:multiple-scales}. The initial data force \(R(0)=1\) and \(\varphi(0)=0\). Write \(\theta=t+\varphi(T)\). Then \(\partial_{t}Y_{0}=-R\sin\theta\) and
\[
\partial_{tT}Y_{0}=-R'\sin\theta-R\varphi'\cos\theta.
\]
Order \(\epsilon^{1}\) is
\[
\partial_{tt}Y_{1}+Y_{1}
=
-2\partial_{tT}Y_{0}-(\partial_{t}Y_{0})^{3}
=
2R'\sin\theta+2R\varphi'\cos\theta+R^{3}\sin^{3}\theta.
\]
The identity \(\sin^{3}\theta=\frac{3\sin\theta-\sin 3\theta}{4}\) puts the right-hand side in the form
\[
\Bigl(2R'+\tfrac34 R^{3}\Bigr)\sin\theta+2R\varphi'\cos\theta-\tfrac14 R^{3}\sin 3\theta.
\]
By \cref{thm:fredholm-osc},
\[
R'=-\frac38 R^{3},
\qquad
\varphi'=0.
\]
With \(R(0)=1\),
\[
R(T)
=
\frac{1}{\sqrt{1+\tfrac34 T}}
=
\frac{1}{\sqrt{1+\tfrac34\epsilon t}},
\qquad
\varphi\equiv 0.
\]
The leading approximation is therefore
\[
y(t)
\sim
\frac{\cos t}{\sqrt{1+\tfrac34\epsilon t}},
\]
uniformly on intervals of length \(O(1/\epsilon)\). That is the accuracy statement 2025 Spring asked for: the error after one multiple-scale step is \(O(\epsilon)\) on \(t\le C/\epsilon\), not \(O(\epsilon)\) for all \(t\ge 0\).
\end{example}

\begin{exercise}[2025 Fall Problem~6, Rayleigh amplitude]
\label{exer:ms-rayleigh}
The Rayleigh oscillator on that paper is
\[
\ddot y+y
=
\epsilon\Bigl(\dot y-\tfrac13(\dot y)^{3}\Bigr),
\qquad
y(0)=0,\quad \dot y(0)=2a,
\]
with \(a>0\) and \(\epsilon>0\) small. Using \cref{con:multiple-scales} with \(Y_{0}=R(T)\cos\bigl(t+\varphi(T)\bigr)\) and \(\theta=t+\varphi(T)\), derive the modulation law
\[
R'
=
\frac{R}{2}\Bigl(1-\frac{R^{2}}{4}\Bigr),
\qquad
\varphi'=0.
\]
Conclude that \(R=0\) is unstable, \(R=2\) is attracting, and every trajectory with \(a>0\) is drawn to the limit cycle of amplitude \(2\) in \(\dot y\) on the time scale \(t\sim 1/\epsilon\). What is \(R(0)\)?
\end{exercise}

\begin{proof}[Solution]
As in \cref{ex:ms-cubic-2025}, \(\partial_{t}Y_{0}=-R\sin\theta\) and \(\partial_{tT}Y_{0}=-R'\sin\theta-R\varphi'\cos\theta\). The \(O(\epsilon)\) equation is
\[
\partial_{tt}Y_{1}+Y_{1}
=
-2\partial_{tT}Y_{0}+\partial_{t}Y_{0}-\tfrac13(\partial_{t}Y_{0})^{3}.
\]
The right-hand side equals
\[
2R'\sin\theta+2R\varphi'\cos\theta-R\sin\theta+\tfrac13 R^{3}\sin^{3}\theta,
\]
because \((\partial_{t}Y_{0})^{3}=-R^{3}\sin^{3}\theta\). Using \(\sin^{3}\theta=\frac{3\sin\theta-\sin 3\theta}{4}\),
\[
\tfrac13 R^{3}\sin^{3}\theta
=
\tfrac14 R^{3}\sin\theta-\tfrac{1}{12}R^{3}\sin 3\theta.
\]
The coefficient of \(\sin\theta\) is \(2R'-R+\frac14 R^{3}\); the coefficient of \(\cos\theta\) is \(2R\varphi'\). Setting both to zero gives \(\varphi'=0\) and
\[
2R'=R-\tfrac14 R^{3}
\qquad\Longleftrightarrow\qquad
R'=\frac{R}{2}\Bigl(1-\frac{R^{2}}{4}\Bigr).
\]
Equilibria: \(R=0\) and \(R=\pm 2\). The linearization of \(R'=\frac R2(1-R^{2}/4)\) has derivative \(\frac12-\frac{3R^{2}}{8}\), equal to \(+\frac12\) at \(R=0\) (unstable) and to \(-1\) at \(R=2\) (attracting on \(R>0\)). The initial velocity \(\dot y(0)=2a\) with \(y(0)=0\) is \(Y_{0}(0,0)=R(0)\cos\varphi(0)=0\) and \(\partial_{t}Y_{0}(0,0)=-R(0)\sin\varphi(0)=2a\). The choice \(\varphi(0)=-\pi/2\), \(R(0)=2a\) (or \(\varphi(0)=\pi/2\), \(R(0)=-2a\), same orbit) matches the data. Thus \(R(T)\to 2\) as \(T\to\infty\) whenever \(a>0\). Lindstedt can describe the cycle itself after translating to a neighbourhood of \(R=2\); it cannot produce the attraction. That is why the 2025 Fall prompt is multiple scales, and why one must not quote a limit-cycle theorem on the present conservative problem.
\end{proof}

\paragraph{Boundary layers and matched asymptotics.}

\begin{theorem}[Distinguished inner balance]
\label{thm:distinguished-balance}
Consider a linear second-order BVP on a fixed interval, with a small parameter \(\epsilon>0\) multiplying the highest derivative, and with one boundary condition at each end. The reduced problem (\(\epsilon=0\)) is of lower order and cannot in general meet both boundary values. An \emph{outer} expansion is a regular expansion of that reduced equation. An \emph{inner} expansion is obtained by stretching \(\xi=(x-x_{*})/\epsilon^{\nu}\) about an endpoint (or an interior turning point) \(x_{*}\), substituting, and choosing the unique \(\nu>0\) for which at least two terms in the inner equation share the lowest power of \(\epsilon\) and all remaining terms are strictly smaller. The inner solution is required to meet the lost boundary condition at \(\xi=0\) and to match the outer solution as \(\xi\to\infty\) (or \(\xi\to-\infty\)). The composite expansion is outer plus inner minus the common part. The layer sits at the endpoint where the outer profile fails the boundary condition; if it fails both, there is a layer at each end. The inner width is \(\epsilon^{\nu}\), not an arbitrary small scale.
\end{theorem}

\begin{proof}
This is the matching recipe of Bender--Orszag \cite[Ch.~9]{bender99}. The algebraic content is the comparison of exponents after stretching: if the two largest inner terms are \(\epsilon^{a}Y''\) and \(\epsilon^{b}Y\) (or \(\epsilon^{c}Y'\)), set \(a=b\) (resp.\ \(a=c\)) to obtain \(\nu\). A wrong \(\nu\) either drops the second derivative (the layer cannot meet a Dirichlet datum) or retains a term that is asymptotically larger than the ones kept (the expansion is not ordered).
\end{proof}

\begin{example}[Model layer of width \(\epsilon\)]
\label{ex:bl-model}
The model that isolates the mechanism is
\[
\epsilon y''+y'=0,
\qquad
y(0)=0,\quad y(1)=1.
\]
Outer (\(\epsilon=0\)): \(y'=0\), so \(y\equiv C\). The right-hand boundary \(y(1)=1\) is compatible with a constant, the left-hand one is not, so the layer is at \(x=0\) (the inflow, given \(+y'\)). Thus \(y_{\mathrm{out}}\equiv 1\). Inner: set \(\xi=x/\epsilon^{\nu}\). Then
\[
\epsilon^{1-2\nu}Y''+\epsilon^{-\nu}Y'=0.
\]
The distinguished balance \(1-2\nu=-\nu\) gives \(\nu=1\), i.e.\ width \(\epsilon\). The inner equation is \(Y''+Y'=0\) with \(Y(0)=0\) and matching \(Y(\infty)=1\), so \(Y=1-e^{-\xi}\). The composite expansion (outer plus inner minus the common part \(1\)) is
\[
y(x)
\sim
1-e^{-x/\epsilon}.
\]
The exact solution is \(\bigl(1-e^{-x/\epsilon}\bigr)/\bigl(1-e^{-1/\epsilon}\bigr)\), which differs from this by an \(O(e^{-1/\epsilon})\) term.
\end{example}

\begin{exercise}[Layer on the other end]
\label{exer:bl-other-end}
For
\[
\epsilon y''-y'=0,
\qquad
y(0)=0,\quad y(1)=1,
\]
locate the layer, determine its width, and construct the leading composite expansion.
\end{exercise}

\begin{proof}[Solution]
Outer: \(-y'=0\), so \(y\equiv C\). Now the left-hand boundary is compatible with a constant and the right-hand one is not, so the layer is at \(x=1\). Thus \(y_{\mathrm{out}}\equiv 0\). Inner: \(\xi=(1-x)/\epsilon^{\nu}\). Then \(\mathrm{d}/\mathrm{d}x=-\epsilon^{-\nu}\mathrm{d}/\mathrm{d}\xi\) and
\[
\epsilon^{1-2\nu}Y''-(-\epsilon^{-\nu}Y')=0,
\]
i.e.\ \(\epsilon^{1-2\nu}Y''+\epsilon^{-\nu}Y'=0\). The same balance \(1-2\nu=-\nu\) gives \(\nu=1\). The inner equation is \(Y''+Y'=0\), now with \(Y(0)=1\) (that is \(y(1)=1\)) and \(Y(\infty)=0\) (match to the outer value \(0\)), so \(Y=e^{-\xi}\). Composite:
\[
y(x)\sim e^{-(1-x)/\epsilon}.
\]
(The sign in front of \(y'\) selects the inflow endpoint. Getting the end wrong is the usual CAM error on this item.)
\end{proof}

\begin{example}[2024 Fall Problem~6, \(\alpha=1\)]
\label{ex:bl-2024fall}
The 2024 Fall operator is different in one structural respect: both \(y''\) and \(y'\) already carry an \(\epsilon\),
\[
\epsilon y''+\epsilon(1+x)^{2} y'-y=x-1,
\qquad
y(0)=\alpha,\quad y(1)=-1.
\]
The reduced equation is algebraic, \(-y=x-1\), hence \(y_{\mathrm{out}}=1-x\). A distinguished inner balance at an endpoint \(x_{*}\) with \(\xi=(x-x_{*})/\epsilon^{\nu}\) compares \(\epsilon\cdot\epsilon^{-2\nu}Y''\) against the undifferentiated \(-Y\): the first-derivative term is \(\epsilon\cdot\epsilon^{-\nu}=\epsilon^{1-\nu}\), which is \(o(1)\) once \(\nu=\frac12\). Thus \(1-2\nu=0\) gives \(\nu=\frac12\), i.e.\ width \(\sqrt{\epsilon}\), not \(\epsilon\). The outer profile already satisfies \(y(0)=1\); it fails \(y(1)=-1\). For \(\alpha=1\) there is a single layer at \(x=1\).

Stretch \(\xi=(1-x)/\sqrt{\epsilon}\). Then
\[
Y''-\sqrt{\epsilon}\,(1+x)^{2} Y'-Y
=
x-1
=
-\sqrt{\epsilon}\,\xi.
\]
The leading inner equation is \(Y''-Y=0\). Boundedness as \(\xi\to+\infty\) (into the outer region) kills \(e^{+\xi}\). Matching \(Y(\infty)=y_{\mathrm{out}}(1)=0\) and the boundary value \(Y(0)=-1\) give \(Y=-e^{-\xi}\). The composite expansion is
\[
y(x)
\sim
(1-x)-e^{-(1-x)/\sqrt{\epsilon}}.
\]
At \(x=1\) this is \(-1\); at \(x=0\) it is \(1-e^{-1/\sqrt{\epsilon}}\), which is \(1\) up to a transcendentally small term, hence consistent with \(\alpha=1\). The first omitted inner terms are \(O(\sqrt{\epsilon})\), which is the formal accuracy asked in part~(iii) of that paper if one stops at this order.
\end{example}

\begin{exercise}[2024 Fall Problem~6, \(\alpha=0\)]
\label{exer:bl-2024-alpha0}
For the same BVP as in \cref{ex:bl-2024fall} with \(\alpha=0\), construct a leading composite expansion, including layers at both ends.
\end{exercise}

\begin{proof}[Solution]
The outer profile \(y_{\mathrm{out}}=1-x\) now fails both boundary conditions. The right-hand layer is as in \cref{ex:bl-2024fall}: \(Y_{\mathrm{right}}=-e^{-(1-x)/\sqrt{\epsilon}}\). At \(x=0\) stretch \(\eta=x/\sqrt{\epsilon}\). The leading inner equation is
\[
Y''-Y
=
-1,
\]
because the original right-hand side equals \(x-1=-1+O(\sqrt{\epsilon})\). A particular solution is \(Y\equiv 1\); the homogeneous solutions are \(e^{\pm\eta}\). Matching \(Y(\infty)=y_{\mathrm{out}}(0)=1\) kills \(e^{+\eta}\) and fixes the particular value \(1\). The boundary condition \(Y(0)=0\) then gives \(Y=1-e^{-\eta}\). The two layers do not overlap at leading order. Subtracting common parts (the outer value \(1\) at the left, \(0\) at the right) produces the composite
\[
y(x)
\sim
(1-x)-e^{-x/\sqrt{\epsilon}}-e^{-(1-x)/\sqrt{\epsilon}}.
\]
At \(x=0\): \(1-1-e^{-1/\sqrt{\epsilon}}\sim 0\). At \(x=1\): \(0-e^{-1/\sqrt{\epsilon}}-1\sim -1\).
\end{proof}

\paragraph{Regular perturbation of an eigenvalue.}

\begin{theorem}[Simple eigenvalue, self-adjoint perturbation]
\label{thm:eig-regular}
Let \(L\) be self-adjoint on a Hilbert space, with a simple eigenvalue \(\lambda_{0}\) and eigenfunction \(u_{0}\neq 0\). Let \(B\) be bounded and self-adjoint, and consider \((L+\epsilon B)u=\lambda u\). There exist expansions \(\lambda=\lambda_{0}+\epsilon\lambda^{(1)}+O(\epsilon^{2})\) and \(u=u_{0}+\epsilon u^{(1)}+O(\epsilon^{2})\) with
\[
\lambda^{(1)}
=
\frac{\langle Bu_{0},u_{0}\rangle}{\langle u_{0},u_{0}\rangle}.
\]
\end{theorem}

\begin{proof}
Substitute and collect \(O(\epsilon)\):
\[
(L-\lambda_{0})u^{(1)}
=
(\lambda^{(1)}-B)u_{0}.
\]
The left-hand side is orthogonal to \(\ker(L-\lambda_{0})=\operatorname{span}\{u_{0}\}\). Taking the inner product against \(u_{0}\) yields the displayed formula. (If \(L\) is a Sturm--Liouville operator, this is the Rayleigh quotient of \(B\) on \(u_{0}\).)
\end{proof}

\begin{example}[2023 Fall Problem~7]
\label{ex:eig-2023}
\[
u''+\bigl(\lambda+\epsilon f(x)\bigr)u=0,
\qquad
u(0)=0,\quad u'(1)=0.
\]
At \(\epsilon=0\) the eigenpairs are \(\lambda_{n}=\mu_{n}^{2}\), \(u_{n}(x)=\sin(\mu_{n}x)\) with \(\mu_{n}=\bigl(n+\tfrac12\bigr)\pi\), \(n=0,1,2,\ldots\). Here \(L=-d^{2}/dx^{2}\) on that domain and \(B=-f\) in the notation of \cref{thm:eig-regular} (the equation is \(Lu=\lambda u+\epsilon f u\)). Write \(\lambda=\lambda_{n}+\epsilon\lambda^{(1)}+O(\epsilon^{2})\) and \(u=u_{n}+\epsilon u^{(1)}+O(\epsilon^{2})\). The \(O(\epsilon)\) equation is
\[
(u^{(1)})''+\lambda_{n}u^{(1)}
=
-\bigl(\lambda^{(1)}+f\bigr)u_{n},
\]
with the same boundary conditions. The homogeneous operator is self-adjoint on this domain, so a solution exists if and only if the right-hand side is orthogonal to \(u_{n}\):
\[
\lambda^{(1)}
=
-\frac{\int_{0}^{1}f(x)u_{n}(x)^{2}\,\mathrm{d}x}{\int_{0}^{1}u_{n}(x)^{2}\,\mathrm{d}x}
=
-2\int_{0}^{1}f(x)\sin^{2}(\mu_{n}x)\,\mathrm{d}x,
\]
using \(\int_{0}^{1}\sin^{2}(\mu_{n}x)\,\mathrm{d}x=\frac12\). That is the \(O(\epsilon)\) expansion of \(\lambda\). The same Fredholm obstruction, with a different homogeneous operator, is the \(\cos s\) coefficient that fixes \(\nu\) in part~(b) of the present problem.
\end{example}

\begin{exercise}[\(f\equiv 1\), and an exact check]
\label{exer:eig-f1}
In the setting of \cref{ex:eig-2023}, take \(f\equiv 1\). Compute \(\lambda^{(1)}\). Then solve the eigenvalue problem exactly and confirm the expansion.
\end{exercise}

\begin{proof}[Solution]
\cref{thm:eig-regular} gives \(\lambda^{(1)}=-1\), because \(B=-1\) and the Rayleigh quotient is \(-1\). Thus \(\lambda=\mu_{n}^{2}-\epsilon+O(\epsilon^{2})\). Exactly, the ODE is \(u''+(\lambda+\epsilon)u=0\) with the same boundary conditions, so the eigenvalues are \(\lambda+\epsilon=\mu_{n}^{2}\), i.e.\ \(\lambda=\mu_{n}^{2}-\epsilon\). The \(O(\epsilon)\) formula is exact, and the remainder is zero.
\end{proof}

\paragraph{Stationary phase.}

This is syllabus item ``stable phase'' in the models box \cite{bender99}. It has not appeared as a P6-type question through 2026 Spring.

\begin{theorem}[Nondegenerate stationary phase]
\label{thm:stat-phase}
Let \(\varphi\in C^{2}(\mathbb{R})\) and let \(a\) be smooth and compactly supported. Suppose \(\varphi'(x_{0})=0\), \(\varphi''(x_{0})\neq 0\), and \(\varphi'\neq 0\) elsewhere on \(\operatorname{supp}a\). Then as \(\lambda\to+\infty\),
\[
\int_{-\infty}^{\infty}a(x)\,e^{\mathrm{i}\lambda\varphi(x)}\,\mathrm{d}x
\sim
a(x_{0})\,
\sqrt{\frac{2\pi}{\lambda\lvert\varphi''(x_{0})\rvert}}\,
\exp\Bigl(\mathrm{i}\lambda\varphi(x_{0})+\mathrm{i}\frac{\pi}{4}\operatorname{sgn}\varphi''(x_{0})\Bigr).
\]
\end{theorem}

\begin{proof}
After translating \(x_{0}\) to the origin, \(\varphi(x)=\varphi(x_{0})+\frac12\varphi''(x_{0})x^{2}+o(x^{2})\). The leading contribution is a Fresnel integral (\cref{ex:fresnel}) against the constant \(a(x_{0})\). Integration by parts on the region where \(\varphi'\neq 0\) shows that the remainder is \(O(\lambda^{-1})\), one order smaller.
\end{proof}

\begin{example}[Fresnel / Gaussian continuation]
\label{ex:fresnel}
For \(\lambda>0\),
\[
\int_{-\infty}^{\infty}e^{\mathrm{i}\lambda x^{2}}\,\mathrm{d}x
=
\sqrt{\frac{\pi}{\lambda}}\,e^{\mathrm{i}\pi/4}.
\]
This is \cref{thm:stat-phase} with \(\varphi(x)=x^{2}\), \(\varphi''(0)=2\), \(a\equiv 1\) in the oscillatory sense (or as a limit of Gaussians \(e^{-\epsilon x^{2}}e^{\mathrm{i}\lambda x^{2}}\) as \(\epsilon\downarrow 0\)). The same formula with \(\lambda\mapsto-\lambda\) replaces \(e^{\mathrm{i}\pi/4}\) by \(e^{-\mathrm{i}\pi/4}\).
\end{example}

\begin{exercise}[Leading oscillation of a Bessel integral]
\label{exer:bessel-phase}
Using \cref{thm:stat-phase}, show that as \(\lambda\to+\infty\),
\[
\int_{0}^{2\pi}e^{\mathrm{i}\lambda\cos\theta}\,\mathrm{d}\theta
\sim
\sqrt{\frac{2\pi}{\lambda}}
\Bigl(
e^{\mathrm{i}(\lambda-\pi/4)}
+
e^{-\mathrm{i}(\lambda-\pi/4)}
\Bigr)
=
2\sqrt{\frac{2\pi}{\lambda}}\cos\Bigl(\lambda-\frac{\pi}{4}\Bigr).
\]
(The left-hand side is \(2\pi J_{0}(\lambda)\).)
\end{exercise}

\begin{proof}[Solution]
Here \(\varphi(\theta)=\cos\theta\) on a circle, so the compact-support hypothesis is replaced by smoothness and periodicity; the same local expansion applies at each stationary point. Now \(\varphi'=-\sin\theta\) vanishes at \(\theta=0\) and \(\theta=\pi\) (and their \(2\pi\)-translates, which are the same points). At \(\theta=0\): \(\varphi=1\), \(\varphi''=-1\). At \(\theta=\pi\): \(\varphi=-1\), \(\varphi''=+1\). Each contributes a term of \cref{thm:stat-phase} with \(a\equiv 1\) and \(\lvert\varphi''\rvert=1\):
\[
\sqrt{\frac{2\pi}{\lambda}}\,e^{\mathrm{i}\lambda-\mathrm{i}\pi/4}
+
\sqrt{\frac{2\pi}{\lambda}}\,e^{-\mathrm{i}\lambda+\mathrm{i}\pi/4},
\]
which is the displayed combination. (Endpoints of \([0,2\pi]\) are identified, so there is no extra boundary contribution.)
\end{proof}

\begin{remark}[WKB, not yet a CAM P6]
\label{rem:wkb}
The same stationary-phase / WKB circle produces the leading approximation \(q^{-1/4}\exp\bigl(\pm\mathrm{i}\lambda\int\sqrt{q}\bigr)\) of \(y''+\lambda^{2} q(x)y=0\) for \(q>0\) and \(\lambda\to\infty\). A turning point \(q(x_{*})=0\) is Airy matching, not a regular WKB step. Neither has appeared as a standalone P6 through 2026 Spring.
\end{remark}

\paragraph{Numerical ODEs.}

This is syllabus item~4, written here as a self-contained CAM chapter. The only standalone ODE-numerics question through 2026 Spring is 2023 Fall Problem~4 (\cref{ex:heun-2023}). Problems~4--5 of the present paper use the same test equation as a lemma inside a PDE von Neumann computation (\cref{ex:ode-test}); they are not this chapter, and the review does not repeat that treatise.

A one-step method for \(y'=f(t,y)\) is
\[
u_{n+1}
=
u_{n}+h F(t_{n},u_{n},h;f).
\]
The increment \(F\) is required to be Lipschitz in \(u\), uniformly for \(0<h\le h_{0}\) and \(t\) in a fixed interval.

\begin{definition}[Consistency and order, one-step]
\label{def:ode-consistency}
Let \(\varphi\) be a smooth solution of the ODE. The local truncation (per unit time) is
\[
\delta(t,h)
=
\frac{\varphi(t+h)-\varphi(t)}{h}-F\bigl(t,\varphi(t),h;f\bigr).
\]
The method is \emph{consistent} if \(\delta(t,h)\to 0\) as \(h\to 0\), and \emph{of order \(p\)} if \(\delta=O(h^{p})\) (equivalently: the step error \(\varphi(t+h)-u_{\mathrm{one\,step}}\) is \(O(h^{p+1})\)).
\end{definition}

The word ``consistency'' is the same word as in Problem~4. The residual is different. Here one feeds a smooth ODE solution to the time stepper. In Problem~4 one feeds a smooth PDE solution to a space-time stencil. Do not quote \(\lvert g(\xi)\rvert\le 1\) as the definition of consistency.

\begin{theorem}[One-step: consistency of order \(p\) \(\Rightarrow\) convergence of order \(p\)]
\label{thm:onestep-conv}
Assume \(F\) is Lipschitz in its second argument, uniformly for \(0<h\le h_{0}\) and \(t\in[0,T]\), with constant \(L\). If the method is of order \(p\) and the starting value satisfies \(u_{0}-\varphi(0)=O(h^{p})\), then
\[
\max_{0\le t_{n}\le T}\lvert u_{n}-\varphi(t_{n})\rvert
\le
C_{T}h^{p}.
\]
In particular, a consistent Lipschitz one-step method is convergent. One-step methods are automatically zero-stable: the homogeneous recurrence \(u_{n+1}=u_{n}\) does not amplify.
\end{theorem}

\begin{proof}
Write \(e_{n}=u_{n}-\varphi(t_{n})\). The exact solution satisfies
\[
\varphi(t_{n+1})
=
\varphi(t_{n})+h F\bigl(t_{n},\varphi(t_{n}),h;f\bigr)+h\delta(t_{n},h),
\]
with \(\lvert\delta\rvert\le C h^{p}\) by \cref{def:ode-consistency}. Subtracting the numerical step gives
\[
e_{n+1}
=
e_{n}+h\Bigl(F(t_{n},u_{n},h;f)-F(t_{n},\varphi(t_{n}),h;f)\Bigr)-h\delta(t_{n},h),
\]
hence \(\lvert e_{n+1}\rvert\le(1+Lh)\lvert e_{n}\rvert+C h^{p+1}\). Unrolling and using \((1+Lh)^{n}\le e^{LT}\) produces
\[
\lvert e_{n}\rvert
\le
e^{LT}\lvert e_{0}\rvert+C h^{p}\frac{e^{LT}-1}{L},
\]
which is \(O(h^{p})\) on \([0,T]\). (If \(L=0\) the second term is \(CTh^{p}\).) This is the ODE counterpart of Lax equivalence, and it is simpler: there is no independent stability hypothesis to check for a Lipschitz one-step method. Linear multistep methods do not enjoy that implication (\cref{thm:dahlquist-eq}).
\end{proof}

\begin{example}[Forward Euler, order \(1\)]
\label{ex:fe-onestep}
Forward Euler is \(u_{n+1}=u_{n}+h f(t_{n},u_{n})\), i.e.\ \(F=f(t,u)\). Then
\[
\delta(t,h)
=
\frac{\varphi(t+h)-\varphi(t)}{h}-\varphi'(t)
=
\frac h2\varphi''(t)+O(h^{2})
=
O(h),
\]
so the method is consistent of order \(1\). By \cref{thm:onestep-conv} it converges with order \(1\) on any fixed interval. On the test equation \(y'=\lambda y\) one has \(u_{n+1}=(1+z)u_{n}\) with \(z=h\lambda\), so \(R(z)=1+z\). The region of absolute stability is the closed disk \(\lvert 1+z\rvert\le 1\), centred at \(-1\) with radius \(1\).
\end{example}

\begin{example}[2023 Fall Problem~4, Heun / explicit RK2]
\label{ex:heun-2023}
The tableau
\[
\begin{array}{c|cc}
0 & 0 \\
1 & 1 & 0 \\
\hline
  & 1/2 & 1/2
\end{array}
\]
is Heun's method. The three examination moves are rewrite, order, and \(R(z)\).

(1) The internal stages are \(k_{1}=f(t_{n},u_{n})\) and \(k_{2}=f(t_{n}+h,u_{n}+h k_{1})\), and
\[
u_{n+1}
=
u_{n}+\frac h2(k_{1}+k_{2}).
\]
Thus
\[
F(t,u,h;f)
=
\frac12\Bigl(f(t,u)+f\bigl(t+h,u+hf(t,u)\bigr)\Bigr).
\]

(2) Let \(\varphi\) solve \(\varphi'=f(t,\varphi)\). Taylor-expand \(f(t+h,\varphi+hf(t,\varphi))\) in \(h\), using \(\varphi''=f_{t}+f_{y}f\):
\[
f\bigl(t+h,\varphi+hf\bigr)
=
f+h(f_{t}+f_{y}f)+O(h^{2})
=
f+h\varphi''+O(h^{2}).
\]
Hence
\[
F\bigl(t,\varphi,h;f\bigr)
=
f+\frac h2\varphi''+O(h^{2})
=
\frac{\varphi(t+h)-\varphi(t)}{h}+O(h^{2}),
\]
because \(\varphi(t+h)=\varphi+h\varphi'+\frac{h^{2}}{2}\varphi''+O(h^{3})\). Thus \(\delta=O(h^{2})\): the method is consistent of order \(2\). Lipschitz continuity of \(f\) in \(y\), uniform in \(t\), implies the same for \(F\). \cref{thm:onestep-conv} gives convergence of order \(2\).

(3) Restrict to \(y'=\lambda y\) with \(z=h\lambda\in\mathbb C\). Then \(k_{1}=\lambda u_{n}\), \(k_{2}=\lambda(1+z)u_{n}\), and
\[
u_{n+1}
=
\Bigl(1+z+\frac{z^{2}}{2}\Bigr)u_{n}.
\]
The stability function is the quadratic Taylor polynomial of \(e^{z}\),
\[
R(z)=1+z+\frac{z^{2}}{2}.
\]
The region of absolute stability is \(\{z\in\mathbb C:\lvert R(z)\rvert\le 1\}\). On the imaginary axis, \(z=\mathrm{i}y\),
\[
\lvert R(\mathrm{i}y)\rvert^{2}
=
\Bigl(1-\frac{y^{2}}{2}\Bigr)^{2}+y^{2}
=
1+\frac{y^{4}}{4}
\ge 1,
\]
with equality only at \(y=0\). The region does not contain the imaginary axis except at the origin, so the method is not \(A\)-stable (and no explicit Runge--Kutta method is: \(R(z)\) is a polynomial, hence \(\lvert R(z)\rvert\to\infty\) as \(\lvert z\rvert\to\infty\) in the left half-plane). It does contain a disk in the left half-plane around the origin, which is enough for non-stiff decaying modes with \(h\lvert\lambda\rvert\) not too large.

For a general Runge--Kutta tableau \((A,b,c)\), the same calculation is
\[
R(z)
=
1+z\,b^{T}(I-zA)^{-1}\mathbf{1}.
\]
For Heun one has \(A=\bigl(\begin{smallmatrix}0&0\\1&0\end{smallmatrix}\bigr)\), \(b=\bigl(\tfrac12,\tfrac12\bigr)\), and the formula recovers \(1+z+z^{2}/2\).
\end{example}

\begin{exercise}[Implicit trapezoid: order \(2\)]
\label{exer:trap-order}
The implicit trapezoidal rule is
\[
u_{n+1}
=
u_{n}+\frac h2\Bigl(f(t_{n},u_{n})+f(t_{n+1},u_{n+1})\Bigr).
\]
Assume \(f\) (hence \(\varphi\)) is smooth. Show that the method is consistent of order \(2\), and conclude convergence of order \(2\) on a fixed interval.
\end{exercise}

\begin{proof}[Solution]
The increment \(F\) is defined implicitly by \(F=\frac12\bigl(f(t,u)+f(t+h,u+hF)\bigr)\). Feeding the exact solution,
\[
\frac{\varphi(t+h)-\varphi(t)}{h}
=
\varphi'+\frac h2\varphi''+\frac{h^{2}}{6}\varphi'''+O(h^{3})
\]
and
\[
\frac12\bigl(\varphi'(t)+\varphi'(t+h)\bigr)
=
\varphi'+\frac h2\varphi''+\frac{h^{2}}{4}\varphi'''+O(h^{3}).
\]
The difference is \(-\frac{h^{2}}{12}\varphi'''+O(h^{3})=O(h^{2})\), so \(\delta=O(h^{2})\). The implicit map \(u\mapsto F\) is Lipschitz for small \(h\) by the implicit-function theorem and Lipschitz continuity of \(f\). \cref{thm:onestep-conv} gives convergence of order \(2\).
\end{proof}

\paragraph{Absolute stability and \(A\)-stability.}

Order on a fixed interval does not control long-time behaviour of decaying modes, nor stepsizes restricted by a fast eigenvalue. That is a different question, asked of the same method.

\begin{definition}[Absolute stability / \(A\)-stability]
\label{def:abs-stab}
A one-step method applied to \(y'=\lambda y\) produces \(u_{n+1}=R(h\lambda)u_{n}\). It is \emph{absolutely stable} at \(z=h\lambda\) if \(\lvert R(z)\rvert\le 1\). It is \emph{\(A\)-stable} if this holds for every \(z\) with \(\operatorname{Re}z\le 0\), i.e.\ if the closed left half-plane lies in the stability region.
\end{definition}

\begin{example}[Stability functions of Euler and trapezoid]
\label{ex:stab-functions}
Forward Euler has \(R(z)=1+z\), stable only in the disk \(\lvert 1+z\rvert\le 1\) (\cref{ex:fe-onestep}). Backward Euler \(u_{n+1}=u_{n}+h f(t_{n+1},u_{n+1})\) has
\[
R(z)=\frac{1}{1-z}.
\]
The implicit trapezoid of \cref{exer:trap-order} has
\[
R(z)=\frac{1+z/2}{1-z/2},
\]
the Cayley transform. Implicit midpoint (the Gauss collocation of order \(2\)) produces the same \(R(z)\) on the linear test equation. Heun is \cref{ex:heun-2023}: \(R(z)=1+z+z^{2}/2\).

\begin{center}
\begin{tikzpicture}[scale=0.95]
\draw[->] (-3.3,0) -- (1.3,0) node[right] {\small \(\operatorname{Re}z\)};
\draw[->] (0,-1.7) -- (0,1.7) node[above] {\small \(\operatorname{Im}z\)};
\fill[blue!12] (-1,0) circle (1);
\draw[blue!55!black, thick] (-1,0) circle (1);
\node[font=\small] at (-1,-1.45) {FE: \(\lvert 1+z\rvert\le 1\)};
\begin{scope}[shift={(5.6,0)}]
\draw[->] (-2.4,0) -- (1.3,0) node[right] {\small \(\operatorname{Re}z\)};
\draw[->] (0,-1.7) -- (0,1.7) node[above] {\small \(\operatorname{Im}z\)};
\fill[blue!12] (-2.35,-1.55) rectangle (0,1.55);
\draw[blue!55!black, thick] (0,-1.55) -- (0,1.55);
\node[font=\small] at (-1.15,-1.45) {BE / trap: \(\operatorname{Re}z\le 0\)};
\end{scope}
\end{tikzpicture}
\end{center}

The left picture is the entire absolute-stability region of forward Euler. The right picture is the \(A\)-stability requirement; both backward Euler and trapezoid fill it (\cref{exer:be-astable}). Trapezoid in fact has \(\lvert R(\mathrm{i}y)\rvert=1\) for all real \(y\): purely imaginary eigenvalues are advanced with no numerical dissipation. Heun's region is a bounded subset of the left half-plane, cut off from the imaginary axis except at \(0\).
\end{example}

\begin{exercise}[Backward Euler is \(A\)-stable]
\label{exer:be-astable}
Prove that \(R(z)=1/(1-z)\) satisfies \(\lvert R(z)\rvert\le 1\) for every \(z\) with \(\operatorname{Re}z\le 0\). Prove that the trapezoidal \(R(z)=(1+z/2)/(1-z/2)\) does as well, and that \(\lvert R(\mathrm{i}y)\rvert=1\).
\end{exercise}

\begin{proof}[Solution]
Let \(z=x+\mathrm{i}y\) with \(x\le 0\). Then \(\lvert 1-z\rvert^{2}=(1-x)^{2}+y^{2}\). Since \(1-x\ge 1\), this is at least \(1\), so \(\lvert 1/(1-z)\rvert\le 1\).

For trapezoid, \(\lvert 1+z/2\rvert\le\lvert 1-z/2\rvert\) if and only if
\[
\bigl(1+\tfrac x2\bigr)^{2}+\bigl(\tfrac y2\bigr)^{2}
\le
\bigl(1-\tfrac x2\bigr)^{2}+\bigl(\tfrac y2\bigr)^{2},
\]
which simplifies to \(x\le -x\), i.e.\ \(x\le 0\). On the imaginary axis \(x=0\) one has equality, so \(\lvert R(\mathrm{i}y)\rvert=1\).
\end{proof}

\begin{exercise}[Classical RK4 is not \(A\)-stable]
\label{exer:rk4-not-astable}
The classical four-stage Runge--Kutta method has
\[
R(z)=1+z+\frac{z^{2}}{2}+\frac{z^{3}}{6}+\frac{z^{4}}{24}.
\]
Show that it is not \(A\)-stable, in two ways: (i) \(R(z)\) is a polynomial; (ii) by evaluating at a concrete point of the left half-plane (e.g.\ \(z=-4\)).
\end{exercise}

\begin{proof}[Solution]
(i) A non-constant polynomial tends to infinity as \(\lvert z\rvert\to\infty\). In particular \(\lvert R(z)\rvert\to\infty\) along the negative real axis, so the left half-plane cannot lie in \(\{\lvert R\rvert\le 1\}\). The same argument shows that no explicit Runge--Kutta method is \(A\)-stable. (ii) Directly: \(R(-4)=1-4+8-64/6+256/24=5-32/3+32/3=5\), whose modulus is \(5>1\). So already \(z=-4\) lies outside the stability region, while \(\operatorname{Re}(-4)<0\).
\end{proof}

\begin{remark}[The PDE amplification factor is this \(R(z)\)]
\label{rem:pde-is-r-of-z}
The amplification factor \(g(\xi)\) of a PDE scheme is \(R(z)\) evaluated at the discrete symbol \(z=\tau\lambda_{h}(\xi)\). Computing \(g\) is the examination move in Problems~4--5; computing \(R(z)\) as a function of a complex variable is the examination move in 2023 Fall Problem~4. They are the same algebra, applied to a Fourier mode or to a scalar ODE. Backward Euler is unconditionally von Neumann stable for heat because every \(\lambda_{h}\) of the discrete Laplacian lies in \(\operatorname{Re}\lambda\le 0\), which is the \(A\)-stability region; it is only conditionally stable for downwind advection because that symbol has \(\operatorname{Re}\lambda_{h}>0\) (\cref{ex:ode-test}).
\end{remark}

\paragraph{Linear multistep methods: zero-stability and convergence.}

\begin{definition}[Linear \(k\)-step method]
\label{def:lmm}
A linear \(k\)-step method is
\[
\sum_{j=0}^{k}\alpha_{j}u_{n+j}
=
h\sum_{j=0}^{k}\beta_{j}f(t_{n+j},u_{n+j}),
\]
with generating polynomials \(\rho(\zeta)=\sum_{j=0}^{k}\alpha_{j}\zeta^{j}\) and \(\sigma(\zeta)=\sum_{j=0}^{k}\beta_{j}\zeta^{j}\), normalized by \(\alpha_{k}=1\) or by any convenient leading coefficient. The method is explicit if \(\beta_{k}=0\) (Adams--Bashforth) and implicit otherwise (Adams--Moulton, BDF).
\end{definition}

\begin{definition}[LMS consistency of order \(p\)]
\label{def:lmm-order}
The method \((\rho,\sigma)\) is consistent of order \(p\) if, for every smooth \(\varphi\),
\[
\sum_{j=0}^{k}\alpha_{j}\varphi(t+jh)-h\sum_{j=0}^{k}\beta_{j}\varphi'(t+jh)
=
O(h^{p+1})
\]
as \(h\to 0\). Equivalently, the order conditions
\[
\sum_{j=0}^{k}\alpha_{j}=0,
\qquad
\sum_{j=0}^{k}j^{q}\alpha_{j}
=
q\sum_{j=0}^{k}j^{q-1}\beta_{j}
\qquad (q=1,\ldots,p)
\]
hold (the \(q=0\) sum is \(\sum\alpha_{j}=0\)). In generating-function form: \(\rho(e^{h})-h\sigma(e^{h})=O(h^{p+1})\).
\end{definition}

\begin{theorem}[Root condition / zero-stability]
\label{thm:root-condition}
The linear multistep method \((\rho,\sigma)\) is \emph{zero-stable} if the roots of \(\rho\) satisfy \(\lvert\zeta\rvert\le 1\), and every root with \(\lvert\zeta\rvert=1\) is simple. Equivalently: the homogeneous recurrence \(\sum\alpha_{j}u_{n+j}=0\) has no exponentially growing solution, and solutions associated with unimodular roots remain bounded (they are polynomials of degree \(0\) only). Without zero-stability the method diverges for \(f\equiv 0\) already, independently of consistency.
\end{theorem}

\begin{proof}
The general solution of the homogeneous recurrence is \(\sum p_{m}(n)\,\zeta_{m}^{n}\), with \(p_{m}\) a polynomial of degree one less than the multiplicity of \(\zeta_{m}\). Boundedness for all \(n\), uniformly in the starting data of size \(O(1)\), forces \(\lvert\zeta_{m}\rvert\le 1\) and \(\deg p_{m}=0\) on the unit circle.
\end{proof}

\begin{theorem}[Dahlquist equivalence]
\label{thm:dahlquist-eq}
A linear multistep method is convergent (on a fixed interval, for \(f\) Lipschitz in \(y\)) if and only if it is consistent and zero-stable. If it is consistent of order \(p\) and zero-stable, and the starting values satisfy \(u_{j}-\varphi(t_{j})=O(h^{p})\) for \(j=0,\ldots,k-1\), then the convergence order is \(p\).
\end{theorem}

\begin{proof}[Proof idea]
Necessity of consistency is Taylor expansion against a smooth solution. Necessity of zero-stability is the \(f\equiv 0\) test. Sufficiency is a discrete Gronwall argument in the companion-form of the recurrence, using that the powers of the companion matrix of \(\rho\) remain bounded precisely when the root condition holds \cite[Ch.~III]{hairer93}. One-step methods fit the theorem with \(k=1\), \(\rho(\zeta)=\zeta-1\): the root condition is automatic, which recovers \cref{thm:onestep-conv}.
\end{proof}

\begin{example}[Leapfrog / explicit midpoint is weakly stable]
\label{ex:leapfrog-weak}
The explicit midpoint rule
\[
u_{n+1}=u_{n-1}+2h f(t_{n},u_{n})
\]
has \(\rho(\zeta)=\zeta^{2}-1\), \(\sigma(\zeta)=2\zeta\). Roots of \(\rho\): \(\pm 1\), both simple, so the method is zero-stable. With \(\alpha_{2}=1\), \(\alpha_{1}=0\), \(\alpha_{0}=-1\), \(\beta_{1}=2\), the order conditions read \(\sum\alpha_{j}=0\), \(\sum j\alpha_{j}=2=\sum\beta_{j}\), \(\sum j^{2}\alpha_{j}=4=2\sum j\beta_{j}\), while \(\sum j^{3}\alpha_{j}=8\neq 3\sum j^{2}\beta_{j}=6\), so the order is \(2\).

Zero-stability is not absolute stability. On \(y'=\lambda y\) with \(z=h\lambda\), the characteristic equation is \(\zeta^{2}-2z\zeta-1=0\), with roots \(\zeta_{\pm}=z\pm\sqrt{z^{2}+1}\). For \(z=0\) these are \(\pm 1\). For \(z\) small and negative (a decaying ODE), the physical root is \(\zeta_{+}=1+z+O(z^{2})\approx e^{z}\), while the parasitic root is \(\zeta_{-}=-1+z+O(z^{2})\), of modulus greater than \(1\). The parasitic mode grows. This is weak stability: the method is zero-stable and convergent on fixed intervals as \(h\to 0\), but it is unusable for decaying modes at fixed \(h\) with \(t\to\infty\). The same phenomenon appears as a computational mode in a two-step PDE scheme (\cref{def:one-two-step}).
\end{example}

\begin{example}[Adams--Bashforth 2]
\label{ex:ab2}
The two-step Adams--Bashforth method is
\[
u_{n+1}-u_{n}
=
\frac h2\bigl(3f(t_{n},u_{n})-f(t_{n-1},u_{n-1})\bigr).
\]
In the indexing of \cref{def:lmm} (shift the labels by one),
\[
\rho(\zeta)=\zeta^{2}-\zeta,
\qquad
\sigma(\zeta)=\tfrac32\zeta-\tfrac12.
\]
Roots of \(\rho\): \(0\) and \(1\), so zero-stable. Order conditions through \(q=2\):
\begin{align*}
\sum\alpha_{j}&=0,\\
\sum j\alpha_{j}&=2-1=1=\sum\beta_{j},\\
\sum j^{2}\alpha_{j}&=4-1=3=2\sum j\beta_{j}=2\cdot\tfrac32,
\end{align*}
and \(\sum j^{3}\alpha_{j}=8-1=7\neq 3\sum j^{2}\beta_{j}=\frac92\), so the order is \(2\). The method is explicit, hence not \(A\)-stable.
\end{example}

\begin{example}[BDF2]
\label{ex:bdf2}
BDF2 is
\[
\frac32 u_{n+1}-2u_{n}+\frac12 u_{n-1}
=
h f(t_{n+1},u_{n+1}),
\]
i.e.\ \(\alpha_{2}=\frac32\), \(\alpha_{1}=-2\), \(\alpha_{0}=\frac12\), \(\beta_{2}=1\), \(\beta_{1}=\beta_{0}=0\). Then
\[
\rho(\zeta)=\tfrac12(3\zeta^{2}-4\zeta+1)=\tfrac12(3\zeta-1)(\zeta-1),
\qquad
\sigma(\zeta)=\zeta^{2}.
\]
Roots of \(\rho\): \(1\) and \(\frac13\), both inside the closed disk, \(1\) simple, so zero-stable. This is the standard stiff-capable two-step method; its \(A\)-stability is recorded with the barriers below.
\end{example}

\begin{exercise}[BDF2: order exactly \(2\)]
\label{exer:bdf2-order}
For BDF2 as in \cref{ex:bdf2}, verify the order conditions of \cref{def:lmm-order} for \(q=0,1,2\), and show that the \(q=3\) condition fails. Conclude that the method is convergent of order \(2\).
\end{exercise}

\begin{proof}[Solution]
\(\sum\alpha_{j}=\frac32-2+\frac12=0\). For \(q=1\), \(\sum j\alpha_{j}=1\cdot(-2)+2\cdot\frac32=1\) and \(1\cdot\sum\beta_{j}=1\). For \(q=2\), \(\sum j^{2}\alpha_{j}=1\cdot(-2)+4\cdot\frac32=4\) and \(2\sum j\beta_{j}=2\cdot 2\cdot 1=4\). For \(q=3\), \(\sum j^{3}\alpha_{j}=1\cdot(-2)+8\cdot\frac32=10\) and \(3\sum j^{2}\beta_{j}=3\cdot 4\cdot 1=12\neq 10\). Thus the order is exactly \(2\). Zero-stability is \cref{ex:bdf2}. Starting values of order \(2\) (e.g.\ one implicit Euler step plus the initial datum) and \cref{thm:dahlquist-eq} give convergence of order \(2\).
\end{proof}

\begin{theorem}[Dahlquist barriers]
\label{thm:dahlquist-barriers}
(First barrier.) A zero-stable linear \(k\)-step method has order at most \(k+1\) if \(k\) is odd, and at most \(k+2\) if \(k\) is even.

(Second barrier.) An \(A\)-stable linear multistep method is implicit and of order at most \(2\). Among order-\(2\) \(A\)-stable LMS methods, the implicit trapezoidal rule has the smallest error constant.
\end{theorem}

\begin{proof}[Proof idea]
The first barrier is an obstruction among the order conditions plus the root condition: one cannot place too many roots of \(\rho\) inside the disk while matching many moments against \(e^{z}\) \cite[Ch.~III]{hairer93}. The second barrier is a maximum-modulus obstruction for the rational function \(r(z)=\rho(z)/\sigma(z)\) as an approximation to \(e^{-z}\) (or \(\sigma/\rho\) to \(e^{z}\)) throughout the left half-plane \cite{hairer96}. Explicit methods have polynomial \(R\) (one-step) or a characteristic polynomial that cannot keep all roots in the disk for all \(\operatorname{Re}z\le 0\).
\end{proof}

\begin{remark}[\(A(\alpha)\)-stability of BDF]
\label{rem:bdf-aalpha}
BDF1 is backward Euler, \(A\)-stable, order \(1\). BDF2 is \(A\)-stable, order \(2\), and saturates the second barrier among BDFs. BDF3 and higher are not \(A\)-stable; they are \(A(\alpha)\)-stable, meaning \(\lvert R(z)\rvert\le 1\) on the sector \(\lvert\arg(-z)\rvert\le\alpha\) with \(\alpha<90^{\circ}\) (\(\alpha\approx 86^{\circ}\) for BDF3, decreasing with \(k\)). For parabolic problems whose symbols lie in a strict sector of the left half-plane, \(A(\alpha)\) is enough. For a CAM prompt that writes ``BDF'' or ``stiff'', BDF2 is the method to write down.
\end{remark}

\paragraph{Stiffness.}

\begin{definition}[Stiff IVP, operational]
\label{def:stiff}
An IVP \(y'=f(t,y)\) on \([0,T]\) is \emph{stiff} when it has decaying modes whose Lipschitz constants (or Jacobian eigenvalues) span several orders of magnitude, so that the stepsize restriction coming from absolute stability of an explicit method is much tighter than the stepsize restriction coming from accuracy on the smooth slow component. There is no single scalar ``stiffness ratio'' that captures every example; the operational test is: forward Euler (or any explicit RK) is forced to take \(h=O(1/\lvert\lambda_{\mathrm{fast}}\rvert)\), while the exact fast transients are already negligible and the slow solution is smooth on a scale \(O(1)\).
\end{definition}

\begin{example}[Prototypical stiff scalar]
\label{ex:stiff-proto}
Let \(\lambda\ll -1\) and
\[
y'=\lambda\bigl(y-\cos t\bigr)-\sin t,
\qquad
y(0)=1.
\]
The exact solution is \(y(t)=\cos t\) for every \(\lambda\): the fast eigenvalue is invisible in the solution, which is \(O(1)\)-smooth. Accuracy of a consistent method, judged on this \(y\), would permit \(h\) of size \(O(1)\). Absolute stability of forward Euler nevertheless requires \(\lvert 1+h\lambda\rvert\le 1\), i.e.\ \(0<h\le 2/\lvert\lambda\rvert\). For \(\lambda=-10^{4}\) this is \(h\le 2\cdot 10^{-4}\), four orders of magnitude below the accuracy scale. Backward Euler and trapezoid, being \(A\)-stable, have no such restriction: one may step with \(h\) chosen for the smoothness of \(\cos t\). That is the ODE original of the remark in Problem~4 that backward Euler is unconditionally stable for heat and only conditionally stable for downwind advection: heat has \(\operatorname{Re}\lambda_{h}\le 0\), downwind does not.

A linear system \(y'=Ay+g(t)\) is stiff when the eigenvalues of \(A\) satisfy \(\operatorname{Re}\lambda_{j}\le 0\) and \(\max\lvert\lambda_{j}\rvert/\min_{\operatorname{Re}\lambda_{j}<0}\lvert\lambda_{j}\rvert\) is large, while the solution itself (after a short transient) varies on the slow scale.
\end{example}

\begin{exercise}[Forward Euler on a stiff decaying mode]
\label{exer:fe-stiff}
For \(y'=\lambda y\) with \(\lambda=-100\) and \(y(0)=1\), (i) write the forward Euler recurrence and the condition for \(\lvert u_{n}\rvert\) to be nonincreasing; (ii) compare the maximal such \(h\) to the time scale on which \(\lvert y(t)\rvert\) has already dropped below \(10^{-8}\); (iii) state the corresponding restriction, or lack of one, for backward Euler.
\end{exercise}

\begin{proof}[Solution]
(i) \(u_{n+1}=(1-100h)u_{n}\). Nonincreasing moduli: \(\lvert 1-100h\rvert\le 1\). For \(h>0\) this is \(0<h\le 2/100=0.02\). (If \(h>0.02\), then \(1-100h<-1\) and the sequence alternates with growing amplitude: the exact solution decays, the numerical one explodes.) (ii) The exact modulus is \(e^{-100t}\), which is below \(10^{-8}\) once \(t>(8\log 10)/100\approx 0.18\). After a fifth of a unit of time the fast mode is gone; an accuracy-based stepper on the remaining (here, zero) slow component would not be constrained by \(\lambda=-100\). Euler still cannot take \(h>0.02\) at any later time, because the stability region does not contain \(z=-100h\) for those \(h\). (iii) Backward Euler has \(u_{n+1}=u_{n}/(1+100h)\), and \(\lvert 1/(1+100h)\rvert<1\) for every \(h>0\). No restriction.
\end{proof}

\begin{examnote}[ODE slot, across papers]
\label{exnote:ode-review}
If the optional question is a conservative oscillator with a period: energy for closedness (\cref{prop:energy-periodic}), Lindstedt for the frequency (\cref{con:lindstedt,thm:fredholm-osc}), never a limit-cycle theorem. If it says ``large \(t\)'' and has damping: multiple scales (\cref{con:multiple-scales}), amplitude ODE, range \(t=O(1/\epsilon)\). If \(\epsilon\) multiplies \(y''\) in a BVP: outer profile, distinguished inner width (\cref{thm:distinguished-balance}), match. If \(\epsilon\) sits in an eigenvalue: Rayleigh quotient (\cref{thm:eig-regular}). If it is an RK tableau: rewrite as \(F\), Taylor for order plus \cref{thm:onestep-conv}, \(R(z)\) on \(y'=\lambda y\) for the stability region. If it writes \(\rho,\sigma\) / Adams / BDF: order conditions, root condition, Dahlquist equivalence; BDF2 if it writes ``stiff''. If it is a PDE scheme: that is Problems~4--5 of this paper, not this review; the only ODE object there is the test-equation lemma \(R(z)\) (\cref{ex:ode-test,rem:pde-is-r-of-z}).
\end{examnote}


\subsection{Problem 7. [20 points]}
The polar of \(C \subseteq \mathbb{R}^{n}\) is defined as the set

\[
C^{\circ} = \bigl\{ y \in \mathbb{R}^{n} \mid y^{T} x \le 1 \text{ for all } x \in C \bigr\}.
\]

\begin{enumerate}[label=(\alph*)]
    \item (2 points) Show that \(C^{\circ}\) is convex (even if \(C\) is not).
    \item (3 points) What is the polar of a cone?
    \item (3 points) What is the polar of the unit ball for a norm \(\|\cdot\|\)?
    \item (4 points) What is the polar of the set \(C = \{ x \mid \mathbf{1}^{T} x = 1,\, x \succeq 0 \}\)?
    \item (8 points) Show that if \(C\) is closed and convex, with \(0 \in C\), then \((C^{\circ})^{\circ} = C\).
\end{enumerate}

\setcounter{section}{7}
\setcounter{theorem}{0}
\setcounter{definition}{0}
\setcounter{remark}{0}
\setcounter{note}{0}

\subsubsection{Theorems used in Problem 7}

This is the convex-analysis slot of the CAM syllabus item ``linear programming and optimization methods''. The object is Boyd--Vandenberghe's set polar \cite{boyd04} (also Rockafellar \cite{rockafellar70}, \S14). There is no graph, no tree, and no Laplacian. The cone in (b) is a cone of rays through the origin. The symbol \(x\succeq 0\) in (d) is Boyd's generalized inequality for the nonnegative orthant: each coordinate \(x_{i}\ge 0\). It is not the matrix statement \(X\succeq 0\).

Write \(h_{C}(y)\coloneqq\sup_{x\in C}y^{T}x\) for the support function of \(C\) (with the convention \(h_{\emptyset}=-\infty\)). Then the polar is the sublevel set
\[
C^{\circ}
=
\bigl\{ y\in\mathbb{R}^{n} \mid h_{C}(y)\le 1 \bigr\}.
\]
Each constraint \(y^{T}x\le 1\) is a closed halfspace in the \(y\)-variable, with outward normal \(x\) and distance \(1/\|x\|_{2}\) from the origin (if \(x\neq 0\)). So \(C^{\circ}\) is the intersection of these halfspaces. Points of \(C\) far from the origin press the walls of \(C^{\circ}\) toward \(0\); recession directions of \(C\) collapse the inequality from \(\le 1\) to \(\le 0\).

\begin{center}
\begin{tikzpicture}[scale=1.05, >=Latex]
  \draw[gray!60] (-2.3,0) -- (2.4,0);
  \draw[gray!60] (0,-1.7) -- (0,2.15);
  \fill[blue!12] (-2.3,-1.7) -- (1.85,-1.7) -- (-0.55,2.15) -- (-2.3,2.15) -- cycle;
  \draw[blue!55!black, thick] (1.85,-1.7) -- (-0.55,2.15);
  \draw[red!70!black, thick, ->] (0,0) -- (1.35,0.85);
  \fill[red!70!black] (1.35,0.85) circle (1.6pt);
  \node[red!70!black, font=\small, anchor=west] at (1.42,0.95) {\(x\in C\)};
  \node[blue!55!black, font=\small] at (-1.15,-0.85) {\(\{y:y^{T}x\le 1\}\)};
  \fill (0,0) circle (1.2pt);
  \node[font=\scriptsize, anchor=north west] at (0.06,-0.08) {\(0\)};
\end{tikzpicture}
\end{center}

\begin{definition}[Cone]
\label{def:p7-cone}
A set \(K\subseteq\mathbb{R}^{n}\) is a \emph{cone} (vertex at the origin) if it is a union of rays through \(0\): equivalently, \(x\in K\) and \(\lambda\ge 0\) imply \(\lambda x\in K\), or as a set identity
\[
K
=
\mathbb{R}_{+}K
\coloneqq
\bigl\{ \lambda x : x\in K,\, \lambda\ge 0 \bigr\}.
\]
Any nonempty such cone contains the origin (take \(\lambda=0\)). Given a generating set, or geometric base, \(S\subseteq\mathbb{R}^{n}\), the same formula produces the (not necessarily convex) cone of rays through \(S\),
\[
\mathbb{R}_{+}S
=
\bigl\{ \lambda x : x\in S,\, \lambda\ge 0 \bigr\}.
\]
The set \(K\) is a \emph{convex cone} if it is a cone and convex, equivalently if \(x,y\in K\) and \(\theta_{1},\theta_{2}\ge 0\) imply \(\theta_{1}x+\theta_{2}y\in K\). The convex cone generated by \(S\) is the conic hull of finite nonnegative combinations
\[
\operatorname{cone}(S)
=
\biggl\{
\sum_{i=1}^{k}\theta_{i}x_{i}
:
k\in\mathbb{N},\ 
x_{i}\in S,\ 
\theta_{i}\ge 0
\biggr\}.
\]
For a finite generator \(S=\{a_{1},\ldots,a_{m}\}\) this is the finitely generated cone \(K=\{A\theta:\theta\succeq 0\}\) with \(A=(a_{1}\mid\cdots\mid a_{m})\). Problem~7(b) uses only the ray identity \(K=\mathbb{R}_{+}K\), not convexity of \(K\).
\end{definition}

\begin{definition}[Dual cone]
\label{def:p7-dual-cone}
The Boyd dual cone of \(K\) is
\[
K^{*}
=
\bigl\{ y\in\mathbb{R}^{n} \mid y^{T}x\ge 0\text{ for all }x\in K \bigr\}.
\]
The polar cone in Rockafellar's sign convention is \(-K^{*}\) \cite{rockafellar70}.
\end{definition}

\begin{definition}[Dual norm]
\label{def:p7-dual-norm}
If \(\|\cdot\|\) is a norm on \(\mathbb{R}^{n}\), the dual norm is
\[
\|y\|_{*}
=
\sup_{\|x\|\le 1} y^{T}x
=
\sup_{x\neq 0}\frac{y^{T}x}{\|x\|}.
\]
The unit ball of \(\|\cdot\|_{*}\) is the polar of the unit ball of \(\|\cdot\|\). In particular \(\|\cdot\|_{2}\) is self-dual, and \(\|\cdot\|_{1}\) is dual to \(\|\cdot\|_{\infty}\).
\end{definition}

\begin{definition}[Probability simplex]
\label{def:p7-simplex}
The set in (d) is the probability simplex
\[
\Delta^{n-1}
=
\bigl\{ x\in\mathbb{R}^{n} \mid \mathbf{1}^{T}x=1,\ x\succeq 0 \bigr\}
=
\operatorname{conv}\{e_{1},\ldots,e_{n}\}.
\]
Here \(x\succeq 0\) means \(x_{i}\ge 0\) for every \(i\), i.e.\ \(x\) lies in the nonnegative orthant \(\mathbb{R}^{n}_{+}\).
\end{definition}

\begin{lemma}[Strict separation in \(\mathbb{R}^{n}\)]
\label{lem:p7-separate}
Let \(C\subseteq\mathbb{R}^{n}\) be nonempty, closed, and convex, and let \(z\notin C\). There exist \(a\in\mathbb{R}^{n}\setminus\{0\}\) and \(\alpha\in\mathbb{R}\) such that
\[
a^{T}x\le\alpha
\quad\text{for all }x\in C,
\qquad
a^{T}z>\alpha.
\]
\end{lemma}

\begin{proof}
Let \(\pi\) be a Euclidean nearest point of \(z\) in \(C\) (a minimizing sequence of \(\|x-z\|_{2}\) is bounded, hence has a convergent subsequence, and the limit lies in \(C\)). Set \(a\coloneqq z-\pi\neq 0\). For every \(x\in C\) and every \(t\in(0,1]\) the point \(\pi+t(x-\pi)\) still lies in \(C\), so
\[
\|\pi-z\|_{2}^{2}
\le
\bigl\|\pi+t(x-\pi)-z\bigr\|_{2}^{2}
=
\|\pi-z\|_{2}^{2}
-2t\,a^{T}(x-\pi)
+t^{2}\|x-\pi\|_{2}^{2}.
\]
Dividing by \(t\) and sending \(t\to 0^{+}\) gives \(a^{T}(x-\pi)\le 0\), i.e.\ \(a^{T}x\le a^{T}\pi\). On the other side,
\[
a^{T}z
=
a^{T}(\pi+a)
=
a^{T}\pi+\|a\|_{2}^{2}
>
a^{T}\pi.
\]
Take \(\alpha=a^{T}\pi\).
\end{proof}

\begin{examnote}[What the script must contain]
\label{exnote:p7}
(a) Intersection of halfspaces, hence convex. (b) Polar of a cone is \(\{y:y^{T}x\le 0\ \forall x\in K\}=-K^{*}\). (c) Dual unit ball \(\{y:\|y\|_{*}\le 1\}\). (d) \(\{y:y\preceq\mathbf{1}\}\), because \(\max_{x\in\Delta^{n-1}}y^{T}x=\max_{i}y_{i}\). (e) Always \(C\subseteq(C^{\circ})^{\circ}\); the reverse is \cref{lem:p7-separate} plus a scaling that uses \(0\in C\). Do not quote Fenchel duality. The eight-point part is the separating-hyperplane argument.
\end{examnote}

\subsubsection{Solution of Problem 7(a)}

\paragraph{Answer.}
\(C^{\circ}\) is an intersection of closed halfspaces, hence closed and convex, whether or not \(C\) is convex.

\begin{proof}
If \(C=\emptyset\), the condition \(y^{T}x\le 1\) for all \(x\in C\) holds vacuously, so \(C^{\circ}=\mathbb{R}^{n}\), which is convex. Now assume \(C\neq\emptyset\). For each \(x\in C\) the set
\[
H_{x}
=
\bigl\{ y\in\mathbb{R}^{n} \mid y^{T}x\le 1 \bigr\}
\]
is a closed halfspace, hence convex. By definition
\[
C^{\circ}
=
\bigcap_{x\in C} H_{x}.
\]
An arbitrary intersection of convex sets is convex, so \(C^{\circ}\) is convex.

The same identity shows why convexity of \(C\) is irrelevant: \(H_{x}\) depends on \(x\) linearly, so
\[
C^{\circ}
=
(\operatorname{conv} C)^{\circ}
=
(\overline{\operatorname{conv}}\,C)^{\circ}.
\]
Nonconvexity of \(C\) is quotiented out before the polar is formed. (Closedness of \(C^{\circ}\) is free, and is used in (e): the bipolar is always closed.)
\end{proof}

\subsubsection{Solution of Problem 7(b)}

\paragraph{Answer.}
If \(K\subseteq\mathbb{R}^{n}\) is a cone in the sense of \cref{def:p7-cone}, then
\[
K^{\circ}
=
\bigl\{ y\in\mathbb{R}^{n} \mid y^{T}x\le 0\text{ for all }x\in K \bigr\}
=
-K^{*}.
\]

\begin{proof}
Let \(y\in K^{\circ}\). Then \(y^{T}(\lambda x)\le 1\) for every \(x\in K\) and every \(\lambda\ge 0\). If \(y^{T}x>0\) for some \(x\in K\), the left-hand side tends to \(+\infty\) as \(\lambda\to\infty\), contradicting the bound \(1\). Hence \(y^{T}x\le 0\) for every \(x\in K\).

Conversely, if \(y^{T}x\le 0\) for every \(x\in K\), then \(y^{T}x\le 1\) holds automatically, so \(y\in K^{\circ}\).

The identification with \(-K^{*}\) is \cref{def:p7-dual-cone}. Convexity of \(K\) is not used. For the nonnegative orthant one has \((\mathbb{R}^{n}_{+})^{\circ}=-\mathbb{R}^{n}_{+}\).
\end{proof}

\begin{center}
\begin{tikzpicture}[scale=0.95, >=Latex]
  \draw[gray!60] (-2.2,0) -- (2.3,0) node[right, font=\small] {\(y_{1}\)};
  \draw[gray!60] (0,-2.2) -- (0,2.3) node[above, font=\small] {\(y_{2}\)};
  \fill[red!15] (0,0) -- (2.2,0) -- (2.2,2.2) -- (0,2.2) -- cycle;
  \fill[blue!15] (0,0) -- (-2.2,0) -- (-2.2,-2.2) -- (0,-2.2) -- cycle;
  \node[red!70!black, font=\small] at (1.15,1.15) {\(K=\mathbb{R}^{2}_{+}\)};
  \node[blue!50!black, font=\small] at (-1.15,-1.15) {\(K^{\circ}=-\mathbb{R}^{2}_{+}\)};
  \fill (0,0) circle (1.2pt);
\end{tikzpicture}
\end{center}

\subsubsection{Solution of Problem 7(c)}

\paragraph{Answer.}
Let \(B=\{x:\|x\|\le 1\}\) be the closed unit ball of a norm \(\|\cdot\|\). Then
\[
B^{\circ}
=
\bigl\{ y\in\mathbb{R}^{n} \mid \|y\|_{*}\le 1 \bigr\},
\]
the closed unit ball of the dual norm.

\begin{proof}
By definition of the dual norm,
\[
h_{B}(y)
=
\sup_{\|x\|\le 1} y^{T}x
=
\|y\|_{*}.
\]
Hence \(y\in B^{\circ}\) if and only if \(\|y\|_{*}\le 1\).

In particular the Euclidean ball is self-polar. The unit ball of \(\|\cdot\|_{\infty}\) is the cube \([-1,1]^{n}\), whose polar is the unit ball of \(\|\cdot\|_{1}\), and conversely.
\end{proof}

\begin{center}
\begin{tikzpicture}[scale=0.85, >=Latex]
  \begin{scope}
    \draw[gray!60] (-1.7,0) -- (1.8,0);
    \draw[gray!60] (0,-1.7) -- (0,1.8);
    \fill[red!12] (-1,-1) rectangle (1,1);
    \draw[red!70!black, thick] (-1,-1) rectangle (1,1);
    \fill[blue!18, opacity=0.85] (1,0) -- (0,1) -- (-1,0) -- (0,-1) -- cycle;
    \draw[blue!50!black, thick] (1,0) -- (0,1) -- (-1,0) -- (0,-1) -- cycle;
    \node[font=\small] at (0,-2.05) {\(\|\cdot\|_{\infty}\) ball and its polar \(\|\cdot\|_{1}\) ball};
  \end{scope}
  \begin{scope}[xshift=5.6cm]
    \draw[gray!60] (-1.7,0) -- (1.8,0);
    \draw[gray!60] (0,-1.7) -- (0,1.8);
    \draw[thick] (0,0) circle (1);
    \node[font=\small] at (0,-2.05) {\(\|\cdot\|_{2}\) ball, self-polar};
  \end{scope}
\end{tikzpicture}
\end{center}

\subsubsection{Solution of Problem 7(d)}

\paragraph{Answer.}
\[
C^{\circ}
=
\bigl\{ y\in\mathbb{R}^{n} \mid y\preceq\mathbf{1} \bigr\}
=
\bigl\{ y\in\mathbb{R}^{n} \mid y_{i}\le 1\text{ for all }i \bigr\}.
\]

\begin{proof}
The set \(C\) is the probability simplex of \cref{def:p7-simplex}. For any \(y\in\mathbb{R}^{n}\),
\[
h_{C}(y)
=
\sup\bigl\{ y^{T}x : \mathbf{1}^{T}x=1,\ x\succeq 0 \bigr\}
=
\max_{1\le i\le n} y_{i},
\]
because a linear function on a simplex attains its maximum at a vertex, and the vertices are the standard basis vectors \(e_{i}\). Therefore \(y\in C^{\circ}\) if and only if \(\max_{i}y_{i}\le 1\), i.e.\ \(y\preceq\mathbf{1}\).

In \(\mathbb{R}^{2}\) the set \(C\) is the segment \(\operatorname{conv}\{e_{1},e_{2}\}\), and \(C^{\circ}\) is the unbounded region to the southwest of the lines \(y_{1}=1\) and \(y_{2}=1\). The origin does not lie in \(C\), so \(C^{\circ}\) is unbounded; this is compatible with (e), whose recovery statement assumes \(0\in C\).
\end{proof}

\begin{center}
\begin{tikzpicture}[scale=0.95, >=Latex]
  \draw[gray!60] (-1.6,0) -- (2.2,0) node[right, font=\small] {\(y_{1}\)};
  \draw[gray!60] (0,-1.6) -- (0,2.2) node[above, font=\small] {\(y_{2}\)};
  \fill[blue!12] (-1.6,-1.6) -- (1,-1.6) -- (1,1) -- (-1.6,1) -- cycle;
  \draw[blue!50!black, thick] (1,-1.6) -- (1,1) -- (-1.6,1);
  \draw[red!70!black, very thick] (1,0) -- (0,1);
  \fill[red!70!black] (1,0) circle (1.7pt);
  \fill[red!70!black] (0,1) circle (1.7pt);
  \node[red!70!black, font=\small, anchor=west] at (1.08,-0.18) {\(e_{1}\)};
  \node[red!70!black, font=\small, anchor=south west] at (0.08,1.08) {\(e_{2}\)};
  \node[red!70!black, font=\small] at (0.62,0.62) {\(C\)};
  \node[blue!50!black, font=\small] at (-0.55,-0.55) {\(C^{\circ}\)};
  \fill (0,0) circle (1.2pt);
\end{tikzpicture}
\end{center}

\subsubsection{Solution of Problem 7(e)}

\paragraph{Answer.}
If \(C\subseteq\mathbb{R}^{n}\) is closed and convex and \(0\in C\), then \((C^{\circ})^{\circ}=C\).

\begin{proof}
\emph{The inclusion \(C\subseteq(C^{\circ})^{\circ}\) holds for an arbitrary set \(C\).}
If \(x\in C\), then \(y^{T}x\le 1\) for every \(y\in C^{\circ}\), which is precisely the membership \(x\in(C^{\circ})^{\circ}\). No closedness, convexity, or \(0\in C\) is used here.

\emph{The reverse inclusion uses all three hypotheses.}
Let \(z\notin C\). By \cref{lem:p7-separate} there exist \(a\neq 0\) and \(\alpha\in\mathbb{R}\) with \(a^{T}x\le\alpha\) for all \(x\in C\) and \(a^{T}z>\alpha\). Since \(0\in C\), one has \(0=a^{T}0\le\alpha\), so \(\alpha\ge 0\).

If \(\alpha>0\), set \(\tilde{a}\coloneqq a/\alpha\). Then \(\tilde{a}^{T}x\le 1\) for all \(x\in C\), hence \(\tilde{a}\in C^{\circ}\), while \(\tilde{a}^{T}z>1\). Thus \(z\notin(C^{\circ})^{\circ}\).

If \(\alpha=0\), then \(a^{T}x\le 0\) for all \(x\in C\) and \(a^{T}z>0\). For every \(t\ge 0\) one has \((ta)^{T}x\le 0\le 1\), so \(ta\in C^{\circ}\). Choosing \(t>1/(a^{T}z)\) gives \((ta)^{T}z>1\), hence again \(z\notin(C^{\circ})^{\circ}\).

Therefore no point outside \(C\) can lie in the bipolar, i.e.\ \((C^{\circ})^{\circ}\subseteq C\). Combined with the first paragraph, \((C^{\circ})^{\circ}=C\).
\end{proof}

\begin{remark}[Why the hypotheses are sharp]
\label{rem:p7-bipolar}
The bipolar \((C^{\circ})^{\circ}\) is always closed, convex, and contains \(0\). Dropping any hypothesis enlarges the right-hand side: in general
\[
(C^{\circ})^{\circ}
=
\overline{\operatorname{conv}}\bigl(C\cup\{0\}\bigr).
\]
If \(0\notin C\), the bipolar adjoins the origin (and then takes closed convex hull). If \(C\) is not convex or not closed, the bipolar returns the missing hull. That is why (e) assumes all three.
\end{remark}

\begin{summary}[Problem 7]
\label{sum:p7}
\(C^{\circ}=\{h_{C}\le 1\}\) is always convex. A cone polarizes to \(\{y^{T}x\le 0\}=-K^{*}\). A norm ball polarizes to the dual ball. The simplex polarizes to \(\{y\preceq\mathbf{1}\}\). Closed convex sets containing \(0\) are recovered by the bipolar.
\end{summary}


\begin{thebibliography}{99}
\bibitem[Ding(2024)]{ding24}
P.~Ding.
\newblock \emph{Linear Model and Extensions}.
\newblock Book manuscript; preprint arXiv:2401.00649. Chinese course title: 应用回归分析.
\newblock Vector calculus (linear / quadratic forms, chain rule, Hessian): Appendix~A. OLS first-order condition / normal equation: Chapter~3. Local shelf: \texttt{qe-review/applied-math/reference/ding-linear-model/}.
\bibitem[GolubVanLoan(1996)]{golub96}
G.~H. Golub and C.~F. Van~Loan.
\newblock \emph{Matrix Computations}, 3rd~ed.
\newblock Johns Hopkins, 1996.
\newblock Official CAM syllabus item~3. Least squares via Householder QR versus the normal equations: Ch.~5. Kronecker structure / not forming \(B^{T}\otimes A\).
\bibitem[TrefethenBau(1997)]{trefethen97}
L.~N. Trefethen and D.~Bau, III.
\newblock \emph{Numerical Linear Algebra}.
\newblock SIAM, 1997.
\newblock Official CAM syllabus item~7. Householder QR: Lectures~10--11. Stability of least squares, and \(\kappa(A^{T}A)=\kappa(A)^{2}\): Lecture~19. SVD and \(A^{+}\): Lectures~31--32.
\bibitem[Qingshu(11th)]{qingshu11}
清疏.
\newblock \emph{大学生数学竞赛班第十一届讲义}.
\newblock Interpolation remainder / \(K\)-method: Theorem~10.1. Integral remainder and Green kernel: Theorem~10.2, \S10.1.4. Midpoint and trapezoid by integrating the remainder: Exercises~10.4--10.5.
\bibitem[LeVeque(2007)]{leveque07}
R.~J. LeVeque.
\newblock \emph{Finite Difference Methods for Ordinary and Partial Differential Equations}.
\newblock SIAM, 2007.
\newblock Difference quotients and truncation: Ch.~1. Backward Euler / absolute stability: Ch.~2. Advection, CFL, von Neumann, Lax equivalence: Chs.~8--10. Local shelf: \texttt{qe-review/applied-math/reference/fdm-stability/}.
\bibitem[Strikwerda(2004)]{strikwerda04}
J.~C. Strikwerda.
\newblock \emph{Finite Difference Schemes and Partial Differential Equations}.
\newblock SIAM, 2004.
\newblock Discrete maximum principle for elliptic schemes: Ch.~2. Hyperbolic Fourier analysis: Chs.~5--6.
\bibitem[GKO(1995)]{gko95}
B.~Gustafsson, H.-O.~Kreiss, and J.~Oliger.
\newblock \emph{Time Dependent Problems and Difference Methods}.
\newblock Wiley, 1995.
\newblock Official CAM syllabus item~5. Fourier multipliers and the \(1+K\tau\) form of von Neumann stability.
\bibitem[BrennerScott(2010)]{brennerscott10}
S.~Brenner and R.~Scott.
\newblock \emph{The Mathematical Theory of Finite Element Methods}.
\newblock Springer, 2010.
\newblock Official CAM syllabus item~8. Weak form, C\'ea, interpolation in Sobolev spaces. Not the finite-difference language of Problems~4--5.
\bibitem[LaxRichtmyer(1956)]{laxrichtmyer56}
P.~D. Lax and R.~D. Richtmyer.
\newblock Survey of the stability of linear finite difference equations.
\newblock \emph{Comm.\ Pure Appl.\ Math.} 9 (1956), 267--293.
\newblock Lax equivalence: a consistent approximation to a linear well-posed IVP converges if and only if it is stable.
\bibitem[Trefethen(2000)]{trefethen00}
L.~N. Trefethen.
\newblock \emph{Spectral Methods in MATLAB}.
\newblock SIAM, 2000.
\newblock Fourier / Chebyshev collocation and differentiation matrices. The remaining PDE-discretization item on the CAM syllabus.
\bibitem[Grade2ODE]{grade2ode}
乐绎华.
\newblock \emph{常微分方程} chapter of \texttt{grade-2.tex} (大二笔记, after 王高雄第三版).
\newblock Local copy: \texttt{qe-review/applied-math/reference/grade2-ode/ode.tex}.
\newblock Reduction of order when \(t\) is absent; equilibria in the phase plane; the rest of the elementary ODE catalogue.
\bibitem[Grade2Period]{grade2period}
乐绎华.
\newblock \emph{常微分方程周期性}.
\newblock Local copy: \texttt{qe-review/applied-math/reference/grade2-ode/ode-periodicity.tex}.
\newblock Energy identity plus uniqueness \(\Rightarrow\) a periodic orbit, written for \(\ddot x+x+x^{3}=0\).
\bibitem[Nayfeh(1973)]{nayfeh73}
A.~H. Nayfeh.
\newblock \emph{Perturbation Methods}.
\newblock Wiley, 1973.
\newblock Lindstedt--Poincar\'e / strained coordinates: Ch.~4. Frequency correction at \(O(\epsilon^{2})\) for a weakly nonlinear center.
\bibitem[BenderOrszag(1999)]{bender99}
C.~M. Bender and S.~A. Orszag.
\newblock \emph{Advanced Mathematical Methods for Scientists and Engineers}.
\newblock Springer, 1999.
\newblock Official CAM syllabus item~1. Regular perturbation of eigenvalues: Ch.~7. Boundary layers / matched asymptotics: Ch.~9. Multiple scales: Ch.~11. Stationary phase is Ch.~6; it has not yet appeared as a P6-type question.
\bibitem[HairerNorsettWanner(1993)]{hairer93}
E.~Hairer, S.~P. N{\o}rsett, and G.~Wanner.
\newblock \emph{Solving Ordinary Differential Equations I: Nonstiff Problems}.
\newblock Springer, 1993.
\newblock Official CAM syllabus item~4. One-step / Runge--Kutta order and \(R(z)\): Ch.~II. Linear multistep, zero-stability, Dahlquist's first barrier: Ch.~III. Absolute stability: Ch.~IV.
\bibitem[HairerWanner(1996)]{hairer96}
E.~Hairer and G.~Wanner.
\newblock \emph{Solving Ordinary Differential Equations II: Stiff and Differential-Algebraic Problems}.
\newblock Springer, 1996.
\newblock Stiffness, \(A\)-stability, the second Dahlquist barrier, and BDF: Chs.~IV--V.
\bibitem[BoydVandenberghe(2004)]{boyd04}
S.~Boyd and L.~Vandenberghe.
\newblock \emph{Convex Optimization}.
\newblock Cambridge University Press, 2004.
\newblock Dual cone: \S2.6. Separating hyperplane: \S2.5.1. Dual norm: \S A.1.6. Generalized inequality \(x\succeq 0\) for \(\mathbb{R}^{n}_{+}\): \S2.4. The polar of a set is Additional Exercise~2.31 of the companion problem book.
\bibitem[Rockafellar(1970)]{rockafellar70}
R.~T.~Rockafellar.
\newblock \emph{Convex Analysis}.
\newblock Princeton University Press, 1970.
\newblock Polar of a set, polar cone, and the bipolar theorem: \S14.
\end{thebibliography}

\end{document}
